\documentclass[12pt,oneside,openany]{book}
% ============================================================
%  Preamble -- Geometric Anatomy of Theoretical Physics (F. P. Schuller); machinery from the Brunton/Mehrling notes
% ============================================================
\usepackage[T1]{fontenc}
\usepackage[utf8]{inputenc}
\usepackage{amsmath,amsthm,mathtools}
\usepackage{newpxtext,newpxmath} % provides the AMS symbols: do not load amssymb
\usepackage{bm}
\usepackage{microtype}
% Screen layout (31.5x18 cm = Acer Aspire E15 screen minus 15 mm (left/right) and 7 mm (top/bottom) viewer margin, PDF viewer in reading mode, not printed): two parallel columns (paracol), the text column
% 18.8 cm and a 10 cm right column (was 15 + 6.5 until 2026-11-05) for notes, answers, feedback,
% examples, proofs, course notes, history and asides. Every right-column item starts level with the
% paragraph it belongs to; if the right column is still busy, the text waits (white space).
% The old A4 print layout was: a4paper,inner=2.4cm,outer=3.6cm,marginparwidth=3.0cm,marginparsep=0.3cm (twoside).
\ifdefined\scrollmode % continuous-scroll build (scroll.py): tall header-free pages, stacked afterwards
  \usepackage[paperwidth=31.5cm,paperheight=100cm,top=0.3cm,bottom=0.3cm,left=0.9cm,textwidth=29.3cm,
    marginparsep=0.5cm,marginparwidth=10cm]{geometry}
\else
\usepackage[paperwidth=31.5cm,paperheight=18cm,includehead,top=0.4cm,bottom=1.5cm,left=0.9cm,textwidth=29.3cm,
  marginparsep=0.5cm,marginparwidth=10cm,headheight=14pt,headsep=0.15cm]{geometry} % one-line header holds the page number: no footer
\fi
\usepackage{paracol}
\usepackage{needspace} % \Needspace*{10.5cm} before a tall unbreakable right-column box (history) that would run off the page
\setcolumnwidth{18.8cm/0.5cm,10cm}
% counters shared by the two columns (equations in a proof, examples, answers ... in the right column)
% (sectioning counters cannot be global with this kernel: they are copied to the right column by
% \synccounter before each switch, see \side@emit)
\globalcounter{equation}\globalcounter{footnote}\globalcounter{figure}
\usepackage{booktabs,array,tabularx,longtable}
\usepackage[dvipsnames]{xcolor}
\usepackage[shortlabels,inline]{enumitem}
\usepackage{tikz}
\usetikzlibrary{arrows.meta,positioning,decorations.pathreplacing,calc}
\usepackage{pgfplots}
\pgfplotsset{compat=1.18}
\usepackage[most]{tcolorbox}
\usepackage{fancyhdr}
\usepackage[colorlinks=true,linkcolor=NavyBlue,urlcolor=NavyBlue,citecolor=NavyBlue]{hyperref}
\usepackage[nameinlink]{cleveref}
\usepackage{listings}

\setlist{itemsep=2pt,topsep=4pt}
\setlength{\parskip}{3pt}
\setlength{\emergencystretch}{2em}
\hfuzz=5pt

% ---------- headers ----------
\pagestyle{fancy}
\fancyhf{}
\fancyhead[L]{\small\leftmark}
\fancyhead[R]{\small\makebox[\marginparwidth][r]{\rightmark\quad\thepage}}
\renewcommand{\headrulewidth}{0.3pt}
\fancypagestyle{plain}{\fancyhf{}\fancyhead[R]{\small\makebox[\marginparwidth][r]{\thepage}}\renewcommand{\headrulewidth}{0pt}}
\ifdefined\scrollmode
  \raggedbottom \pagestyle{empty}\fancypagestyle{plain}{\fancyhf{}\renewcommand{\headrulewidth}{0pt}}
  % first and last page of the \include d chapter, for scroll.py
  \newcounter{scrollabs}\newwrite\scrollpg\immediate\openout\scrollpg=\jobname.scrollpages
  \AddToHook{shipout/before}{\stepcounter{scrollabs}}
  \AddToHook{include/before}{\immediate\write\scrollpg{\the\numexpr\value{scrollabs}+1\relax}}
  \AddToHook{include/after}{\immediate\write\scrollpg{\the\value{scrollabs}}}
  % a chapter longer than ~3 pages of 100 cm starts a new long page at the next \section (clean seam, PDF viewers
  % stop at ~508 cm); the page number of every break goes to the same file, scroll.py stacks accordingly
  \newcounter{scrollsince}
  \newcommand{\scrollbreak}{\clearpage\immediate\write\scrollpg{\the\numexpr\value{scrollabs}+1\relax}\setcounter{scrollsince}{\value{scrollabs}}}
  \AddToHook{include/before}{\setcounter{scrollsince}{\value{scrollabs}}}
  \AddToHook{cmd/section/before}{\ifnum\numexpr\value{scrollabs}-\value{scrollsince}\relax>2 \scrollbreak\fi}
\fi

% ---------- colours (bright on purpose: yellow notes, blue definitions, red theorems, green examples) ----------
\definecolor{defcol}{HTML}{1565C0}  % blue: definitions
\definecolor{thmcol}{HTML}{D32F2F}  % red: theorems, propositions, lemmas, corollaries
\definecolor{intcol}{HTML}{2E7D32}  % green: examples (also the green of the figures)
\definecolor{notecol}{HTML}{FBC02D} % yellow: margin notes and their answers
\definecolor{intucol}{HTML}{8E24AA} % purple: intuition
\definecolor{fincol}{HTML}{EF6C00}  % orange: applications
\definecolor{cscol}{HTML}{00897B}   % teal: case studies
\definecolor{keycol}{HTML}{C2185B}  % magenta: key formulas
\definecolor{exercol}{HTML}{3949AB} % indigo: exercises
\definecolor{srccol}{HTML}{546E7A}  % grey-blue: sources, course notes, remarks
\definecolor{histcol}{HTML}{B8860B} % bronze: history boxes (classical columns)
\definecolor{recallcol}{HTML}{00ACC1} % cyan: recall boxes (material already seen earlier)
\definecolor{doubtcol}{HTML}{C62828}
\definecolor{ideacol}{HTML}{1565C0}
\definecolor{okcol}{HTML}{388E3C}
\colorlet{anscol}{notecol!85!black}

% ---------- table rows: a little air between lines in every table ----------
\AddToHook{env/tabular/begin}{\renewcommand{\arraystretch}{1.35}}
\AddToHook{env/tabularx/begin}{\renewcommand{\arraystretch}{1.35}}
\AddToHook{env/longtable/begin}{\renewcommand{\arraystretch}{1.35}}

% ---------- theorem-like environments (numbered by chapter), coloured headings ----------
\newtheoremstyle{coldef}{}{}{\normalfont}{}{\bfseries\color{defcol}}{.}{ }{}
\newtheoremstyle{colex}{}{}{\normalfont}{}{\bfseries\color{intcol}}{.}{ }{}
\newtheoremstyle{colexer}{}{}{\normalfont}{}{\bfseries\color{exercol}}{.}{ }{}
\newtheoremstyle{colthm}{}{}{\itshape}{}{\bfseries\color{thmcol}}{.}{ }{}
\newtheoremstyle{colrem}{}{}{\normalfont}{}{\itshape\color{srccol}}{.}{ }{}
\theoremstyle{coldef}
\newtheorem{definition}{Definition}[chapter]
\theoremstyle{colex}
\newtheorem{example}[definition]{Example}
\theoremstyle{colexer}
\newtheorem{exercise}{Exercise}[chapter]
\theoremstyle{colthm}
\newtheorem{theorem}[definition]{Theorem}
\newtheorem{proposition}[definition]{Proposition}
\newtheorem{lemma}[definition]{Lemma}
\newtheorem{corollary}[definition]{Corollary}
\theoremstyle{colrem}
\newtheorem{remark}[definition]{Remark}
\crefname{exercise}{Exercise}{Exercises}
\globalcounter{definition}\globalcounter{exercise}

% Right-column items are queued and written into the right column at the next point where the
% text column is between paragraphs at top level (end of the paragraph, end of a box, end of a
% list): there the columns are synchronised, the queued items are typeset in the right column,
% and the text continues from where it was.
\makeatletter
\newif\ifnote@inbox
\newif\ifside@inright
\def\note@queue{}
\def\side@pcname{paracol}
\newcommand{\side@emit}{\ifx\note@queue\@empty\else
  \let\note@tmp\note@queue\gdef\note@queue{}%
  \synccounter{chapter}\synccounter{section}\synccounter{subsection}\synccounter{subsubsection}%
  \switchcolumn*\side@inrighttrue\note@tmp\par\side@inrightfalse\switchcolumn\fi}
\newcommand{\note@flush}{\ifvmode\ifnote@inbox\else\ifside@inright\else
  \ifx\@currenvir\side@pcname\side@emit\fi\fi\fi\fi}
\newcommand{\side@put}[1]{\gappto\note@queue{#1}\note@flush}
\AddToHook{para/after}{\note@flush}
\tcbset{note queue/.code={\appto\kvtcb@afterbox{\note@flush}}}
% every chapter/appendix file is one paracol environment: it starts right after the \chapter
% line of the file (\begin{paracol}{2}) and is closed here (appendix A, only long tables, has none:
% longtable does not work inside paracol)
\AddToHook{include/end}{\par\note@flush\ifx\@currenvir\side@pcname\end{paracol}\fi}
\newcommand{\flushside}{\par\note@flush}
\makeatother
\tcbset{thmbox/.style={enhanced,breakable,colback=#1!7,colframe=#1!7,borderline west={3pt}{0pt}{#1},
  sharp corners,left=6pt,right=4pt,top=2pt,bottom=2pt,before skip=8pt,after skip=8pt,
  before upper={\csname note@inboxtrue\endcsname},note queue}}
\tcolorboxenvironment{definition}{thmbox=defcol}
\tcolorboxenvironment{theorem}{thmbox=thmcol}
\tcolorboxenvironment{proposition}{thmbox=thmcol}
\tcolorboxenvironment{lemma}{thmbox=thmcol}
\tcolorboxenvironment{corollary}{thmbox=thmcol}
% (examples are boxed in the right column, see \begin{example} below)
\tcolorboxenvironment{exercise}{thmbox=exercol}
\tcolorboxenvironment{remark}{thmbox=srccol}

% ---------- content boxes ----------
\tcbset{notebox/.style={enhanced,breakable,colback=#1!7,colframe=#1,
  fonttitle=\bfseries,coltitle=white,colbacktitle=#1,boxrule=0.6pt,arc=2pt,
  left=6pt,right=6pt,top=4pt,bottom=4pt,before skip=10pt,after skip=10pt,
  before upper={\csname note@inboxtrue\endcsname},note queue}}
% Box title: the box type in capitals, then the topic
\newcommand{\boxtitle}[2]{\MakeUppercase{#1}\ifx&#2&\else: #2\fi}
% Intuition: the "physicist's reading" of a result
\newtcolorbox{intuition}[1][]{notebox=intucol,title={\boxtitle{Intuition}{#1}}}
% Application: where a method is used in science and engineering (orange)
\newtcolorbox{application}[1][]{notebox=fincol,title={\boxtitle{Application}{#1}}}
% Case study: a data set worked through in the lectures (cylinder wake, Lorenz, measles, ...)
\newtcolorbox{casestudy}[1][]{notebox=cscol,title={\boxtitle{Case study}{#1}}}
% Summary / key formulas
\newtcolorbox{keyformulas}{notebox=keycol,fontupper=\small,title={\boxtitle{Key formulas}{}}}
\newtcolorbox{keyconcepts}{notebox=keycol,fontupper=\small,title={\boxtitle{Key concepts}{}}}
% Sources: which lectures, videos and files a chapter is built on
\newtcolorbox{sources}{notebox=srccol,fontupper=\footnotesize,title={\boxtitle{Sources for this chapter}{}},before upper=\raggedright}
\newcolumntype{L}{>{\raggedright\arraybackslash}X}
\newcolumntype{V}{>{\raggedright\arraybackslash}p{4.6cm}} % video titles in the Sources box
% Reading: pointers to the other Schuller notes sets and to his papers
\newtcolorbox{reading}[1][]{notebox=srccol,title={\boxtitle{Reading}{#1}}}

% non-floating figure (floats inside paracol columns can leave a page almost empty): same caption
% and numbering as figure
\makeatletter
\newenvironment{nfigure}{\par\medskip\begin{center}\def\@captype{figure}}{\end{center}\medskip}
\makeatother
% Note on course material (errata, missing lectures, differences 2013/2024): right column, see below
% Secondary material ("beyond the core"): can be skipped on a first reading
\newtcolorbox{secondary}[1][]{enhanced,breakable,colback=white,colframe=white,boxrule=0pt,
  borderline west={1.5pt}{0pt}{srccol!60},sharp corners,left=8pt,right=0pt,top=1pt,bottom=1pt,
  before skip=6pt,after skip=6pt,fontupper=\small,
  title={\footnotesize\sffamily$\diamond$\ \boxtitle{Beyond the core}{#1}},
  coltitle=srccol,colbacktitle=white,fonttitle=\bfseries,titlerule=0pt,toptitle=0pt,bottomtitle=0pt,
  before upper={\csname note@inboxtrue\endcsname},note queue}
% ============================================================
%  The right column: the reader's notes/answers/feedback in \scriptsize; examples, proofs, course notes,
%  history and asides in \footnotesize
% ============================================================
%  \note{...}     the reader's comment/doubt, NEW (not yet read by Claude)
%  \note*{...}    a request that has been carried out (green DONE tag)
%  \note+{...}    a comment that has been read and needs no answer (blue SEEN tag)
%  \begin{answer}{question}[topic] ... \end{answer}
%                 a doubt and its answer in one box: the question on top, then a dark yellow
%                 bar "ANSWER n: topic" and the answer; \answerpart{topic} starts a further
%                 answer in the same box (one per question); \label{ans:...} makes it citable
%  \feedback{...} a comment on one of the reader's own solutions (indigo)
%  \begin{history}[topic], \begin{aside}[topic], \begin{coursenote}[topic]
%  example and proof environments are also printed here (sideexample = example)
%  Inside a coloured box of the main text, margin items appear right after the box.
\makeatletter
\tcbset{sidebox/.style={enhanced,colback=#1!25,colframe=#1,boxrule=0.5pt,arc=1pt,
    left=3pt,right=3pt,top=2pt,bottom=2pt,before skip=0pt,after skip=0pt,
    fontupper=\scriptsize\sffamily\mdseries\upshape\raggedright,fonttitle=\scriptsize\sffamily\bfseries,
    before upper={\note@inboxtrue}}}
\newcommand{\side@tag}[2]{\colorbox{#1}{\textcolor{white}{\bfseries\scriptsize #2}}\par}
% Notes and answered notes are numbered per chapter like footnotes: a small dark-yellow number in
% the text marks the point where the note was written, the same number starts its box.
\newcounter{sidenote}[chapter]
\globalcounter{sidenote}
\newcommand{\side@numput}[1]{\stepcounter{sidenote}%
  \ifhmode\ifdim\lastkern=0.01pt\textsuperscript{\sffamily\bfseries\color{anscol!75!black},}\fi
    \textsuperscript{\sffamily\bfseries\color{anscol!75!black}\thesidenote}\kern0.01pt\fi
  \begingroup\edef\x{\endgroup\noexpand\side@put{\def\noexpand\side@num{\thesidenote}\unexpanded{#1}}}\x}
\newcommand{\side@numtag}{{\bfseries\color{anscol!75!black}\side@num}\enspace}
\NewDocumentCommand{\note}{s t+ +m}{\side@numput{\begin{tcolorbox}[sidebox=notecol]%
  \IfBooleanT{#1}{\side@tag{okcol}{DONE}}\IfBooleanT{#2}{\side@tag{defcol}{SEEN}}\side@numtag#3\end{tcolorbox}}}
\newcommand{\feedback}[1]{\side@put{\begin{tcolorbox}[sidebox=exercol]\side@tag{exercol}{FEEDBACK}#1\end{tcolorbox}}}
\newcounter{answer}[chapter]
\renewcommand{\theanswer}{\thechapter.\arabic{answer}}
\crefname{answer}{answer}{answers}
\globalcounter{answer}
\newcommand{\answerpart}[1]{\refstepcounter{answer}%
  \tcbsubtitle[colback=anscol,colframe=anscol,coltext=black,fontupper=\scriptsize\sffamily\bfseries,
    left=3pt,top=0.5pt,bottom=0.5pt,before skip=2pt,after skip=2pt]{ANSWER \theanswer\ifx&#1&\else: #1\fi}}
\NewDocumentEnvironment{answer}{m O{} +b}{\side@numput{\begin{tcolorbox}[sidebox=notecol,colback=notecol!12]%
  \side@numtag#1\answerpart{#2}#3\end{tcolorbox}}}{}
% Content boxes of the right column (examples, proofs, course notes, history, asides): \footnotesize,
% one step above the \scriptsize of the reader's notes
\tcbset{sidecontent/.style={sidebox=#1,colback=#1!6,breakable,
    fontupper=\footnotesize\mdseries\upshape\raggedright,fonttitle=\footnotesize\sffamily\bfseries,coltitle=white,colbacktitle=#1}}
\NewDocumentEnvironment{aside}{O{} +b}{\side@put{\begin{tcolorbox}[sidecontent=histcol,
  title={\boxtitle{Aside}{#1}}]#2\end{tcolorbox}}}{}
% Recall: a repetition of something explained earlier (another chapter, earlier in this one),
% restated briefly with a \cref to the original, so that it is visibly a repetition.
% Longer repeated passages go verbatim to a "Set aside" section at the end of the chapter.
\NewDocumentEnvironment{recall}{O{} +b}{\side@put{\begin{tcolorbox}[sidecontent=recallcol,
  title={\boxtitle{Recall}{#1}}]#2\end{tcolorbox}}}{}
\NewDocumentEnvironment{coursenote}{O{} +b}{\side@put{\begin{tcolorbox}[sidecontent=srccol,
  title={\boxtitle{Course note}{#1}}]#2\end{tcolorbox}}}{}
% Examples and proofs keep their theorem-like look (numbering, labels, QED) but go to the right column
\let\origexample\example \let\endorigexample\endexample
\RenewDocumentEnvironment{example}{o +b}{\side@put{\begin{tcolorbox}[thmbox=intcol,
  fontupper=\footnotesize,before skip=0pt,after skip=0pt]%
  \IfValueTF{#1}{\begin{origexample}[#1]}{\begin{origexample}}#2\end{origexample}\end{tcolorbox}}}{}
\NewDocumentEnvironment{sideexample}{o +b}{\IfValueTF{#1}{\begin{example}[#1]}{\begin{example}}#2\end{example}}{}
\let\origproof\proof \let\endorigproof\endproof
% A proof must start level with the top of its theorem: a theorem directly followed by \begin{proof}
% is held back (\thm@pend) and typeset by the proof, right after the proof has been queued.
\def\thm@pend{}
\newtoks\thm@tk
\newcommand{\thm@typeset}{{\edef\@currenvir{\side@pcname}\begin{tcolorbox}[thmbox=thmcol]\expandafter\IfValueTF\expandafter{\thm@opt}%
  {\begin{orig\thm@cur}[\thm@opt]\the\thm@tk\end{orig\thm@cur}}{\begin{orig\thm@cur}\the\thm@tk\end{orig\thm@cur}}\end{tcolorbox}}}
\newcommand{\thm@peek}[3]{\def\thm@a{proof}\def\thm@b{#3}%
  \ifx\thm@a\thm@b\gdef\thm@pend{\thm@typeset}\else\thm@typeset\fi#1#2{#3}}
\newcommand{\thm@look}[2]{\@ifnextchar\begin{\thm@peek{#1{#2}}}{\thm@typeset#1{#2}}}
\RenewDocumentEnvironment{proof}{o +b}{\par{\edef\@currenvir{\side@pcname}\side@put{\begin{tcolorbox}[thmbox=thmcol,colback=thmcol!3,
  fontupper=\footnotesize,before skip=0pt,after skip=0pt]%
  \IfValueTF{#1}{\begin{origproof}[#1]}{\begin{origproof}}#2\end{origproof}\end{tcolorbox}}%
  \thm@pend}\gdef\thm@pend{}}{}
\@for\thm@n:=theorem,proposition,lemma,corollary\do{%
  \RemoveFromHook{env/\thm@n/before}[tcolorbox]\RemoveFromHook{env/\thm@n/after}[tcolorbox]%
  \expandafter\let\csname orig\thm@n\expandafter\endcsname\csname\thm@n\endcsname
  \expandafter\let\csname endorig\thm@n\expandafter\endcsname\csname end\thm@n\endcsname
  \expandafter\edef\csname thm@def@\thm@n\endcsname{\noexpand\RenewDocumentEnvironment{\thm@n}{o +b}%
    {\noexpand\global\noexpand\thm@tk{##2}\noexpand\gdef\noexpand\thm@cur{\thm@n}\noexpand\gdef\noexpand\thm@opt{##1}\noexpand\thm@look}{}}%
  \csname thm@def@\thm@n\endcsname}
\makeatother
% History: a fact from the history of finance/mathematics, framed like a classical
% temple front: a dentilled entablature (the title) on two Ionic columns (not breakable)
\newcommand{\histcolumn}[2]{% #1 = foot of the column, #2 = top of the capital (coordinate names)
  \fill[histcol] ([xshift=-6pt]#1) rectangle ([xshift=6pt,yshift=2.5pt]#1);
  \fill[histcol] ([xshift=-4.5pt,yshift=2.5pt]#1) rectangle ([xshift=4.5pt,yshift=4pt]#1);
  \fill[histcol!15] ([xshift=-3.5pt,yshift=4pt]#1) rectangle ([xshift=3.5pt,yshift=-5pt]#2);
  \foreach \dx in {-3.5,-1.2,1.2,3.5}
    \draw[histcol,line width=0.4pt] ([xshift=\dx pt,yshift=4pt]#1) -- ([xshift=\dx pt,yshift=-5pt]#2);
  \fill[histcol] ([xshift=-4.5pt,yshift=-5pt]#2) rectangle ([xshift=4.5pt,yshift=-2.5pt]#2);
  \foreach \dx in {-4.5,4.5} {
    \draw[histcol,line width=0.8pt,fill=white] ([xshift=\dx pt,yshift=-4.5pt]#2) circle (1.9pt);
    \fill[histcol] ([xshift=\dx pt,yshift=-4.5pt]#2) circle (0.6pt);}
  \fill[histcol] ([xshift=-6.5pt,yshift=-2.5pt]#2) rectangle ([xshift=6.5pt]#2);}
\makeatletter
\NewDocumentEnvironment{history}{O{} +b}{\side@put{\begin{tcolorbox}[histbox,title={\boxtitle{History}{#1}}]#2\end{tcolorbox}}}{}
\makeatother
\tcbset{histbox/.style={enhanced,colback=histcol!6,colframe=histcol,boxrule=0.8pt,sharp corners,
  fonttitle=\footnotesize\sffamily\bfseries,fontupper=\normalfont\footnotesize\raggedright,coltitle=white,colbacktitle=histcol,
  left=18pt,right=18pt,top=6pt,bottom=3pt,before skip=0pt,after skip=0pt,
  before upper={\csname note@inboxtrue\endcsname},
  overlay={
    \draw[histcol,line width=2.5pt,dash pattern=on 2.5pt off 2pt]
      ([yshift=-1.25pt]title.south west) -- ([yshift=-1.25pt]title.south east);
    \coordinate (hfL) at ([xshift=9pt]frame.south west);
    \coordinate (htL) at ([xshift=9pt,yshift=-2.5pt]title.south west);
    \coordinate (hfR) at ([xshift=-9pt]frame.south east);
    \coordinate (htR) at ([xshift=-9pt,yshift=-2.5pt]title.south east);
    \histcolumn{hfL}{htL}\histcolumn{hfR}{htR}}}}
% The reader's own solution to an exercise, placed right below the exercise
%   \begin{mysolution}[optional topic] ... \end{mysolution}
\newtcolorbox{mysolution}[1][]{notebox=exercol,colback=exercol!5,fontupper=\small,
  title={\boxtitle{My solution}{#1}}}
% A request about the notes themselves (collected in the "Requests" chapter)
%   \request{done|open|noted}{where it was written}{the request}{what was done}
\makeatletter
\newcommand{\request}[4]{\ifstrequal{#1}{done}{\def\req@col{okcol}}{\ifstrequal{#1}{open}{\def\req@col{thmcol}}{\def\req@col{defcol}}}%
  \begin{tcolorbox}[sidebox=notecol,colback=notecol!12,fontupper=\small\sffamily\raggedright,
  before skip=5pt,after skip=5pt]%
  \side@tag{\req@col}{\MakeUppercase{#1}}%
  {\itshape #2.}\ #3%
  \tcbsubtitle[colback=anscol,colframe=anscol,coltext=black,fontupper=\footnotesize\sffamily\bfseries,
    left=3pt,top=0.5pt,bottom=0.5pt,before skip=2pt,after skip=2pt]{WHAT WAS DONE}#4\end{tcolorbox}}
\makeatother



% ---------- video timestamps ----------
%  \ts{4}{1}{3:10}  -> lecture 4 (one video per lecture: the second argument is always 1), at 3:10
%  Each tag is a link to the YouTube video at that second. The IDs (playlist of the 28 lectures)
%  are in videos.tex, generated by ../tools/transcripts.py;
%  an unknown video or a non-numeric time prints the tag without link.
\ExplSyntaxOn
\cs_new_protected:Npn \tsvideo #1#2 { \tl_const:cx { c__ts_#1_tl } { \tl_to_str:n {#2} } }
\int_new:N \l__ts_int
\seq_new:N \l__ts_seq
\cs_new_protected:Npn \tslink #1#2#3#4
  {
    \tl_if_exist:cTF { c__ts_#1@#2_tl }
      { \regex_match:nnTF { \A\d+(:\d+)*\Z } {#3} { \__ts_href:nnnn {#1} {#2} {#3} {#4} } {#4} }
      {#4}
  }
\cs_new_protected:Npn \__ts_href:nnnn #1#2#3#4
  {
    \int_zero:N \l__ts_int
    \exp_args:NNV \seq_set_split:Nnn \l__ts_seq \c_colon_str {#3}
    \seq_map_inline:Nn \l__ts_seq { \int_set:Nn \l__ts_int { 60*\l__ts_int + ##1 } }
    \exp_args:Ne \href { \tl_to_str:n { https://youtu.be/ } \tl_use:c { c__ts_#1@#2_tl } \tl_to_str:n { ?t= } \int_use:N \l__ts_int } {#4}
  }
\ExplSyntaxOff
% Generated by tools/transcripts.py: YouTube ID of every lecture (lecture@1)
\tsvideo{1@1}{V49i_LM8B0E} % Introduction/Logic of propositions and predicates- 01 - Frederic Schuller
\tsvideo{2@1}{AAJB9l-HAZs} % Axioms of set Theory - Lec 02 - Frederic Schuller
\tsvideo{3@1}{6EIWRg-6ftQ} % Classification of sets - Lec 03 - Frederic Schuller
\tsvideo{4@1}{1wyOoLUjUeI} % Topological spaces - construction and purpose - Lec 04 - Frederic Schuller
\tsvideo{5@1}{hiD6Tz06k30} % Topological spaces - some heavily used invariants - Lec 05 - Frederic Schuller
\tsvideo{6@1}{uGEV0Wk0eIk} % Topological manifolds and manifold bundles- Lec 06 - Frederic Schuller
\tsvideo{7@1}{Fa6SRAwY80Y} % Differentiable structures  definition and classification - Lec 07 - Frederic Schuller
\tsvideo{8@1}{4l-qzZOZt50} % Tensor space theory I: over a field - Lec 08 - Frederic P Schuller
\tsvideo{9@1}{UPGoXBfm6Js} % Differential structures: the pivotal concept of tangent vector spaces - Lec 09 - Frederic Schuller
\tsvideo{10@1}{XZcKSoI17r0} % Construction of the tangent bundle - Lec 10 - Frederic Schuller
\tsvideo{11@1}{V0TPgeiyWCo} % Tensor space theory II: over a ring - Lec 11 - Frederic Schuller
\tsvideo{12@1}{QLnzIOGIvfo} % Grassmann algebra and deRham cohomology - Lec 12 - Frederic Schuller
\tsvideo{13@1}{mJ8ZDdA10GY} % Lie groups and their Lie algebras - Lec 13 - Frederic Schuller
\tsvideo{14@1}{Vlbcd_lPNMA} % Classification of Lie algebras and Dynkin diagrams - Lec 14 - Frederic Schuller
\tsvideo{15@1}{H1D09cuFWlM} % The Lie group SL(2,C) and its Lie algebra sl(2,C) - lec 15 - Frederic Schuller
\tsvideo{16@1}{G9uVcit_VwY} % Dynkin diagrams from Lie algebras, and vice versa - Lec 16 - Frederic Schuller
\tsvideo{17@1}{h-d8TUg022A} % Representation theory of Lie groups and Lie algebras - Lec 17 - Frederic Schuller
\tsvideo{18@1}{7qO5y6Es9Ns} % Reconstruction of a Lie group from its algebra - Lec 18 - Frederic Schuller
\tsvideo{19@1}{vYAXjTGr_eM} % Principal fibre bundles - Lec 19 - Frederic Schuller
\tsvideo{20@1}{q2GYZz6q3QI} % Associated fibre bundles - Lec 20 - Frederic Schuller
\tsvideo{21@1}{jFvyeufXyW0} % Conncections and connection 1-forms - Lec 21 - Frederic Schuller
\tsvideo{22@1}{KhagmmNvOvQ} % Local representations of a connection on the base manifold: Yang-Mills fields - Lec 22
\tsvideo{23@1}{jGHaZc2fuX8} % Parallel transport - Lec 23 - Frederic Schuller
\tsvideo{24@1}{j36o4DLLK2k} % Curvature and torsion on principal bundles - Lec 24 - Frederic Schuller
\tsvideo{25@1}{ClIVG7ilm_Q} % Covariant derivatives - Lec 25 - Frederic Schuller
\tsvideo{26@1}{C93KzJ7-Es4} % Application: Quantum mechanics on curved spaces - Lec 26 - Frederic Schuller
\tsvideo{27@1}{Way8FfcMpf0} % Application: Spin structures - lec 27 - Frederic Schuller
\tsvideo{28@1}{F3oGhXNhIDo} % Application: Kinematical and dynamical symmetries - Lec 28 - Frederic Schuller

\makeatletter % consecutive tags may break between them, but never before punctuation
\DeclareRobustCommand{\ts}[3]{\mbox{\footnotesize\sffamily\color{srccol}\tslink{#1}{#2}{#3}{\color{srccol}[$\blacktriangleright$\,L#1\,@\,#3]}}\@ifnextchar\ts{\allowbreak}{}}
% link to a whole video (no time), e.g. in the Sources box: \seg{4}{2}{text}
\newcommand{\seg}[3]{\tslink{#1}{#2}{0}{#3}}
\makeatother

% ---------- maths shortcuts ----------
\newcommand{\R}{\mathbb{R}}
\newcommand{\N}{\mathbb{N}}
\newcommand{\E}{\mathbb{E}}
\newcommand{\Prob}{\mathbb{P}}
\newcommand{\Q}{\mathbb{Q}}
\DeclareMathOperator{\Var}{Var}
\DeclareMathOperator{\Cov}{Cov}
\DeclareMathOperator{\Corr}{Corr}
\DeclareMathOperator{\rank}{rank}
\DeclareMathOperator{\diag}{diag}
\newcommand{\C}{\mathbb{C}}
\newcommand{\K}{\mathcal{K}}  % Koopman operator
\newcommand{\vb}[1]{\mathbf{#1}} % vectors and matrices of data: \vb x, \vb X
\newcommand{\dd}{\,\mathrm{d}}
\newcommand{\ind}{\mathbf{1}}
\newcommand{\eqd}{\stackrel{d}{=}}

% ---------- code: the Matlab demos of the lectures, rewritten in Python ----------
%  \begin{code}[topic] ... \end{code}   (a listing in a grey-blue box; must be in the text column)
\lstdefinestyle{py}{language=Python,basicstyle=\ttfamily\footnotesize,columns=fullflexible,keepspaces,
  keywordstyle=\color{defcol}\bfseries,commentstyle=\color{intcol}\itshape,stringstyle=\color{fincol},
  showstringspaces=false,breaklines=true,upquote=true}
\newtcblisting{code}[1][]{notebox=srccol,colback=srccol!4,listing only,listing options={style=py},
  title={\boxtitle{Python}{#1}}}
%  \codefile{topic}{../code/chNN_name.py}: the same box, reading a script of ../code (run it from the
%  repository root to regenerate the figures in figures/), so notes and scripts never diverge
\newcommand{\codefile}[2]{\tcbinputlisting{notebox=srccol,colback=srccol!4,listing only,
  listing options={style=py},listing file={#2},title={\boxtitle{Python}{#1}~{\normalfont\ttfamily\footnotesize(\detokenize{#2})}}}}

% ---------- file references ----------
\DeclareRobustCommand{\file}[1]{\texttt{\small\detokenize{#1}}}


% Compile only some chapters while working on them, e.g.:
% \includeonly{chapters/ch01-logic}

\begin{document}

\frontmatter
\begin{titlepage}
  \centering
  \vspace*{3cm}
  {\Huge\bfseries The Geometric Anatomy of Theoretical Physics\par}
  \vspace{1cm}
  {\Large Lecture notes\par}
  \vspace{0.5cm}
  {\large based on the 28 video lectures by Frederic P.~Schuller\\
  (Friedrich-Alexander-Universit\"at Erlangen-N\"urnberg, winter semester 2013/14)\par}
  \vspace{2cm}
  {\large Personal study notes of Alessio Martini\par}
  \vfill
  \begin{minipage}{0.85\textwidth}\small
    These notes are a personal, non-commercial study aid derived from the recorded lectures
    \emph{Lectures on the Geometric Anatomy of Theoretical Physics} by Frederic P.~Schuller
    (YouTube and \url{https://www.fau.tv/series/lectures-on-the-geometric-anatomy-of-theoretical-physics-w13}).
    Lecture content is paraphrased from the (automatic) transcripts of the videos, not transcribed;
    the blackboard notation is reconstructed from what is said, and doubtful points are marked~[?].
    Any mistakes are mine. All rights on the original lectures belong to their author.
  \end{minipage}
  \vspace{1cm}
\end{titlepage}

\chapter*{How to use these notes}
\addcontentsline{toc}{chapter}{How to use these notes}
\begin{paracol}{2}

There is one chapter per lecture, and the sections of a chapter follow the lecture in the order in which
Schuller presents them. Every chapter opens with a \textbf{Sources} box (the video, its length, what it
contains) and a \textbf{Reading} box pointing to the matching chapters of the other Schuller notes sets.
Definitions, theorems and proofs are written as on his blackboard; everything else he says (motivation,
remarks, digressions, jokes with a point) is kept, the side remarks in the right column.

\paragraph{Boxes.}
\begin{itemize}
  \item \textcolor{defcol}{\textbf{Definitions}} (blue), \textcolor{thmcol}{\textbf{theorems, propositions, corollaries}} (red),
        \textcolor{exercol}{\textbf{exercises}} and \textcolor{srccol}{\textbf{remarks}} are marked with a
        coloured bar on the left. A definition box titled \emph{New concepts} marks the first time a concept appears.
  \item \textcolor{keycol}{\textbf{Key concepts}} at the end of each chapter: retrieval questions, then the concepts.
  \item \textcolor{intucol}{\textbf{Intuition}} boxes give the picture behind a definition.
  \item \textcolor{srccol}{\textbf{Reading}} boxes point to the other notes sets and to Schuller's papers.
  \item \textbf{Beyond the core} blocks (grey bar on the left, smaller type) hold material that can be
        skipped on a first reading.
\end{itemize}

\paragraph{The right column.} The page is laid out for a screen: the text is always on the left and a
column on the right holds \textcolor{intcol}{\textbf{examples}}, \textcolor{thmcol}{\textbf{proofs}},
\textcolor{srccol}{\textbf{course notes}}, \textcolor{histcol}{\textbf{history}} boxes and
\textcolor{histcol}{\textbf{asides}}, your \textcolor{notecol!80!black}{\textbf{notes}} with their
\textcolor{anscol}{\textbf{answers}}, and \textcolor{exercol}{\textbf{feedback}} on your solutions.

\paragraph{Video timestamps.} A tag such as \ts{1}{1}{28:22}\ means: lecture~1, at minute 28:22 of the
video. Every tag is a link to the YouTube video at that second.

\paragraph{Your annotations.} While reading, write directly in the chapter files. The macro
\verb|\note{...}| prints a comment in the right column next to the text where it is written. Requests
about the notes themselves are collected in the chapter \emph{Requests about these notes}.

\flushside
\end{paracol}

% !TEX root = ../main.tex
\chapter*{Requests about these notes}\label{ch:requests}
\addcontentsline{toc}{chapter}{Requests about these notes}
\markboth{Requests about these notes}{}
\begin{paracol}{2}

Your margin notes are of two kinds: notes on the \emph{content} (a doubt, an idea), which stay next to the
text they refer to, and requests about the \emph{notes themselves} (layout, boxes, conventions), which are
collected here with what was done about them.

\paragraph{How the margin notes are marked.}
\begin{itemize}
  \item \verb|\note{...}|: a \emph{new} note, not yet read. ``Read my notes'' means these, and only these.
  \item A note that has been answered becomes one box with the question on top and a dark yellow
        \textbf{ANSWER} bar below it (\verb|\begin{answer}{question}[topic]| \dots\ \verb|\end{answer}|).
  \item \verb|\note*{...}|: a request that has been carried out (green DONE tag).
  \item \verb|\note+{...}|: a comment that has been read and needs no answer (blue SEEN tag).
  \item \verb|\feedback{...}|: a comment on one of your own solutions (indigo).
\end{itemize}

\section*{The list}
No requests yet.

\flushside
\end{paracol}


\tableofcontents

\mainmatter
\part{Foundations: logic and sets}\label{part:foundations}
% BEGIN CHAPTERS (tools/stubs.py)
% !TEX root = ../main.tex
\chapter{Introduction; logic of propositions and predicates}\label{ch:logic}
\begin{paracol}{2}

\begin{sources}
\begin{tabularx}{\linewidth}{@{}l l L@{}}
\toprule
\textbf{Video} & \textbf{Length} & \textbf{Content}\\
\midrule
\seg{1}{1}{Lecture 1: Introduction / Logic of propositions and predicates} & 1:40:48 & Aim and structure of the
course; propositional logic (unary and binary operators, implication, proof by contradiction, NAND);
predicate logic (quantifiers, order of quantification); axiomatic systems, proofs, consistency;
consistency of propositional logic; G\"odel's incompleteness theorem (statement and idea); outlook on
axiomatic set theory.\\
\bottomrule
\end{tabularx}

\smallskip
\textbf{Transcript} \file{materials/transcripts/C01-V01.txt} (YouTube automatic captions).
Recorded at FAU Erlangen, winter semester 2013/14.
\end{sources}

\begin{reading}[Where this lecture sits in the other notes]
None of the other Schuller courses starts from logic: \emph{Gravity and Light} (Lecture~1) and
\emph{Theoretische Mechanik} (Lecture~2) start directly from topology, which here is Lecture~4.
The overview of the ``building'' in \cref{sec:lo-structure} is the same picture as the bird's-eye view of
\emph{Theoretische Mechanik}, Lecture~1. Links to the chapters will be added when those notes exist.
\end{reading}

% ==================================================================
\section{What the course is for}\label{sec:lo-aims}

Theoretical physics casts our concepts about the real world into rigorous mathematical form, ``for better or
worse''; it does so not for its own sake, but to explore fully what those concepts imply \ts{1}{1}{0:07}.
Schuller puts the spirit of the subject in the words of Wittgenstein: \emph{what we cannot speak about clearly,
we must pass over in silence} \ts{1}{1}{0:51}. If a concept about the world cannot be put into precise
mathematical form, this usually shows that some aspect of the concept has not been understood.
\begin{coursenote}[The two Wittgenstein quotes]
The first is the last proposition (7) of the \emph{Tractatus Logico-Philosophicus} (1921): ``Wovon man nicht
sprechen kann, dar\"uber muss man schweigen.'' The second is close to Tractatus 5.43, where it is said of the
propositions of \emph{logic}: ``Alle S\"atze der Logik sagen aber dasselbe. N\"amlich nichts.''
\end{coursenote}

Mathematics, however, is \emph{just a language}. To extract physical conclusions from a mathematical
formulation one must interpret it, and that is not the task of mathematics but of physicists, which is where
it becomes difficult again \ts{1}{1}{1:32}. Again with Wittgenstein: the theorems of mathematics all say the
same thing, namely nothing \ts{1}{1}{1:52}. This does not mean that mathematics is useless. A theorem of the
form ``$A$ if and only if $B$'' says nothing that $A$ did not already say: logically it is a tautology. But
psychologically, for our understanding of $A$, a reformulation in terms of $B$ may be very useful
\ts{1}{1}{2:13}.

With this understanding, the idea of the course is to provide the proper language for theoretical physics
\ts{1}{1}{2:56}.

\begin{intuition}[The aim in one line]
Provide proper mathematical language for classical mechanics, electromagnetism, quantum mechanics and
statistical physics \ts{1}{1}{3:19}.
\end{intuition}

These are the subjects of the modules TP1 to TP4 at Erlangen \ts{1}{1}{4:04}. The course will not revise all
the mathematics they need; it develops the mathematics from a higher point of view, now that some knowledge
of these subjects can be assumed. So far they have been studied with a fair amount of \emph{analysis}, of
\emph{algebra} (mostly linear, sometimes non-linear) and of \emph{geometry} in one form or another
\ts{1}{1}{4:25}: classical mechanics and electromagnetism appeal to geometric intuition to set up and solve a
problem; quantum mechanics, in its modern form, is done with linear algebra; statistical physics uses some
geometry, maybe without noticing, when it looks at convex cones of states \ts{1}{1}{5:06}.

Drawing a circle around each of the three fields, the physical theories sit roughly in the intersections
\ts{1}{1}{5:26}, as in \cref{fig:lo-venn}. Classical mechanics and electrodynamics use analysis and geometry,
and so does general relativity; quantum mechanics lies in the intersection of analysis and algebra, usually
called \emph{functional analysis}; statistical physics, ``to do it justice'', lies more or less between algebra
and geometry \ts{1}{1}{5:51}. This is not a pure picture, only where the intersections come into play. In the
middle, where all three intersect, sits \emph{differential geometry}, and it is mainly differential geometry
that the course is about \ts{1}{1}{6:33}.

\begin{nfigure}
\begin{tikzpicture}[scale=0.95]
  \draw[thick,defcol] (0,0) circle (2.1) node[above left=1.05cm and 0.6cm] {\textbf{Analysis}};
  \draw[thick,thmcol] (2.6,0) circle (2.1) node[above right=1.05cm and 0.6cm] {\textbf{Algebra}};
  \draw[thick,intcol] (1.3,-2.2) circle (2.1) node[below=1.35cm] {\textbf{Geometry}};
  \node[font=\footnotesize,align=center] at (1.3,1.25) {QM\\(functional\\analysis)};
  \node[font=\footnotesize,align=center] at (-0.45,-1.45) {CM, EM,\\GR};
  \node[font=\footnotesize,align=center] at (3.05,-1.45) {statistical\\physics};
  \node[font=\footnotesize\bfseries,align=center] at (1.3,-0.65) {differential\\geometry};
\end{tikzpicture}
\caption{Where the physical theories sit among the three mathematical fields \ts{1}{1}{5:51}. The course is
about the centre.}\label{fig:lo-venn}
\end{nfigure}

% ==================================================================
\section{The structure of the course: a building}\label{sec:lo-structure}

The development of differential geometry starts with sets \ts{1}{1}{7:17}. Every space of physics (the
configuration space or the phase space of classical mechanics, the space in which electromagnetism works, the
physical space over which the Hilbert space of quantum mechanics is built) is, at its coarsest level, a set of
points \ts{1}{1}{7:38}. But what precisely is a set? Saying ``a collection of elements'' only raises the
questions of what a collection is and what elements are. There is a fundamental problem: if we start writing a
book on mathematics whose pages are still empty, what can the first definition be? A definition needs notions
we already have \ts{1}{1}{7:58}. The way out is to start \emph{axiomatically}; but then, in what language are
the axioms written? So below set theory there is one more building block: \emph{logic}, propositional and
predicate logic, which defines our language \ts{1}{1}{8:40}. The building then goes up as follows
(\cref{fig:lo-building}).

\begin{enumerate}
  \item \textbf{Logic and axiomatic set theory}: the language and the axioms for sets \ts{1}{1}{9:02}.
  \item \textbf{Topology.} The four subjects need a notion of \emph{continuity}: famously, particles do not
        jump; a trajectory does not stop somewhere and start again somewhere else. Continuity is obtained by
        distinguishing certain subsets of the basic set, which leads to topological spaces \ts{1}{1}{9:23}.
  \item \textbf{Topological manifolds.} There are so many topological spaces that they cannot even be
        classified, so we restrict to those that are immediately interesting for physics (topologists may
        disagree that they are the only interesting ones): spaces that locally around each point look like
        some $\R^d$ \ts{1}{1}{10:04}.
  \item \textbf{Differentiable manifolds.} Classical physics needs derivatives, for instance the velocity of a
        curve, which cannot be defined on a topological manifold. Differentiable manifolds are ``an entirely
        new beast'' and need the technology of \emph{charts}, which makes contact with the formalism of all
        the standard courses: these have all been formulated in charts so far \ts{1}{1}{10:46}.
  \item \textbf{Bundles}: to each point of a manifold one attaches another manifold, for instance a vector
        space at each point (the tangent space) \ts{1}{1}{12:52}.
  \item \textbf{Geometry}: tensor fields on these structures \ts{1}{1}{14:58}.
  \item \textbf{Physics} on top: classical mechanics, electromagnetism, quantum mechanics, statistical physics,
        special and general relativity \ts{1}{1}{16:02}.
\end{enumerate}

\begin{intuition}[Charts and atlases]
A chart is what a sailor uses: travelling across the globe, you take out the atlas, turn to the chart of the
place where you are and navigate in it. The chart must be a one-to-one picture of the world locally. If you
move so that you would leave the page, you turn to another page that slightly overlaps with the first and
navigate on \ts{1}{1}{11:30}. Looking for overlapping charts and understanding how the overlaps work is the
whole philosophy of manifolds, and of differentiable manifolds in particular: the charts are pieces of
$\R^d$, where differentiation is understood, and the notion of differentiation is lifted from the charts to
the manifold \ts{1}{1}{12:11}.
\end{intuition}

Bundles are a very important concept in differential geometry, and much of modern physics can be cast in
their language: particle physics, gravity, and certainly quantum mechanics \ts{1}{1}{13:14}. Two warnings come
with them \ts{1}{1}{13:35}:
\begin{itemize}
  \item there is \emph{no such thing as a position vector}, however often it is repeated: a position vector
        is just a collection of coordinates, and these do not constitute a vector;
  \item there is \emph{no such thing as a wave function} in quantum mechanics: it is not a function at all,
        but a section of a complex line bundle over physical space \ts{1}{1}{14:16}.
\end{itemize}
One gets away with both notions as long as one does not look too closely; then they fail. ``You cannot
understand this now, but you will at the end of the course'' \ts{1}{1}{14:37}.

The geometry comes as tensor fields on these structures \ts{1}{1}{14:58}: a \emph{symplectic} geometry, the
geometry of the phase space of classical mechanics, and in fact of virtually all dynamical systems of interest
\ts{1}{1}{15:19}; a \emph{metric} geometry, needed to some extent in electromagnetism (there, really, only a
volume form is needed), and certainly in special and general relativity, which need a so-called
\emph{Lorentzian} metric \ts{1}{1}{15:41}. Up to the geometry everything is ``physical mathematics''; above it
is physics, and all of it can be properly understood by mastering this mathematics \ts{1}{1}{16:23}. The
course develops the mathematics and then shows, mainly by examples, how the insights play out at the level of
physics \ts{1}{1}{16:43}. The picture should not be taken too seriously (bundles, for instance, sometimes come
before or alongside differentiable manifolds) \ts{1}{1}{14:58}.

\begin{nfigure}
\begin{tikzpicture}[every node/.style={draw,minimum width=9.5cm,minimum height=0.62cm,font=\small,align=center}]
  \node[fill=thmcol!10] (ph) at (0,6.2) {\textbf{Physics}: CM, EM, QM, statistical physics, SR, GR};
  \node[fill=intcol!10] (ge) at (0,5.3) {\textbf{Geometry}: tensor fields (symplectic, metric, Lorentzian, volume form)};
  \node[fill=defcol!8]  (bu) at (0,4.4) {\textbf{Bundles} (no position vectors, no wave functions)};
  \node[fill=defcol!8]  (dm) at (0,3.5) {\textbf{Differentiable manifolds} (charts, atlases)};
  \node[fill=defcol!8]  (tm) at (0,2.6) {\textbf{Topological manifolds} (locally $\R^d$)};
  \node[fill=defcol!8]  (to) at (0,1.7) {\textbf{Topological spaces} (continuity)};
  \node[fill=srccol!10] (se) at (0,0.8) {\textbf{Axiomatic set theory}};
  \node[fill=srccol!20] (lo) at (0,-0.1) {\textbf{Logic}: propositional and predicate logic};
\end{tikzpicture}
\caption{The building of the course, from the bottom up \ts{1}{1}{8:40}. Up to geometry it is mathematics;
the top floor is physics.}\label{fig:lo-building}
\end{nfigure}

% ==================================================================
\section{Propositional logic}\label{sec:lo-prop}

Logic is treated as the first part of Chapter~1, \emph{Axiomatic set theory}, because it is the language
in which set theory will be written \ts{1}{1}{17:39}.

\begin{definition}[New concepts: proposition]\label{def:lo-proposition}
A \textbf{proposition} $p$ is a variable that can take the values \emph{true} or \emph{false}, and no others
\ts{1}{1}{18:32}.
\end{definition}
The name comes from the Latin \emph{propositio}, ``a setting forth'', a statement put forward.

That is all a proposition is for propositional logic. In particular, it is not its task to decide whether a
complex statement such as ``there is extraterrestrial life'' is true or not: propositional logic starts from
complete propositions and only assumes that each is true or false \ts{1}{1}{19:14}. Nor is it its task to
decide whether a sentence such as ``in winter it is colder than outside'' is a proposition at all (this one does
not seem particularly reasonable) \ts{1}{1}{19:56}. Propositions are its atomic, elementary notion
\ts{1}{1}{20:16}.

New propositions are built from given ones by \textbf{logical operators} \ts{1}{1}{20:38}.

\paragraph{Unary operators.} A unary operator takes one proposition $p$ and makes a new proposition from it
\ts{1}{1}{21:05}. Since $p$ can only be true or false, there are exactly four of them \ts{1}{1}{21:47}: the
negation $\neg p$, the identity $\mathrm{id}\,p$, the tautology $\top p$ (always true, whatever $p$ is) and the
contradiction $\bot p$ (always false) \ts{1}{1}{22:07}:
\begin{center}
\renewcommand{\arraystretch}{1.35}
\begin{tabular}{c|cccc}
$p$ & $\neg p$ & $\mathrm{id}\,p$ & $\top p$ & $\bot p$\\\hline
T & F & T & T & F\\
F & T & F & T & F
\end{tabular}
\end{center}
These are all the possibilities \ts{1}{1}{22:28}.

\paragraph{Binary operators.} A binary operator takes two given propositions $p,q$; the pair can take the
four values TT, TF, FT, FF \ts{1}{1}{22:48}. For each of the four rows we may choose T or F, so there are
$2^4=16$ binary operators \ts{1}{1}{23:29}. Some of them:
\begin{center}
\renewcommand{\arraystretch}{1.35}
\begin{tabular}{cc|ccccc}
$p$ & $q$ & $p\wedge q$ & $p\vee q$ & $p\veebar q$ & $p\Rightarrow q$ & $p\Leftrightarrow q$\\\hline
T & T & T & T & F & T & T\\
T & F & F & T & T & F & F\\
F & T & F & T & T & T & F\\
F & F & F & F & F & T & T
\end{tabular}
\end{center}
with $\wedge$ ``and'' \ts{1}{1}{23:09}, $\vee$ ``or'' (inclusive) and $\veebar$ the exclusive or
\ts{1}{1}{25:33}.

\begin{coursenote}[Why not ``maps'' from $\{T,F\}\times\{T,F\}$?]
Asked whether a binary operator is a map from the Cartesian product of $\{\text{true},\text{false}\}$ with
itself to $\{\text{true},\text{false}\}$, Schuller answers that this is the idea, but he cannot say it that way:
we do not yet know what a set or a Cartesian product is \ts{1}{1}{24:10}. Even more abstractly, one can
write down \emph{deduction rules} and never write a table. He does not, because for standard logic it would
serve no purpose \ts{1}{1}{24:31}. For non-standard logics, such as intuitionist logic, one removes some
deduction rules and gets a logic that can no longer be expressed by tables: ``very esoteric'', very
interesting, but every course has to define its bottom line, and this one presents tables instead of
deduction rules and axioms, which is sufficient for what follows \ts{1}{1}{24:52}.
\end{coursenote}

\begin{history}[``Or'': inclusive or exclusive?]
As told in class \ts{1}{1}{25:54}: in the Middle Ages people discussed whether the ``or'' in some sentence meant
the inclusive or the exclusive one, and appealed to the Bible to see how God used it. Around the turn of the
20th century Frege, Wittgenstein and others said: it is just a table, a definition; there are 16 binary
operators, each of which may or may not deserve a name.\par
\smallskip
\emph{References.} The medieval anecdote is reported as said in class (not checked). Truth tables in their
modern form: L.~Wittgenstein, \emph{Tractatus Logico-Philosophicus} (1921), 4.31 and 5.101; E.~L.~Post,
``Introduction to a general theory of elementary propositions'', \emph{Amer.\ J.\ Math.} 43 (1921) 163--185.
G.~Frege, \emph{Begriffsschrift} (1879).
\end{history}

This is all ``child's play'', but one binary operator is sometimes ill understood: the \textbf{implication}
$\Rightarrow$ \ts{1}{1}{26:35}. From school we use the arrow to mean ``I wrote down something true, did something
valid, and arrived here''. But the implication is a binary operator like $\wedge$: it takes two propositions
and makes one new proposition, $p\Rightarrow q$, which as a whole is true or false \ts{1}{1}{26:57}. Its table
is the column above: from true follows true, true; from true follows false, false (``you don't pass the
test''); from false follows true, true; and from false follows false, also true \ts{1}{1}{27:18}. This is the
definition \ts{1}{1}{28:01}.

\begin{definition}[New concepts: ex falso quodlibet]\label{def:lo-exfalso}
The principle that from a false assumption anything can be concluded, the implication as a whole being true,
is called \textbf{ex falso quodlibet} (Latin: ``from the false, whatever you like'') \ts{1}{1}{28:22}.
\end{definition}

The first two rows are intuitive; why define the last two like this? The answer is hidden in a little
theorem \ts{1}{1}{28:42}. To state it we need the \textbf{equivalence} $p\Leftrightarrow q$, true when both
propositions have the same value and false otherwise (table above) \ts{1}{1}{29:05}.

\begin{theorem}[Contraposition]\label{thm:lo-contra}
$(p\Rightarrow q)\ \Leftrightarrow\ \bigl((\neg q)\Rightarrow(\neg p)\bigr)$ \ts{1}{1}{29:25}.
\end{theorem}
\begin{proof}
Write out the truth table in full \ts{1}{1}{32:10}:
\begin{center}
\renewcommand{\arraystretch}{1.3}
\begin{tabular}{cc|cc|cc}
$p$ & $q$ & $\neg p$ & $\neg q$ & $p\Rightarrow q$ & $\neg q\Rightarrow\neg p$\\\hline
T & T & F & F & T & T\\
T & F & F & T & F & F\\
F & T & T & F & T & T\\
F & F & T & T & T & T
\end{tabular}
\end{center}
The last two columns agree in every row \ts{1}{1}{33:12}, so the equivalence is true whatever $p$ and $q$ are.
\end{proof}
Note that putting $\neg$ in front exchanges the order of $p$ and $q$; since binding strengths have not been
fixed yet, brackets are used \ts{1}{1}{29:46}. Proving such a statement by writing out the truth table is ``the
first vanilla problem in any exam'' \ts{1}{1}{32:10}.

\begin{corollary}[Proof by contradiction]\label{cor:lo-contradiction}
Assertions can be proven by way of contradiction \ts{1}{1}{30:06}: to prove that $q$ follows from the true
assumption $p$, one may equivalently assume that $q$ is false and show that then $p$ cannot be true.
\end{corollary}
This is a very elementary principle and a very sharp tool of mathematics \ts{1}{1}{30:48}.

\begin{secondary}[Intuitionist logic]
If logic is formulated by deduction rules and axioms instead of tables, one can ``surgically'' remove a rule
so that \cref{thm:lo-contra} no longer holds; the result cannot be written with tables any more, and it gives a
mathematics in which proofs by contradiction are not allowed \ts{1}{1}{30:48}. In every good mathematics
department there is one professor who follows this \emph{intuitionist} logic, and ``it's a pain'': everything
must be proven constructively \ts{1}{1}{31:29}. Some make the point that a constructive result is more
meaningful, and an even deeper one: that proofs by contradiction cannot be trusted. The course does not go into
this: its logic allows proofs by contradiction, essentially by the definition of $\Rightarrow$
\ts{1}{1}{31:49}.
\end{secondary}

\begin{exercise}[A tautology]\label{ex:lo-taut}
Write the truth table of $p\vee\neg p$ and of $(p\wedge(p\Rightarrow q))\Rightarrow q$. Which of the two is
true for every choice of $p,q$?
\end{exercise}

\begin{remark}[Binding strength]\label{rem:lo-binding}
We agree on the decreasing binding strength $\neg,\ \wedge,\ \vee,\ \Rightarrow,\ \Leftrightarrow$: $\neg$ binds
most strongly, $\Leftrightarrow$ least \ts{1}{1}{33:37}. With more symbols the sequence must be extended
\ts{1}{1}{33:59}.
\end{remark}

\begin{remark}[Higher-order operators and NAND]\label{rem:lo-nand}
An operator $\heartsuit(p_1,\dots,p_n)$ that takes $n$ propositions is an \textbf{operator of order $n$}: unary
operators have order 1, binary ones order 2 \ts{1}{1}{34:41}. All operators of any order can be constructed from
one single binary operator \ts{1}{1}{35:02}, the \textbf{NAND} operator (``not and''), $p\uparrow q:\Leftrightarrow
\neg(p\wedge q)$ \ts{1}{1}{35:24}:
\begin{center}
\renewcommand{\arraystretch}{1.35}
\begin{tabular}{cc|c}
$p$ & $q$ & $p\uparrow q$\\\hline
T&T&F\\ T&F&T\\ F&T&T\\ F&F&T
\end{tabular}
\end{center}
So the binding strengths are not really needed: they are a convenience \ts{1}{1}{36:07}.
\end{remark}

\begin{exercise}[First steps with NAND]\label{ex:lo-nand}
Express $\neg p$ and $p\wedge q$ using only $\uparrow$ (``finger practice'', as on the first problem sheet
\ts{1}{1}{36:27}). The full statement is \cref{ex:lo-nand-full} in the appendix.
\end{exercise}

Asked again about the equivalence operator: it is true exactly when both propositions take the same value;
that is its definition \ts{1}{1}{36:48}.

% ==================================================================
\section{Predicate logic}\label{sec:lo-pred}

\begin{definition}[New concepts: predicate]\label{def:lo-predicate}
A \textbf{predicate} is, informally, a proposition-valued function of some variable or variables
\ts{1}{1}{37:38}. For example $P(x)$ is true or false depending on $x$ \ts{1}{1}{37:58}; $Q(x,y)$ is true or false
depending on the combination of $x$ and $y$ \ts{1}{1}{38:21}.
\end{definition}
The name comes from the Latin \emph{praedicare}, ``to proclaim'': what is asserted of a subject.

As propositional logic does not study how complex propositions arise, predicate logic, strictly speaking, does
not say how predicates are built \ts{1}{1}{38:41}. One could think of $Q(x,y)$ being true if $x>y$ for real
numbers $x,y$; but constructing such a predicate needs a further language for combining objects $x$ and $y$
\ts{1}{1}{39:02}. And from which set do $x$ and $y$ come? \emph{No such thing}: what $x$ and $y$ are, where they
come from and what their nature is, is left entirely open; they are just variables \ts{1}{1}{39:44}. This is
weird, because since elementary school we have been conditioned always to say which set a variable comes from;
but it must be left open, because the notion of set is to be defined later \emph{in this language}
\ts{1}{1}{40:05}.

\begin{coursenote}[Outlook: the element predicate]
Set theory will introduce a fundamental predicate of two variables, $x\in y$, which for two objects is either
true or false \ts{1}{1}{40:25}. If $x$ and $y$ are sets, $x\in y$ is a predicate of two variables. This is only
an outlook: in predicate logic the symbol $\in$ does not exist yet \ts{1}{1}{41:08}. ``It is always important to
know what a subject is not talking about'': here, how predicates are built \ts{1}{1}{41:28}.
\end{coursenote}

\paragraph{New predicates from old.} One can build new predicates from given ones \ts{1}{1}{41:48}, for instance
\[
  Q(x,y,z):\Leftrightarrow P(x)\wedge R(y,z),
\]
where $:\Leftrightarrow$ means ``is defined to be equivalent to'' (the analogue of $:=$): once $x,y,z$ are
supplied, $P(x)$ and $R(y,z)$ are propositions, which can be joined by $\wedge$ into a new predicate, now of three
variables \ts{1}{1}{42:10}.

\paragraph{Quantifiers.} More interesting: a predicate of one variable can be converted into a proposition
\ts{1}{1}{42:51}.

\begin{definition}[New concepts: universal quantifier]\label{def:lo-forall}
For a predicate $P$ of one variable, $\forall x\colon P(x)$ is the proposition, read ``for all $x$, $P(x)$ is
true'' \ts{1}{1}{44:14}, which is \emph{defined} to be true if $P(x)$ is true independently of $x$, that is,
whatever is substituted for $x$ \ts{1}{1}{44:38}.
\end{definition}
The whole expression is a proposition constructed from a predicate of one variable \ts{1}{1}{43:54}.

\begin{example}[A ``feel-good'' predicate]
Predicate logic does not provide specific predicates, but as a ``feel-good example'' (and a feel-good example
always involves some imprecision) \ts{1}{1}{45:26}: let
$P(x):\Leftrightarrow (x\text{ is a human being})\Rightarrow(x\text{ has been created})$, which ``includes all
faiths'' \ts{1}{1}{46:10}. The whole proposition is the implication; where the $x$ come from is hidden in its
first part. This is certainly always true, so $\forall x\colon P(x)$ is true \ts{1}{1}{46:36}.
\end{example}

A student asks whether a predicate can be made into a proposition in other ways, ``for some $x$''
\ts{1}{1}{46:57}. Indeed: the brother of the universal quantifier \ts{1}{1}{47:18}.

\begin{definition}[New concepts: existential quantifier]\label{def:lo-exists}
$\exists x\colon P(x)$, read ``there exists an $x$ such that $P(x)$'' \ts{1}{1}{47:59}, is the proposition
\[
  \exists x\colon P(x)\ :\Leftrightarrow\ \neg\bigl(\forall x\colon \neg P(x)\bigr),
\]
the negated universal quantification of the negated predicate \ts{1}{1}{48:19}.
\end{definition}

\begin{corollary}\label{cor:lo-neg}
$\forall x\colon\neg P(x)\ \Leftrightarrow\ \neg\bigl(\exists x\colon P(x)\bigr)$ \ts{1}{1}{48:40}: if something is
not true for any $x$, then it is not true that there exists an $x$ for which it is true \ts{1}{1}{49:03}.
\end{corollary}

\begin{aside}[The hairdresser's door]
Simple once written down, but careful with negations and quantifiers \ts{1}{1}{49:28}. On the doors of many
barber shops it says: ``What hairdressers can do, only hairdressers can do.'' It sounds good, but it means that
hairdressers can do nothing a non-hairdresser can do: if a non-hairdresser can draw nice pictures, a hairdresser
cannot \ts{1}{1}{49:48}.
\end{aside}

\begin{remark}[Quantifying predicates of several variables]\label{rem:lo-free}
Quantification also applies to predicates of more than one variable \ts{1}{1}{50:30}. $\forall x\colon P(x,y)$ is a
predicate of the one variable $y$ that was not quantified \ts{1}{1}{51:14}: $x$ is the \textbf{bound} variable,
$y$, which survives the quantification, the \textbf{free} variable \ts{1}{1}{51:35}. The remaining variable can
be quantified again \ts{1}{1}{51:56}.
\end{remark}

\begin{remark}[The order of quantification matters]\label{rem:lo-order}
In general \ts{1}{1}{52:17}
\[
  \forall x\ \exists y\colon P(x,y)\qquad\text{and}\qquad \exists y\ \forall x\colon P(x,y)
\]
are different propositions \ts{1}{1}{52:37}. One should therefore seriously avoid, although it is tricky in the
heat of a calculation, writing the quantifiers \emph{after} the expression (``\dots for all $x,y,z$''): when
written informally behind, it is unclear where they should be pulled to. Write all quantifiers in front, in
order \ts{1}{1}{52:58}.
\end{remark}

\begin{example}[Neutral and inverse elements]
For the additive neutral element of the real numbers: \emph{there exists} $0\in\R$ such that \emph{for all}
$x\in\R$, $x+0=x$; the same $0$ works for every $x$ \ts{1}{1}{53:40}. For inverses: \emph{for every} $x$ \emph{there
exists} a $y$ with $x+y=0$; $y$ may be chosen depending on $x$ \ts{1}{1}{54:01}. (Schuller first stated the
neutral element the wrong way round, and corrected himself: ``you see, already a mistake again''.)
\end{example}

\begin{exercise}[Negating quantifiers]\label{ex:lo-negq}
Using \cref{def:lo-exists}, show that $\neg(\forall x\ \exists y\colon P(x,y))\Leftrightarrow\exists x\ \forall y\colon
\neg P(x,y)$. Then negate the statement ``every real number has an additive inverse''.
\end{exercise}

% ==================================================================
\section{Axiomatic systems and the theory of proofs}\label{sec:lo-proofs}

Proofs are met in physics and, above all, in mathematics lectures; but what \emph{is} a proof
\ts{1}{1}{55:17}? A way of arguing why something is true given that something else was assumed true; but what are
the valid rules for writing one down? This can be defined, and there are subtleties which make some proofs
that are given really non-proofs \ts{1}{1}{55:38}. All of this prepares the axiomatic system of set theory
\ts{1}{1}{55:58}.

\begin{definition}[New concepts: axiomatic system]\label{def:lo-axsys}
An \textbf{axiomatic system} is a finite sequence of propositions $a_1,a_2,\dots,a_N$, which are called the
\textbf{axioms} \ts{1}{1}{56:19}\ts{1}{1}{58:27}.
\end{definition}

\begin{aside}[Pre-mathematical numbers]
A smart objection: what do ``finite'' and the numbers $1,2,\dots,N$ mean before sets are defined? One could
write $a_{/},a_{//},a_{///},\dots$ with strokes instead; logicians call these \emph{pre-mathematical numbers}:
``it's what you're doing in Bavarian pubs'' (counting beers by strokes on the beer mat) \ts{1}{1}{56:44}.
\end{aside}

\begin{coursenote}[Disclaimer]
A mathematical logician in the audience would shout ``this is a little too na\"ive'': the course does not use the
full formal approach of first-order predicate logic (this already showed in the use of tables) \ts{1}{1}{57:24}.
It gives the right idea, conceptually exactly what logicians do, without the formal nitty-gritty, which for
what follows would be fully redundant. It is not what a professional logician would write in a research
monograph, but nothing of it is wrong \ts{1}{1}{57:46}.
\end{coursenote}

The idea: first the axioms are given, then a homework problem ``prove $p$ starting from the axioms''
\ts{1}{1}{59:08}.

\begin{definition}[New concepts: proof]\label{def:lo-proof}
A \textbf{proof} of a proposition $p$ within an axiomatic system $a_1,\dots,a_N$ is a finite sequence of
propositions (or propositional schemes) $q_1,\dots,q_M$ (the steps of the proof) \ts{1}{1}{59:30} such that
$q_M=p$ (the last step is the proposition to be proven) \ts{1}{1}{1:00:15} and, for every $1\le j\le M$, one of the
following holds \ts{1}{1}{1:02:41}\ts{1}{1}{1:03:03}:
\begin{itemize}
  \item[(A)] $q_j$ is a proposition from the list of axioms \ts{1}{1}{1:03:23};
  \item[(T)] $q_j$ is a tautology \ts{1}{1}{1:03:44};
  \item[(M)] there exist $1\le m,n<j$ such that $(q_m\wedge q_n)\Rightarrow q_j$ is true \ts{1}{1}{1:05:28}\ts{1}{1}{1:06:10}.
\end{itemize}
The letters stand for \emph{axiom}, \emph{tautology} and \emph{modus ponens} \ts{1}{1}{1:05:28}.
\end{definition}
So at any step one may write down an axiom (A), a plain tautology (T), or a proposition deduced from two
\emph{previous} steps (M) \ts{1}{1}{1:07:12}. \emph{Modus ponens} (Latin, ``the mode that affirms'') is a very old
deduction rule.

\begin{definition}[New concepts: tautology]\label{def:lo-tautology}
A \textbf{tautology} is a proposition that is true independently of the elementary propositions it is made of
\ts{1}{1}{1:04:05}; for example $p\vee\neg p$. Generically the value of an expression containing $p$ depends on $p$,
but this particular combination is always true, as the table shows \ts{1}{1}{1:04:25}.
\end{definition}

\begin{remark}[Propositional schemes]
One single axiom can have the form ``$\forall x\colon P(x)$''. It is just one proposition, but one may then use
the individual instances $P(x)$, which are true in maybe infinitely many cases; this is called a
\textbf{propositional scheme}. Still, the steps of a proof must be \emph{finitely many}, and this is very
important \ts{1}{1}{1:00:38}\ts{1}{1}{1:01:18}.
\end{remark}

\begin{aside}[``This proof is not finite'']
In his first semester Schuller was, as he says today, fortunate that his analysis professor was an intuitionist:
no proofs by contradiction, and he was sometimes really mad at textbook proofs of analysis theorems, often saying
``this proof is not finite'' [?: the transcript has ``not cyanide'', probably German \emph{nicht finit}]. The
students did not understand what he meant; with the formal theory of proofs, where a proof is a finite sequence
of propositions, they did \ts{1}{1}{1:01:39}.
\end{aside}

\paragraph{Is rule (M) needed?} A student asks whether (M) does not already follow from (A) and (T)
\ts{1}{1}{1:07:54}. It is needed, ``definitely, seriously, desperately, fully'': the intuition ``if this is true,
then I can do this'' is kindergarten intuition; $\Rightarrow$ is just a binary operator, with no built-in
direction \ts{1}{1}{1:08:15}. There is not yet enough technology for a meaningful example, but the course will
write out \emph{one} proof in full according to this scheme: the uniqueness of the empty set. An axiom will
stipulate that there exists an empty set; showing that there is only one is intuitively clear, but written out in
full it takes about 16 lines \ts{1}{1}{1:08:36}\ts{1}{1}{1:08:56}.

\begin{intuition}[Machine code and C++]
All other proofs are given the usual way, by quick arguments. Not even the most celebrated theorems in the best
journals are proven in the formal fashion, which would take ``millions and billions'' of lines for comparatively
simple statements \ts{1}{1}{1:09:37}. The formal proof is like machine code, zeros and ones; the proofs we write
are in a higher language, like C++ \ts{1}{1}{1:09:58}. The understanding is that a written proof is a proof only if
fellow mathematicians agree that it could in principle be broken down into a proof of the formal type
\ts{1}{1}{1:10:18}. If somebody suspects that it would not (for instance, not into finitely many steps), they object;
and as long as somebody says ``I'm not convinced'', the proof has to be broken down further and further
\ts{1}{1}{1:10:39}. ``Don't tell anyone: you can obstruct any lecture in mathematics'' this way \ts{1}{1}{1:10:59}.
\end{intuition}

Why may (M) use any two previous steps, even with others in between, but no later ones? Everything before the
current step is an axiom or a tautology or was reached from them by implication arrows; going beyond the current
step would amount to using equivalence arrows (``and later I show\dots''), and apart from circularity it is the
directedness of the implication that allows exactly the previous steps \ts{1}{1}{1:11:21}\ts{1}{1}{1:11:41}.

\begin{remark}[Notation]
If $p$ can be proven from the axiomatic system $a_1,\dots,a_N$, one writes $a_1,\dots,a_N\vdash p$ (``proves
$p$'') \ts{1}{1}{1:12:23}.
\end{remark}

\begin{remark}[Recognising versus finding]
The definition makes it easy to \emph{recognise} a proof: a computer can check whether a given sequence is a proof
of $p$ \ts{1}{1}{1:13:05}\ts{1}{1}{1:13:28}. \emph{Finding} a proof is an altogether different matter: even for set
problems, which can be solved, it is not easy, and sometimes it is so hard that whoever finds a proof first
becomes a celebrated mathematician \ts{1}{1}{1:13:49}. The theory of proofs gives no hint on how to prove things;
it tells you whether what you have written is a proof, which can be very useful \ts{1}{1}{1:14:30}.
\end{remark}

\begin{remark}[Tautologies are redundant axioms]\label{rem:lo-tautax}
A tautology among the axioms can be removed from the list without impairing the power of the axiomatic system:
the same propositions can still be proven, because the tautology can be pulled into any proof by (T)
\ts{1}{1}{1:14:56}\ts{1}{1}{1:15:41}. The extreme case: the axiomatic system of \emph{propositional logic} is the
\emph{empty sequence}, since everything one wants to show there is a tautology \ts{1}{1}{1:16:22}\ts{1}{1}{1:16:46}.
\end{remark}

\begin{definition}[New concepts: consistency]\label{def:lo-consistent}
An axiomatic system $a_1,\dots,a_N$ is \textbf{consistent} if there exists a proposition $q$ that cannot be proven
from it \ts{1}{1}{1:17:34}\ts{1}{1}{1:17:57}:
\[
  \exists q\colon\ \neg\,(a_1,\dots,a_N\vdash q).
\]
\end{definition}

The idea behind this maybe surprising definition \ts{1}{1}{1:18:40}: consider an axiomatic system containing
contradictory propositions, say $s$ and, later in the list, $\neg s$ \ts{1}{1}{1:19:28}\ts{1}{1}{1:19:54}. By (A) both
$s$ and $\neg s$ can be pulled into a proof; then by (M), for an arbitrary $q$, it suffices that
$(s\wedge\neg s)\Rightarrow q$ is true \ts{1}{1}{1:20:20}. It is a tautology: $s\wedge\neg s$ is always false, and by
the definition of the implication arrow, ex falso quodlibet, the whole statement is always true
\ts{1}{1}{1:21:02}\ts{1}{1}{1:21:23}. So any statement $q$ can be proven, and also $\neg q$ \ts{1}{1}{1:21:44}. Not being
able to prove every statement is a sign of not having contradictory axioms; hence the definition excludes exactly
this \ts{1}{1}{1:22:05}.

\begin{theorem}\label{thm:lo-propconsistent}
Propositional logic is consistent \ts{1}{1}{1:22:45}.
\end{theorem}
\begin{proof}
It suffices to exhibit a proposition that cannot be proven within propositional logic \ts{1}{1}{1:23:11}. Its
axiomatic system is the empty sequence (\cref{rem:lo-tautax}), so only (T) and (M) can be used in a proof
\ts{1}{1}{1:23:56}. The only way to introduce a new step is to write down a tautology, and (M) only allows, from
previous tautologies, a $q_j$ for which the whole implication is a tautology; so only tautologies can be proven
\ts{1}{1}{1:24:17}\ts{1}{1}{1:24:43}. But $q\wedge\neg q$ is not a tautology, so it cannot be proven
\ts{1}{1}{1:25:03}. Hence propositional logic is consistent \ts{1}{1}{1:25:26}.
\end{proof}
This is really trivial \ts{1}{1}{1:25:46}. The important remark is that while it is perfectly clear how to
\emph{define} consistency, it is perfectly difficult to \emph{prove} it for a given axiomatic system, propositional
logic being a big exception. For a more powerful system, like the axioms of set theory, it is very difficult; and
famous results, first obtained by G\"odel, say that under certain circumstances it is not even possible
\ts{1}{1}{1:26:07}\ts{1}{1}{1:26:27}. Only a rough taste of one of them \ts{1}{1}{1:26:48}:

\begin{theorem}[G\"odel, roughly and non-technically]\label{thm:lo-goedel}
Any axiomatic system that is powerful enough to encode the elementary arithmetic of the natural numbers is either
inconsistent or contains a proposition that can be neither proven nor disproven \ts{1}{1}{1:27:09}\ts{1}{1}{1:28:19}.
\end{theorem}
``Elementary arithmetic'' is imprecise language: roughly elementary-school mathematics (prime numbers, products);
what exactly is needed one sees from the proof \ts{1}{1}{1:27:57}. Variants of the formulation say that the
system contains true statements that cannot be proven \ts{1}{1}{1:29:10}.

\begin{history}[G\"odel's incompleteness theorems]
This sent shockwaves through the mathematical world \ts{1}{1}{1:29:10}: in mathematics, of all the arts, the notion
of truth had seemed clear: something was true if it could be proven from a system of assumptions
\ts{1}{1}{1:29:32}. Then G\"odel came along and, ironically, \emph{proved} this and variants of it
\ts{1}{1}{1:29:52}.\par
\smallskip
\emph{References.} K.~G\"odel, ``\"Uber formal unentscheidbare S\"atze der Principia Mathematica und verwandter
Systeme I'', \emph{Monatshefte f\"ur Mathematik und Physik} 38 (1931) 173--198 (the first and second
incompleteness theorems; announced at K\"onigsberg in September 1930).
\end{history}

\begin{secondary}[The idea of G\"odel's proof]
The proof is involved; its basic idea is the following \ts{1}{1}{1:29:52}.
\begin{enumerate}
  \item Assign to each mathematical, and also \emph{meta}-mathematical, statement a number, now called its
        \textbf{G\"odel number} \ts{1}{1}{1:30:15}. In a mathematical statement such as $a^2+b^2=c^2$, or one with
        quantifiers, roughly the plus sign gets a number, the exponent gets one, every variable a different one,
        the equal sign another \ts{1}{1}{1:31:04}. One can then talk \emph{about} the equation: ``the left-hand
        side has more symbols than the right-hand side'' is a meta-mathematical statement, and G\"odel showed how
        to assign a number even to such statements \ts{1}{1}{1:31:24}\ts{1}{1}{1:31:44}.
  \item The construction: every elementary symbol is assigned a prime; one takes powers of primes and multiplies
        them, getting big numbers, huge ones for meta-mathematical statements \ts{1}{1}{1:32:07}. Since numbers are
        multiplied, elementary arithmetic is needed: this explains the mysterious condition on the axiomatic
        system \ts{1}{1}{1:32:29}.
  \item Then a barber-type argument (``the barber shaves all people in his village who do not shave themselves'')
        identifies a proposition that is neither provable nor disprovable \ts{1}{1}{1:32:50}\ts{1}{1}{1:33:14}.
\end{enumerate}
The construction is reflexive, like Cantor's proof that there are uncountably many real numbers
\ts{1}{1}{1:34:08}: one assumes that every number written as an integer, a decimal point and decimal places
$d_1d_2d_3\dots$ can be listed, and uses the numbering to construct a new real number that is not in the list
\ts{1}{1}{1:34:29}. In G\"odel's proof one assigns G\"odel numbers to statements and produces a statement that uses
the G\"odel number of the very same statement, which leads to the contradiction \ts{1}{1}{1:34:49}.
\end{secondary}

With the definitions of this lecture (axiomatic system, provable, hence disprovable: the negation can be proven,
and consistent) the \emph{statement} of the theorem is fully understood; the proof is a different matter
\ts{1}{1}{1:35:10}.

% ==================================================================
\section{Outlook: axiomatic set theory}\label{sec:lo-outlook}

This completes the language in which axiomatic set theory will be defined, starting next time
\ts{1}{1}{1:35:30}.
\begin{itemize}
  \item The \textbf{element symbol} $\in$ is nothing but a predicate of two variables, written in the middle
        ($x\in y$ instead of ${\in}(x,y)$) to save two brackets and a comma, as is often done for predicates of
        two variables \ts{1}{1}{1:35:51}\ts{1}{1}{1:36:11}.
  \item It is the \emph{only} elementary predicate of set theory: even the equality of sets will be defined in
        terms of it \ts{1}{1}{1:36:32}. Nobody will explain what it means that $x$ is an element of $y$; and both
        $x$ and $y$ will be sets: there is no hierarchy of ``grown-up sets'' and mere elements, sets can be
        contained in sets \ts{1}{1}{1:36:53}.
  \item A priori there is no reason to exclude a set being an element of itself \ts{1}{1}{1:37:13}. If one
        restricts to the sets that do not contain themselves, one proves in two lines that the collection of
        all such sets is not a set: asking whether it contains itself, both answers lead to a contradiction
        \ts{1}{1}{1:37:34}\ts{1}{1}{1:37:54}. So na\"ive set theory is tricky, and axioms are needed to regulate the
        use of the element predicate, a kind of \emph{operational} definition \ts{1}{1}{1:38:15}.
  \item Some axioms say that certain sets exist: an empty set; for a given set, its power set, the set of all its
        subsets \ts{1}{1}{1:38:35}. In fact the only set one really needs to postulate is the empty set: every set
        is ultimately constructed from it (the set containing the empty set, and so on) \ts{1}{1}{1:38:56}.
  \item There are about nine axioms (the number varies with how they are written) that reproduce the set theory
        on which all of modern mathematics is built \ts{1}{1}{1:39:16}.
  \item One of them is the \textbf{axiom of choice} (Zorn's lemma is an equivalent form). Why be so picky? Without
        it one cannot prove that every vector space has a basis, which is used very often
        \ts{1}{1}{1:39:16}\ts{1}{1}{1:39:37}.
\end{itemize}
Already at the bottom of the building there are screws that must be justified for the whole building to stand as
we want; the purpose of this lecture and the next is to look at the foundations and to have a justifiably good
feeling about the constructions to come \ts{1}{1}{1:39:58}\ts{1}{1}{1:40:18}.

\begin{coursenote}[Names used]
The paradox in the third point is Russell's paradox (B.~Russell, 1901; letter to Frege, 1902); the axioms are those
of Zermelo--Fraenkel set theory with choice (ZFC); both names come in the next lecture.
\end{coursenote}

% ==================================================================
\section*{Key concepts}
\begin{keyconcepts}
\textbf{Retrieval questions.}
\begin{enumerate}
  \item Why is $(F\Rightarrow F)$ true, and which theorem justifies this definition?
  \item How many unary and binary operators are there, and from which single operator can all be built?
  \item Define $\exists$ in terms of $\forall$. Why does the order of quantifiers matter?
  \item State the three rules (A), (T), (M) for a step of a proof.
  \item What does it mean for an axiomatic system to be consistent, and why does a contradiction among the axioms
        destroy it?
  \item Why is propositional logic consistent? What does G\"odel's theorem say?
\end{enumerate}
\textbf{Concepts.}
\begin{itemize}
  \item Proposition: a variable taking the values true or false (\cref{def:lo-proposition}).
  \item Implication: $p\Rightarrow q$ false only for $p$ true, $q$ false; ex falso quodlibet (\cref{def:lo-exfalso}).
  \item Contraposition: $(p\Rightarrow q)\Leftrightarrow(\neg q\Rightarrow\neg p)$; proof by contradiction
        (\cref{thm:lo-contra}, \cref{cor:lo-contradiction}).
  \item NAND: $p\uparrow q:\Leftrightarrow\neg(p\wedge q)$ generates all operators (\cref{rem:lo-nand}).
  \item Predicate; $\forall x\colon P(x)$; $\exists x\colon P(x):\Leftrightarrow\neg\forall x\colon\neg P(x)$
        (\cref{def:lo-forall,def:lo-exists}); bound and free variables.
  \item Axiomatic system, proof by (A), (T), (M), $a_1,\dots,a_N\vdash p$ (\cref{def:lo-axsys,def:lo-proof}).
  \item Consistency: $\exists q\colon\neg(a_1,\dots,a_N\vdash q)$ (\cref{def:lo-consistent}); propositional logic is
        consistent (\cref{thm:lo-propconsistent}); G\"odel (\cref{thm:lo-goedel}).
\end{itemize}
\end{keyconcepts}

\flushside
\end{paracol}

% !TEX root = ../main.tex
\chapter{Axioms of set theory}\label{ch:l02}
\begin{paracol}{2}

\begin{sources}
\begin{tabularx}{\linewidth}{@{}l l L@{}}
\toprule
\textbf{Video} & \textbf{Length} & \textbf{Content}\\
\midrule
\seg{2}{1}{Lecture 2: Axioms of set theory} & 1:51:56 & The element relation; $\notin$, $\subseteq$, $=$;
the nine Zermelo--Fraenkel axioms with choice (ZFC): $\in$, empty set (with a fully formal proof of its
uniqueness), pair sets, union sets, replacement (restricted comprehension, complements), power sets,
infinity (the non-negative integers), choice, foundation. From 1:35:51 the video continues with the beginning of
the next session (choice and foundation).\\
\bottomrule
\end{tabularx}

\smallskip
\textbf{Transcript} \file{materials/transcripts/C02-V01.txt} (local whisper transcription: the video has no
YouTube captions).
\end{sources}

\begin{reading}[Where this lecture sits in the other notes]
Not treated in the other Schuller courses, which take sets for granted.
\end{reading}

% ==================================================================
\section{The element relation}\label{sec:se-in}

Last time propositional and predicate logic were introduced as the language in which set theory is to be
formulated; today that is what we do \ts{2}{1}{0:08}. Axiomatic set theory is built on one postulate: there is a
fundamental \emph{relation} called $\in$ \ts{2}{1}{0:08}.

\begin{definition}[New concepts: relation]\label{def:se-relation}
A \textbf{relation} is a predicate of two variables; ``relation'' is just shorter to say \ts{2}{1}{1:03}.
\end{definition}

Obviously $\in$ will be the element symbol we expect in set theory, but no preconceived notions: it is simply a
relation, called $\in$ \ts{2}{1}{1:28}. There will be \emph{no definition}, in the strict sense, of what $\in$ is,
apart from being a predicate of two variables, nor of what a set is \ts{2}{1}{1:53}\ts{2}{1}{2:16}. Instead we write
down nine axioms that speak of $\in$ and of sets; they teach us how to use $\in$ and what constitutes a set. Only
the interplay between $\in$ and what we call sets defines both \ts{2}{1}{2:38}\ts{2}{1}{3:03}. As said last time, this
is necessary if we want to write down mathematics from scratch, without any prior notions in terms of which a set
could be defined \ts{2}{1}{3:25}.

\paragraph{Overview of the axioms.} Schuller memorises them ``like a poem'' by their initial letters
\ts{2}{1}{3:46}\ts{2}{1}{4:11}:
\[
  \in\ \ \emptyset\qquad P\ \ U\ \ R\ \ P\qquad I\ \ C\ \ F
\]
\begin{itemize}
  \item $\in$, $\emptyset$: \textbf{basic existence axioms}: a property of the element relation, and the existence of
        an empty set \ts{2}{1}{4:33};
  \item P, U, R, P: heavily used \textbf{construction axioms}, which say how to build new sets from given ones
        (a recurring theme: given a structure, how to build new structures of the same type): \textbf{p}air sets,
        \textbf{u}nion sets, \textbf{r}eplacement (``a powerful thing''), \textbf{p}ower sets \ts{2}{1}{4:55}\ts{2}{1}{5:18};
  \item I, C: \textbf{further existence axioms} (infinity, choice), which could also be called more advanced
        construction axioms, and are historically newer \ts{2}{1}{5:41};
  \item F: the axiom of \textbf{foundation}, rather a \emph{non-existence} axiom: ``don't use sets of the following
        kind'' \ts{2}{1}{6:04}.
\end{itemize}
They are discussed in this order, partly because later axioms rely on earlier ones, partly to give an immediate
application of each \ts{2}{1}{6:30}\ts{2}{1}{6:53}.
\begin{coursenote}[The mnemonic]
The transcript has ``E E perp I C F''. Read as ``$\in$, $\emptyset$, PURP, ICF'': the two Es are the element relation
and the empty set, ``perp'' is P-U-R-P.
\end{coursenote}

\paragraph{Relations defined from $\in$.} A relation eats two variables; by tradition its symbol is written between
them, which saves two brackets and a comma \ts{2}{1}{7:39}\ts{2}{1}{8:02}.

\begin{definition}[New concepts: $\notin$, subset, equality of sets]\label{def:se-subset}
\begin{align*}
  x\notin y &:\Leftrightarrow \neg(x\in y), \\
  x\subseteq y &:\Leftrightarrow \forall a\colon (a\in x\Rightarrow a\in y), \\
  x=y &:\Leftrightarrow (x\subseteq y)\wedge(y\subseteq x)
\end{align*}
\ts{2}{1}{8:24}\ts{2}{1}{8:48}\ts{2}{1}{9:12}.
\end{definition}
So the equality sign, which one would normally take without questioning, is \emph{defined} in terms of the element
relation \ts{2}{1}{9:12}\ts{2}{1}{9:34}.

% ==================================================================
\section{The axiom on the element relation}\label{sec:se-ax-in}

The Zermelo--Fraenkel axioms of set theory \ts{2}{1}{9:57}. The first declares the relation between $\in$ and sets
\ts{2}{1}{10:21}.

\begin{definition}[Axiom on the $\in$-relation]\label{ax:se-in}
$x\in y$ is a proposition if and only if $x$ and $y$ are both sets \ts{2}{1}{10:48}\ts{2}{1}{11:09}. Formally:
$\forall x\,\forall y\colon (x\in y)\veebar\neg(x\in y)$ (\emph{tertium non datur}) \ts{2}{1}{23:04}\ts{2}{1}{23:33}.
\end{definition}
In predicate logic the nature of the variables was left entirely open; this axiom says that $x\in y$ is a
proposition, and $\in$ a relation, only on what we are going to call sets. It is the first thing we know about sets
\ts{2}{1}{11:37}. (The formal version is ``a little bit unnecessary'': we know what it means to be a proposition
\ts{2}{1}{23:59}.)

This seemed so trivial that for a long time it was not considered much of a condition; but it is. A famous
counterexample of something that is \emph{not} a set \ts{2}{1}{12:04}:

\begin{proposition}[Russell]\label{prop:se-russell}
Assume there is an object $u$ that contains all sets that do not contain themselves:
$\exists u\ \forall z\colon (z\in u\Leftrightarrow z\notin z)$, with $z$ ranging over sets
\ts{2}{1}{12:32}\ts{2}{1}{13:35}\ts{2}{1}{13:58}. Then $u$ is not a set.
\end{proposition}
\begin{proof}
If $u$ were a set, then by \cref{ax:se-in} $u\in u$ would be a proposition, true or false \ts{2}{1}{15:39}\ts{2}{1}{16:09}.
If $u\in u$ is true, $u$ is one particular $z$, so by the definition of $u$, $u\notin u$: contradiction; so the
assumption cannot be true \ts{2}{1}{16:33}\ts{2}{1}{17:03}. Then it must be false, since it is a proposition: $u\notin u$.
But $u$ contains all sets that do not contain themselves, so $u\in u$: contradiction again \ts{2}{1}{17:29}\ts{2}{1}{17:54}.
Both cases are contradictory, so $u$ is not a set \ts{2}{1}{18:17}.
\end{proof}
One may be disturbed by a set being, or not being, an element of itself, having been trained to distinguish sets
from their elements; the axiom of foundation will eventually say ``don't even look at such stuff''. But at this
point it is perfectly fine: if $z$ is a set, $z\in z$ is a proposition \ts{2}{1}{14:21}\ts{2}{1}{14:49}\ts{2}{1}{15:15}.
This is the \textbf{Russell paradox} (``Russell drama'') \ts{2}{1}{18:43}.

In \emph{naive set theory}, taught at school and in undergraduate courses, one easily writes
$u=\{z\mid z\notin z\}$, conditioning the elements of a set on a property, and so constructs objects that are not
sets \ts{2}{1}{19:07}\ts{2}{1}{19:38}\ts{2}{1}{20:02}. The contradictions of naive set theory were finally resolved
``in a big struggle by a number of people''; today we honour Zermelo and Fraenkel in particular. The first axiom
already shows the problem; eight more are needed not to run into difficulties \ts{2}{1}{20:31}\ts{2}{1}{21:00}.

\begin{history}[Russell's drama]
As told in class \ts{2}{1}{18:43}\ts{2}{1}{20:31}: Russell and Whitehead, Cambridge philosophers, logicians and
mathematicians, wrote the thick \emph{Principia Mathematica} on logic as the basic language of mathematics, their
life's work, built on naive set theory; when Russell discovered the contradiction he could not sleep, and had
nightmares of people in the Cambridge library throwing the volumes out of the window.\par
\smallskip
\emph{Checked.} Russell found the paradox in 1901 and wrote it to Frege in June 1902, as the second volume of
Frege's \emph{Grundgesetze der Arithmetik} (whose system it destroyed) was in press. \emph{Principia Mathematica}
(three volumes, 1910--1913) came \emph{after}, and avoids the paradox with the theory of types. The nightmare is
Russell's dream, told by G.~H.~Hardy in \emph{A Mathematician's Apology} (1940), of a librarian in the year 2100
deciding whether to throw away the last copy of \emph{Principia}. Zermelo's axioms: 1908; Fraenkel's replacement: 1922.
\end{history}

% ==================================================================
\section{The empty set}\label{sec:se-empty}

\begin{definition}[Axiom on the existence of an empty set]\label{ax:se-empty}
There exists a set that contains no elements \ts{2}{1}{21:57}. Formally: $\exists x\ \forall y\colon y\notin x$
\ts{2}{1}{22:29}.
\end{definition}

\begin{theorem}\label{thm:se-empty-unique}
There is only one empty set; it is therefore given a name, $\emptyset$ \ts{2}{1}{24:25}.
\end{theorem}
\begin{proof}[Standard textbook-style proof]
Let $x$ and $x'$ be empty sets; we show $x=x'$ \ts{2}{1}{24:59}\ts{2}{1}{25:38}. For any $y$, $y\in x$ is false, since
$x$ has no elements. So the implication $y\in x\Rightarrow y\in x'$ is true, whatever is written after the arrow (ex
falso quodlibet), independently of $y$ \ts{2}{1}{26:11}. By the definition of $\forall$,
$\forall y\colon(y\in x\Rightarrow y\in x')$ is true, i.e.\ $x\subseteq x'$ \ts{2}{1}{26:39}. Conversely, $y\in x'$ is
always false, so $\forall y\colon (y\in x'\Rightarrow y\in x)$, i.e.\ $x'\subseteq x$ \ts{2}{1}{27:20}\ts{2}{1}{27:47}.
Together, $x=x'$ \ts{2}{1}{28:13}.
\end{proof}

Didn't we define what a proof is? This textbook proof does not have that form \ts{2}{1}{28:37}. So here, once and only
once in the course, is the formal version \ts{2}{1}{28:37}\ts{2}{1}{29:49}. The assumptions become the axioms of the system
(propositional logic alone would have an empty list) \ts{2}{1}{30:57}\ts{2}{1}{31:19}:
\[
  a_1:\ \forall y\colon y\notin x,\qquad a_2:\ \forall y\colon y\notin x'.
\]
We need steps $q_1,\dots,q_M$ with $q_M$ ``$x=x'$'', each justified by (T), (A) or (M) of \cref{def:lo-proof}
\ts{2}{1}{30:13}\ts{2}{1}{31:41}:
\begin{enumerate}
  \item[$q_1$] $\bigl(\forall y\colon y\notin x\bigr)\Rightarrow\bigl(\forall y\colon (y\in x\Rightarrow y\in x')\bigr)$
        \hfill (T): a tautology, by ex falso quodlibet \ts{2}{1}{32:04}\ts{2}{1}{32:49};
  \item[$q_2$] $\forall y\colon y\notin x$ \hfill (A): axiom $a_1$ \ts{2}{1}{33:12}\ts{2}{1}{33:34};
  \item[$q_3$] $\forall y\colon (y\in x\Rightarrow y\in x')$, i.e.\ $x\subseteq x'$ \hfill (M) with $q_1,q_2$:
        $(q_1\wedge q_2)\Rightarrow q_3$ is a tautology \ts{2}{1}{33:56}\ts{2}{1}{34:40}\ts{2}{1}{35:24};
  \item[$q_4$] $\bigl(\forall y\colon y\notin x'\bigr)\Rightarrow\bigl(\forall y\colon (y\in x'\Rightarrow y\in x)\bigr)$
        \hfill (T) \ts{2}{1}{36:08};
  \item[$q_5$] $\forall y\colon y\notin x'$ \hfill (A): axiom $a_2$ \ts{2}{1}{36:30};
  \item[$q_6$] $\forall y\colon (y\in x'\Rightarrow y\in x)$, i.e.\ $x'\subseteq x$ \hfill (M) with $q_4,q_5$
        \ts{2}{1}{36:52}\ts{2}{1}{37:37};
  \item[$q_7$] $x=x'$ \hfill (M) with $q_3,q_6$ \ts{2}{1}{37:59}.
\end{enumerate}
\begin{coursenote}[Which steps the modus ponens uses]
On the board Schuller first said that $q_6$ used steps 5 and 6, then corrected himself: 4 and 5
\ts{2}{1}{37:37}. Strictly, propositions are not \emph{equal} but \emph{equivalent}, so the steps should be written with
$\Leftrightarrow$ rather than $=$ if one is picky \ts{2}{1}{30:35}.
\end{coursenote}
``Something you probably want to have a look at again at home'' \ts{2}{1}{35:01}. The formal proof is nothing more
than the textbook proof broken down pedantically, but now we are certain it is a proof \ts{2}{1}{38:21}. This is the
spirit of axiomatic set theory: two levels, the usual language and the precise first-order version that mirrors it
\ts{2}{1}{38:43}\ts{2}{1}{39:05}. We will not follow the strict technique further: proofs get longer and longer, and no
book spells out all its proofs this way; but it must be possible in principle \ts{2}{1}{39:27}\ts{2}{1}{39:50}.

A student asks whether $x\subseteq y$ and $x\in y$ imply one another. No, in neither direction; finding
counterexamples for both is easy \ts{2}{1}{29:02}\ts{2}{1}{29:25}.
\begin{exercise}[$\subseteq$ versus $\in$]\label{ex:se-subset-in}
Find sets with $x\subseteq y$ but $x\notin y$, and with $x\in y$ but $x\not\subseteq y$ (use $\emptyset$ and the sets of
the next sections).
\end{exercise}

% ==================================================================
\section{Pair sets}\label{sec:se-pair}

``A bit tiring'' to go through all nine axioms, but this is the foundation of all of mathematics. A field such as
$\R$ also needs nine axioms (commutativity, associativity, neutral and inverse element for $+$, the same for
$\cdot$, one distributive law), so building all of mathematics on nine axioms is not too bad
\ts{2}{1}{40:14}\ts{2}{1}{40:36}\ts{2}{1}{40:58}.

\begin{definition}[Axiom on pair sets]\label{ax:se-pair}
Let $x$ and $y$ be sets. Then there exists a set that contains as its elements precisely $x$ and $y$
\ts{2}{1}{40:58}\ts{2}{1}{41:20}. Formally: $\forall x\,\forall y\ \exists m\ \forall u\colon\bigl(u\in m\Leftrightarrow
(u=x\vee u=y)\bigr)$ \ts{2}{1}{42:09}\ts{2}{1}{42:32}. Notation: $m=\{x,y\}$ \ts{2}{1}{42:56}\ts{2}{1}{43:23}.
\end{definition}
The notation is just a symbol (``I could put a little flower in the upper right corner as well''); what the axiom
ensures is that, if $x$ and $y$ are sets, this object \emph{is a set} \ts{2}{1}{43:23}\ts{2}{1}{43:46}.

\paragraph{Does the order matter?} Is $\{x,y\}=\{y,x\}$? Yes: $a\in\{x,y\}\Rightarrow a\in\{y,x\}$ and conversely, so
each is a subset of the other \ts{2}{1}{44:09}\ts{2}{1}{44:30}\ts{2}{1}{44:52}. ``We don't agree on it, we prove it'';
that the whole thing is a set is guaranteed by the axiom \ts{2}{1}{45:13}.

\begin{definition}[New concepts: one-element set]\label{def:se-singleton}
For a set $x$, $\{x\}:=\{x,x\}$, a set by the pair set axiom \ts{2}{1}{45:13}\ts{2}{1}{45:38}. Since equality is defined
through $\in$, writing $x$ twice adds no information \ts{2}{1}{46:00}.
\end{definition}
So far we have one- and two-element sets, nothing more: at this stage there is no justification for writing
$\{x,y,z\}$ \ts{2}{1}{46:22}. The most demanding task now is to get rid of old prejudices about sets, and to see how
what we knew follows from the axioms (some things will not be allowed), much as one learns that commutativity of
the real numbers is something special that must be stated as a law \ts{2}{1}{46:44}\ts{2}{1}{47:08}\ts{2}{1}{47:33}.

% ==================================================================
\section{Union sets}\label{sec:se-union}

\begin{definition}[Axiom on union sets]\label{ax:se-union}
Let $x$ be a set. Then there exists a set whose elements are precisely the elements of the elements of $x$
\ts{2}{1}{47:57}\ts{2}{1}{48:20}\ts{2}{1}{48:46}. Notation: $\bigcup x$ \ts{2}{1}{49:58}\ts{2}{1}{50:22}.
\end{definition}

\begin{exercise}[Axioms in first-order language]\label{ex:se-formal}
Homework for every further axiom: convert the verbal statement into first-order predicate logic, as done above for
the empty set and the pair sets: ``if you think you understand it, put it in first-order predicate logic''
\ts{2}{1}{49:13}\ts{2}{1}{49:35}.
\end{exercise}

Don't we need at least two sets to build a union? No: the set $x$ already collects all the sets to be put into the
union \ts{2}{1}{49:35}\ts{2}{1}{49:58}.

\begin{example}[Union of two one-element sets]
Let $a,b$ be sets. $\{a\}$ and $\{b\}$ are sets (pair set axiom), hence so is $x=\{\{a\},\{b\}\}$ (pair set axiom
again) \ts{2}{1}{52:14}. Its union contains the elements of the elements of $x$: $\bigcup x=\{a,b\}$. Here we already
knew this is a set; the union set axiom says the result is a set whatever $x$ \ts{2}{1}{52:39}\ts{2}{1}{53:03}.
\end{example}

\begin{example}[A three-element set]
For sets $a,b,c$, $x=\{\{a,b\},\{c\}\}$ is a set, so $\bigcup x$ is a set by the union set axiom; it deserves a name:
$\{a,b,c\}:=\bigcup x$ \ts{2}{1}{53:03}\ts{2}{1}{53:29}\ts{2}{1}{53:56}. ``In elementary school you learn this is a set,
but how do you know? It's just claimed. Now the axiom tells you'' \ts{2}{1}{54:21}.
\end{example}

\begin{definition}[New concepts: finite sets with specified elements]\label{def:se-finite}
Let $a_1,\dots,a_N$ be sets (the indices are again pre-mathematical stroke numbers) \ts{2}{1}{54:45}. For $N\ge3$ define
recursively
\[
  \{a_1,\dots,a_N\}:=\bigcup\bigl\{\{a_1,\dots,a_{N-1}\},\{a_N\}\bigr\},
\]
the recursion being started at $N=1,2$ by the pair set axiom \ts{2}{1}{55:27}\ts{2}{1}{55:56}.
\end{definition}
This gives everything we would informally call a finite set with specified elements \ts{2}{1}{55:56}.

\begin{intuition}[As many sets as fit into a set]
The big thing about the axiom is the assumption that the sets to be united are collected \emph{in a set} $x$.
Russell's $u$ is not a set, so it can never stand in place of $x$: one can only build unions of ``as many sets as
still fit into a set'' \ts{2}{1}{56:20}\ts{2}{1}{56:52}\ts{2}{1}{57:19}. This is much more than taking finitely many
sets (it contains, roughly, taking infinitely many), but not more sets than fit into a set: subtle, but very
precise \ts{2}{1}{57:44}. The elementary-school notion ``a set is a collection of elements'' comes back, now with a
justification from the first four axioms \ts{2}{1}{57:44}\ts{2}{1}{58:09}.
\end{intuition}

% ==================================================================
\section{Replacement and restricted comprehension}\label{sec:se-replacement}

Pedagogically, replacement comes before power sets, although it is not logically necessary
\ts{2}{1}{58:32}\ts{2}{1}{59:03}. ``A very innocent-sounding name; one of the most powerful axioms'' \ts{2}{1}{59:41}.

\begin{definition}[New concepts: functional relation, image]\label{def:se-functional}
A relation $R$ is \textbf{functional} if for every $x$ there exists precisely one $y$ such that $R(x,y)$
\ts{2}{1}{1:01:03}\ts{2}{1}{1:01:38}. The \textbf{image} of a set $m$ under a functional relation $R$ ``consists of'' all
those $y$ for which there is an $x\in m$ with $R(x,y)$ \ts{2}{1}{1:02:52}\ts{2}{1}{1:03:15}\ts{2}{1}{1:03:37}.
\end{definition}
This is what one calls a function: to every $x$ precisely one $y$ is assigned (to another $x$ maybe another $y$) but
at this stage there is no talk of $x$ or $y$ being sets: it is just a relation \ts{2}{1}{1:02:03}\ts{2}{1}{1:02:29}.
And ``your alarm bell should ring'' at the harmless word \emph{consists}: who says that all those $y$ are the
elements of a set \ts{2}{1}{1:03:58}\ts{2}{1}{1:04:19}?

\begin{exercise}[``There exists precisely one'']\label{ex:se-exists-one}
Define $\exists!\,y\colon P(y)$ (``there exists precisely one $y$'') in terms of $\forall$, $\exists$ and $=$
(homework in class, ``instead of an examples class'') \ts{2}{1}{1:01:38}\ts{2}{1}{1:02:03}.
\end{exercise}

\begin{definition}[Axiom of replacement]\label{ax:se-replacement}
Let $R$ be a functional relation and $m$ a set. Then the image of $m$ under $R$ is again a set
\ts{2}{1}{1:00:10}\ts{2}{1}{1:00:39}.
\end{definition}
A construction axiom: start with a set, construct the image under a functional relation, and the axiom guarantees
the result is a set \ts{2}{1}{1:00:39}\ts{2}{1}{1:04:42}. You have met it already, even at school, but in a weaker form,
restricted comprehension, which follows from it. The full axiom is needed to cover all of modern mathematics, but
in this course only its weaker consequence will be used \ts{2}{1}{1:05:03}\ts{2}{1}{1:05:23}\ts{2}{1}{1:05:45}.

\begin{theorem}[Principle of restricted comprehension]\label{thm:se-comprehension}
Let $P$ be a predicate of one variable and $m$ a set. Then the elements $y\in m$ for which $P(y)$ holds constitute a
set, denoted $\{y\in m\mid P(y)\}$ \ts{2}{1}{1:06:20}\ts{2}{1}{1:07:28}\ts{2}{1}{1:07:52}.
\end{theorem}
\begin{proof}
Two cases \ts{2}{1}{1:11:10}.
\begin{enumerate}
  \item There is no $y\in m$ with $P(y)$. Then the symbol is defined as $\emptyset$, which is a set; replacement is not
        needed \ts{2}{1}{1:11:39}\ts{2}{1}{1:12:42}.
  \item There is a $y\in m$ with $P(y)$; pick one, $\hat y$ \ts{2}{1}{1:13:19}. Define the relation
        \[
          R(x,y):\Leftrightarrow\bigl(P(x)\wedge y=x\bigr)\vee\bigl(\neg P(x)\wedge y=\hat y\bigr)
        \]
        \ts{2}{1}{1:13:49}\ts{2}{1}{1:14:32}\ts{2}{1}{1:14:52}. It is functional: if $P(x)$ holds, $x$ is sent to itself;
        if not, to the one fixed $\hat y$ \ts{2}{1}{1:15:15}\ts{2}{1}{1:15:36}. Define $\{y\in m\mid P(y)\}$ as the image of
        $m$ under $R$, a set by replacement \ts{2}{1}{1:15:59}\ts{2}{1}{1:16:21}.
\end{enumerate}
``You need to stare at this at home for five minutes'' \ts{2}{1}{1:15:59}\ts{2}{1}{1:16:43}.
\end{proof}
Replacement implies restricted comprehension, but not conversely \ts{2}{1}{1:06:20}.

\begin{remark}[Restricted quantifiers]\label{rem:se-restricted-q}
The proof uses quantifiers restricted to a set, which first-order logic does not have; they are short forms
\ts{2}{1}{1:11:39}\ts{2}{1}{1:12:12}:
\[
  \forall y\in m\colon P(y):\Leftrightarrow\forall y\colon(y\in m\Rightarrow P(y)),\qquad
  \exists y\in m\colon P(y):\Leftrightarrow\neg\bigl(\forall y\in m\colon\neg P(y)\bigr),
\]
so that $\exists y\in m\colon P(y)\Leftrightarrow\exists y\colon(y\in m\wedge P(y))$. Schuller first wrote the
existential version with $\Rightarrow$; a student objected, and the correct form is the one above: for $\forall$ an
implication, so that a $y$ outside $m$ does not affect the truth of the statement, and $\exists$ derived from it by
the double negation, which turns the implication into a conjunction \ts{2}{1}{1:20:07}\ts{2}{1}{1:20:47}\ts{2}{1}{1:21:26}\ts{2}{1}{1:21:57}.
``Very good observation. The laces of my shoes came undone because of this mistake'' \ts{2}{1}{1:22:28}.
\end{remark}

\begin{intuition}[Restricted versus universal comprehension]
Not to be confused with the inconsistent \textbf{principle of universal comprehension}, Russell's problem: ``the $y$
for which $P(y)$ holds form a set'' \ts{2}{1}{1:08:48}\ts{2}{1}{1:09:27}. That could be all $y$ with $y\notin y$
\ts{2}{1}{1:10:05}. The difference is that a set $m$ must be given from the start: the selection can only kick out some
elements of $m$, never make anything bigger, so it is not unbounded. The ``restricted'' is exactly ``$\in m$''
\ts{2}{1}{1:10:05}\ts{2}{1}{1:10:39}\ts{2}{1}{1:11:10}.
\end{intuition}

\begin{definition}[New concepts: complement]\label{def:se-complement}
For sets $u\subseteq m$, $m\setminus u:=\{x\in m\mid x\notin u\}$, a set by restricted comprehension, ultimately by
replacement \ts{2}{1}{1:17:27}\ts{2}{1}{1:17:49}\ts{2}{1}{1:18:11}.
\end{definition}

\begin{exercise}[Intersections]\label{ex:se-intersection}
Let $x$ be a set. Define the intersection of all the elements of $x$, using restricted comprehension (homework in
class) \ts{2}{1}{1:18:37}\ts{2}{1}{1:18:58}. Note the asymmetry: the existence of unions had to be postulated; intersections
follow from replacement \ts{2}{1}{1:19:22}.
\end{exercise}

% ==================================================================
\section{Power sets}\label{sec:se-power}

The power set is the collection of all subsets of a set \ts{2}{1}{1:22:28}\ts{2}{1}{1:23:36}. Historically, in naive set
theory, the principle of universal comprehension was thought to be needed to define it: ``all the $y$ that are
subsets of $m$'' \ts{2}{1}{1:23:36}\ts{2}{1}{1:24:28}. Given only $m$, we cannot say from which bigger set these subsets
come, so restricted comprehension does not help \ts{2}{1}{1:24:56}. And one wants to apply power sets repeatedly,
the power set of the power set\dots, going through the whole universe of sets; to \emph{generate} sets one cannot know
beforehand where they come from \ts{2}{1}{1:24:56}\ts{2}{1}{1:25:22}. There is no other choice than to postulate
\ts{2}{1}{1:25:53}:

\begin{definition}[Axiom on the existence of power sets]\label{ax:se-power}
Let $m$ be a set. Then there exists a set, denoted $\mathcal P(m)$, whose elements are precisely the subsets of $m$
\ts{2}{1}{1:25:53}\ts{2}{1}{1:26:37}.
\end{definition}

\begin{example}[Power set of a two-element set]
For $m=\{a,b\}$: $\mathcal P(m)=\bigl\{\emptyset,\{a\},\{b\},\{a,b\}\bigr\}$: the empty set (a subset of any set), the
one-element subsets, and $m$ itself \ts{2}{1}{1:26:37}\ts{2}{1}{1:27:37}.
\end{example}
Easy to explain how to build it; but is the result a set? Neither the non-existent universal principle nor
restricted comprehension can say so: an extra axiom is needed \ts{2}{1}{1:28:04}\ts{2}{1}{1:28:29}. ``Not very
economical'', but the power sets we want \ts{2}{1}{1:28:29}.

% ==================================================================
\section{Infinity}\label{sec:se-infinity}

\begin{definition}[Axiom of infinity]\label{ax:se-infinity}
There exists a set that contains the empty set and, with every one of its elements $y$, also contains $\{y\}$
\ts{2}{1}{1:29:23}\ts{2}{1}{1:29:48}.
\end{definition}
Without it we do not get standard mathematics \ts{2}{1}{1:29:23}. (``With every one of its elements'' refers to the
elements of the set being postulated, not of $\emptyset$, which has none \ts{2}{1}{1:29:48}.)

\begin{remark}[The simplest such set]\label{rem:se-naturals}
Informally, the simplest such set is $\bigl\{\emptyset,\{\emptyset\},\{\{\emptyset\}\},\{\{\{\emptyset\}\}\},\dots\bigr\}$
\ts{2}{1}{1:30:39}\ts{2}{1}{1:31:22}. Too many brackets for elementary school, so trivial names are introduced:
$0:=\emptyset$, $1:=\{\emptyset\}$, $2:=\{\{\emptyset\}\}$, and so on \ts{2}{1}{1:31:55}\ts{2}{1}{1:32:21}. The axiom guarantees
other sets too; this is one particular of them \ts{2}{1}{1:32:21}.
\end{remark}

\begin{corollary}\label{cor:se-N}
The non-negative integers form a set according to axiomatic set theory \ts{2}{1}{1:32:50}.
\end{corollary}
Without the axiom of infinity one could not prove it \ts{2}{1}{1:33:17}.
\begin{coursenote}[The precise construction]
``Informally speaking'': the set of \cref{rem:se-naturals} is obtained from any set given by the axiom, by
restricted comprehension, as the smallest one (the intersection of all subsets with the property). These are
Zermelo's numerals, $n+1=\{n\}$; the now more common von Neumann numerals use $n+1=n\cup\{n\}$.
\end{coursenote}

What next: do we require the real numbers to exist? As a \emph{set} (nothing further) they can be understood as, or
defined as, the power set of the integers \ts{2}{1}{1:33:17}\ts{2}{1}{1:33:47}. Then the power set of the reals is a bigger
set, its power set bigger still, and so on through the universe of sets \ts{2}{1}{1:33:47}\ts{2}{1}{1:34:19}. And the only set
whose existence was explicitly demanded so far is the empty set: from it we already arrived at the reals
\ts{2}{1}{1:34:19}.
\begin{coursenote}[``The reals are the power set of the integers'']
As a statement about sets this means: $\mathcal P(\mathbb N)$ has the cardinality of $\R$, so it can serve as the
underlying set; $\R$ as the field of real numbers is built differently (e.g.\ Dedekind cuts of $\mathbb Q$).
\end{coursenote}

The two remaining axioms are announced: C, which one may put in or leave out (ZF with or without choice), and F,
which forbids ``bootstrapping'' sets that contain themselves; with all axioms, essentially every set of standard
mathematics is constructed from the empty set alone. ``So set theory is very simple: it's just the empty set''
\ts{2}{1}{1:34:48}\ts{2}{1}{1:35:19}.

% ==================================================================
\section{Choice}\label{sec:se-choice}

The video continues with the next session: ``welcome back''; seven axioms discussed, two to go
\ts{2}{1}{1:35:51}.

\begin{definition}[Axiom of choice]\label{ax:se-choice}
Let $x$ be a set whose elements are non-empty and mutually disjoint (the intersection of any two elements of $x$ is
$\emptyset$). Then there exists a set $y$ that contains exactly one element of each element of $x$
\ts{2}{1}{1:36:28}\ts{2}{1}{1:37:08}.
\end{definition}
Intuitively: from each of the elements of $x$, which are all sets, pick one element, and collect them all in $y$
\ts{2}{1}{1:37:08}\ts{2}{1}{1:38:12}. Sometimes $y$ is called a ``dark'' set: given $x$, there is no prescription of how to
pick the elements \ts{2}{1}{1:38:41}.

\begin{intuition}[Shoes and socks]
An intuition, ``I think by Russell'' \ts{2}{1}{1:38:41}\ts{2}{1}{1:39:14}: if $x$ consists of pairs of shoes,
$\{\{\text{left}_1,\text{right}_1\},\{\text{left}_2,\text{right}_2\},\dots\}$, no axiom of choice is needed: there is an
algorithm, ``always pick the left shoe'' \ts{2}{1}{1:39:57}. For pairs of socks, two different individual socks that
cannot be distinguished, one must invoke the axiom of choice: there is no further means of telling which one is
picked \ts{2}{1}{1:40:37}\ts{2}{1}{1:41:15}. It is just an analogy \ts{2}{1}{1:41:52}.
\end{intuition}

A student asks whether it should say ``one element of each \emph{set} of $x$''. The point is that everything that
appears is a set: ``element'' only emphasises that something lies in something else \ts{2}{1}{1:41:52}\ts{2}{1}{1:42:22}.

\begin{remark}[Independence and uses]\label{rem:se-choice}
The axiom of choice is independent of the other axioms: set theory can be done with or without it; standard
mathematics uses it, and so will we \ts{2}{1}{1:42:22}\ts{2}{1}{1:43:27}. Two theorems of this course need it
\ts{2}{1}{1:43:27}\ts{2}{1}{1:44:12}:
\begin{enumerate}
  \item every vector space has a basis;
  \item for every equivalence relation there exists a complete system of representatives (equivalence relations
        come later).
\end{enumerate}
\end{remark}

\begin{history}[Zermelo and the axiom of choice]
As said in class (``I think, I'm not entirely sure'') \ts{2}{1}{1:44:52}\ts{2}{1}{1:45:19}: the axiom was introduced by
Zermelo to prove a theorem that had long been unproven, so it was first suspected of being made up just for that;
today it is used throughout, and nothing about it is suspicious.\par
\smallskip
\emph{Checked.} E.~Zermelo, ``Beweis, da\ss{} jede Menge wohlgeordnet werden kann'', \emph{Math.\ Ann.}\ 59 (1904)
514--516: the well-ordering theorem, conjectured by Cantor. Independence: G\"odel (1938) showed choice consistent with
ZF, P.~Cohen (1963) that its negation is too. The shoes-and-socks image is indeed Russell's (\emph{Introduction to
Mathematical Philosophy}, 1919).
\end{history}

% ==================================================================
\section{Foundation}\label{sec:se-foundation}

The last axiom is not an existence or construction axiom but a non-existence axiom; notwithstanding its name it was
one of the last to be introduced \ts{2}{1}{1:45:19}\ts{2}{1}{1:46:05}.

\begin{definition}[Axiom of foundation]\label{ax:se-foundation}
Every non-empty set $x$ contains an element $y$ that has none of its elements in common with $x$:
$x\ne\emptyset\Rightarrow\exists y\in x\colon y\cap x=\emptyset$ \ts{2}{1}{1:46:31}\ts{2}{1}{1:47:01}.
\end{definition}

\begin{corollary}\label{cor:se-no-self}
There is no set $x$ with $x\in x$ \ts{2}{1}{1:47:25}.
\end{corollary}
\begin{proof}
(Not written out in class.) Let $x$ be a set; $\{x\}$ is a non-empty set, so by foundation it has an element with no
element in common with $\{x\}$. Its only element is $x$, so $x\cap\{x\}=\emptyset$, i.e.\ $x\notin x$.
\end{proof}
Recall the attempt to build the set of all sets that do not contain themselves: that is a non-set, since it violates
the first axiom; and self-containing sets could not be excluded by a principle of universal comprehension, since
there is none, only a restricted one. So the axiom of foundation has to take care of it
\ts{2}{1}{1:48:06}\ts{2}{1}{1:48:30}\ts{2}{1}{1:48:55}. It also excludes defining sets self-referentially; without such sets
every set must ultimately be built from the empty set, as the non-negative integers were by the axiom of infinity
\ts{2}{1}{1:49:20}\ts{2}{1}{1:49:50}.

\paragraph{Summary.} The nine axioms $\in,\emptyset$, P, U, R, P, I, C, F should be memorised; every time a set is
constructed, ask whether it can actually be constructed from them \ts{2}{1}{1:49:50}\ts{2}{1}{1:50:17}. Their totality is
\textbf{Zermelo--Fraenkel set theory with the axiom of choice}, \textbf{ZFC}; without C, ZF \ts{2}{1}{1:50:17}\ts{2}{1}{1:50:56}.
There are other, weaker axiomatisations (not really stronger ones, as far as Schuller knows); as far as anybody knows,
all of standard mathematics is built on ZFC and can, in principle, be constructed from it
\ts{2}{1}{1:50:56}\ts{2}{1}{1:51:19}.

% ==================================================================
\section*{Key concepts}
\begin{keyconcepts}
\textbf{Retrieval questions.}
\begin{enumerate}
  \item How are $\subseteq$ and $=$ defined from $\in$?
  \item Why is $\{z\mid z\notin z\}$ not a set, and which axiom shows it?
  \item List the nine axioms of ZFC with their character (existence, construction, non-existence).
  \item Why is $\{a,b,c\}$ a set? Why is $\{x\}$ one?
  \item State replacement and derive restricted comprehension from it. Why is universal comprehension inconsistent?
  \item Why must the existence of power sets be postulated? Which set does infinity give?
  \item Where will the axiom of choice be needed? What does foundation exclude?
\end{enumerate}
\textbf{Concepts.}
\begin{itemize}
  \item $x\subseteq y:\Leftrightarrow\forall a\colon(a\in x\Rightarrow a\in y)$; $x=y:\Leftrightarrow x\subseteq y\wedge y\subseteq x$ (\cref{def:se-subset}).
  \item Russell's paradox (\cref{prop:se-russell}); uniqueness of $\emptyset$ with its formal proof (\cref{thm:se-empty-unique}).
  \item Pair sets $\{x,y\}$, $\{x\}$; unions $\bigcup x$; finite sets $\{a_1,\dots,a_N\}$ (\cref{ax:se-pair,ax:se-union,def:se-finite}).
  \item Functional relation, image; replacement; $\{y\in m\mid P(y)\}$ (\cref{ax:se-replacement,thm:se-comprehension}).
  \item $\mathcal P(m)$; infinity and $0=\emptyset$, $1=\{\emptyset\}$, \dots (\cref{ax:se-power,ax:se-infinity}).
  \item Choice (\cref{ax:se-choice}); foundation, no $x\in x$ (\cref{ax:se-foundation,cor:se-no-self}); ZFC.
\end{itemize}
\end{keyconcepts}

\flushside
\end{paracol}

% !TEX root = ../main.tex
% STUB: generated by tools/stubs.py; write the chapter from materials/transcripts/C03-V01.txt
\chapter{Classification of sets}\label{ch:l03}
\begin{paracol}{2}

\begin{sources}
\begin{tabularx}{\linewidth}{@{}l l L@{}}
\toprule
\textbf{Video} & \textbf{Length} & \textbf{Content}\\
\midrule
\seg{3}{1}{Lecture 3: Classification of sets} & 1:34:42 & (to be written)\\
\bottomrule
\end{tabularx}

\smallskip
\textbf{Transcript} \file{materials/transcripts/C03-V01.txt}.
\end{sources}

\begin{reading}[Where this lecture sits in the other notes]
Not treated in the other courses (only here). (From the lecture titles; to be checked when the chapter is written.)
\end{reading}

\emph{This chapter has not been written yet.}

\flushside
\end{paracol}

% !TEX root = ../main.tex
\chapter{Topological spaces - construction and purpose}\label{ch:l04}
\begin{paracol}{2}

\begin{sources}
\begin{tabularx}{\linewidth}{@{}l l L@{}}
\toprule
\textbf{Video} & \textbf{Length} & \textbf{Content}\\
\midrule
\seg{4}{1}{Lecture 4: Topological spaces - construction and purpose} & 1:38:42 & Definition of a topology and of a
topological space; how many topologies a finite set has; chaotic, discrete, a three-point and the standard topology on
$\R^d$ (open balls, proof that it is a topology); induced (subset) topology, closed sets, open versus closed, the circle
as a subset and as a quotient; product topology; Furstenberg's topological proof that there are infinitely many primes
(remark); convergence of sequences; a digression on non-transitive dice and ``progress'' in research proposals;
continuity, with examples; homeomorphisms, homeomorphic spaces, and the open classification problem.\\
\bottomrule
\end{tabularx}

\smallskip
\textbf{Transcript} \file{materials/transcripts/C04-V01.txt} (YouTube automatic captions).
\end{sources}

\begin{reading}[Where this lecture sits in the other notes]
\emph{Gravity and Light}, Lecture~1 (Topology) and Lecture~2 (Topological manifolds); \emph{Theoretische Mechanik},
Lecture~2 (Topological manifolds). In both courses topology is the first step after the introduction, and sets,
maps and equivalence relations are taken for granted; here they come from \cref{ch:l02,ch:l03}. (From the lecture
titles; those chapters are not written yet, so the overlap is not checked.)
\end{reading}

% ==================================================================
\section{Topologies and topological spaces}\label{sec:to-def}

Today the lecture starts looking closely at topological spaces, on the basis of the set theory of the previous
lectures \ts{4}{1}{0:04}. As announced last time, establishing a so-called \emph{topology} on a set is the weakest
structure one can put on it in order to have two very important notions: \emph{convergence} of sequences of points of
a set, and \emph{continuity} of maps between two topological spaces \ts{4}{1}{0:25}\ts{4}{1}{0:46}. The definition is at
first sight rather abstract, which is its downside; its upside is that it is extremely simple. It is the result of a
historical development (commented on occasionally), and is to be taken as the simplest definition that mathematicians
in the end found useful \ts{4}{1}{0:46}\ts{4}{1}{1:07}.

\begin{definition}[New concepts: topology, open set]\label{def:to-topology}
Let $M$ be a set (we know what that is from axiomatic set theory). A \textbf{choice} $\mathcal O\subseteq\mathcal P(M)$
of subsets of the power set of $M$ (possibly the entire power set) is called a \textbf{topology} on $M$ if three
conditions hold \ts{4}{1}{1:07}\ts{4}{1}{1:29}\ts{4}{1}{1:50}\ts{4}{1}{2:11}:
\begin{enumerate}
  \item $\emptyset\in\mathcal O$ and $M\in\mathcal O$ \ts{4}{1}{2:31};
  \item $U,V\in\mathcal O\ \Rightarrow\ U\cap V\in\mathcal O$ \ts{4}{1}{2:53}\ts{4}{1}{3:15}\ts{4}{1}{3:35};
  \item for every collection $C\subseteq\mathcal O$ of open sets (arbitrary, not only finite), $\bigcup C\in\mathcal O$
        \ts{4}{1}{4:15}\ts{4}{1}{4:36}.
\end{enumerate}
The elements of $\mathcal O$ are called \textbf{open sets} \ts{4}{1}{21:05}. (The word ``open'' is just a way of speaking
that makes things simple; it gets its meaning in \cref{sec:to-standard}.)
\end{definition}
That is all \ts{4}{1}{5:17}. You are not entirely free in the choice: that is the idea of the three conditions
\ts{4}{1}{2:11}. Condition~2 requires that if you choose $U$ and $V$ you also choose their intersection
\ts{4}{1}{3:35}. The intersection is written with the notation of \cref{ex:se-intersection}: one intersects the elements
of a set, so one can intersect infinitely many sets but never more sets than fit into a set \ts{4}{1}{3:15}\ts{4}{1}{3:55}\ts{4}{1}{4:15}
(``in kindergarten you wrote $U\cap V$'', which is still fine for two sets). Condition~3 could be written in the same
parallel fashion, $\{U,V\}\subseteq\mathcal O\Rightarrow\bigcup\{U,V\}\in\mathcal O$, if one puts the two sets
into a set \ts{4}{1}{4:56}.
\begin{coursenote}[Finite versus arbitrary]
Note the asymmetry of the conditions: \emph{finite} intersections (it suffices to ask for two, then induct), but
\emph{arbitrary} unions. Not said in class.
\end{coursenote}

\begin{definition}[New concepts: topological space]\label{def:to-space}
The pair $(M,\mathcal O)$ of a set and a topology on it is called a \textbf{topological space} \ts{4}{1}{11:17}\ts{4}{1}{11:38}.
\end{definition}
This is the first real instance of what was described very abstractly last time as a \emph{space}: a set with some
additional information, here a choice of subsets \ts{4}{1}{11:38}\ts{4}{1}{12:01}.

\begin{remark}[A real choice: how many topologies are there?]\label{rem:to-count}
Unless $M$ has only one element, there are many different topologies on one and the same set \ts{4}{1}{5:38}\ts{4}{1}{5:59}.
For finite $M$, with $|M|$ the number of elements, the number of topologies is \ts{4}{1}{6:20}\ts{4}{1}{6:42}\ts{4}{1}{7:03}\ts{4}{1}{7:24}:
\begin{center}
\begin{tabular}{l|ccccccc}
$|M|$ & 1 & 2 & 3 & 4 & 5 & 6 & 7\\\hline
topologies & 1 & 4 & 29 & 355 & 6\,942 & 209\,527 & 9\,535\,241
\end{tabular}
\end{center}
For $|M|=1$ the only choice is $\{\emptyset,M\}$. Already for a seven-element set the number is extreme, and an infinite
set has infinitely many topologies \ts{4}{1}{7:45}. So it is a real choice, and depending on it the notions of
continuity of maps and convergence of sequences, which will be defined solely in terms of the topology, change
\ts{4}{1}{8:06}.
\end{remark}
\begin{coursenote}[The numbers]
The captions give the last two numbers as ``2 95 2 7'' and ``9,535,000 241''; the counts of topologies on an
$n$-element set (labelled points) are $1,4,29,355,6942,209527,9535241$, as in the table. The values for 6 and 7 are
as in the standard integer-sequence tables (A000798, checked), not recomputed here.
\end{coursenote}

\subsection*{Examples}

To get familiar with the concept, first the basic examples \ts{4}{1}{8:26}.

\begin{example}[The chaotic topology]\label{ex:to-chaotic}
On any set $M$ take $\mathcal O=\{\emptyset,M\}$ \ts{4}{1}{8:26}. Condition~1 holds trivially. Condition~2: the
intersection of any two elements is $\emptyset\cap\emptyset=\emptyset$, $\emptyset\cap M=\emptyset$, $M\cap M=M$, all
already in $\mathcal O$. Condition~3: the only possible unions are $\emptyset$ or $M$, again in $\mathcal O$
\ts{4}{1}{8:47}\ts{4}{1}{9:08}. It is called the \textbf{chaotic topology} \ts{4}{1}{9:29}. It will not be very useful,
only as an extreme example, ``but sometimes the extreme examples teach you a lot'' \ts{4}{1}{9:52}.
\end{example}

\begin{example}[The discrete topology]\label{ex:to-discrete}
On any set $M$ take $\mathcal O=\mathcal P(M)$, every subset \ts{4}{1}{9:52}\ts{4}{1}{10:14}. The empty set and $M$ are
in the power set; the intersection or the union of any collection of subsets is a subset \ts{4}{1}{10:14}\ts{4}{1}{10:34}.
This is the \textbf{discrete topology}: ``discrete'' in the sense of separated, not in the moral sense
\ts{4}{1}{10:56}.
\end{example}

\begin{example}[A three-point set]\label{ex:to-three}
Let $M=\{1,2,3\}$, which has 29 topologies; choose
$\mathcal O=\bigl\{\emptyset,\{1\},\{2\},\{1,2\},\{1,2,3\}\bigr\}$ \ts{4}{1}{12:21}\ts{4}{1}{12:42}. (The one-element set $\{3\}$ and
the sets containing $3$ other than $M$ are missing.) Condition~1 holds as $\emptyset$ and $M$ are in. Any intersection of
two elements, and any union of elements, again lies in $\mathcal O$: you can check it \ts{4}{1}{13:04}\ts{4}{1}{13:24}.
It has no name, ``not that interesting'', but it is a topology; the chaotic and the discrete ones are two others on the same
set \ts{4}{1}{13:24}\ts{4}{1}{13:46}.
\end{example}

\subsection{The standard topology on $\R^d$}\label{sec:to-standard}

The most heavily used example, which underlies the usual constructions of convergence and continuity on the real
numbers \ts{4}{1}{13:46}\ts{4}{1}{14:06}. Take $M=\R^d=\R\times\dots\times\R$ ($d$ factors; the Cartesian product is known
from set theory) \ts{4}{1}{14:06}\ts{4}{1}{14:27}\ts{4}{1}{14:48}. One could put the chaotic or the discrete topology on it, but there is
one very useful one, the \textbf{standard topology} $\mathcal O_{\rm standard}$ \ts{4}{1}{14:48}\ts{4}{1}{15:09}. In contrast to the
examples above, where the open sets were listed, here they are defined implicitly, as there are uncountably many: the
construction proceeds in three steps \ts{4}{1}{15:09}\ts{4}{1}{15:30}\ts{4}{1}{15:50}.

\paragraph{Step 1: the balls.} For every $x\in\R^d$ and every $r\in\R^+$ define
\[
  B_r(x):=\Bigl\{y\in\R^d\ \Bigm|\ \sqrt{\textstyle\sum_{i=1}^d (y_i-x_i)^2}<r\Bigr\}
\]
(restricted comprehension; $y=(y_1,\dots,y_d)$ is an ordered tuple) \ts{4}{1}{15:50}\ts{4}{1}{16:12}\ts{4}{1}{16:33}\ts{4}{1}{16:54}\ts{4}{1}{17:15}.
Usually this is written as $\|y-x\|<r$, the $2$-norm: squares of the differences, then the square root
\ts{4}{1}{17:36}. At this point it is just a set, nothing to do with topology yet \ts{4}{1}{19:19}; intuitively it is the
\emph{open ball} of radius $r$ around $x$ \ts{4}{1}{18:59}.

One could use the $2n$-norm instead, with the power $2n$ and the $2n$-th root, for some $n$; for the standard topology it
makes no difference. In fact for \emph{any} norm the resulting topology is the same: a famous theorem
\ts{4}{1}{17:56}\ts{4}{1}{18:17}\ts{4}{1}{18:37}. Not to be too general, take $n=1$ \ts{4}{1}{18:37}.

\paragraph{Step 2: the open sets.} A subset $U\subseteq\R^d$ lies in $\mathcal O_{\rm standard}$ by definition if
\[
  \forall p\in U\colon\ \exists r\in\R^+\colon\ B_r(p)\subseteq U
\]
\ts{4}{1}{19:19}\ts{4}{1}{19:40}\ts{4}{1}{20:02}. That is the definition. Remember that ``open'' is just the name of the elements
of a topology \ts{4}{1}{21:05}.

\begin{example}[The picture in $\R^2$]\label{ex:to-picture}
Let the blackboard plane be $\R^2$ and $U$ the interior of a dotted curve, the points on the curve not belonging to $U$
\ts{4}{1}{20:23}\ts{4}{1}{20:44}. Is $U$ open? Pick a point of $U$ blindly; there is a radius such that the ball
around it, drawn dotted too, lies inside $U$ \ts{4}{1}{21:05}\ts{4}{1}{21:25}. The ball consists of all points with distance
\emph{strictly} less than $r$: the boundary does not belong to it \ts{4}{1}{21:45}. Even a point very close to the
boundary has some (small) ball inside $U$ \ts{4}{1}{22:06}. If instead part of the boundary belongs to $U$, take $p$ on
it: however small the positive $r$, part of the ball lies outside $U$, so $U$ is not open \ts{4}{1}{22:27}\ts{4}{1}{22:48}.
\end{example}

\begin{intuition}[Other norms: TV screens and boxes]
With the $2n$-norm for $n=2$ (fourth power, fourth root; one can check it is a norm) the open ``balls'' look a little
like the screens of old TV sets: not round, not square, rounded corners \ts{4}{1}{23:08}\ts{4}{1}{23:29}. They give the same
topology: around a point with a TV set around it one finds a smaller disc, around which a smaller TV set, and so on, so
that what is open for the discs is open for the TV sets and the other way around \ts{4}{1}{23:49}\ts{4}{1}{24:10}. For
$n=10\,000$ the corners are only a little rounded \ts{4}{1}{24:10}\ts{4}{1}{24:31}; for $n\to\infty$ one arrives at
boxes, the \textbf{maximum norm} (``the infinity norm''), writable with absolute values and no infinity
\ts{4}{1}{24:31}\ts{4}{1}{24:52}\ts{4}{1}{25:12}. Every box contains a disc and every disc contains a box, so the topology is always the same
\ts{4}{1}{24:52}\ts{4}{1}{25:12}.
\end{intuition}

\paragraph{Step 3: it is a topology.} It remains to prove that this construction really yields a topology
\ts{4}{1}{25:12}\ts{4}{1}{25:33}.
\begin{proposition}\label{prop:to-standard}
$\mathcal O_{\rm standard}$ on $\R^d$ is a topology.
\end{proposition}
\begin{proof}
Condition~1 \ts{4}{1}{25:57}.
\begin{enumerate}
  \item $\emptyset\in\mathcal O_{\rm standard}$: we must check $\forall p\in\emptyset\colon(\dots)$, with the rest of the
        condition irrelevant. The statement is $\forall p\colon (p\in\emptyset\Rightarrow\dots)$ and $p\in\emptyset$ is always
        false, so the implication is true by ex falso quodlibet, whatever follows; the predicate in $p$ is true for every
        $p$, so the $\forall$-proposition is true \ts{4}{1}{25:57}\ts{4}{1}{26:17}\ts{4}{1}{26:37}\ts{4}{1}{26:58}. The empty set is open
        ultimately because of the definition of the implication.
  \item $\R^d\in\mathcal O_{\rm standard}$: any ball of any radius around any point of $\R^d$ is by definition a subset of
        $\R^d$ \ts{4}{1}{27:18}\ts{4}{1}{27:39}.
\end{enumerate}
Condition~2 \ts{4}{1}{27:39}\ts{4}{1}{28:00}. Let $U,V$ be open; take an arbitrary $p\in U\cap V$ \ts{4}{1}{28:00}\ts{4}{1}{28:20}. Then $p\in U$
and $p\in V$. As $U$ is open there is $r\in\R^+$ with $B_r(p)\subseteq U$; as $V$ is open there is $s\in\R^+$ with
$B_s(p)\subseteq V$ \ts{4}{1}{28:41}\ts{4}{1}{29:01}. Then $B_{\min(r,s)}(p)$ (the smaller of the two, or either one if equal)
lies in both $U$ and $V$, hence in $U\cap V$ \ts{4}{1}{29:23}\ts{4}{1}{29:44}. We found a ball in $U\cap V$ around an arbitrary
point of it, so $U\cap V\in\mathcal O_{\rm standard}$ \ts{4}{1}{30:04}.

Condition~3 \ts{4}{1}{30:04}\ts{4}{1}{30:25}. (``We are not always going to go through all steps; just to demonstrate this once.'')
Let $C\subseteq\mathcal O_{\rm standard}$ be a collection of open sets and $p\in\bigcup C$. Then $p\in U$ for some $U\in C$,
and as $U$ is open there is $r\in\R^+$ with $B_r(p)\subseteq U$. But $U\subseteq\bigcup C$, so $B_r(p)\subseteq\bigcup C$
\ts{4}{1}{30:25}\ts{4}{1}{30:47}\ts{4}{1}{31:08}\ts{4}{1}{31:28}. Hence $\bigcup C$ is open \ts{4}{1}{31:28}\ts{4}{1}{31:49}.
\end{proof}
\begin{recall}[Ex falso quodlibet]
That an implication with false premise is true, so a $\forall$ over an empty range holds, is the argument of the
formal proof that the empty set is unique, \cref{thm:se-empty-unique}.
\end{recall}

% ==================================================================
\section{Constructing new topologies from given ones}\label{sec:to-construct}

Once there is a toolbox of topologies, new ones are mostly custom designed for one's problem and built from the
building blocks one has; it is an important technique to know how to construct new topologies from given ones
\ts{4}{1}{32:14}\ts{4}{1}{32:38}\ts{4}{1}{32:59}.

\subsection{The induced topology on a subset}

\begin{definition}[New concepts: induced (subset) topology]\label{def:to-induced}
Let $(M,\mathcal O)$ be a topological space and $N\subseteq M$ a (proper) subset. Define
\[
  \mathcal O|_N:=\{\,U\cap N\mid U\in\mathcal O\,\}
\]
\ts{4}{1}{32:59}\ts{4}{1}{33:19}\ts{4}{1}{33:40}. Then $\mathcal O|_N$ is a topology on $N$, called the \textbf{induced}
(or \textbf{subset}) \textbf{topology}: ``a definition slash theorem'', since the claim that it is a topology is part of
it \ts{4}{1}{33:40}\ts{4}{1}{34:00}\ts{4}{1}{34:23}.
\end{definition}
It is constructed from the topology of the bigger set by intersecting each open set of it with the subset; every element
is a subset of $N$, so $\mathcal O|_N\subseteq\mathcal P(N)$ and there is a chance that it is a topology
\ts{4}{1}{34:23}\ts{4}{1}{34:45}.

\begin{proof}
It can be boring to go through such proofs, but for the first two or three definitions it is useful for a feeling of how
this works; one always asks the same questions \ts{4}{1}{35:07}.
\begin{enumerate}
  \item $\emptyset\in\mathcal O|_N$: $\emptyset=\emptyset\cap N$ and $\emptyset\in\mathcal O$, so it has the required form
        \ts{4}{1}{35:07}\ts{4}{1}{35:29}.
  \item $N\in\mathcal O|_N$: $N=M\cap N$ and $M\in\mathcal O$ by condition~1 \ts{4}{1}{35:49}\ts{4}{1}{36:10}.
  \item Let $S,T\in\mathcal O|_N$. (Schuller first tried to write it shorter, ran into trouble because he had not used the
        definition first, and repeated more carefully \ts{4}{1}{36:30}\ts{4}{1}{37:15}\ts{4}{1}{37:35}\ts{4}{1}{37:56}.)
        By definition there are $U,V\in\mathcal O$ with $S=U\cap N$ and $T=V\cap N$
        \ts{4}{1}{38:16}\ts{4}{1}{38:36}\ts{4}{1}{38:57}. Then
        \[
          S\cap T=(U\cap N)\cap(V\cap N)=(U\cap V)\cap N,
        \]
        and $U\cap V\in\mathcal O$ as $\mathcal O$ is a topology, so $S\cap T\in\mathcal O|_N$
        \ts{4}{1}{39:18}\ts{4}{1}{39:38}\ts{4}{1}{40:00}.
  \item Arbitrary unions: checked in the same way, ``I am not going to write this down now'' \ts{4}{1}{41:01}
        (compare \cref{ex:to-induced-union}).
\end{enumerate}
\end{proof}
A student asks whether the step $(U\cap N)\cap(V\cap N)=(U\cap V)\cap N$ has been shown: no, it was left for the problem
sheet; it is clear from ``kindergarten set theory'', and with the definition of the intersection from the last problem
sheet the finite case reduces to what one learns there \ts{4}{1}{40:00}\ts{4}{1}{40:20}.
\begin{aside}[Set theory in school]
A digression in class: in the 1970s set theory was a fashion in elementary schools, and parents got unnerved because
``it's mathematics, I don't understand'': red cars, blue cars, and the intersection of red cars and blue cars being empty
\ts{4}{1}{40:20}\ts{4}{1}{40:40}.
\end{aside}
\begin{exercise}[Induced topology]\label{ex:to-induced-union}
Prove the identity $(U\cap N)\cap(V\cap N)=(U\cap V)\cap N$ from the definition of the intersection, and check
condition~3 for $\mathcal O|_N$ (arbitrary unions) \ts{4}{1}{41:01}.
\end{exercise}

\begin{example}[Open in the subset, not open in the big space]\label{ex:to-interval}
Equip $\R$ with the standard topology and take $N=[-1,1]=\{x\in\R\mid -1\le x\le 1\}$ with the induced topology
\ts{4}{1}{41:01}\ts{4}{1}{41:21}\ts{4}{1}{41:43}\ts{4}{1}{42:04}. The set $(0,1]$ is a subset of $\R$; it is \emph{not} open in the
standard topology, because of the point $1$ on its boundary (the argument above) \ts{4}{1}{42:04}\ts{4}{1}{42:26}. But it is
also a subset of $N$, and in the induced topology it \emph{is} open, as
\[
  (0,1]=(0,2)\cap[-1,1],
\]
with $(0,2)$ open in $\R$ \ts{4}{1}{42:26}\ts{4}{1}{42:48}\ts{4}{1}{43:10}\ts{4}{1}{43:31}.
\end{example}
Any set that can be written as an open set of the bigger space intersected with the subspace is open in the induced
topology: a set can be non-open in the big space and nevertheless open in the induced topology of a subset it happens
to lie in \ts{4}{1}{43:31}\ts{4}{1}{43:53}. So the question is always ``open with respect to which topology?'' and one must not
draw quick conclusions from seeing subsets \ts{4}{1}{43:53}\ts{4}{1}{44:14}.

\subsection{Closed sets}

Since ``non-open'' is tempting to call ``closed'', one should say that this is \emph{not} the convention
\ts{4}{1}{44:14}\ts{4}{1}{44:35}.
\begin{definition}[New concepts: closed set]\label{def:to-closed}
Let $(M,\mathcal O)$ be a topological space. A subset $C\subseteq M$ is called \textbf{closed} if its complement
$M\setminus C$ is open \ts{4}{1}{44:35}\ts{4}{1}{44:58}\ts{4}{1}{45:21}.
\end{definition}
\begin{example}[A closed interval]
$[0,1]$ is closed in $\R$ with the standard topology. It has its end points, which only means it is not open; but it is
closed because
\[
  \R\setminus[0,1]=(-\infty,0)\cup(1,\infty)
\]
is a union of two open sets (around each point of either a small ball lies in it), hence open
\ts{4}{1}{45:21}\ts{4}{1}{45:43}\ts{4}{1}{46:04}\ts{4}{1}{46:24}.
\end{example}
\begin{coursenote}[Board notation]
The transcript has ``the interval 0 1'' without brackets; the closed interval $[0,1]$ is meant, as the complement has
open ends at $0$ and $1$ \ts{4}{1}{45:21}\ts{4}{1}{46:04}.
\end{coursenote}

A warning \ts{4}{1}{46:24}\ts{4}{1}{46:46}. For a generic subset of a generic topological space, with nothing known about the
space, six statements are possible: it can be (a) open, (b) closed, (c) open and closed, (d) open and not closed, (e)
not open and closed, (f) not open and not closed \ts{4}{1}{46:46}\ts{4}{1}{47:07}\ts{4}{1}{47:27}\ts{4}{1}{47:48}. Being open and being closed have
nothing to do with each other, except via the complement; from openness one cannot conclude anything about closedness, nor
the other way around \ts{4}{1}{48:09}.
\begin{coursenote}[The six cases]
As listed, (a) and (b) are ``can be open'' and ``can be closed''; the four mutually exclusive cases are
(c) to (f). Schuller's point is that all four combinations of ``open'' and ``closed'' can occur.
\end{coursenote}

\begin{proposition}[Observation: $\emptyset$ and $M$ are open and closed]\label{prop:to-clopen}
In every topological space $(M,\mathcal O)$ the sets $\emptyset$ and $M$ are both open and closed
\ts{4}{1}{48:09}\ts{4}{1}{48:32}.
\end{proposition}
\begin{proof}
$\emptyset=M\setminus M$ and $M\in\mathcal O$, so $\emptyset$ is closed. $M=M\setminus\emptyset$ and
$\emptyset\in\mathcal O$, so $M$ is closed. Both are open by condition~1 \ts{4}{1}{48:32}\ts{4}{1}{48:53}\ts{4}{1}{49:14}.
\end{proof}
There is no contradiction \ts{4}{1}{49:14}\ts{4}{1}{49:36}. Later a \emph{connected} topological space will be defined; one can prove
that in such a space these are the \emph{only} sets that are both open and closed. If one finds this funny, one can look at
connected spaces; but no topology avoids the fact that $\emptyset$ and $M$ are both, as it is a general property
\ts{4}{1}{49:36}\ts{4}{1}{49:56}\ts{4}{1}{50:16}.

\subsection{An application: a topology on the circle}

A real-life example where one uses the induced topology \ts{4}{1}{50:16}. Consider the set $S^1$, the circle, which one
thinks of only as a set: it has no shape; a shape will come only with geometry, at the very end of the building of
structures (``even this is already too much'': it is a \emph{continuous} circle, to come) \ts{4}{1}{50:39}\ts{4}{1}{50:59}\ts{4}{1}{51:19}.

\paragraph{As a subset of $\R^2$.} One can define it as a subset of $\R^2$:
$S^1:=\{(x,y)\in\R^2\mid x^2+y^2=1\}$ \ts{4}{1}{51:19}\ts{4}{1}{51:41}. To get a topology without starting from scratch, let
$\mathcal O_{S^1}:=\mathcal O_{\rm standard}|_{S^1}$ \ts{4}{1}{51:41}\ts{4}{1}{52:04}\ts{4}{1}{52:25}.
\begin{itemize}
  \item An arc of the circle without its two end points is open in this topology: it is $S^1\cap U$ for a suitable open
        $U\in\mathcal O_{\rm standard}(\R^2)$ \ts{4}{1}{52:25}\ts{4}{1}{52:46}\ts{4}{1}{53:07}.
  \item The empty set is open, as always: $\emptyset=\emptyset\cap S^1$ \ts{4}{1}{53:27}.
  \item A set consisting of one point is \emph{not} open: it cannot be written as the circle intersected with an open set of
        $\R^2$. An open set that merely touches the circle has no ball around the point (the boundary is not in it) and
        generates nothing; one that covers the point covers more than one point of the circle
        \ts{4}{1}{53:48}\ts{4}{1}{54:09}\ts{4}{1}{54:29}.
\end{itemize}

\paragraph{As a quotient set.} A different route \ts{4}{1}{54:29}\ts{4}{1}{54:49}: take the real line and identify points under
the equivalence relation $x\sim y:\Leftrightarrow y=x+2\pi$ [?: as heard, ``$y=x+2\pi$''; for $\sim$ to be an equivalence
relation it must be $y=x+2\pi n$ for some $n\in\mathbb Z$, which is what is meant] \ts{4}{1}{54:49}\ts{4}{1}{55:11}. Two points whose
difference is $2\pi$ are identified: starting at $0$, the point $2\pi$ is identified with $0$ and one starts again; this way
the circle is obtained as a quotient set (\cref{def:cl-quotient}) \ts{4}{1}{55:11}\ts{4}{1}{55:31}. One can then study how a topology
is inherited by a quotient set from the topology of the original set; the construction was not shown in class
\ts{4}{1}{55:31}\ts{4}{1}{55:51}.
\begin{exercise}[Quotient topology]\label{ex:to-quotient}
Problem sheet: the topology inherited by a quotient set \ts{4}{1}{55:51}\ts{4}{1}{56:13}. The full statement is in
\cref{ex:to-quotient-full}, \cref{exam:ch04}.
\end{exercise}
So there are sometimes several ways to do the same, and it is convenient to be able to build subsets, quotient sets and the
topologies they inherit \ts{4}{1}{56:13}.

\subsection{The product topology}

One more very important construction \ts{4}{1}{56:13}\ts{4}{1}{56:35}. Let $(A,\mathcal O_A)$ and $(B,\mathcal O_B)$ be topological
spaces (not the induced topology, so no stroke). Very often, and implicitly, one equips the Cartesian product $A\times B$ with
the \textbf{product topology} $\mathcal O_{A\times B}$ \ts{4}{1}{56:35}\ts{4}{1}{56:55}\ts{4}{1}{57:17}.

\begin{definition}[New concepts: product topology]\label{def:to-product}
A set $U\subseteq A\times B$ is an element of $\mathcal O_{A\times B}$ if and only if for every point $p=(a,b)\in U$
there exist $S\in\mathcal O_A$ with $a\in S$ and $T\in\mathcal O_B$ with $b\in T$ such that
$S\times T\subseteq U$ \ts{4}{1}{57:17}\ts{4}{1}{57:40}\ts{4}{1}{58:01}\ts{4}{1}{58:22}\ts{4}{1}{58:42}.
\end{definition}
Intuitively: represent $A$ and $B$ as the two axes (a big assumption, as topological spaces can look very crazy; there is no
notion of dimension yet) \ts{4}{1}{59:04}\ts{4}{1}{59:24}. To check whether a set $U$ is open, take a point $(a,b)$ and look for
an open set $S$ around $a$ in $A$ and $T$ around $b$ in $B$; the box $S\times T$ must lie entirely inside $U$. If it does not
work, choose smaller $S$ and $T$ around the point; whether it then works depends on the topologies of $A$ and $B$
\ts{4}{1}{59:47}\ts{4}{1}{1:00:08}\ts{4}{1}{1:00:28}\ts{4}{1}{1:00:49}. This is inspired by $A,B$ being the real lines with the standard topology; that is the
picture to remember \ts{4}{1}{1:00:49}\ts{4}{1}{1:01:10}.
\begin{exercise}[Product topology]\label{ex:to-product}
Check that $\mathcal O_{A\times B}$ is a topology on $A\times B$ (problem sheet) \ts{4}{1}{1:01:10}\ts{4}{1}{1:01:31}.
\end{exercise}
\begin{remark}[Finite products; $\R^d$]\label{rem:to-product}
Two remarks \ts{4}{1}{1:01:31}\ts{4}{1}{1:01:53}.
\begin{enumerate}
  \item The same works for any finite product $A_1\times A_2\times\dots\times A_n$.
  \item The standard topology $\mathcal O_{\rm standard}$ on $\R^d$ coincides with the product topology of $d$ copies of
        $\R$, each with the standard topology: the argument is the one with the maximum norm, boxes always contain discs and
        discs always contain boxes \ts{4}{1}{1:01:53}\ts{4}{1}{1:02:15}\ts{4}{1}{1:02:36}.
\end{enumerate}
One could have been more minimalistic, defining the standard topology only on $\R$ and then on $\R^d$ via the product topology.
Both ways coincide; there is no a priori reason why they should, but they do \ts{4}{1}{1:02:36}\ts{4}{1}{1:02:57}\ts{4}{1}{1:03:17}.
\end{remark}

\begin{secondary}[Furstenberg: topology in number theory]
A remark before continuing \ts{4}{1}{1:03:17}\ts{4}{1}{1:03:38}. In this course topologies give structure to sets on the way to
manifolds and physical space. But topology is not only useful in this ``roughly geometric'' context: it is a very versatile
subject and deserves its own status \ts{4}{1}{1:03:38}\ts{4}{1}{1:03:58}. For instance one can equip the integers with a (cleverly
chosen) topology: Furstenberg discovered that with a suitable topology on $\mathbb Z$ one obtains a purely topological proof, by
just thinking about open and closed sets, that there are infinitely many primes \ts{4}{1}{1:03:58}\ts{4}{1}{1:04:20}\ts{4}{1}{1:04:42}. So
even in number theory topology is a tool. The example is famous as it is unexpected, and it explains why the definition has this
abstract form: it covers many situations, and the brain work goes into choosing the right topology to solve one's
problem, which sometimes takes ingenuity \ts{4}{1}{1:04:42}\ts{4}{1}{1:05:03}\ts{4}{1}{1:05:23}.\par
\smallskip
\emph{Not said in class.} Furstenberg's topology on $\mathbb Z$ has the arithmetic progressions as a basis; the proof is a
page long.
\end{secondary}
\begin{history}[Hausdorff's neighbourhoods]
\emph{Not said in class} (the lecture only says that the definition came from a historical development, \ts{4}{1}{0:46}\ts{4}{1}{1:07}).\par
\smallskip
\emph{Checked.} The notion of an abstract topological space with neighbourhood axioms is due to F.~Hausdorff,
\emph{Grundz\"uge der Mengenlehre} (1914); the definition by open sets, as above, became standard later. H.~Furstenberg,
``On the infinitude of primes'', \emph{Amer.\ Math.\ Monthly} 62 (1955) 353.
\end{history}

% ==================================================================
\section{Convergence}\label{sec:to-conv}

The whole structure is only useful if one can do something with it: convergence of sequences and continuity of maps
\ts{4}{1}{1:05:23}\ts{4}{1}{1:05:45}.

\begin{definition}[New concepts: sequence, convergence, open neighbourhood]\label{def:to-conv}
A \textbf{sequence} on a topological space $(M,\mathcal O)$ is a map $q\colon\N\to M$, from the non-negative integers into the
set underlying the space (Schuller first said ``integers'' and corrected to the naturals) \ts{4}{1}{1:06:09}\ts{4}{1}{1:06:33}\ts{4}{1}{1:08:00}. A sequence
$q$ \textbf{converges} against a point $a\in M$, the \textbf{limit point}, if for every open set $U\in\mathcal O$ with $a\in U$
there is $N\in\N$ such that
\[
  \forall n>N\colon\ q(n)\in U
\]
\ts{4}{1}{1:06:33}\ts{4}{1}{1:06:58}\ts{4}{1}{1:07:19}\ts{4}{1}{1:08:00}\ts{4}{1}{1:08:21}.
\end{definition}
An open set $U$ with $a\in U$ is called an \textbf{open neighbourhood} of $a$: it lies around the point (what lies around you is
your neighbourhood; this is just topology lingo) \ts{4}{1}{1:07:19}\ts{4}{1}{1:07:39}. The condition must hold for \emph{every} open neighbourhood,
which is very important \ts{4}{1}{1:08:00}: from some $N$ on, all the images of the sequence lie in it
\ts{4}{1}{1:08:00}\ts{4}{1}{1:08:21}.

\begin{example}[Chaotic topology: everything converges to everything]\label{ex:to-conv-chaotic}
Let $M$ carry the chaotic topology \ts{4}{1}{1:08:43}. Being a sequence has nothing to do with the topology; it is a map into the
set \ts{4}{1}{1:09:05}. Claim: \emph{every sequence converges against every point} \ts{4}{1}{1:09:25}. Indeed let $q$ be arbitrary and
$a\in M$ arbitrary \ts{4}{1}{1:09:48}. Of the two open sets only $M$ contains $a$, so we must check $\exists N\ \forall n>N\colon q(n)\in M$,
which holds for every $n$ as $q$ maps into $M$, independently of the point and of the sequence
\ts{4}{1}{1:10:09}\ts{4}{1}{1:10:29}\ts{4}{1}{1:10:50}.
\end{example}
Hence ``chaotic'', and not useful \ts{4}{1}{1:10:50}.
\begin{aside}[A quick exam solution]
Schuller's joke: if an examiner asks to show that a sequence converges but is not careful enough to specify the topology, you
can claim that it converges against every point, since you choose the chaotic topology and were not asked to choose another one
\ts{4}{1}{1:10:50}\ts{4}{1}{1:11:11}.
\end{aside}

\begin{example}[Discrete topology: only almost constant sequences]\label{ex:to-conv-discrete}
Equip $M$ with the discrete topology \ts{4}{1}{1:11:33}. Every set is open, in particular every $\{a\}$ is an open
neighbourhood of $a$, so a sequence converges against $a$ only if $q(n)=a$ for all large $n$. So only \emph{almost constant}
sequences converge, and they converge against that constant; ``almost constant'' means constant everywhere apart from a finite
exception set \ts{4}{1}{1:11:33}\ts{4}{1}{1:11:53}\ts{4}{1}{1:12:14}.
\end{example}

\begin{example}[$\R^d$ with the standard topology]\label{ex:to-conv-standard}
Let $M=\R^d$ with $\mathcal O_{\rm standard}$ \ts{4}{1}{1:12:37}. One can prove that a sequence $q$ converges against $a\in\R^d$
if and only if
\[
  \forall\varepsilon>0\ \exists N\in\N\ \forall n>N\colon\ \|q(n)-a\|<\varepsilon,
\]
for some (any) norm \ts{4}{1}{1:12:37}\ts{4}{1}{1:13:00}\ts{4}{1}{1:13:21}\ts{4}{1}{1:13:42}.
\end{example}
\begin{exercise}[The $\varepsilon$ criterion]\label{ex:to-conv-eps}
Prove the equivalence in \cref{ex:to-conv-standard} from \cref{def:to-conv} and the definition of $\mathcal O_{\rm standard}$
(stated in class without proof).
\end{exercise}
A (garbled) question from a student, apparently about the relation with the analysis way of saying ``for all\dots'' [?], is answered
as follows \ts{4}{1}{1:13:42}\ts{4}{1}{1:14:03}. The discrete topology is the \emph{finest} one: everything convergent in it is convergent
in any other topology, but in topologies that are not that fine, with fewer subsets chosen, there are more convergent
sequences, which do not converge in the finer one \ts{4}{1}{1:13:42}\ts{4}{1}{1:14:03}\ts{4}{1}{1:14:23}.

\begin{example}[$q(n)=1-1/n$]\label{ex:to-conv-1n}
In $M=\R$ take $q(n)=1-\frac1n$ for $n\ge1$ (start at $n=1$ not to run into trouble at zero) \ts{4}{1}{1:14:23}\ts{4}{1}{1:14:45}. On the real line the
points are $0,\frac12,\frac23,\dots$, closer and closer to $1$ \ts{4}{1}{1:14:45}\ts{4}{1}{1:15:05}. Every value is different, so the sequence is not almost constant, and
it does \emph{not} converge in $\R$ with the discrete topology; but it converges against $1$ in $\R$ with the standard
topology \ts{4}{1}{1:15:05}\ts{4}{1}{1:15:34}\ts{4}{1}{1:15:57}.
\end{example}
The intuition of ``getting closer and closer'' is that of the standard topology, which uses the norm, i.e.\ the distance
\ts{4}{1}{1:15:57}\ts{4}{1}{1:16:18}. It generalises to every \emph{metric space}: from a metric (a distance function) one can construct a
topology; but by far not every topology is induced by a distance map \ts{4}{1}{1:16:18}\ts{4}{1}{1:16:39}.

\begin{secondary}[Digression: research proposals, non-transitive dice and topologies of theories]
A half-serious side remark \ts{4}{1}{1:16:39}\ts{4}{1}{1:17:00}. To get research money one writes a proposal of the type ``we do this
and that, because then we will know more''; in physics: ``we come closer to the truth'', or ``in recent years we have converged
against the true theory''. After the (spiritual) approval of the office for gender and diversity it goes to experts, who have no time
and read only the introduction, which must say that we are getting better and making progress
\ts{4}{1}{1:17:00}\ts{4}{1}{1:17:20}\ts{4}{1}{1:17:40}.

What would this mean? One way to say that theories make progress is that they get better and better; for that one needs a measure
of whether one theory is better than another, a quality function $Q$ \ts{4}{1}{1:17:40}\ts{4}{1}{1:18:01}. Consider \textbf{non-transitive dice}
(Schuller said ``non-commutative'' and corrected himself): three dice $a,b,c$ with different distributions of numbers on their
faces (the distribution on the board was only illustrative) \ts{4}{1}{1:18:01}\ts{4}{1}{1:18:27}\ts{4}{1}{1:18:48}. Two dice play
against each other: both are thrown, the one with the higher number gets a point, and after a long run the one with more points
wins \ts{4}{1}{1:19:08}\ts{4}{1}{1:19:29}\ts{4}{1}{1:19:51}. One can calculate with which probability each wins, so one could say that
the quality of $a$ was greater than that of $b$ if $a$ wins against $b$ \ts{4}{1}{1:19:51}\ts{4}{1}{1:20:12}. Similarly $b$ may win
against $c$. Then one would conclude that $a$ wins against $c$, but remarkably there is a distribution of numbers on the three dice such
that $c$ wins against $a$: from $Q(a)>Q(b)$ and $Q(b)>Q(c)$ one cannot conclude $Q(a)>Q(c)$ \ts{4}{1}{1:20:12}\ts{4}{1}{1:20:34}\ts{4}{1}{1:20:54}.
So no such real-valued quality measure exists \ts{4}{1}{1:20:54}\ts{4}{1}{1:21:14}.

This is remarkable: anybody who assigns to complex situations, such as research proposals or the abilities of a girl or a boy, a
real-valued function measuring them on one scale, and compares on it, makes an assumption on which economics, for instance, is built
\ts{4}{1}{1:21:14}\ts{4}{1}{1:21:35}. Already the little complexity of numbers on dice has no quality function; so nobody should say
``this theory is better than that theory'' without being asked immediately for the quality function, which is generically
non-transitive; only special distributions are transitive \ts{4}{1}{1:21:35}\ts{4}{1}{1:21:57}.

One might say: even without a quality function, the theories come \emph{closer to the true theory}. That would be a very weak but
still impressive statement; but to make it one needs to know a \emph{topology on the space of all theories}, which is far too
complex; Schuller does not know how to write it down, and assumes, unless somebody convinces him otherwise, that statements of the
type ``theories converge somehow'' are not meaningful \ts{4}{1}{1:21:57}\ts{4}{1}{1:22:17}\ts{4}{1}{1:22:38}\ts{4}{1}{1:22:58}. Since no
structure making fewer assumptions than a topology is known by which to define convergence, topology also reveals nonsensical or
ill-defined statements that sound good; it raises at least sufficient suspicion about any such claim, and applies to other
circumstances too \ts{4}{1}{1:22:58}\ts{4}{1}{1:23:20}\ts{4}{1}{1:23:41}.
\end{secondary}
\begin{aside}[Hairdressers]
A passing joke: ``what hairdressers can do, only hairdressers can do; in logic, in topology, maybe you can conclude\dots''
\ts{4}{1}{1:23:20}. (Captions garbled; the point is the one of the digression above.) [?]
\end{aside}
\begin{history}[Efron's dice]
\emph{Not said in class.} \emph{Checked.} Sets of non-transitive dice were found by the statistician Bradley Efron and were
popularised by Martin Gardner in \emph{Scientific American} (Dec.\ 1970).
\end{history}

% ==================================================================
\section{Continuity}\label{sec:to-cont}

Convergence was interesting, but continuity is at the heart of topology, as it will finally be used to define the structure
preserving maps \ts{4}{1}{1:23:41}\ts{4}{1}{1:24:02}.

\begin{definition}[New concepts: continuous map]\label{def:to-cont}
Let $(M,\mathcal O_M)$ and $(N,\mathcal O_N)$ be topological spaces and $f\colon M\to N$ a map between the underlying sets. Then $f$
is called \textbf{continuous} if for every open set $V\in\mathcal O_N$ in the target space the pre-image $\operatorname{preim}_f(V)$
(\cref{def:cl-preimage}) is an open set in the domain, $\operatorname{preim}_f(V)\in\mathcal O_M$
\ts{4}{1}{1:24:02}\ts{4}{1}{1:24:23}\ts{4}{1}{1:24:43}\ts{4}{1}{1:25:07}\ts{4}{1}{1:25:29}\ts{4}{1}{1:25:50}.
\end{definition}
Extremely simple: ``pre-images of open sets are open''. That hides the fact that the first openness is in the target space and the second in the domain, so more precisely
$\forall V\in\mathcal O_N\colon\operatorname{preim}_f(V)\in\mathcal O_M$ \ts{4}{1}{1:25:50}\ts{4}{1}{1:26:14}\ts{4}{1}{1:26:39}. It is at the heart of topology, as will be
seen \ts{4}{1}{1:26:39}.
\begin{recall}[Pre-image]
The pre-image of $V\subseteq N$ is the set of all points $m\in M$ with $f(m)\in V$, \cref{def:cl-preimage}
\ts{4}{1}{1:25:29}.
\end{recall}

\begin{example}[Discrete domain: every map is continuous]\label{ex:to-cont-discrete}
Take any map $f\colon M\to N$ between arbitrary sets, with the discrete topology on $M$ and any topology on $N$ \ts{4}{1}{1:26:39}\ts{4}{1}{1:27:00}.
Then $f$ is continuous \ts{4}{1}{1:27:20}: for any open $V\subseteq N$ the pre-image is a subset of $M$, and every subset of $M$ is open (for $V=\emptyset$ the
pre-image is the empty set; no need to distinguish cases) \ts{4}{1}{1:27:20}\ts{4}{1}{1:27:40}\ts{4}{1}{1:28:01}. No property of $f$ was used.
\end{example}
If every map is continuous, continuity is a useless property; so one \emph{never} equips the domain with the discrete topology
if one still wants to work with continuity \ts{4}{1}{1:28:01}\ts{4}{1}{1:28:21}.

\begin{example}[Chaotic target: every map is continuous]\label{ex:to-cont-chaotic}
Take any map $M\to N$ with any topology on $M$ and the chaotic topology on $N$: again every map is continuous, as one can quickly
think about (the only non-empty open set of $N$ is $N$, with pre-image $M$, and $\emptyset$ has pre-image $\emptyset$)
\ts{4}{1}{1:28:21}\ts{4}{1}{1:28:42}\ts{4}{1}{1:29:02}.
\end{example}

\begin{example}[$\R^d\to\R^f$: the $\varepsilon$--$\delta$ criterion]\label{ex:to-cont-eps}
For $f\colon\R^d\to\R^{f'}$ with both spaces carrying the standard topology one recovers the usual definition of continuity: for
every $\varepsilon>0$ there is $\delta>0$ such that for all $x,y$ with $\|x-y\|<\delta$ one has $\|f(x)-f(y)\|<\varepsilon$ \ts{4}{1}{1:29:02}\ts{4}{1}{1:29:24}. The
$\varepsilon$--$\delta$ criterion is recovered from the general definition, which makes sense for any pair of topological spaces
\ts{4}{1}{1:29:24}\ts{4}{1}{1:29:44}.
\end{example}
\begin{coursenote}[A correction on the board]
A student points at the target dimension (written with the letter of the map): it must be an $n$ (here $f'$ to avoid a clash with the map $f$)
\ts{4}{1}{1:29:44}\ts{4}{1}{1:30:04}. ``Thank you very much, yes of course.''
\end{coursenote}
\begin{exercise}[$\varepsilon$--$\delta$]\label{ex:to-cont-eps-proof}
Derive the $\varepsilon$--$\delta$ criterion from \cref{def:to-cont} for $\mathcal O_{\rm standard}$ on $\R^d$ and $\R^{f'}$ (stated in
class without proof).
\end{exercise}

% ==================================================================
\section{Homeomorphisms: structure preserving maps of topological spaces}\label{sec:to-homeo}

\begin{definition}[New concepts: homeomorphism, homeomorphic]\label{def:to-homeo}
Let $f\colon M\to N$ be a bijection (\cref{def:cl-bij}) between two sets and make topological spaces out of them with $\mathcal O_M$ and
$\mathcal O_N$. Then $f$ is called a \textbf{homeomorphism} if \ts{4}{1}{1:30:04}\ts{4}{1}{1:30:26}\ts{4}{1}{1:30:47}
\begin{enumerate}
  \item $f\colon M\to N$ is continuous, and
  \item the inverse $f^{-1}\colon N\to M$ (it exists and is unique since $f$ is a bijection, \cref{def:cl-inverse}) is continuous
        \ts{4}{1}{1:31:08}\ts{4}{1}{1:31:30}.
\end{enumerate}
Two topological spaces between which a homeomorphism exists are called \textbf{homeomorphic} \ts{4}{1}{1:35:07}, and one writes
$(M,\mathcal O_M)\cong_{\rm top}(N,\mathcal O_N)$ [?: the same isomorphism sign as in set theory is used].
\end{definition}
The word ``homeomorphism'' is long, so one says ``homeo'' and suppresses ``morphism'' \ts{4}{1}{1:31:30}\ts{4}{1}{1:31:51}. The origin of the name:
\emph{homoios} (Greek), similar, and \emph{morphe}, shape (not said in class).

Homeomorphisms are \emph{the structure preserving maps} in topology, the central notion of map \ts{4}{1}{1:30:47}\ts{4}{1}{1:31:51}. Why? $f$ takes
one from $M$ to $N$ as a set, and $f^{-1}$ back; since $f$ is continuous the pre-image of any open set is open, and the
pre-image under $f^{-1}$ is the image under $f$; so, if a homeomorphism exists, it provides a one-to-one pairing of the open sets of $M$ with those of $N$
according to the chosen topologies \ts{4}{1}{1:31:51}\ts{4}{1}{1:32:15}\ts{4}{1}{1:32:36}\ts{4}{1}{1:32:57}\ts{4}{1}{1:33:21}. One can say that the topologies
have the same structure \ts{4}{1}{1:33:21}\ts{4}{1}{1:33:44}.

\begin{proposition}[Homeomorphisms pair up open sets]\label{prop:to-homeo-pair}
Let $f\colon M\to N$ be a homeomorphism. Then $U\mapsto f(U)$ is a bijection $\mathcal O_M\to\mathcal O_N$ with inverse $V\mapsto f^{-1}(V)$.
\end{proposition}
\begin{proof}
(Not written out in class; the content of the argument above.) For $U\in\mathcal O_M$, $f(U)$ is the pre-image of $U$ under the continuous map
$f^{-1}$, hence open in $N$; for $V\in\mathcal O_N$, $f^{-1}(V)$ is the pre-image under the continuous map $f$, so open in $M$. The two assignments are inverse to each other
because $f$ is a bijection.
\end{proof}

\begin{intuition}[Isomorphic as sets, different as topological spaces]\label{int:to-hair}
A bijection exists, so the two sets are isomorphic as sets, ignoring the topology; if they are also homeomorphic they are isomorphic as topological
spaces, which is more \ts{4}{1}{1:33:44}\ts{4}{1}{1:34:05}\ts{4}{1}{1:34:26}\ts{4}{1}{1:34:47}. The same isomorphism sign as in set theory is used, but the bijection
now has to satisfy the additional homeomorphism condition. Conversely, topological isomorphism implies set-theoretic isomorphism; but two sets that are
isomorphic can, once equipped with topologies, become non-isomorphic (different topologically) \ts{4}{1}{1:34:47}\ts{4}{1}{1:35:07}\ts{4}{1}{1:35:30}\ts{4}{1}{1:35:51}.
Things that are the same up to reshuffling can become non-isomorphic at a higher structural level by establishing different structures
\ts{4}{1}{1:35:51}. ``Two children, both blonde; now I dye their hair, one blonde, the other red: now they are different. One finds other distinguishing
properties, but fine'' \ts{4}{1}{1:36:13}.
\end{intuition}

\paragraph{Classification.} Recall the recurrent theme of mathematics, \cref{int:cl-theme}: classify spaces by the structure preserving maps between them
\ts{4}{1}{1:36:13}\ts{4}{1}{1:36:33}. In set theory one classifies by cardinality (not carried far, there is a very big theory behind it); for finite sets, by the number of
elements: two finite sets with the same number of elements are isomorphic as sets \ts{4}{1}{1:36:33}\ts{4}{1}{1:36:53}.
For topological spaces one would like properties that distinguish them in the same way; such properties will be defined in the next lecture
\ts{4}{1}{1:36:53}\ts{4}{1}{1:37:13}. But nobody knows a \emph{complete} set of properties, like the number of elements for sets, such that spaces with the same
values of all of them are already homeomorphic \ts{4}{1}{1:37:13}\ts{4}{1}{1:37:34}\ts{4}{1}{1:37:55}. There is no classification of topological spaces as far as Schuller
knows (``there are just too many, or we haven't thought carefully enough''); this is an open classification problem. One can prove that, ``and then make donations
to the university'': it is important to know what the open problems of a field are, to solve one of them one had better know they are open
\ts{4}{1}{1:37:55}\ts{4}{1}{1:38:15}.
\begin{coursenote}[On the open problem]
Not said in class: the homeomorphism problem is undecidable for manifolds of dimension $\ge4$ (Markov, 1958), so no complete, algorithmic list
of invariants exists there; for general topological spaces the situation is wilder still. Stated here from memory, not checked.
\end{coursenote}

% ==================================================================
\section*{Key concepts}
\begin{keyconcepts}
\textbf{Retrieval questions.}
\begin{enumerate}
  \item State the three conditions of a topology. Which are finite, which arbitrary?
  \item Why are the chaotic, the discrete and the $\{\emptyset,\{1\},\{2\},\{1,2\},\{1,2,3\}\}$ examples topologies?
  \item How is the standard topology on $\R^d$ built in three steps, and why does the choice of norm not matter?
  \item Why is $(0,1]$ open in $[-1,1]$ but not in $\R$? Why is $\{p\}$ not open in $S^1\subseteq\R^2$?
  \item Define closed sets; why are $\emptyset$ and $M$ both open and closed?
  \item Which sequences converge in the chaotic, which in the discrete topology? Does $1-\frac1n$ converge in the standard one?
  \item Define continuity. Why is every map continuous out of a discrete space and into a chaotic one?
  \item Define homeomorphism; why does the classification of topological spaces remain open?
\end{enumerate}
\textbf{Concepts.}
\begin{itemize}
  \item Topology, open sets, topological space (\cref{def:to-topology,def:to-space}); standard topology on $\R^d$ (\cref{prop:to-standard}).
  \item Induced topology $\mathcal O|_N=\{U\cap N\}$ (\cref{def:to-induced}); closed set $=$ open complement (\cref{def:to-closed,prop:to-clopen}).
  \item Product topology (\cref{def:to-product}); circle as subset and as quotient.
  \item Convergence: every open neighbourhood contains the tail (\cref{def:to-conv}).
  \item Continuity: pre-images of open sets are open (\cref{def:to-cont}); homeomorphism (\cref{def:to-homeo}).
\end{itemize}
\end{keyconcepts}

\flushside
\end{paracol}

% !TEX root = ../main.tex
% STUB: generated by tools/stubs.py; write the chapter from materials/transcripts/C05-V01.txt
\chapter{Topological spaces - some heavily used invariants}\label{ch:l05}
\begin{paracol}{2}

\begin{sources}
\begin{tabularx}{\linewidth}{@{}l l L@{}}
\toprule
\textbf{Video} & \textbf{Length} & \textbf{Content}\\
\midrule
\seg{5}{1}{Lecture 5: Topological spaces - some heavily used invariants} & 1:56:44 & (to be written)\\
\bottomrule
\end{tabularx}

\smallskip
\textbf{Transcript} \file{materials/transcripts/C05-V01.txt}.
\end{sources}

\begin{reading}[Where this lecture sits in the other notes]
\emph{Gravity and Light}, Lecture 1 and Tutorial 1. (From the lecture titles; to be checked when the chapter is written.)
\end{reading}

\emph{This chapter has not been written yet.}

\flushside
\end{paracol}

% !TEX root = ../main.tex
% STUB: generated by tools/stubs.py; write the chapter from materials/transcripts/C06-V01.txt
\chapter{Topological manifolds and manifold bundles}\label{ch:l06}
\begin{paracol}{2}

\begin{sources}
\begin{tabularx}{\linewidth}{@{}l l L@{}}
\toprule
\textbf{Video} & \textbf{Length} & \textbf{Content}\\
\midrule
\seg{6}{1}{Lecture 6: Topological manifolds and manifold bundles} & 1:49:18 & (to be written)\\
\bottomrule
\end{tabularx}

\smallskip
\textbf{Transcript} \file{materials/transcripts/C06-V01.txt}.
\end{sources}

\begin{reading}[Where this lecture sits in the other notes]
\emph{Gravity and Light}, Lecture 2 (Topological manifolds); \emph{Theoretische Mechanik}, Lecture 2. (From the lecture titles; to be checked when the chapter is written.)
\end{reading}

\emph{This chapter has not been written yet.}

\flushside
\end{paracol}

% !TEX root = ../main.tex
\chapter{Differentiable structures: definition and classification}\label{ch:l07}
\begin{paracol}{2}

\begin{sources}
\begin{tabularx}{\linewidth}{@{}l l L@{}}
\toprule
\textbf{Video} & \textbf{Length} & \textbf{Content}\\
\midrule
\seg{7}{1}{Lecture 7: Differentiable structures: definition and classification} & 1:14:34 & Refining the maximal $C^0$ atlas: $C^k$, smooth, analytic and complex atlases;
the theorem that a maximal $C^k$ atlas ($k\ge1$) contains a unique smooth one; differentiable manifolds; two incompatible smooth atlases of the real line; differentiability of maps between manifolds and its
well-definedness; diffeomorphisms; levels of structure on a set (and student questions); how many smooth structures a topological manifold carries (dimensions 1--3, above 4, exactly 4);
tangent spaces without embedding; outlook to tensor spaces and modules.\\
\bottomrule
\end{tabularx}

\smallskip
\textbf{Transcript} \file{materials/transcripts/C07-V01.txt} (YouTube automatic captions).
\end{sources}

\begin{reading}[Where this lecture sits in the other notes]
\emph{Gravity and Light}, Lecture~4 (Differentiable manifolds); \emph{Theoretische Mechanik}, Lecture~3 (Differentiable manifolds). Both introduce the $C^k$ atlas and the
differentiability of maps through charts; the classification results of this lecture (unique smooth structure up to dimension~3, uncountably many on $\R^4$) are, as far as the titles say, specific to the present course.
(From the lecture titles and the plan of the course; not checked in detail.)
\end{reading}

% ==================================================================
\section{Adding structure: refining the atlas}\label{sec:df-atlas}

Today differentiable manifolds \ts{7}{1}{0:08}. One passes from topological to differentiable manifolds by \emph{removing charts} from the topological atlas: this is the first step, ``adding structure by refining the maximal topological atlas'', the maximal $C^0$ atlas
\ts{7}{1}{0:08}\ts{7}{1}{0:34}\ts{7}{1}{0:55}.
Recall (\cref{def:mf-c0atlas,prop:mf-c0}): on a topological manifold $M_0$ there is essentially only one maximal topological atlas, because all charts one can choose are automatically $C^0$ compatible by the very definition of a topological manifold; for a topological manifold an atlas was a fully
redundant notion. Now something non-redundant is done \ts{7}{1}{0:55}\ts{7}{1}{1:16}\ts{7}{1}{1:37}.
\begin{recall}[$C^0$ compatibility]
Two charts $(U,x)$, $(V,y)$ are $C^0$ compatible if $U\cap V=\emptyset$ or the transition map $y\circ x^{-1}\colon x(U\cap V)\to y(U\cap V)$ is continuous (\cref{def:mf-compat}); for a topological manifold this always holds (\cref{prop:mf-c0}):
$x$ and $y$ are homeomorphisms, hence so is $y\circ x^{-1}$.
\end{recall}

\begin{definition}[New concepts: $C^\star$ atlas]\label{def:df-star}
Let $M_0=(M,\mathcal O)$ be a topological manifold of dimension $D$. An atlas $\mathcal A$ of $M_0$ is called a \textbf{$C^\star$ atlas} (where $C^\star$ stands for one of the compatibility classes below; the symbol is heard in the transcript as ``flower'' and written
as a placeholder, so I use $\star$ [?]) if any two charts $(U,x)$ and $(V,y)$ in $\mathcal A$ are $C^\star$ compatible, that is: either
\begin{enumerate}
  \item $U\cap V=\emptyset$ (the two chart regions, open subsets of the manifold, do not meet: then they are already compatible), or
  \item $U\cap V\ne\emptyset$ and the chart transition map
        \[
          y\circ x^{-1}\colon x(U\cap V)\to y(U\cap V)
        \]
        (both $x(U\cap V)$ and $y(U\cap V)$ are subsets of $\R^D$, $D=\dim M$) is $C^\star$ as a map $\R^D\supseteq x(U\cap V)\to y(U\cap V)\subseteq\R^D$
\end{enumerate}
\ts{7}{1}{1:37}\ts{7}{1}{2:02}\ts{7}{1}{2:28}\ts{7}{1}{2:48}\ts{7}{1}{3:09}\ts{7}{1}{3:30}\ts{7}{1}{3:51}\ts{7}{1}{4:12}.
\end{definition}
The chart transition map is not a map between manifolds but a map from $\R^{\dim M}$ to $\R^{\dim M}$. In a topological manifold it is always a homeomorphism, because $x$ and $y$ are; but now one does not want the transition maps to be pairwise $C^0$ compatible, one wants them
to be $C^\star$: ``before you think I finally turned nuts, what is $C^\star$?'' \ts{7}{1}{4:12}\ts{7}{1}{4:33}\ts{7}{1}{4:54}\ts{7}{1}{5:15}. There are several possibilities for what $C^\star$ can be:

\begin{definition}[New concepts: $C^0$, $C^k$, smooth, analytic and complex atlases]\label{def:df-ck}
For the manifold $M_0$ and an atlas $\mathcal A$:
\begin{itemize}
  \item \textbf{$C^0$:} the transition maps are continuous as maps $\R^D\to\R^D$ (as last time) \ts{7}{1}{5:15}\ts{7}{1}{5:36}.
  \item \textbf{$C^k$:} the transition maps are $k$ times continuously differentiable as maps $\R^D\to\R^D$. This is a condition pushed down to the transition maps, which is meaningful since differentiability on $\R^D$ is understood; the definition of $C^k$ \emph{for manifolds} is what
        is about to be constructed. The atlas is called a $C^k$ atlas, and one speaks of a manifold with $k$ times continuously differentiable transition functions \ts{7}{1}{5:36}\ts{7}{1}{5:57}\ts{7}{1}{6:19}\ts{7}{1}{6:40}.
  \item \textbf{$C^\infty$:} one can differentiate arbitrarily often; such atlases (and manifolds) are called \textbf{smooth}. (Very often used, as will be seen) \ts{7}{1}{6:40}\ts{7}{1}{7:01}.
  \item \textbf{$C^\omega$:} (an omega, looking a little like an infinity) \textbf{real analytic}: the transition functions are real analytic, i.e.\ can be Taylor expanded around every point of their domain; the atlas gives an \textbf{analytic manifold}
        \ts{7}{1}{7:22}\ts{7}{1}{7:43}\ts{7}{1}{8:04}.
  \item \textbf{Complex:} one speaks of a \textbf{complex atlas} and a \textbf{complex manifold} (details below) \ts{7}{1}{8:04}\ts{7}{1}{8:24}.
\end{itemize}
\end{definition}
(One argues in a second, after a few more definitions, that for $k\ge1$ there is no big difference between $C^k$ and $C^\infty$ \ts{7}{1}{7:01}\ts{7}{1}{7:22}.)
Analyticity is a stronger condition than $C^\infty$: every analytic function is $C^\infty$, but not every $C^\infty$ function is analytic, and the analysis class gives examples \ts{7}{1}{8:04}\ts{7}{1}{8:24}.

\paragraph{Complex manifolds.} So far one always mapped into some $\R^D$; how does this become complex? There are different possibilities. One can keep the definition as it is, but then one must take an even-dimensional real manifold, which is a \emph{complex} manifold if the transition functions,
being continuous anyway (``I must emphasise this, otherwise they could not come from an atlas as defined''), satisfy the \textbf{Cauchy--Riemann equations} \ts{7}{1}{8:24}\ts{7}{1}{8:46}\ts{7}{1}{9:07}\ts{7}{1}{9:27}.
\begin{definition}[Cauchy--Riemann equations]\label{def:df-cr}
Look at a map $f\colon\R^2\to\R^2$; as a set the complex plane is $\C\cong\R^2$, and a complex number is a pair $(x,y)\leftrightarrow x+iy$. Decompose the dependent variable too,
\[
  f(x+iy)=u(x,y)+i\,v(x,y),
\]
with two real functions $u,v$ of two real variables. The \textbf{Cauchy--Riemann equations} are
\[
  \frac{\partial u}{\partial x}=\frac{\partial v}{\partial y},\qquad \frac{\partial u}{\partial y}=-\frac{\partial v}{\partial x}.
\]
A \emph{continuous} map $\R^2\to\R^2$ satisfying them is as good as a complex differentiable function $\C\to\C$ \ts{7}{1}{9:27}\ts{7}{1}{9:48}\ts{7}{1}{10:08}\ts{7}{1}{10:29}\ts{7}{1}{10:49}\ts{7}{1}{11:10}\ts{7}{1}{11:32}.
\end{definition}
(A student's remark at \ts{7}{1}{12:35}: the map is differentiable from $\C$ to $\C$, not $\R^2\to\R^2$; Schuller agrees: $\R^2$ is the complex plane. The lecturer first asked whether the Cauchy--Riemann equations were known; ``no? okay, look it up''
\ts{7}{1}{9:27}\ts{7}{1}{9:48}.)
Differentiability in the complex plane is a much stronger requirement than over the reals, because very roughly one can approach the limit point in so many different ways; the field studying it is function theory (complex analysis) \ts{7}{1}{11:32}\ts{7}{1}{11:52}. With more variables one
gets indices, $\partial u^a/\partial x^b=\ldots$, and one imposes the Cauchy--Riemann equations on the transition maps; since one always has to split into two, this requires $\dim M$ to be even as a real manifold \ts{7}{1}{11:52}\ts{7}{1}{12:13}\ts{7}{1}{12:35}.
\begin{coursenote}[The higher-dimensional equations]
For $z^a=x^a+iy^a$, $a=1,\dots,n$ and a map $f^a=u^a+iv^a$ the equations read $\partial u^a/\partial x^b=\partial v^a/\partial y^b$, $\partial u^a/\partial y^b=-\partial v^a/\partial x^b$; the dimension of the real manifold is then $2n$. Not written out in class (``you get indices here'').
\end{coursenote}
So one can add additional or different structure and pick which one one wants; by far the most important are the smooth differentiable structures and the complex ones, for reasons explained in a moment \ts{7}{1}{12:35}\ts{7}{1}{12:55}.

% ==================================================================
\section{Smooth structure is not a big choice}\label{sec:df-whitney}

\begin{theorem}[Whitney]\label{thm:df-whitney}
Let $k\ge1$ (very important: this does \emph{not} include $k=0$). Any maximal $C^k$ atlas contains a $C^\infty$ atlas; and two maximal $C^k$ atlases that contain the same $C^\infty$ atlas are already identical
\ts{7}{1}{12:55}\ts{7}{1}{13:17}\ts{7}{1}{13:39}\ts{7}{1}{14:01}\ts{7}{1}{14:26}.
\end{theorem}
(The theorem is stated, not proved, in class; the name is given as ``I think Whitney''.)
\begin{coursenote}[On the attribution and the wording]
In the transcript the name is ``Whitney'' with ``I think'', and the statement is a loose one; the usual form, due to H.~Whitney (1936), is that every $C^k$ structure, $k\ge1$, contains a compatible smooth one, unique up to diffeomorphism.
Not said in class; ``contains'' is as in class.
\end{coursenote}
\emph{Implication.} Once one has found a $C^1$ atlas for a manifold, it is not a stronger condition to say that the manifold is already $C^\infty$: it is already smooth; a $C^1$ atlas can be turned into a $C^\infty$ one by removing more charts \ts{7}{1}{14:26}\ts{7}{1}{14:49}\ts{7}{1}{15:10}.
The converse is the content of the next sentence: one could have a topological manifold with \emph{no} differentiable atlas at all; but once one sees one, one can go all the way down to $C^\infty$ \ts{7}{1}{15:10}\ts{7}{1}{15:30}. So it may be clever to say that a theorem valid for $C^2$ manifolds is automatically valid for
$C^3,C^4,C^5,\dots$ since they are contained, and one can construct from a $C^2$ (even a $C^1$) manifold a smooth one. When talking about differentiable manifolds or the existence of a differentiable structure one need not distinguish between the cases $C^k$ with $k\ge1$ and $C^\infty$, so
smooth in the above sense \ts{7}{1}{15:30}\ts{7}{1}{15:51}\ts{7}{1}{16:11}\ts{7}{1}{16:37}.

% ==================================================================
\section{Differentiable manifolds}\label{sec:df-manifold}

What is done with the differentiable structure? First one defines the notion of a differentiable manifold, then differentiable maps between two manifolds; these are the structure preserving maps if they are also invertible, which gives a \emph{classification} of smooth manifolds, with the first surprises, even
from the point of view of fundamental physics \ts{7}{1}{16:37}\ts{7}{1}{17:00}\ts{7}{1}{17:22}.

\begin{definition}[New concepts: $C^k$ manifold]\label{def:df-manifold}
A \textbf{$C^k$ manifold} is a triple $(M,\mathcal O,\mathcal A)$ where $(M,\mathcal O)$ is a topological manifold (``of course that stays'') and $\mathcal A$ is a \emph{maximal} $C^k$ atlas \ts{7}{1}{17:22}\ts{7}{1}{17:45}\ts{7}{1}{18:09}.
\end{definition}
One has a new additional layer of structure: the $C^k$ atlas is always a subset of a $C^0$ atlas, one takes out strategically certain charts \ts{7}{1}{18:09}\ts{7}{1}{18:31}.

\begin{remark}[Incompatible atlases]\label{rem:df-incompat}
A given topological manifold can carry different \emph{incompatible} atlases. Only compatibility of charts was defined, but what is meant is clear: technically two atlases are compatible if their union is again an atlas of the same type; or: if between any two charts of the two atlases the
compatibility of type $C^\star$ holds \ts{7}{1}{18:31}\ts{7}{1}{18:58}\ts{7}{1}{19:18}.
\end{remark}

\begin{example}[Two incompatible smooth atlases of the real line]\label{ex:df-line}
Let the underlying topological manifold be $\R$ with the standard topology \ts{7}{1}{19:18}\ts{7}{1}{19:39}.
\begin{itemize}
  \item \emph{Atlas one:} $\mathcal A_1=\{(\R,\mathrm{id}_\R)\}$, a single chart covering all of $\R$ (as before, \cref{ex:mf-atlases}). It is trivially a $C^\infty$ atlas, even $C^1$, everything: the chart overlaps only with itself and $\mathrm{id}\circ\mathrm{id}^{-1}=\mathrm{id}$ is $C^\infty$
        \ts{7}{1}{19:39}\ts{7}{1}{20:00}\ts{7}{1}{20:21}.
  \item \emph{Atlas two:} $\mathcal A_2=\{(\R,x)\}$ again with the entire real line as chart region, but with chart map
        \[
          x\colon\R\to\R,\qquad a\mapsto a^{1/3}
        \]
        (the third root) \ts{7}{1}{20:21}\ts{7}{1}{20:41}. First check that this is a chart map at all: it is one-to-one, continuous, invertible with the third power as inverse, which is continuous too: a homeomorphism from the manifold $\R$ down to the chart $\R^1$. Its transition map
        with itself, $x\circ x^{-1}=\mathrm{id}$, is $C^\infty$, so this is also a $C^\infty$ atlas \ts{7}{1}{20:41}\ts{7}{1}{21:03}\ts{7}{1}{21:25}\ts{7}{1}{21:46}\ts{7}{1}{22:06}.
\end{itemize}
Both can be made into \emph{maximal} $C^\infty$ atlases by adding all the other charts whose transition functions with the chart(s) already there are $C^\infty$ (call them $\mathcal A_1^{\max}$ and $\mathcal A_2^{\max}$) \ts{7}{1}{22:06}\ts{7}{1}{22:27}.
But the two atlases are \emph{incompatible}: were they compatible, their union would be another $C^\infty$ atlas, but the transition map from $(\R,\mathrm{id})$ to $(\R,x)$, $x\circ\mathrm{id}^{-1}=a\mapsto a^{1/3}$, though continuous (it must be), is not $C^\infty$, not even $C^1$: its derivative at $0$ is infinite, not defined
\ts{7}{1}{22:27}\ts{7}{1}{22:48}\ts{7}{1}{23:10}\ts{7}{1}{23:31}\ts{7}{1}{23:53}. So the two charts are not $C^1$ compatible: as long as they are parts of different atlases that is no problem, but they could not be put into the same atlas and give a $C^\infty$ atlas \ts{7}{1}{23:31}\ts{7}{1}{23:53}.
One can therefore equip the real line with at least two different incompatible $C^\infty$ structures (just one example) \ts{7}{1}{23:53}\ts{7}{1}{24:13}.
\end{example}
That looks bad: one seems to have to decide which smooth structure to establish; it is not so much a choice (one would have to do it for all the $C^k$ levels, which would be even worse), but even looking only at smooth structures, a topological manifold does not have a unique one. If one wants to do physics, and physics
wants differentiability, ``what a disaster this is: not even the real line has a unique differentiable structure'' \ts{7}{1}{24:13}\ts{7}{1}{24:34}\ts{7}{1}{24:56}. The situation is not so bad, but to explain this one needs the next definition \ts{7}{1}{24:56}\ts{7}{1}{25:17}.

% ==================================================================
\section{Differentiable maps and diffeomorphisms}\label{sec:df-maps}

A given $C^k$ manifold is equipped with one specific choice of atlas; assume this pick is made (which pick? \ts{7}{1}{25:17}) \ts{7}{1}{25:17}\ts{7}{1}{25:38}.
\begin{definition}[New concepts: differentiable map between manifolds]\label{def:df-map}
Let $\phi\colon M\to N$ be a map, where $(M,\mathcal O_M,\mathcal A_M)$ and $(N,\mathcal O_N,\mathcal A_N)$ are $C^k$ manifolds (for a map it would suffice that $M,N$ are sets; for a continuous map topological spaces would be needed; now both are differentiable manifolds)
\ts{7}{1}{25:17}\ts{7}{1}{25:38}\ts{7}{1}{25:59}. Then $\phi$ is called \textbf{differentiable at the point $p\in M$} (of class $C^k$ at $p$) if for \emph{some} chart $(U,x)$ in the atlas $\mathcal A_M$ with $p\in U$ and \emph{some} chart $(V,y)$ in the atlas $\mathcal A_N$ with $\phi(p)\in V$ the map
\[
  y\circ\phi\circ x^{-1}\colon\ \R^{\dim M}\supseteq x(U)\longrightarrow y(V)\subseteq\R^{\dim N}
\]
is $C^k$ as a map from $\R^{\dim M}$ to $\R^{\dim N}$ \ts{7}{1}{25:59}\ts{7}{1}{26:24}\ts{7}{1}{26:45}\ts{7}{1}{27:10}\ts{7}{1}{27:32}\ts{7}{1}{27:52}\ts{7}{1}{28:13}\ts{7}{1}{28:39}\ts{7}{1}{29:00}.
\end{definition}
(One can only define differentiability of a map if a differentiable atlas has been chosen: for topological manifolds or spaces there is no such notion. The definition uses the chosen atlas in an essential way; which one does not matter, one chooses one \ts{7}{1}{26:24}\ts{7}{1}{26:45}.)
Why is $y\circ\phi\circ x^{-1}$ a map $\R^{\dim M}\to\R^{\dim N}$? Look at $M$, but only at the region $U$ that contains the point of interest, and at $N$, but only at the region $V$ containing the image of that point under $\phi$; the chart maps bring one down to $x(U)\subseteq\R^{\dim M}$
and $y(V)\subseteq\R^{\dim N}$ ($\dim M$ and $\dim N$ may differ, they are different manifolds). One wants to speak of a map between differentiable manifolds, but because an atlas is available one does so \emph{only} through the available atlases: pick charts that cover the point of interest and its
image, and consider the map in its \emph{chart representation}: mimic $\phi$ by the path $x^{-1}$, then $\phi$, then $y$. This is a map from $\R^{\dim M}$ to $\R^{\dim N}$, and for these one knows what differentiability means; so the differentiability of $\phi$ is \emph{defined} by that of its chart representation
\ts{7}{1}{29:00}\ts{7}{1}{29:21}\ts{7}{1}{29:42}\ts{7}{1}{30:02}\ts{7}{1}{30:22}\ts{7}{1}{30:43}\ts{7}{1}{31:04}.

The definition says ``for some chart''; one must check that the notion does not depend on the charts chosen: well-definedness \ts{7}{1}{31:04}\ts{7}{1}{31:25}\ts{7}{1}{31:46}.
\begin{proposition}[Well-definedness]\label{prop:df-welldef}
The notion of differentiability at $p$ of \cref{def:df-map} does not depend on the chosen charts: if $y\circ\phi\circ x^{-1}$ is $C^k$ for one choice of charts $(U,x)\in\mathcal A_M$, $(V,y)\in\mathcal A_N$, then $\tilde y\circ\phi\circ\tilde x^{-1}$ is $C^k$ for every other choice
$(\tilde U,\tilde x)\in\mathcal A_M$, $(\tilde V,\tilde y)\in\mathcal A_N$ with $p\in\tilde U$ and $\phi(p)\in\tilde V$ \ts{7}{1}{31:46}\ts{7}{1}{32:11}.
\end{proposition}
\begin{proof}
Take $U,V$ the subsets of $M,N$ containing $p$ and $\phi(p)$, with charts $x\colon U\to x(U)$ and $y\colon V\to y(V)$ from the respective atlases (very important: from the $C^k$ atlases), and the map $y\circ\phi\circ x^{-1}$ that is $C^k$ as a map $\R^{\dim M}\to\R^{\dim N}$
\ts{7}{1}{32:33}\ts{7}{1}{32:53}\ts{7}{1}{33:15}\ts{7}{1}{33:35}\ts{7}{1}{33:56}. Now let different charts $\tilde x$, $\tilde y$ be chosen, being extra picky: they have their own chart domains $\tilde U\ni p$, $\tilde V\ni\phi(p)$; replacing $U$ by $U\cap\tilde U$ and $V$ by $V\cap\tilde V$ (``otherwise the picture gets
messy'', without loss of generality) one may assume the domains are the same. One considers the map $\tilde y\circ\phi\circ\tilde x^{-1}$; a priori nothing guarantees that it is also $C^k$ \ts{7}{1}{33:56}\ts{7}{1}{34:16}\ts{7}{1}{34:36}\ts{7}{1}{34:57}\ts{7}{1}{35:17}\ts{7}{1}{35:38}.
But it can be written as a composition
\[
  \tilde y\circ\phi\circ\tilde x^{-1}=\bigl(\tilde y\circ y^{-1}\bigr)\circ\bigl(y\circ\phi\circ x^{-1}\bigr)\circ\bigl(x\circ\tilde x^{-1}\bigr),
\]
reading off the diagram \ts{7}{1}{35:38}\ts{7}{1}{35:58}. Here $\tilde y\circ y^{-1}$ is a chart transition map between two charts of the same $C^k$ atlas $\mathcal A_N$, hence $C^k$ as a map $\R^{\dim N}\to\R^{\dim N}$ (twice the same dimension, because it is a transition map: for $M$ and $N$ the dimensions could differ), and
$x\circ\tilde x^{-1}$ is a chart transition map between two charts of $\mathcal A_M$, hence $C^k$ as a map $\R^{\dim M}\to\R^{\dim M}$ \ts{7}{1}{35:58}\ts{7}{1}{36:18}\ts{7}{1}{36:39}\ts{7}{1}{37:00}\ts{7}{1}{37:20}.
The middle factor is $C^k$ by assumption, and a composition of $C^k$ maps (over $\R^n$, $\R^m$ and so on) is again $C^k$; hence once it is checked in one chart it is true in any other chart \ts{7}{1}{37:20}\ts{7}{1}{37:41}\ts{7}{1}{38:02}.
\end{proof}
Now one sees why it was so important to restrict the charts to those of a $C^k$ atlas: otherwise the compatibility condition would not be there; if the transition maps were only $C^0$, as they are for any topological manifold, one could check $C^k$ in one chart, but $C^k$ composed with $C^0$ maps
possibly destroys the $C^k$ property and leaves a merely continuous map. The chart transition maps must be at least of the same differentiability class: hence \emph{on a $C^k$ manifold one can define $C^k$ differentiability for maps between $C^k$ manifolds, but no higher}, because it might be
destroyed \ts{7}{1}{38:02}\ts{7}{1}{38:23}\ts{7}{1}{38:44}\ts{7}{1}{39:06}.

\begin{intuition}[The manifold philosophy]\label{int:df-philosophy}
You want to define something for an object on the manifold, and you define it by looking at its representative in a chart, or in several charts (here one goes from one manifold to the other, so two charts at least); then you show that it is independent of the choice of chart, and for
well-definedness to work you must suitably restrict the transition functions. For a complex manifold, one could check complex differentiability in one representation, but if the transition maps did not satisfy the higher-dimensional Cauchy--Riemann equations on top of being continuous anyway, one could not conclude that what was
checked in one chart is true in any other, hence could not attribute it as a property to the abstract object. For topological manifolds the chart picture was redundant, since continuity can be defined in its own right between topological spaces; once one goes beyond continuity, the chart picture becomes essential: one
steals from the knowledge of differentiability on $\R^n\to\R^m$ and lifts the notion to the manifold level \ts{7}{1}{39:06}\ts{7}{1}{39:26}\ts{7}{1}{39:47}\ts{7}{1}{40:08}\ts{7}{1}{40:30}\ts{7}{1}{40:54}.
\end{intuition}

\begin{definition}[New concepts: diffeomorphism, diffeomorphic manifolds]\label{def:df-diffeo}
A map $\phi\colon M\to N$ between smooth manifolds is called a \textbf{diffeomorphism} if it is bijective as a map between sets and both $\phi$ and its inverse $\phi^{-1}\colon N\to M$ are $C^\infty$ (one could say $C^k$, but let us say $C^\infty$)
\ts{7}{1}{40:54}\ts{7}{1}{41:16}\ts{7}{1}{41:38}. Two smooth manifolds $(M,\mathcal O_M,\mathcal A_M)$ and $(N,\mathcal O_N,\mathcal A_N)$ are called \textbf{diffeomorphic}, written $M\simeq_{\mathrm{diff}}N$, if there exists a diffeomorphism between them
\ts{7}{1}{41:38}\ts{7}{1}{42:00}\ts{7}{1}{42:22}.
\end{definition}
The diffeomorphisms are the isomorphisms, the structure preserving maps between differentiable (smooth) manifolds: ``always the same pattern'' (\cref{int:cl-theme}). Strictly, one should write that the triple $(M,\mathcal O_M,\mathcal A_M)$ is diffeomorphic to the triple $(N,\mathcal O_N,\mathcal A_N)$, as that is the full name of a smooth manifold:
a set, a topological structure and a pick of a smooth atlas. In future the notation is lightened and only the sets $M$ and $N$ are written, but ``if I say these $M$ and $N$ be smooth manifolds you know that secretly they carry with them a choice of topology and a choice of atlas'', and another manifold carries its own possibly
different choices; this must not be suppressed in the mind \ts{7}{1}{42:22}\ts{7}{1}{42:42}\ts{7}{1}{43:03}\ts{7}{1}{43:23}. Being diffeomorphic is an equivalence relation (every manifold is diffeomorphic to itself, as the identity is a diffeomorphism) \ts{7}{1}{43:23}\ts{7}{1}{43:44}, and it is custom not to distinguish
diffeomorphic manifolds: they are called the same, seen through the glasses of differentiable structures \ts{7}{1}{43:44}\ts{7}{1}{44:04}\ts{7}{1}{44:25}.

\begin{exercise}[The two atlases of the real line are diffeomorphic]\label{ex:df-line-diffeo}
\emph{Not said in class.} Show that $(\R,\mathcal A_1^{\max})$ and $(\R,\mathcal A_2^{\max})$ of \cref{ex:df-line} are diffeomorphic: take $\phi\colon(\R,\mathcal A_2^{\max})\to(\R,\mathcal A_1^{\max})$, $\phi(p)=p^{1/3}$, and compute its chart representation with the charts $x=(\cdot)^{1/3}$ and $\mathrm{id}$,
and likewise for $\phi^{-1}$. This anticipates \cref{sec:df-classification}.
\end{exercise}

% ==================================================================
\section{Levels of structure}\label{sec:df-levels}

Remark: it is custom to consider diffeomorphic smooth manifolds to be the same, and this is all a question of the \emph{level of structure} \ts{7}{1}{44:04}\ts{7}{1}{44:25}.
Two sets $M$ and $N$ may be the same seen through the glasses of set theory (there is a bijection); equip each with a topology, and unless there exists a homeomorphism between the extended structures they are no longer the same seen through the glasses of topology. One can push this on: even if they are
homeomorphic as topological spaces, one can equip them with different differentiable structures, and then they may be different as differentiable manifolds. This was seen with the real line with the standard topology: trivially homeomorphic to itself, but with one smooth atlas on one copy and another on the other
copy they are not the same as differentiable manifolds (\cref{ex:df-line}; but see below, they are in fact diffeomorphic) \ts{7}{1}{44:25}\ts{7}{1}{44:46}\ts{7}{1}{45:07}\ts{7}{1}{45:27}\ts{7}{1}{45:47}.
\begin{coursenote}[A slip in the example]
The lecture says that with two different atlases the two copies of the real line are ``no longer the same as differentiable manifolds''; what is true is that the atlases are \emph{incompatible}, i.e.\ different structures on the same set; whether the resulting manifolds are diffeomorphic is
a different question, answered in the next section (they are diffeomorphic, \cref{ex:df-line-diffeo}). Not said in class.
\end{coursenote}
``As soon as I start painting people's hair red and green they are no longer the same''; it always depends on context what ``these manifolds are the same'' means \ts{7}{1}{45:47}\ts{7}{1}{46:09}.

\paragraph{A student asks about group structures.} A student: one had said that some sets are equal or isomorphic as a group structure; is it a linear stacking? Schuller goes back to page one (the picture of the first lecture): a set can be equipped with a topology (modify the picture), or instead of a topology
with a map $\diamond\colon M\times M\to M$ (``some blob or diamond''); if $\diamond$ is associative, has a neutral element in $M$ and inverse elements, it is a \textbf{group}. A set can be equipped with a topology \emph{and} such a structure: one obtains a \textbf{topological group}, also called a Lie group
(see the coursenote). On top of the topology one can say it is a manifold and then put a differentiable structure: a differentiable manifold, which is the path followed so far \ts{7}{1}{46:09}\ts{7}{1}{46:29}\ts{7}{1}{46:49}\ts{7}{1}{47:12}\ts{7}{1}{47:32}\ts{7}{1}{47:54}.
\begin{coursenote}[Topological versus Lie groups]
The lecture says that a set with a topology and a group structure is ``a topological group, also called a Lie group'', which is a slip: a Lie group is a group that is also a \emph{smooth manifold}, with smooth group operations (\cref{ch:l13}). What is true, and is the content of the following paragraph, is that a
locally Euclidean topological group is already a Lie group. Not said in class.
\end{coursenote}
\emph{Hilbert's problem.} Once there, with a set with a topology and a group structure, one has a so-called Lie group; one can continue. Around 1900 Hilbert gave a big speech and said that there are 22 problems to be solved in the twentieth century; one of them was whether every topological group (called Lie group)
is already a differentiable Lie group, and \emph{the answer is yes}. Hence the analysis of Lie groups is best done by already employing a differentiable structure, as in the course: the differentiable structure gives rise to the so-called Lie algebras, and by Lie algebras one can study Lie groups very efficiently, at least some
aspects \ts{7}{1}{47:54}\ts{7}{1}{48:14}\ts{7}{1}{48:35}\ts{7}{1}{48:55}.
\begin{history}[Hilbert's fifth problem]
\emph{Checked.} At the Paris congress of 1900 Hilbert listed 23 problems (published version; the lecture says 22, as heard); the fifth asks whether a continuous group is automatically a differentiable one. It was solved in the affirmative, for locally Euclidean groups, by Gleason and by
Montgomery and Zippin in 1952.\par
\emph{References.} D.~Hilbert, ``Mathematische Probleme'', \emph{G\"ottinger Nachrichten} (1900); A.~Gleason, \emph{Ann.\ Math.}\ 56 (1952); D.~Montgomery, L.~Zippin, \emph{Ann.\ Math.}\ 56 (1952).
\end{history}
So what this picture should convey is that one can have one structure and not the other, but not a differentiable structure without a topological space, and not a topological space without a set; but one can combine the structures in various ways. Another thing: an additional structure like a scalar multiplication
$S\colon K\times M\to M$ (``an S multiplication''), and if the right eight axioms hold, it is a \textbf{vector space}. A vector space can be equipped with a topology and so on; ``you can combine the structures in almost arbitrary ways, but not completely arbitrary ways, because some structures require the other structure to underlie''.
It is always very important whether one talks about a vector space or about a topological vector space \ts{7}{1}{48:55}\ts{7}{1}{49:15}\ts{7}{1}{49:36}\ts{7}{1}{49:57}\ts{7}{1}{50:18}.
\begin{secondary}[Not a linear stacking of structures]
A second student question: for the fundamental groups one says the groups are isomorphic as groups, while the torus and the cylinder are not isomorphic topologically; could one then say that they cannot be isomorphic in the group sense? No: the picture talks about \emph{the same set} being equipped with more
and more structure of this and that type, in that combination. If one looks at a topological space, one can construct from it the fundamental group (\cref{sec:iv-pi1}); but the group has its own underlying set (for the torus, $\mathbb Z\times\mathbb Z$, a different set) which carries a group structure, and one can build another
``building like this'' on that set. It does not mean that in one mathematics problem things always come in this one linear stacking; the picture, which ``looks pretty, maybe more confusing than enlightening'', was not meant to convey that \ts{7}{1}{50:18}\ts{7}{1}{50:38}\ts{7}{1}{50:59}\ts{7}{1}{51:19}\ts{7}{1}{51:40}\ts{7}{1}{52:00}.
\end{secondary}

% ==================================================================
\section{How many smooth structures? The classification}\label{sec:df-classification}

Now that diffeomorphic manifolds are essentially the same, return to the question \ts{7}{1}{52:00}\ts{7}{1}{52:21}: how many differentiable structures are there up to isomorphism? For instance, is the real line with the standard topology and the maximal extension of the atlas $\mathcal A_1$ of \cref{ex:df-line} at
least diffeomorphic to the real line with the standard topology and (the maximal extension of) $\mathcal A_2$? They were clearly not the same atlas, $\mathcal A_1^{\max}\ne\mathcal A_2^{\max}$: one chart in each, but once one has a chart one adds all compatible charts and gets the maximisation; and one started worrying which
structure one would use as \emph{the} differentiable structure \ts{7}{1}{52:21}\ts{7}{1}{52:43}\ts{7}{1}{53:06}\ts{7}{1}{53:27}. But one should not be so precise and look only at the atlas; one should ask whether the resulting differentiable manifold differs \emph{essentially}: a non-essential difference
is one for which they are diffeomorphic. Then from the point of view of differentiability they are the same, since every function differentiable on one can be pushed over to the other and is differentiable there \ts{7}{1}{53:27}\ts{7}{1}{53:47}\ts{7}{1}{54:08}.
There are rather deep results, not proved here, some of them big achievements. The question:
\begin{center}\emph{How many different differentiable structures can one put on a given topological manifold, up to diffeomorphism?}\end{center}
(I hand you a topological manifold, you have to decide which differentiable structure to put on it) \ts{7}{1}{54:08}\ts{7}{1}{54:28}\ts{7}{1}{54:49}\ts{7}{1}{55:10}. The surprising answer depends on the dimension \ts{7}{1}{55:10}\ts{7}{1}{55:31}.

\begin{theorem}[Dimensions $1,2,3$ (Rad\'o--Moise)]\label{thm:df-low}
For $\dim M=1,2$ or $3$ every topological manifold carries, up to diffeomorphism, a \emph{unique} smooth manifold structure \ts{7}{1}{55:31}\ts{7}{1}{55:52}\ts{7}{1}{56:12}\ts{7}{1}{56:34}. (In the transcript: ``Radon Mois theorems''.)
\end{theorem}
Although one could choose different smooth charts already for the real line, and the atlases are not compatible, the resulting manifolds are diffeomorphic; so everything one does with differentiability is unique. One now has a notion of differentiability on the torus, say: ``see what an achievement this is''. At high school one learns
differential calculus of one variable over the real line, at university in an engineering course over $\R^n$ or $\R^3$, and now one can do it on arbitrary manifolds, with holes and whatnot; up to dimension three everything is unique \ts{7}{1}{56:34}\ts{7}{1}{56:58}\ts{7}{1}{57:19}\ts{7}{1}{57:39}.

\begin{theorem}[Dimension greater than four: finitely many]\label{thm:df-high}
For $\dim M>4$ there are only finitely many different smooth manifolds (up to diffeomorphism) that can be made of a given topological manifold \ts{7}{1}{58:01}\ts{7}{1}{1:00:04}\ts{7}{1}{1:00:24}\ts{7}{1}{1:00:46}.
\end{theorem}
Dimension four is left out for now: ``dimension of $M$ greater than four, but \emph{not} four: decidedly not four'' \ts{7}{1}{57:39}\ts{7}{1}{58:01}. There is a toolbox called \emph{surgery theory}: take, say, a sphere and operate on it like a surgeon: cut out two pieces with a scalpel; a colleague gives
a cylinder, which is pushed in and sewn at both ends; what one gets is the torus. The idea (a big field): the sphere and the cylinder are well understood, and if one performs surgery in such a way that one controls the topological invariants one cares about (the fundamental group, later homology), one understands the torus;
that is a trivial example, the aim being to understand all manifolds of dimension higher than four by reducing them to elements and surgery procedures one understands. Done systematically, in the 1960s, this gave the theorem \ts{7}{1}{58:01}\ts{7}{1}{58:22}\ts{7}{1}{58:43}\ts{7}{1}{59:03}\ts{7}{1}{59:23}\ts{7}{1}{59:44}\ts{7}{1}{1:00:04}.
(In the lecture it is first said ``dim $M\ge4$'' and then ``greater than four''; see the coursenote.)
\begin{coursenote}[Dimension four and the wording]
The lecture first says ``greater or equal to four'' and then ``greater than four'' and ``not four''; the correct statement is $\dim M\ge5$, with $\dim M=4$ treated separately, as done next. Not said in class.
\end{coursenote}
That is not as neat as dimensions one, two, three, but the structures are finitely many, so there is a chance, in principle, to enumerate them and write them all down. Take say the four-sphere (``let us find all the differentiable structures''; for a given manifold it may be very difficult in practice, but the theorem
says that they are finitely many) \ts{7}{1}{1:00:46}\ts{7}{1}{1:01:07}.
\begin{application}[If space-time were higher dimensional]\label{app:df-physics}
Some physicists believe that space-time is, say, 10-dimensional (``not proposing or condemning that''); then one must worry which smooth structure to equip the manifold with. Physics can make experiments: if different structures lead to different observations, one can at least enumerate them and make finitely many
observations to distinguish among them, in principle. So higher dimension is in a sense good news: one could determine by experiment which of the structures nature has chosen, if any (if space-time is a differentiable manifold in the first place). ``This has immediate impact on physics, if you take it seriously, and we want to do that''
\ts{7}{1}{1:01:07}\ts{7}{1}{1:01:29}\ts{7}{1}{1:01:49}\ts{7}{1}{1:02:10}\ts{7}{1}{1:02:30}.
\end{application}

\begin{theorem}[Dimension four]\label{thm:df-four}
A $4$-dimensional smooth manifold, like $\R^4$: up to diffeomorphism there are \emph{uncountably many} different smooth manifolds that can be made of the same $4$-dimensional topological manifold; in particular for $\R^4$ with the standard topology, a topological manifold of dimension four.
This holds at least in the non-compact case (c.1)
\ts{7}{1}{1:02:30}\ts{7}{1}{1:02:51}\ts{7}{1}{1:03:14}\ts{7}{1}{1:03:37}\ts{7}{1}{1:03:58}.
\end{theorem}
The \emph{compact case} (c.2) has only partial results: if the second \textbf{Betti number} $b_2$, a topological invariant, is greater than $18$ [?: transcript ``greater than 18''], then there are only (countably) infinitely many different smooth manifolds; not much better, so the compact case is a little
better than the non-compact one but still infinite \ts{7}{1}{1:03:58}\ts{7}{1}{1:04:21}\ts{7}{1}{1:04:45}\ts{7}{1}{1:05:05}\ts{7}{1}{1:05:25}\ts{7}{1}{1:05:46}.
\begin{coursenote}[Which compact case]
The transcript has ``countably infinitely many'' for $b_2>18$; the lecture's statement is only a partial result and I could not check the exact hypothesis; marked [?]. Not said in class.
\end{coursenote}
\begin{definition}[New concepts: Betti numbers (as explained informally)]\label{def:df-betti}
The Betti numbers are defined in terms of homology groups (to come later); intuitively the zeroth Betti number $b_0$ is the number of connected pieces of the space, $b_1$ the number of circular (one-dimensional) holes, $b_2$ the number of two-dimensional holes, and so on; ``the pretzel but with 19 holes or so'' has a large number of holes
\ts{7}{1}{1:04:21}\ts{7}{1}{1:04:45}\ts{7}{1}{1:05:05}\ts{7}{1}{1:05:25}. (Name: after E.~Betti, who introduced the numbers in 1871; checked.)
\end{definition}
\begin{application}[Special relativity on $\R^4$]\label{app:df-r4}
When one says space-time is $\R^4$ with some additional structure (the light cone or the metric, in special relativity) and one finds tangent vectors to world lines, the question is which smooth structure to put on $\R^4$ as a topological manifold in order to do special relativity: there are uncountably many to choose
from. Can one do experiments to distinguish between them? One would need, with luck, uncountably many experiments to find the one: ``it's not so good, right? But it's a result''. ``Unfortunately I have nothing better to say about this, it's just an observation, and this is what the analysis yields''. From a philosophical point of view it is
serious (``I do not know whether from a scientific practical point of view''): in principle our theories could fail on $\R^4$ because the wrong differentiable structure was chosen, and might be repaired by choosing another; ``I don't think this is a point where you could start doing research, but maybe you can. As a mathematical result it stays
and should be noted'' \ts{7}{1}{1:05:46}\ts{7}{1}{1:06:07}\ts{7}{1}{1:06:29}\ts{7}{1}{1:06:53}\ts{7}{1}{1:07:14}\ts{7}{1}{1:07:34}.
\end{application}
\begin{history}[Exotic spheres and exotic $\R^4$]
\emph{Not said in class.} \emph{Checked.} J.~Milnor found in 1956 smooth structures on the 7-sphere that are not diffeomorphic to the standard one (exotic spheres); for $\R^4$, exotic structures follow from the work of Donaldson and Freedman (1982), and C.~Taubes (1987) showed that there are uncountably many
non-diffeomorphic ones.\par
\emph{References.} J.~Milnor, \emph{Ann.\ Math.}\ 64 (1956); S.~Donaldson, \emph{J.\ Diff.\ Geom.}\ 18 (1983); M.~Freedman, \emph{J.\ Diff.\ Geom.}\ 17 (1982); C.~Taubes, \emph{J.\ Diff.\ Geom.}\ 25 (1987).
\end{history}

% ==================================================================
\section{Outlook: tangent spaces without embedding}\label{sec:df-outlook}

So a differentiable structure is defined, one can classify; but what is really new about differentiable manifolds is that, roughly speaking, \emph{they have tangent spaces} \ts{7}{1}{1:07:14}\ts{7}{1}{1:07:34}\ts{7}{1}{1:07:54}. The key feature: at each point of the manifold there exists a tangent space (the words are in quotes: one has an idea of what a tangent space looks like,
but that idea will not be used; it will be done in a different way) \ts{7}{1}{1:07:54}\ts{7}{1}{1:08:16}\ts{7}{1}{1:08:36}.

\begin{intuition}[Why not a plane in the ambient space]\label{int:df-embedding}
Look at the sphere $S^2$ as a topological manifold; one always thinks of it as embedded in three-dimensional space, but as a topological manifold it does not care: the set and the topology make no reference to what lies outside the surface; the same holds when one adds a smooth atlas. Now explain tangent spaces in the
following way: take one point of the sphere, rotate it ``for dramatic purposes'' to the horizon, and imagine the tangent space as all points or vectors in the tangent plane. This makes reference to the space in which the sphere is embedded, which has nothing to do with the smooth structure, independent of the embedding. So tangent spaces
will \emph{not} be defined as planes in an exterior space in which the manifold is embedded. (Due to various theorems, one could do it that way without loss of generality \ts{7}{1}{1:08:36}\ts{7}{1}{1:08:56}\ts{7}{1}{1:09:17}\ts{7}{1}{1:09:37}\ts{7}{1}{1:09:58}\ts{7}{1}{1:10:19}\ts{7}{1}{1:10:39}.)
Nevertheless it is not done, because if the manifold is not two- but four-dimensional and represents our universe, space-time, it would be funny to think of the universe as embedded, physically speaking, in a bigger space in which the tangent spaces lie, since that bigger space would obviously be what should be called the universe,
the thing that comprehends everything. Even from the point of view of theory building in physics it is nice to construct everything \emph{intrinsically}, without recourse to structures outside the object: one must define tangent spaces solely with the structure one is handed, notwithstanding intuition (``you may think of tangent spaces
this way, but it is better not to define them this way'') \ts{7}{1}{1:10:39}\ts{7}{1}{1:11:01}\ts{7}{1}{1:11:22}\ts{7}{1}{1:11:42}\ts{7}{1}{1:12:03}.
\end{intuition}

In the next subsections one fixes a point of the manifold and constructs the tangent space at this point, its dual, and the tensor spaces over this point; later all tangent spaces are looked at together: tangent vectors at one point form a vector space, but taking at every point a different vector gives a
\emph{vector field} on the manifold \ts{7}{1}{1:12:03}\ts{7}{1}{1:12:23}\ts{7}{1}{1:12:44}. Vector spaces are linear algebra: a few linear algebra notions will be reviewed quickly, in particular the construction of tensors, because ``there are some standard confusions that we will now try to eliminate once and for all''
\ts{7}{1}{1:12:44}\ts{7}{1}{1:13:04}.
However, vector fields still have a linear structure: two vector fields can be added, but also scaled, and at different points of the manifold one can scale differently, with a \emph{function}. It turns out that the strict notion of a vector space is too restrictive to deal in a clever way with vector fields, and one goes to another
structure called a \textbf{module}. A show of hands (``who of you saw modules in linear algebra? Oh, fantastic'' --- some had): it will be new for most, ``it's my responsibility, all your future misconceptions will be my responsibility''; normally one meets modules in linear algebra two, a very useful structure, which also in physics
``we use all the time, we just choose to ignore that they are not really vector spaces anymore'' \ts{7}{1}{1:13:04}\ts{7}{1}{1:13:25}\ts{7}{1}{1:13:46}\ts{7}{1}{1:14:07}. This is the plan of the next two lectures: tensor space theory over a field (\cref{ch:l08}) and then over a ring.

% ==================================================================
\section*{Key concepts}
\begin{keyconcepts}
\textbf{Retrieval questions.}
\begin{enumerate}
  \item What does it mean for an atlas to be $C^k$, $C^\infty$, $C^\omega$, complex? Which condition do the transition maps satisfy in each case?
  \item State the Whitney theorem on maximal $C^k$ atlases. Why does it allow one to speak of ``smooth'' only? Why must $k\ge1$?
  \item Define a $C^k$ manifold. Why can a topological manifold carry incompatible atlases? Give the example of $\R$ with the atlases of the identity and of $a\mapsto a^{1/3}$.
  \item Define the differentiability of a map between manifolds at a point. Prove that it does not depend on the charts chosen: where does the compatibility of the atlases enter?
  \item Why can one define $C^k$ differentiability on a $C^k$ manifold but no higher?
  \item Define diffeomorphism and diffeomorphic manifolds. Why is the real line with the two atlases above ``the same'' manifold?
  \item How many smooth structures does a topological manifold carry in dimension $1,2,3$, above $4$, and equal to $4$?
  \item Why are tangent spaces not defined as planes in an ambient space?
\end{enumerate}
\textbf{Concepts.}
\begin{itemize}
  \item $C^\star$ compatibility, $C^k$, smooth, analytic and complex atlases (\cref{def:df-star,def:df-ck,def:df-cr}); the theorem on smooth sub-atlases (\cref{thm:df-whitney}).
  \item $C^k$ manifold (\cref{def:df-manifold}); incompatible atlases (\cref{rem:df-incompat,ex:df-line}).
  \item Differentiable maps between manifolds and their well-definedness (\cref{def:df-map,prop:df-welldef}); the manifold philosophy (\cref{int:df-philosophy}).
  \item Diffeomorphism (\cref{def:df-diffeo,ex:df-line-diffeo}); levels of structure (\cref{sec:df-levels}).
  \item Classification: dimensions $1,2,3$, above $4$, $4$ (\cref{thm:df-low,thm:df-high,thm:df-four}); intrinsic tangent spaces (\cref{int:df-embedding}).
\end{itemize}
\end{keyconcepts}

\flushside
\end{paracol}

% !TEX root = ../main.tex
\chapter{Tensor space theory I: over a field}\label{ch:l08}
\begin{paracol}{2}

\begin{sources}
\begin{tabularx}{\linewidth}{@{}l l L@{}}
\toprule
\textbf{Video} & \textbf{Length} & \textbf{Content}\\
\midrule
\seg{8}{1}{Lecture 8: Tensor space theory I: over a field} & 2:22:59 & Review of fields, rings and vector spaces; linear maps, $\mathrm{Hom}$, $\mathrm{End}$, $\mathrm{Aut}$, the dual space; why modules over a non-commutative ring fail; $(p,q)$ tensors, the tensor space, pointwise operations,
the tensor product; $T^0_1V=V^*$ and $T^1_1V\cong\mathrm{End}(V^*)$; Hamel bases and dimension; $V^{**}\cong V$ in finite dimension; basis-free construction; components, dual basis and reconstruction of a tensor; change of basis; the
summation convention and the up/down convention; matrices as ``witchcraft''; transformation of components; the determinant via top forms and volume forms; the determinant of a bilinear form is not well defined; outlook to
principal bundles and tensor densities.\\
\bottomrule
\end{tabularx}

\smallskip
\textbf{Transcript} \file{materials/transcripts/C08-V01.txt} (YouTube automatic captions).
\end{sources}

\begin{reading}[Where this lecture sits in the other notes]
\emph{Gravity and Light}, Lecture~3 (Multilinear algebra); \emph{Theoretische Mechanik}, Lecture~5 (Multilinear algebra). Both cover the tensor space, the dual basis and the change of basis; the extended
treatment of the determinant through top forms and the discussion of matrices as a bad notation are, as far as the titles say, specific to this course; the module generalisation announced here is the subject of \cref{ch:l11}.
(From the lecture titles and the plan of the course; not checked in detail.)
\end{reading}

% ==================================================================
\section{Fields and rings}\label{sec:te-field}

A quick review of vector spaces \ts{8}{1}{0:04}. A vector space is built on a structure called a \emph{field}: an algebraic field, what a mathematician calls a field, not what a physicist calls a field (a function on space-time) \ts{8}{1}{0:25}.

\begin{definition}[New concepts: field]\label{def:te-field}
A \textbf{field} is a triple $(K,+,\cdot)$ where $K$ is a set and
\[
  +\colon K\times K\to K,\qquad \cdot\colon K\times K\to K
\]
are two maps (addition and multiplication) satisfying
\begin{itemize}
  \item for $+$: it is commutative, associative, there is a neutral element (called $0$, whatever element of the set it is), and every element has an inverse with respect to $+$: adding it gives the neutral element
        (the neutral element is to be written down before the inverse, since the inverse refers to it);
  \item for $\cdot$: the same four laws, but only on the set $K\setminus\{0\}$, where the neutral element of addition has been removed: commutativity, associativity, a neutral element (named $1$) and inverses
\end{itemize}
\ts{8}{1}{0:25}\ts{8}{1}{0:48}\ts{8}{1}{1:09}\ts{8}{1}{1:32}\ts{8}{1}{1:53}\ts{8}{1}{2:14}\ts{8}{1}{2:34}.
\end{definition}
(The set of the four laws is written ``CANI'' in the lecture: Commutativity, Associativity, Neutral element, Inverse \ts{8}{1}{1:32}\ts{8}{1}{1:53}.) The multiplicative laws are required only for $K\setminus\{0\}$ because $0$ has no multiplicative inverse. (A student:
could one not just ignore the zero? ``How do you ignore it, you close your eyes when it comes along?'' --- taking it out \emph{is} the meaning of ignoring it \ts{8}{1}{2:34}\ts{8}{1}{2:55}\ts{8}{1}{3:16}.) These are the tools learned in high school \ts{8}{1}{3:16}.
\begin{coursenote}[The distributive law]
The lecture does not mention the law linking $+$ and $\cdot$ ($a\cdot(b+c)=a\cdot b+a\cdot c$), which is also part of the definition of a field (and of a ring); at the board it is left implicit. Not said in class.
\end{coursenote}
\begin{history}[Field and ring]
\emph{Not said in class.} \emph{Checked.} The word \emph{Körper} (field) for a set of numbers closed under the four operations is due to R.~Dedekind (1871); the word \emph{Ring} for the structure without inverses was introduced by D.~Hilbert (\emph{Zahlbericht}, 1897).\par
\emph{References.} R.~Dedekind, supplement to Dirichlet's \emph{Vorlesungen über Zahlentheorie} (1871); D.~Hilbert, ``Die Theorie der algebraischen Zahlkörper'', \emph{Jahresber.\ DMV} 4 (1897).
\end{history}

A weaker notion, ``wildly important'' later when modules are considered, is the following \ts{8}{1}{3:16}\ts{8}{1}{3:39}.
\begin{definition}[New concepts: ring]\label{def:te-ring}
A \textbf{ring} is a set $R$ with two operations $+$ and $\cdot$ of the same shape, but satisfying only: $(R,+)$ is a commutative (abelian) group (CANI for $+$); for $\cdot$ the commutativity is gone, the associativity and a neutral element need to stay (there are rings without neutral
element; one says a \textbf{unital} ring, a ring ``with unit''), and, most importantly, the multiplicative inverses are gone \ts{8}{1}{3:39}\ts{8}{1}{4:01}\ts{8}{1}{4:22}\ts{8}{1}{4:43}.
\end{definition}
Every field is a ring, but there are rings that are not fields \ts{8}{1}{4:43}\ts{8}{1}{5:03}.
\begin{example}[Rings that are not fields]\label{ex:te-rings}
\begin{itemize}
  \item The integers $\mathbb Z$ with $+$ and $\cdot$ (as in first grade): a \emph{commutative} ring, a nice one; but the multiplicative inverse of $2$ is $\tfrac12$, which is not in $\mathbb Z$. One learned to calculate with rings before fields \ts{8}{1}{5:03}\ts{8}{1}{5:24}\ts{8}{1}{5:45}.
  \item $\R$, $\mathbb Q$ and so on, with $+,\cdot$, are fields, ``very special rings'' \ts{8}{1}{5:45}\ts{8}{1}{6:06}.
  \item A \emph{non-commutative} ring: the $M\times M$ matrices over $\R$, which can be added and multiplied by matrix multiplication; the multiplication is not commutative and not every matrix has an inverse (if the determinant is zero)
        \ts{8}{1}{6:06}\ts{8}{1}{6:26}\ts{8}{1}{6:47}. A very natural structure; rings are seen everywhere; function rings later \ts{8}{1}{6:47}\ts{8}{1}{7:08}.
\end{itemize}
\end{example}

% ==================================================================
\section{Vector spaces and linear maps}\label{sec:te-vs}

Now back to fields \ts{8}{1}{7:08}.
\begin{definition}[New concepts: $K$-vector space]\label{def:te-vs}
Let $K$ be a field (important). A \textbf{$K$-vector space} is a triple $(V,\oplus,\odot)$ (the operations are circled because they are \emph{not} the operations of $K$: the $+$ of $K$ is not the $\oplus$ of $V$, and the $\cdot$ of $K$ is not the $\odot$: totally different structure) with
\[
  \oplus\colon V\times V\to V,\qquad \odot\colon K\times V\to V,
\]
the second taking an element of the underlying field and a vector and producing a vector (not two vectors producing a vector). The standard example is a real vector space, over $\R$: vectors can be added and there is a scalar (``$S$'') multiplication. The axioms:
\begin{itemize}
  \item for $\oplus$ the same CANI laws as for a field's addition, now applied to the vector space elements;
  \item for $\odot$ there is no CANI (it does not take two vectors); there are four other laws, linking the operations: an associative law, two distributive laws (one distributes over the $+$ of $K$, the other over the $\oplus$ of $V$), and the unitary law $1\odot v=v$
\end{itemize}
\ts{8}{1}{7:08}\ts{8}{1}{7:29}\ts{8}{1}{7:49}\ts{8}{1}{8:10}\ts{8}{1}{8:30}\ts{8}{1}{8:50}\ts{8}{1}{9:10}\ts{8}{1}{9:31}\ts{8}{1}{37:08}.
\end{definition}
(The axioms are not spelled out at the board: ``if you don't remember this precisely, go back to your linear algebra book and study this now'' \ts{8}{1}{9:31}\ts{8}{1}{9:52}. The list of eight axioms is mentioned again at \ts{8}{1}{37:08}: commutative, associative, neutral, inverse, an associative law, two distributive laws and the unitary law.)

A subset $U\subseteq V$ is a \textbf{vector subspace} if for all $u_1,u_2\in U$ the sum $u_1\oplus u_2$ is again in $U$ and for all $\lambda\in K$ also $\lambda\odot u_1\in U$: $U$ is \emph{closed} under the two operations, they do not lead out of $U$
\ts{8}{1}{9:52}\ts{8}{1}{10:13}\ts{8}{1}{10:34}.

\paragraph{Dropping the circles.} From now on the circles are dropped, $+$ between vectors and $v$ scaled by $\lambda$ as $\lambda v$; one must always know whether it is the plus between elements of the field or between vectors, the product of two elements of $K$ or the scalar multiplication of a $K$ element with a $V$ element:
the expression always tells which, though they are different operations. E.g.\ $\lambda\mu v$ with $\lambda,\mu\in K$ and $v\in V$: the product between $\lambda$ and $\mu$ is the ordinary one in $K$, the others are scalar multiplications \ts{8}{1}{10:34}\ts{8}{1}{10:55}\ts{8}{1}{11:16}\ts{8}{1}{11:38}\ts{8}{1}{12:02}.
The circles are dropped because $\diamond$-type symbols are needed in a moment for another operation \ts{8}{1}{12:02}\ts{8}{1}{12:22}.

\paragraph{Structure preserving maps.} As always one classifies by the structure preserving maps (\cref{int:cl-theme}); for vector spaces the homomorphisms are the \emph{linear maps} \ts{8}{1}{12:22}\ts{8}{1}{12:43}.
\begin{definition}[New concepts: linear map, vector space isomorphism]\label{def:te-linear}
Let $(V,\oplus,\odot)$ and $(W,\boxplus,\boxdot)$ be $K$-vector spaces (different spaces may carry different pluses and scalar multiplications, so they are distinguished here). A map $f\colon V\to W$ is \textbf{linear} if for all $v_1,v_2\in V$ and all $\lambda\in K$
\[
  f(v_1\oplus v_2)=f(v_1)\boxplus f(v_2),\qquad f(\lambda\odot v)=\lambda\boxdot f(v);
\]
one can add first in $V$ and then map, or map the individual vectors and add in the target: the same result (additivity and homogeneity). A bijective linear map is a \textbf{vector space isomorphism}; its inverse is then automatically linear \ts{8}{1}{12:43}\ts{8}{1}{13:04}\ts{8}{1}{13:25}\ts{8}{1}{13:46}\ts{8}{1}{14:07}\ts{8}{1}{14:48}\ts{8}{1}{15:09}.
\end{definition}
(The domain and the target could be the same vector space, operations included, in which case one need not distinguish: but this is the degenerate case, not the rule \ts{8}{1}{14:07}\ts{8}{1}{14:27}.) Two vector spaces $V,W$ (operations suppressed in the notation) are called \textbf{isomorphic as vector spaces} if there is an isomorphism $V\to W$;
``and so on and so on'', it is again an equivalence relation: one considers vector spaces the same if there is a bijective linear map between them, the same idea as always \ts{8}{1}{15:09}\ts{8}{1}{15:30}\ts{8}{1}{15:51}.

% ==================================================================
\section{The vector space of linear maps, End, Aut and the dual}\label{sec:te-hom}

A very important construction \ts{8}{1}{15:51}\ts{8}{1}{16:12}.
\begin{definition}[New concepts: $\mathrm{Hom}(V,W)$]\label{def:te-hom}
$\mathrm{Hom}(V,W)$ is the set of all linear maps from the vector space $V$ to the vector space $W$. In future a tilde on the arrow, $f\colon V\xrightarrow{\ \sim\ }W$, means that the map is linear (``everything will be linear from now on'', it saves chalk). It can again be made into a vector space, by
the operations
\[
  \diamond\!+\ \colon\ (f,g)\mapsto f\mathbin{\diamond\!+}g,\quad (f\mathbin{\diamond\!+}g)(v):=f(v)+g(v),\qquad (\lambda\mathbin{\diamond\!\cdot}f)(v):=\lambda\,f(v),
\]
where the $+$ on the right is the plus of $W$ (the results are added up; first written ``on $V$'', corrected at the board) and the product on the right is the scalar multiplication of $W$: $\mathrm{Hom}(V,W)$ inherits its operations from the target space
\ts{8}{1}{16:12}\ts{8}{1}{16:32}\ts{8}{1}{16:54}\ts{8}{1}{17:16}\ts{8}{1}{17:37}\ts{8}{1}{17:57}\ts{8}{1}{18:19}\ts{8}{1}{18:42}.
\end{definition}
One has to check that $f\mathbin{\diamond\!+}g$ is again a linear map (it is additive; and likewise the definition of $\diamond\cdot$ satisfies the linearity axiom), so that it lies in $\mathrm{Hom}(V,W)$ \ts{8}{1}{18:42}\ts{8}{1}{19:02}\ts{8}{1}{19:23}.

\begin{secondary}[Why a field: modules over a non-commutative ring]
If one had not a field but a ring (removing commutativity and inverses) one could still write down the vector space axioms; this is a wildly different structure which deserves its own name: a \textbf{module over a ring} (to be looked at later in the course, \cref{ch:l11}), nothing but a vector space structure whose underlying
``field'' is a ring. One can do everything above also for modules until one tries to show that $\lambda\diamond\cdot f$, defined by $(\lambda\diamond\cdot f)(v)=\lambda\,f(v)$, is again a linear map: for that one scales, pulls the scaling out, but one cannot pull it in front of the $\lambda$, because
$\lambda$ is a ring element and they no longer commute. So for a module over a non-commutative ring one cannot make $\mathrm{Hom}(V,W)$ again into a module; one can only do it for a module over a \emph{commutative} ring: a very small subtlety that plays a big role
\ts{8}{1}{19:23}\ts{8}{1}{19:44}\ts{8}{1}{20:05}\ts{8}{1}{20:25}\ts{8}{1}{20:46}\ts{8}{1}{21:07}\ts{8}{1}{21:28}.
\end{secondary}
Here everything is fine, and $\mathrm{Hom}(V,W)$ with $\diamond+$ and $\diamond\cdot$ is a vector space in its own right \ts{8}{1}{21:28}\ts{8}{1}{21:49}.

\begin{coursenote}[Vector spaces over different fields]
A student asks whether $V$ and $W$ could be vector spaces over different fields; Schuller first answers: a $K$-vector space $V$ and a $K'$-vector space $W$ can be made into a $K'$-vector space $\mathrm{Hom}(V,W)$, since the scalar multiplication of $\mathrm{Hom}$ is taken
from the target \ts{8}{1}{21:49}\ts{8}{1}{22:10}\ts{8}{1}{22:35}\ts{8}{1}{22:57}\ts{8}{1}{23:18}\ts{8}{1}{23:39}. The same student comes back at \ts{8}{1}{31:42}\ts{8}{1}{32:12}: ``what does linear mean, then?'' and Schuller agrees: the definition of linearity uses the \emph{same} $\lambda$ on both sides, so
\emph{linearity forces $K=K'$}; the first answer fails exactly there (``it was very clever, the question made me consider it'') \ts{8}{1}{32:12}\ts{8}{1}{32:32}\ts{8}{1}{32:53}. So the first answer is wrong as stated and the second one is the correct one.
\end{coursenote}

\paragraph{Terminology.} ``Terrible if you do not know it, trivial once you do'' \ts{8}{1}{23:39}\ts{8}{1}{24:00}.
\begin{definition}[New concepts: endomorphisms, automorphisms, dual space]\label{def:te-end}
\begin{itemize}
  \item $\mathrm{End}(V):=\mathrm{Hom}(V,V)$, the \textbf{endomorphisms} of $V$: the linear maps from $V$ to itself \ts{8}{1}{24:00}\ts{8}{1}{24:21}.
  \item $\mathrm{Aut}(V)$, the \textbf{automorphisms}: the invertible linear maps $V\to V$ (the isomorphisms of $V$ with itself, called isomorphisms before) \ts{8}{1}{24:21}\ts{8}{1}{24:42}.
  \item $V^*:=\mathrm{Hom}(V,K)$, the \textbf{dual vector space} (dual space) of $V$, where $K$ is considered as a $K$-vector space in itself: its addition is the addition of $K$ and its scalar multiplication the multiplication of $K$ \ts{8}{1}{25:44}\ts{8}{1}{26:04}\ts{8}{1}{26:25}.
\end{itemize}
\end{definition}
$\mathrm{Aut}(V)$ is not a vector subspace of $\mathrm{End}(V)$ (the sum of two invertible maps need not be invertible; the lecture says ``vector subspaces'' and then ``usually proper'': there are fewer invertible linear maps than linear maps, and the inclusion $\mathrm{Aut}(V)\subseteq\mathrm{End}(V)$ is proper unless the dimension is zero) \ts{8}{1}{24:42}\ts{8}{1}{25:04}\ts{8}{1}{25:24}\ts{8}{1}{25:44}.
\begin{coursenote}[Aut is not a subspace]
The lecture calls $\mathrm{Aut}(V)$ a vector subspace of $\mathrm{End}(V)$; it is a subset, indeed a group under composition, but not a subspace: it does not contain $0$ and is not closed under addition ($\mathrm{id}+(-\mathrm{id})=0$). Not said in class.
\end{coursenote}
Since $\mathrm{End}(V)$, $\mathrm{Aut}(V)$ and $V^*$ are special cases of $\mathrm{Hom}$, they inherit the operations just constructed (for $V^*$ the operations are $\diamond+$ and $\diamond\cdot$) \ts{8}{1}{26:25}\ts{8}{1}{26:45}\ts{8}{1}{27:06}. The dual space will play an important role: from a vector space and its dual all tensors are constructed \ts{8}{1}{27:06}.

% ==================================================================
\section{Tensors and the tensor space}\label{sec:te-tensors}

\begin{definition}[New concepts: $(p,q)$ tensor]\label{def:te-tensor}
Let $V$ be a $K$-vector space. A \textbf{$(p,q)$ tensor} $T$ over $V$ is a multilinear map
\[
  T\colon \underbrace{V^*\times\dots\times V^*}_{p}\times\underbrace{V\times\dots\times V}_{q}\ \xrightarrow{\ \sim\ }\ K,
\]
a map from $p$ copies of $V^*$ and (Cartesian product) $q$ copies of $V$ to $K$ that is linear in each slot, independently of the other slots: for the linear map one had additivity and scaling, now this applies to each slot
\ts{8}{1}{27:06}\ts{8}{1}{27:28}\ts{8}{1}{27:49}\ts{8}{1}{28:10}.
The pair $(p,q)$ is called the \textbf{valence} of the tensor; sometimes $p+q$ is called its \textbf{rank} \ts{8}{1}{37:29}\ts{8}{1}{37:50}.
\end{definition}
(The name \emph{tensor}: not explained in the lecture; see the coursenote at the end of the section.)

\begin{definition}[New concepts: the tensor space $T^p_qV$]\label{def:te-tspace}
One defines
\[
  T^p_qV:=\underbrace{V\otimes\dots\otimes V}_{p}\otimes\underbrace{V^*\otimes\dots\otimes V^*}_{q}\ :=\ \{\,T\mid T\text{ is a }(p,q)\text{ tensor over }V\,\},
\]
the set of all $(p,q)$ tensors (note that the order of $V$ and $V^*$ in the symbol is the opposite of the arguments in the maps). This is ``just a name'': ``if you ask how this symbol is defined, I don't need to tell you, it is a name; I could also call it $T^p_qV$ identically, an equally good name''
\ts{8}{1}{28:10}\ts{8}{1}{28:32}\ts{8}{1}{28:52}\ts{8}{1}{29:13}\ts{8}{1}{30:38}\ts{8}{1}{30:59}.
It is equipped with the \textbf{pointwise} addition and scalar multiplication: for example for two $(1,1)$ tensors $T,S$ (each eats an element $\sigma\in V^*$ and a vector $v\in V$)
\[
  (T+S)(\sigma,v):=T(\sigma,v)+S(\sigma,v),\qquad(\lambda T)(\sigma,v):=\lambda\,T(\sigma,v),
\]
where the $+$ is the plus in the target $K$, and the product the multiplication in $K$ \ts{8}{1}{29:13}\ts{8}{1}{29:36}\ts{8}{1}{29:57}\ts{8}{1}{30:18}\ts{8}{1}{30:38}.
This space with these operations is the \textbf{$(p,q)$ tensor space over the vector space $V$} \ts{8}{1}{30:59}\ts{8}{1}{31:20}\ts{8}{1}{36:24}\ts{8}{1}{36:44}.
\end{definition}
It is a vector space: one checks the eight axioms (commutative, associative, neutral, inverse, an associative law, two distributive laws and the unitary law) \ts{8}{1}{36:44}\ts{8}{1}{37:08}.

\begin{secondary}[That the tensor space is a set]
The condition that the maps be multilinear is a condition imposed on the maps; to be sure that this is a set one needs the principle of restricted comprehension, so one has to say from which set the elements are taken. Recall that the notation of a map hides that a map is a relation (a functional relation), which is a subset of the Cartesian
product of the domain and the target: so a map $f\colon A\to B$ is an element of the power set $\mathcal P(A\times B)$, and the multilinear maps are selected from the power set of $(V^*)^p\times V^q\times K$. That is how restricted comprehension shows that the $(p,q)$ tensor space is indeed a set (``a general question that I will not
write down every time'') \ts{8}{1}{33:14}\ts{8}{1}{33:37}\ts{8}{1}{33:58}\ts{8}{1}{34:18}\ts{8}{1}{34:39}\ts{8}{1}{35:00}\ts{8}{1}{35:20}\ts{8}{1}{35:41}\ts{8}{1}{36:02}.
\end{secondary}
(See \cref{thm:se-comprehension} for the principle, and \cref{def:cl-map} for maps as relations.)

\begin{definition}[New concepts: tensor product of tensors]\label{def:te-tensorprod}
There are several ways to construct new tensors from old ones; the most important is the \textbf{tensor product}. Unlike the addition, which takes two tensors of the same valence, it takes $T\in T^p_qV$ and $S\in T^r_sV$ of an entirely different valence and produces a tensor of valence $(p+r,q+s)$ over the same $V$:
\[
  (T\otimes S)(\omega_1,\dots,\omega_{p+r},v_1,\dots,v_{q+s}):=T(\omega_1,\dots,\omega_p,v_1,\dots,v_q)\cdot S(\omega_{p+1},\dots,\omega_{p+r},v_{q+1},\dots,v_{q+s}),
\]
where $T$ eats as many co-vectors and vectors as it can, in that order, the product on the right is the multiplication in $K$, and $S$ eats the rest
\ts{8}{1}{37:08}\ts{8}{1}{37:29}\ts{8}{1}{37:50}\ts{8}{1}{38:11}\ts{8}{1}{38:32}\ts{8}{1}{38:52}\ts{8}{1}{39:13}\ts{8}{1}{39:34}\ts{8}{1}{39:55}\ts{8}{1}{40:16}\ts{8}{1}{40:36}.
\end{definition}
\begin{remark}[Two meanings of the symbol $\otimes$]\label{rem:te-two-otimes}
The symbol $\otimes$ is used twice: in $V\otimes\dots\otimes V^*$ as mere notation (the ``orange'' tensor product, a name for the set of \cref{def:te-tspace}), and between two \emph{elements} of tensor spaces, $T\otimes S$, defined as above (the ``blue'' tensor product)
\ts{8}{1}{38:11}\ts{8}{1}{38:32}. The vectors $v_1,\dots,v_p$ here are individual vectors, not components of a tensor or of a vector (components have not been discussed so far) \ts{8}{1}{39:13}\ts{8}{1}{39:34}.
\end{remark}
\begin{coursenote}[Valence indices]
In the lecture the argument of $T$ is called $(V^*)^p\times V^q$ and the tensor space ``$V\otimes\dots\otimes V^*$'' with the order of $V$ and $V^*$ reversed with respect to the arguments; it is the lecture's own remark (\ts{8}{1}{28:32}\ts{8}{1}{28:52}) that this order is different, I keep it so. The name ``tensor''
(Latin \emph{tendere}, to stretch) is not used in class; the derivation is [?] not said in class.
\end{coursenote}

% ==================================================================
\section{First examples: $T^0_1V$ and $T^1_1V$}\label{sec:te-examples}

Examples are in order \ts{8}{1}{40:58}. Start with $T^0_1V$ over a vector space \ts{8}{1}{40:58}\ts{8}{1}{41:19}.
\begin{example}[$(0,1)$ tensors are co-vectors]\label{ex:te-01}
By the notation $T^0_1V=V^*$ (no tensor product appears) and by definition it consists of all multilinear (linear) maps $V\to K$: indeed the definition of $V^*$. The notation is consistent with what was defined before. So the $(0,1)$ tensors are the dual vectors, the whole space is the dual space,
and a $(0,1)$ tensor is called a \textbf{co-vector}, which in physics is its common name \ts{8}{1}{41:19}\ts{8}{1}{41:40}\ts{8}{1}{42:00}\ts{8}{1}{42:21}. What is called a vector and a co-vector depends on what one calls a vector: if the elements of $V$ are called vectors, those of $V^*$ are co-vectors; in finite dimension one can turn this around (below)
\ts{8}{1}{42:21}\ts{8}{1}{42:42}.
\end{example}

\begin{proposition}[$T^1_1V\cong\mathrm{End}(V^*)$]\label{prop:te-11}
The $(1,1)$ tensors over $V$, i.e.\ the multilinear maps $V^*\times V\to K$ (here, as in the lecture, the argument order is $V^*$ first), are isomorphic as a vector space to $\mathrm{End}(V^*)$
\ts{8}{1}{42:42}\ts{8}{1}{43:03}\ts{8}{1}{43:24}\ts{8}{1}{43:44}.
\end{proposition}
\begin{proof}
For every bilinear map construct an element of $\mathrm{End}(V^*)$, and show how to recover precisely the map one started from; then the two contain the same information, one can map one to the other one to one; one must also check that the correspondence is linear, in both directions, which is very simple
\ts{8}{1}{43:44}\ts{8}{1}{44:04}.

Given $T\in T^1_1V$ construct $\hat T\in\mathrm{End}(V^*)$ as follows. $\hat T$ must assign to every $\omega\in V^*$ another element of $V^*$. Use $T$ with its first slot filled with $\omega$: what remains, $v\mapsto T(\omega,v)$, is an object into which one can plug a vector and get a number,
a map $V\to K$, linear since $T$ is linear in the second slot; such a linear map is an element of $V^*$. Define
\[
  \hat T(\omega):=T(\omega,\,\cdot\,)\in V^*
\]
\ts{8}{1}{44:04}\ts{8}{1}{44:26}\ts{8}{1}{44:46}\ts{8}{1}{45:07}\ts{8}{1}{45:28}\ts{8}{1}{45:49}.
No information is lost: $T$ can be reconstructed, since $T(\omega,v)=\bigl(\hat T(\omega)\bigr)(v)$ (the number obtained by letting the co-vector $\hat T(\omega)$ act on $v$), by the very definition of $\hat T$ \ts{8}{1}{45:49}\ts{8}{1}{46:09}\ts{8}{1}{46:29}.
Hence $T\mapsto\hat T$ is a bijection; one checks that it is also linear \ts{8}{1}{46:29}\ts{8}{1}{46:49}. (Linearity of $\hat T(\omega)$ in $\omega$ comes from linearity of $T$ in the first slot, hence $\hat T\in\mathrm{End}(V^*)$; not spelled out in class.)
\end{proof}
One might also feel that $T^1_1V$ is the same as $\mathrm{End}(V)$, but this is generically not true \ts{8}{1}{46:49}\ts{8}{1}{47:10}.

\paragraph{The list of questions.} Continue a list \ts{8}{1}{47:10}\ts{8}{1}{47:30}:
\begin{itemize}
  \item $T^0_1V=V^*$ is the dual.
  \item Any physics textbook says the $(1,0)$ tensors are, as vector space, the same as the vector space $V$ itself, while the $(0,1)$ ones are the duals: is this really true? In general \emph{not} \ts{8}{1}{47:30}\ts{8}{1}{47:51}.
  \item Are the $(1,1)$ tensors the same as $\mathrm{End}(V)$? Yes, but not generically \ts{8}{1}{47:51}\ts{8}{1}{48:11}.
  \item The fundamental question that links the two and answers them: if one takes the dual and then the dual of the dual, does one get $V$ back? Again no, not always, but yes if the vector space is \emph{finite dimensional}
        \ts{8}{1}{48:11}\ts{8}{1}{48:32}.
\end{itemize}
The dimension has not been defined yet: that is what is done next, and only with a notion of dimension can one restrict to the finite dimensional case \ts{8}{1}{48:32}\ts{8}{1}{48:53}.

% ==================================================================
\section{Bases and dimension}\label{sec:te-basis}

Some care is needed in the definition of dimension: here one deals with \emph{bare} vector spaces, no additional structure, no norm. With no further structure than the vector space structure the only notion of basis one can define is the so-called \textbf{Hamel basis}
\ts{8}{1}{48:53}\ts{8}{1}{49:13}\ts{8}{1}{49:35}.
\begin{definition}[New concepts: Hamel basis, dimension]\label{def:te-basis}
Let $(V,+,\cdot)$ be a vector space. A subset $B\subseteq V$ is called a \textbf{Hamel basis} (to avoid confusion with other notions of basis) if
\begin{enumerate}
  \item every finite subset $\{b_1,\dots,b_N\}\subseteq B$ is \emph{linearly independent}, that is
        \[
          \sum_{i=1}^N\lambda^ib_i=0\ \Longrightarrow\ \lambda^1=\dots=\lambda^N=0
        \]
        (the index is put up for future use); and
  \item $B$ is a \emph{generating system}: for every $v\in V$ there exist a finite number $M$ of numbers $v^1,\dots,v^M\in K$ and finitely many elements $b_1,\dots,b_M\in B$ such that $v=\sum_{i=1}^Mv^ib_i$.
\end{enumerate}
The \textbf{dimension} of $V$ is $\dim V:=$ the cardinality of a Hamel basis; it may be finite or infinite (``various infinities'')
\ts{8}{1}{49:35}\ts{8}{1}{49:57}\ts{8}{1}{50:21}\ts{8}{1}{50:45}\ts{8}{1}{51:06}\ts{8}{1}{53:10}\ts{8}{1}{53:31}\ts{8}{1}{53:52}\ts{8}{1}{55:16}\ts{8}{1}{55:37}.
\end{definition}
(That the dimension is well defined, i.e.\ independent of the basis, ``there is something to prove, but we are not going to do that'' \ts{8}{1}{55:16}\ts{8}{1}{55:37}.)
The condition in the first item is a proposition, an implication, either true or false, and its truth or falseness is equivalent to the linear independence of the finite subset; ``this is something that in high school causes students usually great pain to evaluate, mainly because they have not had a proper explanation
of the implication arrow'' (\cref{sec:lo-prop}) \ts{8}{1}{50:45}\ts{8}{1}{51:06}\ts{8}{1}{51:27}\ts{8}{1}{51:47}.

\emph{Why finiteness.} One only looks at \emph{finite} subsets and finite sums. If one required that all the elements of $B$ be linearly independent in the sum sense there would be two problems: $B$ could be uncountable, and then one cannot even write down a sum, not even a series (a series requires partial sums);
and even for a countable basis the sum would become a series up to infinity, but to know what a series does one needs a notion of convergence, a limit; on a vector space without additional structure, not even a topology, one cannot talk about convergence. So on a bare vector space the only thing one can do is to check that for every finite
subset the condition holds \ts{8}{1}{51:47}\ts{8}{1}{52:08}\ts{8}{1}{52:28}\ts{8}{1}{52:49}\ts{8}{1}{53:10}.
\begin{example}[Fourier analysis]\label{ex:te-fourier}
In Fourier analysis the cosines and sines over the real numbers are countably many and form a \emph{countable} basis, but only at the price of having Fourier series, infinite sums: so they do not constitute a Hamel basis; they constitute what is called a \textbf{Schauder basis} [?: ``sha basis'' in the transcript], and for that one needs
more structure, a notion of convergence. A Hamel basis for the functions that are Fourier expandable would be \emph{uncountable}, whereas the Schauder basis is countable; and one uses the Schauder bases in practice
\ts{8}{1}{54:12}\ts{8}{1}{54:35}\ts{8}{1}{54:55}\ts{8}{1}{55:16}.
\end{example}
\begin{history}[Hamel]
\emph{Not said in class.} \emph{Checked.} The notion is named after G.~Hamel, who in 1905 used a basis of $\R$ as a vector space over $\mathbb Q$ to construct discontinuous additive functions; the existence of a Hamel basis of every vector space is equivalent to the axiom of choice (\cref{sec:se-choice}).\par
\emph{References.} G.~Hamel, ``Eine Basis aller Zahlen und die unstetigen L\"osungen der Funktionalgleichung $f(x+y)=f(x)+f(y)$'', \emph{Math.\ Ann.}\ 60 (1905).
\end{history}

\begin{theorem}[Double dual]\label{thm:te-dd}
If $\dim V<\infty$, then $V^{**}:=(V^*)^*$ is isomorphic as a vector space to $V$
\ts{8}{1}{55:57}\ts{8}{1}{56:20}.
\end{theorem}
The proof is asked on the upcoming problem sheet (a vector space isomorphism has to be \emph{constructed}; see the exercise below and \cref{exam:ch08}) \ts{8}{1}{55:57}\ts{8}{1}{56:20}\ts{8}{1}{1:52:57}.
\begin{exercise}[The canonical isomorphism $V\to V^{**}$]\label{ex:te-dd}
Construct a vector space isomorphism $V\to V^{**}$ for finite dimensional $V$, without choosing a basis. (Problem sheet.) The long version is in \cref{exam:ch08}.
\end{exercise}

\paragraph{Vectors and co-vectors are interchangeable in finite dimension.} Hence if one has $V$ and $V^*$, one can identify $V$ with $V^{**}$ for $\dim V<\infty$; one could call the elements of $V$ vectors and those of $V^*$ co-vectors, but $V^*$ is a vector space, so its elements could also be called vectors
and then the elements of $V$ would be the duals of those, the co-vectors. So in finite dimension what is a vector and what is a co-vector can be turned around; but in manifold theory one thinks of tangent spaces as the vector space and over the tangent space one constructs the co-vector space, so there is
the \emph{conventional asymmetry}: one calls these vectors and those co-vectors; but one needs to know that one starts with the space $V$. ``No problem in practice, but once you start thinking about this you might get confused, so I tried to confuse you beforehand; I do not know where
that works pedagogically, but what I said is true'' \ts{8}{1}{56:20}\ts{8}{1}{56:41}\ts{8}{1}{57:01}\ts{8}{1}{57:21}\ts{8}{1}{57:42}.
Once $V^{**}\cong V$ one can redo the construction of \cref{prop:te-11} from $V$ to $V$ and quickly see that $T^1_1V\cong\mathrm{End}(V)$ as well, and likewise the $(1,0)$ tensors are $V^{**}\cong V$: but only for finite dimensional $V$ \ts{8}{1}{57:42}\ts{8}{1}{58:02}\ts{8}{1}{58:25}.
(The lecture comes back to this point in answer to a student's question, \cref{int:te-10}, in the part on transformation of components.)

% ==================================================================
\section{Components, the dual basis and reconstruction}\label{sec:te-components}

So far bases were used only to classify vector spaces into finite and infinite dimensional ones, but bases have another property \ts{8}{1}{58:25}\ts{8}{1}{58:45}. Instead of talking about \emph{the vector}, a rather abstract concept, if one wants to hand over a concrete vector one agrees on a basis and gives its
\textbf{components} $v^1,\dots,v^M$ with respect to that basis; another vector by its components. Bases are used to communicate which precise vector one means, always subject to the choice of basis \ts{8}{1}{58:45}\ts{8}{1}{59:06}\ts{8}{1}{59:26}\ts{8}{1}{59:47}.

\begin{intuition}[Construct without bases]\label{int:te-basisfree}
There is a warning: if one has the components of a vector and starts defining new objects in terms of components, there is the danger of doing \emph{illegal constructions} that are not good in vector space theory, in particular constructions that would take a very different shape if one changed bases. The general rule: always construct
everything \emph{without choice of basis}; only once it is constructed, one may introduce a basis to study which array of numbers is needed to provide the new object. ``Some people say a gentleman only chooses a basis if he really must''; if one has a concrete reason to choose a basis in order to calculate something it is fine,
but constructions should be basis free \ts{8}{1}{59:47}\ts{8}{1}{1:00:07}\ts{8}{1}{1:00:28}\ts{8}{1}{1:00:49}.
\end{intuition}

\begin{definition}[New concepts: components of a tensor]\label{def:te-components}
From now on only \emph{finite dimensional} vector spaces. Let $T$ be a $(p,q)$ tensor over $V$ and $e_1,\dots,e_{\dim V}$ a basis of $V$ (that induces a basis on the tensor space, see below), and $\varepsilon^1,\dots,\varepsilon^{\dim V}$ a basis of $V^*$ (it has the same size as $\dim V^*=\dim V$). The \textbf{components} of $T$ with
respect to these bases are the numbers (in the underlying field, say the real numbers)
\[
  T^{a_1\dots a_p}{}_{b_1\dots b_q}:=T\bigl(\varepsilon^{a_1},\dots,\varepsilon^{a_p},e_{b_1},\dots,e_{b_q}\bigr)\in K.
\]
The tensor is an abstract thing, a multilinear map; its components are numbers, a whole lot of numbers \ts{8}{1}{1:00:49}\ts{8}{1}{1:01:09}\ts{8}{1}{1:01:31}\ts{8}{1}{1:01:53}\ts{8}{1}{1:02:13}\ts{8}{1}{1:02:36}\ts{8}{1}{1:02:57}\ts{8}{1}{1:03:17}.
\end{definition}
(The lecturer writes the basis of $V$ as $B_1,\dots,B_{\dim V}$, then renames it $e$, ``more standard'', and forgets for a moment the second basis \ts{8}{1}{1:02:13}\ts{8}{1}{1:02:36}.)

\begin{definition}[New concepts: dual basis]\label{def:te-dualbasis}
Usually one does not choose independent bases on $V$ and $V^*$ (one could, they are two vector spaces); one chooses on $V^*$ a basis \emph{induced} by the previous choice on $V$, by the condition
\[
  \varepsilon^a(e_b)=\delta^a_b,\qquad\delta^a_b:=\begin{cases}1&a=b\\0&a\ne b\end{cases}
\]
(the $a$-th basis co-vector eats the $b$-th basis vector). This is a linear relation with a unique solution: for every choice of basis on $V$ and this condition one gets a unique basis of $V^*$, called the \textbf{dual basis} (of the dual space); it is ``dual'' because it is a basis of a dual space and
because of this condition. One could choose just any basis of the dual space; ``we choose the dual basis of the dual space'', and from now on bases are always dual bases unless said otherwise \ts{8}{1}{1:03:17}\ts{8}{1}{1:03:37}\ts{8}{1}{1:03:58}\ts{8}{1}{1:04:19}\ts{8}{1}{1:04:40}\ts{8}{1}{1:05:43}\ts{8}{1}{1:06:03}\ts{8}{1}{1:06:24}\ts{8}{1}{1:26:15}.
\end{definition}
Why impose this condition and not another? No particular reason: but ``if you have to make a choice, make sure that you do not make choices all over the place; if one choice can already be used to imply another it is usually most efficient to keep the number of choices minimal''; one could argue that this
choice is made once and for all. It is a very convenient induction principle, since many calculations simplify because there are so many zeros around \ts{8}{1}{1:04:40}\ts{8}{1}{1:05:01}\ts{8}{1}{1:05:21}.

The most important point: once one has a basis, and a tensor with its components, one can \emph{reconstruct the abstract tensor from its components} \ts{8}{1}{1:06:24}\ts{8}{1}{1:06:44}.
\begin{theorem}[Reconstruction of a tensor from its components]\label{thm:te-recon}
Let $V$ be finite dimensional, $e_a$ a basis, $\varepsilon^a$ its dual basis, and $T\in T^p_qV$ with components $T^{a_1\dots a_p}{}_{b_1\dots b_q}$. Then
\[
  T=T^{a_1\dots a_p}{}_{b_1\dots b_q}\cdot e_{a_1}\otimes\dots\otimes e_{a_p}\otimes\varepsilon^{b_1}\otimes\dots\otimes\varepsilon^{b_q},
\]
where (summation convention, see below) the dot is the scalar multiplication in the tensor space $T^p_qV$ (number in the field times a given tensor), the basis vectors $e_a$ are $(1,0)$ tensors and the $\varepsilon^b$ $(0,1)$ tensors, and the $\otimes$ is the tensor product of \cref{def:te-tensorprod} (the ``blue'' one)
\ts{8}{1}{1:06:44}\ts{8}{1}{1:07:07}\ts{8}{1}{1:07:28}\ts{8}{1}{1:07:49}\ts{8}{1}{1:08:10}\ts{8}{1}{1:08:34}\ts{8}{1}{1:08:55}\ts{8}{1}{1:09:16}.
\end{theorem}
Reading: the string $e_{a_1}\otimes\dots\otimes\varepsilon^{b_q}$ is a $(p,q)$ tensor, as is $T$; there are many such tensors, one for each particular choice of $a_1=7,\dots,a_p=28,b_1=1,\dots,b_q=1$ and so on, and one multiplies each by the number $T^{a_1\dots}{}_{b_1\dots}$
\ts{8}{1}{1:08:34}\ts{8}{1}{1:08:55}\ts{8}{1}{1:09:16}.
\begin{proof}
Not carried out in class: one evaluates the right-hand side on an arbitrary collection of co-vectors and vectors, expanding them in the bases, and uses how the $e$'s and the $\varepsilon$'s act on arbitrary vectors and co-vectors; ``it is very basic'' and a finger exercise on the next problem sheet (\cref{exam:ch08});
``this equal sign here is correct because it is in every textbook, and it has practical importance'' \ts{8}{1}{1:09:16}\ts{8}{1}{1:09:38}\ts{8}{1}{1:09:59}\ts{8}{1}{1:10:20}.
\end{proof}
\begin{exercise}[Reconstruction]\label{ex:te-recon}
Verify \cref{thm:te-recon} for a $(1,1)$ tensor by evaluating both sides on $(\varepsilon^c,e_d)$ and on $(\omega,v)=(\omega_a\varepsilon^a,\,v^be_b)$. (Problem sheet; \cref{exam:ch08}.)
\end{exercise}

% ==================================================================
\section{Change of basis and the summation convention}\label{sec:te-change}

Finite dimensional spaces; let $e_1,\dots,e_d$ ($d:=\dim V$) be a basis of whatever vector space \ts{8}{1}{1:10:20}\ts{8}{1}{1:10:41}.
\begin{definition}[Change of basis]\label{def:te-change}
A new basis $\tilde e_1,\dots,\tilde e_d$ (the same number of elements) consists of vectors, each expressible as a linear combination of the old basis (that is what a basis is):
\[
  \tilde e_a=A^b{}_a\,e_b\qquad(\text{sum over }b=1,\dots,d),
\]
with $A^b{}_a\in K$ (two indices: there is a set of components for every new vector $a$), the scalar multiplication being the one of $V$. Not every such combination gives a basis: one could destroy the linear independence and the generating property. The condition is that one can also recover the old basis from the new one,
\[
  e_a=B^b{}_a\,\tilde e_b,
\]
and if both are possible one checks easily that $B$, as a matrix, is the \emph{inverse} of $A$, $B^a{}_mA^m{}_b=\delta^a_b$
\ts{8}{1}{1:10:41}\ts{8}{1}{1:11:02}\ts{8}{1}{1:11:23}\ts{8}{1}{1:11:43}\ts{8}{1}{1:12:05}\ts{8}{1}{1:12:25}\ts{8}{1}{1:12:45}\ts{8}{1}{1:13:07}.
\end{definition}
(Matrices ``in a second''. The index placement of $A$ and $B$ in the transcript is ambiguous; I have chosen the placement that makes the formulas below consistent [?].)

\begin{definition}[New concepts: summation convention]\label{def:te-einstein}
In the reconstruction \cref{thm:te-recon} there are, actually, a succession of $p+q$ sum symbols, $\sum_{a_1=1}^d\dots\sum_{b_q=1}^d$, each running from $1$ to $\dim V$ (the sums over indices of $V$ and of the co-vector space), which were not written; they must be there, since this is the expansion; but by convention they are forgotten. In the future they are not written:
whenever an index appears once up and once down, in the same product (``product in a generalised sense; not if it is in a sum''), it is summed over, unless something else is explicitly said. This is the \textbf{Einstein summation convention}: omit the $\sum_{a=1}^d$ if $a$ appears once up and once down
\ts{8}{1}{1:13:07}\ts{8}{1}{1:13:28}\ts{8}{1}{1:13:49}\ts{8}{1}{1:14:11}\ts{8}{1}{1:14:31}\ts{8}{1}{1:14:52}\ts{8}{1}{1:17:41}\ts{8}{1}{1:18:02}\ts{8}{1}{1:18:23}\ts{8}{1}{1:18:46}.
\end{definition}
\begin{definition}[The up/down convention]\label{def:te-updown}
How does an index deserve to be up or down? There is only one fundamental choice, and everything else follows: the \emph{basis vectors} of $V$ are labelled by a \emph{downstairs} index, $e_1,\dots,e_d$ (each is an element of $V$, not a component: the first vector, the second vector, the $d$-th), and bases of $V^*$ by an \emph{upstairs}
index, $\varepsilon^1,\dots,\varepsilon^d$. Why does everything else follow? In the definition of the components the tensor first eats the co-vectors, which have their index up, so the resulting component has the index up (an equation should match upstairs and downstairs on both sides); the vector slots take $e_b$ with the index down,
so the component has the index down: ``nothing you have to memorise on top; you just have to memorise the fundamental convention''. When one reconstructs, the indices are already in position \ts{8}{1}{1:15:13}\ts{8}{1}{1:15:35}\ts{8}{1}{1:15:57}\ts{8}{1}{1:16:18}\ts{8}{1}{1:16:40}\ts{8}{1}{1:17:01}\ts{8}{1}{1:17:21}\ts{8}{1}{1:17:41}.
\end{definition}
\begin{history}[The summation convention]
\emph{Not said in class.} \emph{Checked.} The convention of omitting the sum sign over a repeated index was introduced by A.~Einstein in his review of general relativity.\par
\emph{References.} A.~Einstein, ``Die Grundlage der allgemeinen Relativit\"atstheorie'', \emph{Ann.\ Phys.}\ 49 (1916).
\end{history}

\begin{remark}[The convention works only with linear maps]\label{rem:te-linear-only}
A remark: one could think ``let us always do this in all branches of mathematics'', a good convention which saves chalk; but it only works with \emph{linear} spaces, in particular linear maps \ts{8}{1}{1:18:46}\ts{8}{1}{1:19:07}\ts{8}{1}{1:19:30}.
Example: let $f$ be a $(1,1)$ tensor, $f\colon V^*\times V\xrightarrow{\sim}K$. Plug in a co-vector $\omega=\omega_a\varepsilon^a$ (basis of $V^*$ with up index, so the components $\omega_a$ have a down index) and a vector $v=v^be_b$:
\[
  f(\omega,v)=f\bigl(\omega_a\varepsilon^a,\,v^be_b\bigr)=\omega_a\,v^b\,f(\varepsilon^a,e_b)=\omega_a\,v^b\,f^a{}_b,
\]
where one used the multilinearity of $f$: additivity pulls out the (invisible) finite sums over $a$ and over $b$, and linearity pulls out the numbers $\omega_a,v^b\in K$ \ts{8}{1}{1:19:30}\ts{8}{1}{1:19:50}\ts{8}{1}{1:20:11}\ts{8}{1}{1:20:31}\ts{8}{1}{1:20:52}\ts{8}{1}{1:21:13}\ts{8}{1}{1:21:34}\ts{8}{1}{1:21:54}\ts{8}{1}{1:22:15}\ts{8}{1}{1:22:35}\ts{8}{1}{1:22:56}.
So the abstract objects $f,\omega,v$ give a number, and one obtains it by building a big sum and product out of their components: one forgets the abstract objects and uses the components with respect to a basis and its dual basis \ts{8}{1}{1:22:56}\ts{8}{1}{1:23:17}\ts{8}{1}{1:23:38}.
\end{remark}
Imagine now one does not see orange (the chalk in which the sums were written in): what is the difference between the first and the second line? Nothing; but the equal sign is not for nothing: it is because of the linearity of $f$ that one can pull out the sum. If one chooses to omit sum signs, an invisible sum can be invisible here or invisible there: ``that is the point
of invisibility, you do not see where it is'', and so the map better be linear, since otherwise it would make a difference where the sum sits, and one would not see it. The summation convention only works when it plays with linear maps; with polynomial maps, as in algebraic geometry, ``and you use the summation convention at least in unclever ways, you are in trouble''
\ts{8}{1}{1:23:38}\ts{8}{1}{1:23:58}\ts{8}{1}{1:24:19}\ts{8}{1}{1:24:39}\ts{8}{1}{1:25:01}\ts{8}{1}{1:25:21}\ts{8}{1}{1:25:42}.

% ==================================================================
\section{Matrices as ``witchcraft mathematics''}\label{sec:te-matrices}

A few polemical remarks about matrices \ts{8}{1}{1:25:42}\ts{8}{1}{1:26:15}. Having chosen a basis in $V$ and the dual basis in $V^*$ it is \emph{tempting} to think of a co-vector $\omega=\omega_a\varepsilon^a$ as the collection of its components and write them as a \emph{row}, $\omega\ \hat=\ (\omega_1,\dots,\omega_d)$, and the vector $v=v^be_b$ as a \emph{column},
$v\ \hat=\ (v^1,\dots,v^d)^{\mathsf T}$; nothing wrong in itself, as the bases are fixed, one has the bases in mind \ts{8}{1}{1:26:15}\ts{8}{1}{1:26:35}\ts{8}{1}{1:26:56}\ts{8}{1}{1:27:17}\ts{8}{1}{1:27:37}\ts{8}{1}{1:27:58}.
Further, $\phi\in\mathrm{End}(V)$ is, as shown, a $(1,1)$ tensor over $V$, $\phi=\phi^a{}_b\,e_a\otimes\varepsilon^b$, and one is tempted, as is the basis of any engineering mathematics, to think of the linear map as a square array with $\phi^1{}_1$ in the top left, the first row containing $\phi^1{}_1,\phi^1{}_2,\dots,\phi^1{}_d$,
then $\phi^2{}_1$, and so on down to $\phi^d{}_d$ in the last line: the \emph{row-column rule}, the first index the row and the second the column \ts{8}{1}{1:27:58}\ts{8}{1}{1:28:18}\ts{8}{1}{1:28:40}\ts{8}{1}{1:29:01}\ts{8}{1}{1:29:21}\ts{8}{1}{1:29:42}\ts{8}{1}{1:30:03}.
(The word ``tempting'' in a mathematics lecture means: be careful \ts{8}{1}{1:26:15}\ts{8}{1}{1:26:35}.)

\begin{intuition}[A matrix is one way to arrange numbers]\label{int:te-matrix}
This is just \emph{one way} to arrange the information, the $d^2$ numbers, as much as a row or column is one way to arrange $d$ numbers: pure convention. A joke on arrangements: if you look for your great-great-grandfather in the cemetery (``I choose this because you have no emotional attachment to him''), he lies in row 7, column 5; if
another scheme were used he would lie elsewhere \ts{8}{1}{1:30:03}\ts{8}{1}{1:30:26}\ts{8}{1}{1:30:46}\ts{8}{1}{1:31:07}. (``Symmetry'' in the transcript is the speech recogniser's rendering of \emph{cemetery}; I read it so throughout.)
\end{intuition}
What does the convention imply? Take two endomorphisms $\phi,\psi\in\mathrm{End}(V)$ and the composition $\phi\circ\psi$, again an endomorphism (the composition of linear maps is linear) \ts{8}{1}{1:31:07}\ts{8}{1}{1:31:28}\ts{8}{1}{1:31:48}.
\begin{proposition}[Matrix multiplication is the representation of composition]\label{prop:te-matmult}
The components of $\phi\circ\psi$ are $(\phi\circ\psi)^a{}_b=\phi^a{}_m\,\psi^m{}_b$, so that in the row-column arrangement they are the matrix product (row of the first times column of the second).
\end{proposition}
\begin{proof}
The component of an endomorphism: $(\phi\circ\psi)^a{}_b=\varepsilon^a\bigl(\phi(\psi(e_b))\bigr)$ (the tensor ``eats'' the result, then $\varepsilon^a$ acts) \ts{8}{1}{1:31:48}\ts{8}{1}{1:32:09}\ts{8}{1}{1:32:30}.
Now $\psi(e_b)=\psi^m{}_b\,e_m$ (quickly checked from what was said before); $\phi$ is linear, so the invisible sum and the numbers $\psi^m{}_b$ are pulled out, leaving $\psi^m{}_b\,\varepsilon^a\bigl(\phi(e_m)\bigr)$, and again $\varepsilon^a$ is linear, so the numbers can be pulled out: $\varepsilon^a(\phi(e_m))=\phi^a{}_m$ is the component of $\phi$. Hence
\[
  (\phi\circ\psi)^a{}_b=\psi^m{}_b\,\phi^a{}_m=\phi^a{}_m\,\psi^m{}_b,
\]
the last step because the product is in $K$, a field, so commutative \ts{8}{1}{1:32:30}\ts{8}{1}{1:32:50}\ts{8}{1}{1:33:10}\ts{8}{1}{1:33:30}\ts{8}{1}{1:33:51}\ts{8}{1}{1:34:12}\ts{8}{1}{1:34:32}.
\end{proof}
Looking through ``cemetery glasses'' one prefers the second order: the indices to be summed over are close together, so that for the first endomorphism the row index and the column index, and for the second again row and column; one sums over the column index of the first and the row index of the second, and gets the standard rule
\emph{row times column}: the entry in row $a$ and column $b$ of $\phi\circ\psi$ is obtained by running over the columns of the first matrix with the row fixed and over the rows of the second with the column fixed \ts{8}{1}{1:34:32}\ts{8}{1}{1:34:53}\ts{8}{1}{1:35:14}\ts{8}{1}{1:35:36}\ts{8}{1}{1:35:57}\ts{8}{1}{1:36:18}\ts{8}{1}{1:36:38}.
Then one only writes the arrays, defines a new rule (the row-times-column rule, \emph{matrix multiplication}) and writes $\hat\phi\,\hat\psi$ for the arrays: ``the whole point of matrix multiplication is that it is the cemetery rule for finding the representing array of the composition. This is the deep reason behind row times column, but it relies on this
additional construction'' \ts{8}{1}{1:36:38}\ts{8}{1}{1:36:59}\ts{8}{1}{1:37:20}.

The rule extends to $\mathrm{Hom}(V,W)$ with corresponding dimensions, and to the action of a co-vector from the left on a vector: without coordinates and bases it is clear that a co-vector acts on a vector to yield a number (this is the definition), and in components
\[
  \omega(v)=\omega_a\varepsilon^a\bigl(v^be_b\bigr)=\omega_a\,v^b\,\varepsilon^a(e_b)=\omega_a\,v^b\,\delta^a_b=\omega_m\,v^m,
\]
where one pulls out the (invisible) sum and the numbers and uses the dual basis; written as a row and a column this is again a matrix product \ts{8}{1}{1:37:20}\ts{8}{1}{1:37:42}\ts{8}{1}{1:38:02}\ts{8}{1}{1:38:23}\ts{8}{1}{1:38:44}.
Without the summation convention $f(\omega,v)=\omega_a\,f^a{}_b\,v^b$: these are all numbers, so they can be written in a different order; one writes $\omega$ as a row, $f$ as a matrix, $v$ as a column, and ends with the expression one has certainly seen: $\omega^{\mathsf T}\,f\,v$, where one writes a \emph{transpose} to emphasise that one comes from the left \ts{8}{1}{1:38:44}\ts{8}{1}{1:39:06}\ts{8}{1}{1:39:27}\ts{8}{1}{1:39:47}.

\begin{remark}[Resist the temptation]\label{rem:te-resist}
The transpose is ``not a good object at all'': very bad notation, almost pretending to be basis independent but not at all. One might think ``$\omega$ comes from the left and $v$ comes from the right, I calculate like in the introductory course''; but that is not necessary, it requires a lot of additional convention and is not worth it. For higher
rank it gets absurd: a $(1,2)$ tensor would have to be a cube of numbers, a $(2,2)$ tensor a four-dimensional cube, ``what is the point?'' It is a singular possibility for $(1,1)$ and $(0,2)$ tensors to come up with such witchcraft notation. One should think in terms of components and the summation rule, and always try to understand
these objects from the component-free, basis-free notation: the task is to gain facility in switching between basis notation and basis-free notation, as done above by plugging in the basis, and to forget the rest, which is confusing. It is mentioned because it is ingrained in our thinking, since it was used often in the beginning; but the lesson
is: do not do this. ``A gentleman only chooses a basis if he must'', and, the stronger principle (the widow principle), ``a lady must never go to the cemetery'': one should enjoy the fortunes of the basis-free notation
\ts{8}{1}{1:39:47}\ts{8}{1}{1:40:08}\ts{8}{1}{1:40:28}\ts{8}{1}{1:40:49}\ts{8}{1}{1:41:09}\ts{8}{1}{1:41:30}\ts{8}{1}{1:41:51}\ts{8}{1}{1:42:12}.
\end{remark}

% ==================================================================
\section{Transformation of components under a change of basis}\label{sec:te-transform}

After this longish aside, the change of basis: define a new basis $\tilde e_a=A^b{}_a\,e_b$ (with the inverse $B$ so that $e_a=B^b{}_a\tilde e_b$) and study how the components of vectors and co-vectors change; the components are basis dependent
\ts{8}{1}{1:42:12}\ts{8}{1}{1:42:32}\ts{8}{1}{1:42:54}.
\begin{proposition}[Transformation of components]\label{prop:te-transform}
Under $\tilde e_a=A^b{}_a\,e_b$, $e_a=B^b{}_a\,\tilde e_b$ (with $B=A^{-1}$), the components of a co-vector $\omega$ and a vector $v$ transform as
\[
  \tilde\omega_a=A^m{}_a\,\omega_m,\qquad \tilde v^a=B^a{}_m\,v^m,
\]
i.e.\ vector components transform with the \emph{inverse} matrix of the one for co-vector components (``vectors transform this way, co-vectors the other way''), and for a $(p,q)$ tensor
\[
  \tilde T^{a_1\dots a_p}{}_{b_1\dots b_q}=B^{a_1}{}_{m_1}\cdots B^{a_p}{}_{m_p}\;A^{n_1}{}_{b_1}\cdots A^{n_q}{}_{b_q}\;T^{m_1\dots m_p}{}_{n_1\dots n_q}:
\]
an upstairs index transforms like a vector index (with $B$), a downstairs index like a co-vector index (with $A$)
\ts{8}{1}{1:42:54}\ts{8}{1}{1:43:15}\ts{8}{1}{1:43:36}\ts{8}{1}{1:43:57}\ts{8}{1}{1:44:18}\ts{8}{1}{1:44:38}\ts{8}{1}{1:44:59}\ts{8}{1}{1:48:26}\ts{8}{1}{1:48:47}\ts{8}{1}{1:49:07}\ts{8}{1}{1:49:28}.
\end{proposition}
\begin{proof}
\emph{Co-vector.} The components are $\tilde\omega_a=\omega(\tilde e_a)$ (applying $\omega$ to the basis vector defines the components: abstract object, plug in the bases, get the components). Using $\tilde e_a=A^m{}_a\,e_m$, linearity (pull out the numbers and the invisible sum), $\tilde\omega_a=A^m{}_a\,\omega(e_m)=A^m{}_a\,\omega_m$
\ts{8}{1}{1:42:54}\ts{8}{1}{1:43:15}\ts{8}{1}{1:43:36}\ts{8}{1}{1:43:57}\ts{8}{1}{1:44:18}.

\emph{Vector.} The $a$-th component of $v$ is the vector applied to the $a$-th co-vector of the dual basis, $v^a=\varepsilon^a(v)$ (the vector considered as a $(1,0)$ tensor eating a co-vector: the identification of $V$ with $V^{**}$ shown above for finite dimension). After a false start (see the coursenote) the lecturer expands $v$ in the
\emph{new} basis, $v=\tilde v^b\tilde e_b$, and expresses the new basis vectors by the old ones, $\tilde e_b=A^m{}_b\,e_m$; then, using the dual basis relation $\varepsilon^a(e_m)=\delta^a_m$ (this is where the dual basis helps) and linearity of $\varepsilon^a$:
\[
  v^a=\varepsilon^a(v)=\tilde v^b\,A^m{}_b\,\varepsilon^a(e_m)=A^a{}_b\,\tilde v^b,
\]
that is $v=A\tilde v$ in components, so $\tilde v^a=B^a{}_m\,v^m$: ``it goes with the opposite matrix, with the inverse of this matrix'' \ts{8}{1}{1:44:59}\ts{8}{1}{1:45:20}\ts{8}{1}{1:45:40}\ts{8}{1}{1:46:01}\ts{8}{1}{1:46:21}\ts{8}{1}{1:46:42}\ts{8}{1}{1:47:03}\ts{8}{1}{1:47:24}\ts{8}{1}{1:47:44}\ts{8}{1}{1:48:05}\ts{8}{1}{1:48:26}.

\emph{Tensors.} Same by multilinearity, slot by slot \ts{8}{1}{1:48:26}\ts{8}{1}{1:48:47}.
\end{proof}
\begin{coursenote}[The first attempt at the vector rule]
At \ts{8}{1}{1:44:59}--\ts{8}{1}{1:47:44} the derivation for the vector components goes wrong at first (an expansion in the wrong basis leads to $\varepsilon^a(\tilde e_m)$, ``I cannot expand that, something is going wrong here''); it is corrected by expanding $v$ in the twiddled basis, as above. The lecturer's names for the matrices
($A$, $B$, the placement of the indices) vary in the transcript and I fixed them as above [?]. The final rule is as stated in class: vectors transform with the inverse matrix with respect to co-vectors.
\end{coursenote}
The matrix $B$ is the inverse of $A$ as a linear map $V\to V$ (``well defined, we know what that is''); $A^a{}_b$ and $B^a{}_b$ are the components of $A$ and $B$ as endomorphisms; ``the inverse matrix'' either uses the matrix multiplication witchcraft or the notion of inverse of a linear map, which is free of coordinates: the relation $A^a{}_mB^m{}_b=\delta^a_b$ \emph{follows} from
$B$ being the inverse, it is not the definition \ts{8}{1}{1:49:28}\ts{8}{1}{1:49:48}\ts{8}{1}{1:50:09}\ts{8}{1}{1:50:29}. This is the transformation rule for a tensor under a change of basis; one has to see where the twiddles, the $A$'s and the $B$'s are, and ``it is best every time you need this, if you want to know explicitly how an upstairs index
transforms, to quickly rederive it, otherwise you get probably confused'' \ts{8}{1}{1:50:29}\ts{8}{1}{1:50:49}\ts{8}{1}{1:51:10}.

A student asks how one goes ``from here to there'' in the step of the vector derivation where $v$ is treated as a $(1,0)$ tensor eating a co-vector \ts{8}{1}{1:51:10}\ts{8}{1}{1:51:31}. The answer:
\begin{intuition}[The $(1,0)$ tensors are not immediately the vectors]\label{int:te-10}
It is immediate that the $(0,1)$ tensors are $V^*$ (all linear maps $V\to K$, by definition). That the $(1,0)$ tensors are $V$ is \emph{not} immediate, and holds only for finite dimension: an element of $T^1_0V$ is a linear map $V^*\to K$, i.e.\ an element of $V^{**}$; to understand a vector as acting on a co-vector one needs
the vector space isomorphism $V\to V^{**}$ of the problem sheet (\cref{thm:te-dd,ex:te-dd}); once one interprets $v\in V$ as an element of $V^{**}$ via this isomorphism, it can act on a co-vector: it is the co-vector acting on the vector, just the other way around. In physics this subtlety is usually not mentioned (``usually this is just presented as: let us say this acts like this'')
\ts{8}{1}{1:51:31}\ts{8}{1}{1:51:55}\ts{8}{1}{1:52:15}\ts{8}{1}{1:52:36}\ts{8}{1}{1:52:57}\ts{8}{1}{1:53:17}\ts{8}{1}{1:53:38}\ts{8}{1}{1:53:59}\ts{8}{1}{1:54:20}\ts{8}{1}{1:54:40}\ts{8}{1}{1:55:01}.
\end{intuition}
The lecturer thanks the student: since the construction of the isomorphism was outsourced to the problem sheet, it could not be understood technically from what was written, ``but it is outsourced there precisely because it is so useful, so you should deal with that'' (a sub-question on the sheet will make it explicit); the construction is not hard, but needs a little more study than the
immediate isomorphism of the $(0,1)$ tensors with the co-vectors \ts{8}{1}{1:53:59}\ts{8}{1}{1:54:20}.

% ==================================================================
\section{Forms, volume and the determinant}\label{sec:te-det}

Something that will be used and that often comes with a lot of confusion: the \textbf{determinant} \ts{8}{1}{1:55:01}\ts{8}{1}{1:55:22}. Trust the order of presentation: first without the basis, then the basis to get the dimension, then back to the basis-free notation, and again introducing components in that precise order; the whole subject must be understood in this order. If one thinks from
the start of a vector as a collection of numbers, its components, one is already on the wrong track: one has to think abstractly and only after introducing a basis is one allowed to do so \ts{8}{1}{1:55:22}\ts{8}{1}{1:55:42}\ts{8}{1}{1:56:02}\ts{8}{1}{1:56:23}.
In ``witchcraft mathematics'' one takes an array of numbers and calculates the determinant according to some mysterious rule, then shows with heavy work that exchanging two rows or two columns gives a minus sign, ``a very mysterious rule acting on a cemetery'' \ts{8}{1}{1:56:23}\ts{8}{1}{1:56:43}.

\paragraph{Two arrays that look alike.} The array also hides a fact: if $f$ is a $(1,1)$ tensor with components $f^a{}_b$ one may arrange them in rows and columns; a $(0,2)$ tensor $g$ has components $g_{ab}$ with both indices down (it eats two vectors, in particular two basis vectors, which have the index down), and these too can be arranged in an array. But an endomorphism
(a $(1,1)$ tensor) and a bilinear form (a $(0,2)$ tensor) are \emph{entirely different beasts}: $f^a{}_b$ has one index up and one down and transforms with an $A$ and a $B$ (the up index with one, the down index with the other) under a change of basis; $g_{ab}$ transforms with two $A$'s. In mathematics one writes such transformations as
$A^{-1}fA$ for the endomorphism and $A^{\mathsf T}gA$ for the bilinear map (``the transpose, as shown before, a concept from witchcraft mathematics''): two totally different objects transform totally differently. If one does not change bases one may delude oneself into thinking that both are matrices and that what can be done with one
can be done with the other, e.g.\ define the determinant ``according to some witchcraft rule'' also for $g_{ab}$: no; \emph{determinants are only defined for endomorphisms}, as seen when done basis free
\ts{8}{1}{1:56:43}\ts{8}{1}{1:57:04}\ts{8}{1}{1:57:26}\ts{8}{1}{1:57:46}\ts{8}{1}{1:58:07}\ts{8}{1}{1:58:28}\ts{8}{1}{1:58:48}\ts{8}{1}{1:59:11}\ts{8}{1}{1:59:31}\ts{8}{1}{1:59:51}\ts{8}{1}{2:00:13}\ts{8}{1}{2:00:33}.

\begin{definition}[New concepts: $n$-form]\label{def:te-form}
Let $V$ be a finite dimensional vector space, $\dim V=d$. An \textbf{$n$-form} is a $(0,n)$ tensor $\omega$ that is \emph{totally anti-symmetric}: for $0<n\le d$,
\[
  \omega(v_1,\dots,v_n)=\operatorname{sgn}(\pi)\;\omega\bigl(v_{\pi(1)},\dots,v_{\pi(n)}\bigr)\qquad\text{for every permutation }\pi\in S_n,
\]
where $S_n$ is the permutation group (``not the sphere'') and $\operatorname{sgn}(\pi)=+1$ for an even and $-1$ for an odd permutation; more simply: if one exchanges any two arguments one picks up a minus sign, exchanging the same or another two picks up yet another minus sign. For $n=0$: a $(0,0)$ tensor is by convention
the underlying number field, a scalar, so a $0$-form is just a number; for $n=1$ a $(0,1)$ tensor the condition has no effect (nothing to exchange), so every co-vector is a $1$-form. $n$ must lie between $0$ and the dimension of the vector space \ts{8}{1}{2:00:33}\ts{8}{1}{2:00:53}\ts{8}{1}{2:01:16}\ts{8}{1}{2:01:38}\ts{8}{1}{2:02:01}\ts{8}{1}{2:02:23}\ts{8}{1}{2:02:44}\ts{8}{1}{2:03:05}\ts{8}{1}{2:03:25}\ts{8}{1}{2:03:46}\ts{8}{1}{2:04:06}.
\end{definition}
(``This anti-symmetry is wildly important.'' For $n>d$ there are anti-symmetric $(0,n)$ tensors in the sense of the definition, but they are zero, since the arguments are then linearly dependent and the anti-symmetry kills everything; so they are not considered \ts{8}{1}{2:04:27}\ts{8}{1}{2:04:47}\ts{8}{1}{2:05:08}.)

\begin{definition}[New concepts: top form, volume form]\label{def:te-top}
A \textbf{top form} on $V$ is a form of the top rank, with as many entries as the vector space has dimensions: a $d$-form \ts{8}{1}{2:04:06}\ts{8}{1}{2:04:27}\ts{8}{1}{2:04:47}\ts{8}{1}{2:05:08}\ts{8}{1}{2:05:28}.
\end{definition}
\begin{proposition}[Top forms are unique up to a factor]\label{prop:te-topuniq}
If $\omega$ and $\omega'$ are two non-vanishing top forms, there is a number $c\in K$ such that $\omega'=c\,\omega$: essentially there is, up to a scalar multiple, only one top form on a vector space \ts{8}{1}{2:05:28}\ts{8}{1}{2:05:49}.
\end{proposition}
\begin{proof}
``Very easy to see because it is totally anti-symmetric: it can only differ by a factor'' \ts{8}{1}{2:05:49}\ts{8}{1}{2:06:10}. Explicitly (not said in class): expanding $d$ vectors $v_i=v_i^{\,a}e_a$ in a basis and using multilinearity and anti-symmetry, a top form is determined by its value on $(e_1,\dots,e_d)$, namely
$\omega(v_1,\dots,v_d)=\bigl(\sum_{\pi\in S_d}\operatorname{sgn}(\pi)\,v_1^{\pi(1)}\cdots v_d^{\pi(d)}\bigr)\,\omega(e_1,\dots,e_d)$, so two top forms are proportional.
\end{proof}
The choice of a top form $\omega$ is a \textbf{choice of volume form} (``choice of volume'') on the vector space: ``the choice is not big, one can choose one, another one is five times the one you chose''. A vector space together with a choice of top form can be called a vector space \emph{with volume}, an additional structure one chooses, all of whose choice lies in the factor
\ts{8}{1}{2:06:10}\ts{8}{1}{2:06:31}\ts{8}{1}{2:06:55}\ts{8}{1}{2:08:39}\ts{8}{1}{2:08:59}.
\begin{definition}[New concepts: volume spanned by vectors]\label{def:te-vol}
Let $v_1,\dots,v_d$ be $d$ vectors in a $d$-dimensional vector space $V$ with a choice of volume form $\omega$. The \textbf{volume spanned} by them (the order matters) is
\[
  \mathrm{vol}(v_1,\dots,v_d):=\omega(v_1,\dots,v_d)
\]
\ts{8}{1}{2:06:55}\ts{8}{1}{2:07:16}\ts{8}{1}{2:07:36}.
\end{definition}
In a three-dimensional space with three vectors, plugging them into the volume form gives a volume; if the three are linearly dependent, if they lie in one plane, the anti-symmetry kills it: the spanned volume is zero; they must be linearly independent to span a non-zero volume, and the same in higher dimensions \ts{8}{1}{2:07:36}\ts{8}{1}{2:07:57}\ts{8}{1}{2:08:18}.
One might have thought that volume needs an inner product, because one needs an inner product to have angles and lengths; but \emph{volume is the choice of a top form}. This will come back later: it is one place where one needs the volume form \ts{8}{1}{2:08:18}\ts{8}{1}{2:08:39}.

\begin{definition}[New concepts: determinant of an endomorphism]\label{def:te-det}
There is no determinant on a vector space without a volume form, ``but it is not that bad'' \ts{8}{1}{2:08:59}\ts{8}{1}{2:09:24}. Let $\phi\in\mathrm{End}(V)$, $V$ finite dimensional of dimension $d$ (emphasised: an endomorphism, a $(1,1)$ tensor, a linear map $V\to V$). Its \textbf{determinant} is the number
\[
  \det\phi:=\frac{\omega\bigl(\phi(e_1),\dots,\phi(e_d)\bigr)}{\omega(e_1,\dots,e_d)}
\]
for some top form $\omega\ne0$ and some basis $e_1,\dots,e_d$ of $V$; ``it is not the determinant of the components, it is a determinant of the abstract object'' \ts{8}{1}{2:09:24}\ts{8}{1}{2:09:44}\ts{8}{1}{2:10:05}\ts{8}{1}{2:10:27}.
\end{definition}
\begin{proposition}[Well-definedness of the determinant]\label{prop:te-detwd}
$\det\phi$ does not depend on the choice of $\omega$ nor on the choice of the basis \ts{8}{1}{2:10:27}\ts{8}{1}{2:10:49}\ts{8}{1}{2:11:12}.
\end{proposition}
\begin{proof}
\emph{Independence of $\omega$} (immediately seen, as the lecturer asks the class): another top form gives a factor $c$ (\cref{prop:te-topuniq}) in the numerator and the same factor $c$ in the denominator, which cancel \ts{8}{1}{2:11:12}\ts{8}{1}{2:11:35}.
\emph{Independence of the basis:} ``you show by plugging in the basis change we had before: you plug it in here, you plug it in there, you pull it out and so on, and you find it is independent of that too, because essentially you get the $\omega$ of the basis change matrix here, and those factors cancel as well; check, not going to do this right here''
\ts{8}{1}{2:11:35}\ts{8}{1}{2:11:55}\ts{8}{1}{2:12:18}.
Explicitly (not said in class): for a new basis $\tilde e_a=A^b{}_a\,e_b$, multilinearity and anti-symmetry give $\omega(\tilde e_1,\dots,\tilde e_d)=\kappa(A)\,\omega(e_1,\dots,e_d)$ with $\kappa(A)=\sum_\pi\operatorname{sgn}(\pi)A^{\pi(1)}{}_1\cdots A^{\pi(d)}{}_d$, and by linearity of $\phi$ likewise
$\omega(\phi(\tilde e_1),\dots,\phi(\tilde e_d))=\kappa(A)\,\omega(\phi(e_1),\dots,\phi(e_d))$; the common factor cancels in the quotient.
\end{proof}
One sees why an endomorphism is needed: the vector $\phi(e_a)$ must again be a vector, to be fed to $\omega$ \ts{8}{1}{2:12:18}.

\paragraph{Reading off the standard rules.} (i) The determinant of $\phi$ coincides with the determinant according to the witchcraft rule applied to the array of its components: ``what you learned is correct, the determinant rule will coincide'' \ts{8}{1}{2:12:18}\ts{8}{1}{2:12:39}. (ii) To what power of the components of $\phi$ does the determinant contribute: each $\phi(e_a)$ gives one
component and $\omega$ eats $d$ of them, hence the products of $d$ components that appear in the witchcraft rule (for a $2\times2$ matrix $ad-bc$, for $3\times3$ ``$a\,d\,m$ minus the other products'' --- the power is the dimension of the vector space) \ts{8}{1}{2:12:39}\ts{8}{1}{2:13:00}\ts{8}{1}{2:13:20}\ts{8}{1}{2:13:41}\ts{8}{1}{2:14:02}.
(iii) Exchanging two rows or two columns gives a minus sign: this can be read off from the definition since $\omega$ is totally anti-symmetric (exchanging in the argument list of $\omega$ the two vectors $\phi(e_a)$ gives a minus sign); all these rules follow directly from the abstract definition \ts{8}{1}{2:14:02}\ts{8}{1}{2:14:23}\ts{8}{1}{2:14:43}\ts{8}{1}{2:15:04}.
(iv) Under a change of basis nothing happens, as already said (independence of the basis): in witchcraft rules $\det(A^{-1}\hat\phi A)=\det\hat\phi$ \ts{8}{1}{2:15:04}\ts{8}{1}{2:15:24}. (v) $\det(\phi\circ\psi)=\det\phi\cdot\det\psi$, the product being in $K$: this is not a matrix product but a composition, the determinant eating an endomorphism, totally basis free; it
follows from the definition by looking at compositions \ts{8}{1}{2:15:24}\ts{8}{1}{2:15:44}\ts{8}{1}{2:16:05}\ts{8}{1}{2:16:26}. ``All the standard rules come up in this way.''

\begin{secondary}[The witchcraft check of independence of the basis]
The lecture also shows the check in matrix language: under a change of basis the representing array changes as $\hat\phi\mapsto A^{-1}\hat\phi A$; the determinant of the product is the product of the determinants, these are numbers, so one can push $\det A^{-1}$ next to $\det A$ and use the rule backwards: $\det A^{-1}\cdot\det A=\det(A^{-1}A)=\det\mathbb 1=1$ (the unit endomorphism, which also gives $1$ from the abstract
definition); so $\det\hat\phi$ does not depend on the basis, ``they better not, we defined them this way''. Not a big result, an important consistency check \ts{8}{1}{2:17:28}\ts{8}{1}{2:17:48}\ts{8}{1}{2:18:08}\ts{8}{1}{2:18:29}\ts{8}{1}{2:18:50}.
\end{secondary}

\paragraph{The determinant of a bilinear form.} In physics, and also in mathematics, one has the opportunity to use the determinant of a bilinear form; this will also be on the problem sheet (\cref{exam:ch08}). As defined the determinant can only be defined for an endomorphism, as one needs to produce vectors again
\ts{8}{1}{2:16:26}\ts{8}{1}{2:16:47}\ts{8}{1}{2:17:07}\ts{8}{1}{2:17:28}. The witchcraft version, to be replaced by a better one on the problem sheet, is to take the determinant of the array $g_{ab}$ of a $(0,2)$ tensor. This cannot coincide with the proper definition, because the object needs to be an endomorphism: it is ``deeply illegal'' from an abstract point of view
\ts{8}{1}{2:18:50}\ts{8}{1}{2:19:10}\ts{8}{1}{2:19:32}.
\begin{proposition}[The determinant of a bilinear form is not well defined]\label{prop:te-detg}
Under a change of basis the components of a $(0,2)$ tensor $g$ transform as $\tilde g=A^{\mathsf T}gA$ (arrays), so that, with the witchcraft determinant,
\[
  \det\tilde g=(\det A)^2\,\det g,
\]
and the determinant of a bilinear map is \emph{not} independent of the choice of basis: it picks up the factor $(\det A)^2$, the square of the determinant of the change of basis matrix
\ts{8}{1}{2:19:32}\ts{8}{1}{2:19:53}\ts{8}{1}{2:20:13}\ts{8}{1}{2:20:34}\ts{8}{1}{2:20:55}.
\end{proposition}
\begin{proof}
$\det(A^{\mathsf T}gA)=\det(A^{\mathsf T})\det g\det A$ and $\det A^{\mathsf T}=\det A$ \ts{8}{1}{2:19:53}\ts{8}{1}{2:20:13}\ts{8}{1}{2:20:34}.
\end{proof}
Hence one would not use it; but it turns out that one does something ill defined from an abstract point of view, and in other places there will be factors coming up in the denominator which are also not well defined under a change of basis, but the \emph{product} of something not well defined and something picking up the opposite factor is independent of the
change of basis: one starts considering these. ``You commit a mistake and later on the opposite mistake and the mistakes cancel: how can we cast this into proper mathematics?'' The answer: by considering \emph{bundles} \ts{8}{1}{2:20:55}\ts{8}{1}{2:21:16}\ts{8}{1}{2:21:37}.
\begin{application}[Principal bundles and tensor densities]\label{app:te-density}
One considers basis changes by taking the base manifold of interest and attaching at every point, as a fibre, the transformation group of the change of bases: this is called a \textbf{principal fibre bundle}; with it one will be able to give a \emph{bundle definition} of what a tensor is, and of what a \textbf{tensor density} is: the density picks up such a factor. All this can be done without hand-waving in the bundle formalism, ``and
that is one of the big reasons why we need bundles'' \ts{8}{1}{2:21:37}\ts{8}{1}{2:21:59}\ts{8}{1}{2:22:19}\ts{8}{1}{2:22:40}. (Principal bundles are treated in \cref{ch:l19}.)
\end{application}

% ==================================================================
\section*{Key concepts}
\begin{keyconcepts}
\textbf{Retrieval questions.}
\begin{enumerate}
  \item State the definition of a field and of a ring. Which axioms does a ring lose, and which example (a) is a ring but not a field, (b) is a non-commutative ring?
  \item Define a $K$-vector space and a linear map. Why can $\mathrm{Hom}(V,W)$ be made into a vector space, and why does the construction fail for modules over a non-commutative ring? Why does linearity force $V$ and $W$ to be over the same field?
  \item Define $\mathrm{End}(V)$, $\mathrm{Aut}(V)$ and $V^*$. Why is $\mathrm{Aut}(V)$ not a subspace?
  \item Define a $(p,q)$ tensor, the tensor space $T^p_qV$ and the tensor product. What are the two meanings of the symbol $\otimes$?
  \item Prove that $T^1_1V\cong\mathrm{End}(V^*)$. Which statements about $T^1_0V$ and $T^1_1V$ hold only in finite dimension?
  \item Define a Hamel basis and the dimension. Why only finite subsets? Why are sines and cosines not a Hamel basis?
  \item Define the components of a tensor, the dual basis, and state the reconstruction theorem.
  \item State the summation convention and the up/down convention. Why does the summation convention need linearity?
  \item How do the components of a vector, a co-vector and a $(p,q)$ tensor change under a change of basis?
  \item Define $n$-forms, top forms and the volume spanned by $d$ vectors; define $\det\phi$ for an endomorphism and show it is well defined. Why is the determinant of a bilinear form not well defined?
\end{enumerate}
\textbf{Concepts.}
\begin{itemize}
  \item Field, ring, vector space, linear map (\cref{def:te-field,def:te-ring,def:te-vs,def:te-linear}); $\mathrm{Hom}$, $\mathrm{End}$, $\mathrm{Aut}$, dual (\cref{def:te-hom,def:te-end}).
  \item $(p,q)$ tensor, tensor space, tensor product (\cref{def:te-tensor,def:te-tspace,def:te-tensorprod,rem:te-two-otimes}); $T^1_1V\cong\mathrm{End}(V^*)$ (\cref{prop:te-11}).
  \item Hamel basis and dimension, double dual (\cref{def:te-basis,thm:te-dd}).
  \item Components, dual basis, reconstruction (\cref{def:te-components,def:te-dualbasis,thm:te-recon}).
  \item Change of basis, summation and up/down conventions, transformation of components (\cref{def:te-change,def:te-einstein,def:te-updown,prop:te-transform}).
  \item Forms, top forms, volume, determinant (\cref{def:te-form,def:te-top,def:te-vol,def:te-det,prop:te-detwd,prop:te-detg}).
\end{itemize}
\end{keyconcepts}

\flushside
\end{paracol}

% !TEX root = ../main.tex
\chapter{Differential structures: the pivotal concept of tangent vector spaces}\label{ch:l09}
\begin{paracol}{2}

\begin{sources}
\begin{tabularx}{\linewidth}{@{}l l L@{}}
\toprule
\textbf{Video} & \textbf{Length} & \textbf{Content}\\
\midrule
\seg{9}{1}{Lecture 9: Differential structures: the pivotal concept of tangent vector spaces} & 1:44:15 & The vector space $C^\infty(M)$; tangent vectors as directional derivative operators along curves;
the tangent space $T_pM$ and the proof that addition and scalar multiplication close; an aside on algebras, derivations, Lie algebras and the Poisson bracket; the theorem $\dim T_pM=\dim M$ with the chart-induced basis
$\partial/\partial x^a$; the difference between $\partial_a$ and $\partial/\partial x^a$; legality of acting on coordinate functions; components, change of chart and why there are no position vectors.\\
\bottomrule
\end{tabularx}

\smallskip
\textbf{Transcript} \file{materials/transcripts/C09-V01.txt} (YouTube automatic captions).
\end{sources}

\begin{reading}[Where this lecture sits in the other notes]
\emph{Gravity and Light}, Lecture~5 (Tangent spaces); \emph{Theoretische Mechanik}, Lecture~4. Both construct the tangent space of a smooth manifold; the aside on algebras and derivations, the careful distinction between the partial
derivative $\partial_a$ on $\R^D$ and the tangent vector $\partial/\partial x^a$, and the remarks on position vectors are, as far as the titles say, emphasised in this course.
(From the lecture titles and the plan of the course; not checked in detail.)
\end{reading}

% ==================================================================
\section{The vector space of smooth functions}\label{sec:tv-cinf}

Today the tangent spaces to a manifold: the definitions made here are ``probably the single most important definitions in differential geometry'', that of a tangent vector to a manifold and everything built from it \ts{9}{1}{0:09}. Having reviewed vector spaces last time
(\cref{ch:l08}), including infinite dimensional ones, one makes use of that now: the tangent space will be a finite dimensional vector space if the manifold is finite dimensional, with the same dimension as the manifold has as a topological manifold. But first an \emph{infinite
dimensional} vector space is constructed from a smooth manifold, an object that will return again and again \ts{9}{1}{0:30} \ts{9}{1}{0:51}.

Let $M$ be a smooth manifold. Throughout this chapter ``smooth manifold'' comes with a topology and a smooth atlas (\cref{def:df-manifold}); this is suppressed in the notation, and made explicit again only when several manifolds with several topologies and atlases are compared. ``From now on we are grown up'' \ts{9}{1}{1:11} \ts{9}{1}{1:32}.

\begin{definition}[New concepts: the vector space $C^\infty(M)$]\label{def:tv-cinf}
Let $M$ be a smooth manifold. The set
\[
  C^\infty(M):=\{f\colon M\to\R\mid f\text{ smooth}\}
\]
is a vector space over $\R$ (the manifolds are real) with the \textbf{pointwise} operations: for $f,g\in C^\infty(M)$ and $\lambda\in\R$
\[
  (f+g)(p):=f(p)+g(p),\qquad (\lambda\cdot f)(p):=\lambda\,f(p)\qquad\text{for all }p\in M,
\]
where on the right-hand sides are the addition and the multiplication of the real numbers, since $f$ and $g$ send points to real numbers \ts{9}{1}{1:52} \ts{9}{1}{2:13} \ts{9}{1}{2:54} \ts{9}{1}{3:36}.
\end{definition}
That this is a vector space is quickly checked (``pretty obvious''), and it is an infinite dimensional vector space which will play an important role \ts{9}{1}{3:56} \ts{9}{1}{4:17}.
(Smoothness of $f$ is meant in the sense of \cref{def:df-map}, with $\R$ carrying its single standard chart.)

% ==================================================================
\section{Tangent vectors as directional derivative operators}\label{sec:tv-vector}

The central concept follows \ts{9}{1}{4:17}. A curve is a map from the real line; to avoid worrying about intervals (``one could be very picky and say it is an interval from $a$ to $b$ and then one must worry about the domains''; one can always reparametrise) one is a little cavalier:
for the purposes here a curve has the full real line as parameter range; more precision will be added when it matters \ts{9}{1}{4:39} \ts{9}{1}{4:59} \ts{9}{1}{5:20}.

\begin{definition}[New concepts: tangent vector, directional derivative operator]\label{def:tv-vector}
Let $\gamma\colon\R\to M$ be a smooth curve through a point $p\in M$. Without loss of generality (shift the parameter) $\gamma(0)=p$; this could be dropped but makes the formulae simpler \ts{9}{1}{5:40} \ts{9}{1}{6:02} \ts{9}{1}{6:23}.
The \textbf{directional derivative operator at $p$ along $\gamma$} is the linear map
\[
  X_{\gamma,p}\colon C^\infty(M)\longrightarrow\R,\qquad f\longmapsto X_{\gamma,p}(f):=(f\circ\gamma)'(0).
\]
Here $f\circ\gamma\colon\R\to\R$ is a map from the real line to itself, for which the derivative is known (the prime), evaluated at the parameter value $0$ at which $\gamma$ passes through $p$: the result is a real number. The target $\R$ is regarded as a real vector space
over itself, $\R^1$, so that linearity makes sense on both sides. It is linear essentially because it is built with a derivative \ts{9}{1}{6:43} \ts{9}{1}{7:26} \ts{9}{1}{8:07} \ts{9}{1}{8:28} \ts{9}{1}{9:09} \ts{9}{1}{9:50}.
In differential geometry $X_{\gamma,p}$ is called the \textbf{tangent vector to the curve $\gamma$ at the point $p\in M$} \ts{9}{1}{10:11} \ts{9}{1}{10:32}.
\end{definition}
(The name \emph{tangent}: from Latin \emph{tangere}, to touch; not said in class.)
The terminology should not mislead one into thinking that this is anything like a collection of numbers; ``finally it will be'', but at this point the tangent vector, formally and abstractly, \emph{is} a directional derivative operator: tangent vectors are operators, derivative operators
acting on smooth functions. That is how the concept of a vector survives in differential geometry \ts{9}{1}{10:53} \ts{9}{1}{11:14}.

\begin{intuition}[Direction and derivative are married]\label{int:tv-married}
From undergraduate physics one has a direction $\vec X$ and the $\nabla$ operator, and the dot product $\vec X\cdot\nabla$ gives a directional derivative: in vector calculus one has an independent notion of direction and an independent notion of derivative and combines them. In
differential geometry the two concepts are married and \emph{cannot be separated}: there is only the combination. This is the key point in the technical, hence the practical, understanding of tangent vectors, and the consequences will be worked out today \ts{9}{1}{11:14} \ts{9}{1}{11:35} \ts{9}{1}{11:55} \ts{9}{1}{12:15}.
\end{intuition}

\begin{intuition}[$X_{\gamma,p}$ is the velocity of $\gamma$ at $p$]\label{int:tv-velocity}
A precise intuition: $X_{\gamma,p}$ is the \emph{velocity} of the curve $\gamma$ at $p$. Recall that tangent spaces are to be constructed not by referring to an outside space in which the manifold is embedded (\cref{int:df-embedding}), but intrinsically: given a smooth
manifold and a curve through a point, with no reference to the bulk around it, one of the tangent vectors at the point is the velocity of that curve, which is a vector, by definition, namely $X_{\gamma,p}$ \ts{9}{1}{12:38} \ts{9}{1}{12:58} \ts{9}{1}{13:20} \ts{9}{1}{13:40}.
\end{intuition}
\emph{Consistency check.} Let $\delta(\lambda):=\gamma(2\lambda)$, the curve that runs through the same points as $\gamma$ at double the speed (also $\delta(0)=p$). Acting on an arbitrary smooth $f$, the chain rule gives an inner derivative:
\[
  X_{\delta,p}(f)=(f\circ\delta)'(0)=2\,(f\circ\gamma)'(0)=2\,X_{\gamma,p}(f),
\]
so $X_{\delta,p}=2X_{\gamma,p}$, and any other factor (three times, twenty times $\pi$) pulls out in the same way. Running through the curve at double speed gives a vector twice as long: this is how a velocity must scale. This is only an intuition; the definition is the technical one above
\ts{9}{1}{14:01} \ts{9}{1}{14:22} \ts{9}{1}{14:43} \ts{9}{1}{15:03} \ts{9}{1}{15:23} \ts{9}{1}{15:44}.

% ==================================================================
\section{The tangent space $T_pM$ and why it is closed}\label{sec:tv-space}

\begin{definition}[New concepts: tangent space]\label{def:tv-space}
The \textbf{tangent vector space to $M$ at $p\in M$}, $T_pM$, has as underlying set
\[
  T_pM:=\{X_{\gamma,p}\mid\gamma\text{ a smooth curve through }p\},
\]
the set of \emph{all} tangent vectors at $p$ to all smooth curves through $p$. It is equipped with two operations $\oplus$ and $\odot$ (called so for the moment, since it is not yet clear that they close on the set), defined \textbf{pointwise} as maps into $\R$: for
$X_{\gamma,p},X_{\delta,p}\in T_pM$, $\lambda\in\R$ and $f\in C^\infty(M)$
\[
  (X_{\gamma,p}\oplus X_{\delta,p})(f):=X_{\gamma,p}(f)+X_{\delta,p}(f),\qquad(\lambda\odot X_{\gamma,p})(f):=\lambda\,X_{\gamma,p}(f),
\]
using on the right the addition and multiplication in $\R$ \ts{9}{1}{16:08} \ts{9}{1}{16:31} \ts{9}{1}{18:37} \ts{9}{1}{19:41} \ts{9}{1}{20:24} \ts{9}{1}{21:05}.
\end{definition}
The intuition: pick $p$ and a curve; look also at the same curve run through at double speed, at $\mu$ times the speed, backwards (a tangent vector in the opposite direction); then take any other smooth curve through $p$ at any speed (technically a different map
$\R\to M$): the collection of all the corresponding directional derivative operators is the set underlying the tangent space. ``It's very simple'' \ts{9}{1}{16:51} \ts{9}{1}{17:12} \ts{9}{1}{17:33} \ts{9}{1}{17:54}.

\paragraph{The question of closure.} The sum of two linear maps $C^\infty(M)\to\R$ is again a linear map $C^\infty(M)\to\R$; that is not in doubt. But an element of $T_pM$ is not just any linear map: it is one constructed in the specific way $f\mapsto(f\circ\sigma)'(0)$ from a curve.
So one must find, for the sum, a curve $\sigma$ with
\[
  X_{\gamma,p}(f)+X_{\delta,p}(f)=X_{\sigma,p}(f)\qquad\text{for all }f\in C^\infty(M);
\]
only then can one be confident that the addition closes on $T_pM$. The same for the multiplication by scalars \ts{9}{1}{21:46} \ts{9}{1}{22:06} \ts{9}{1}{22:27} \ts{9}{1}{22:48} \ts{9}{1}{23:08}.

\begin{proposition}[$\oplus$ and $\odot$ close on $T_pM$]\label{prop:tv-closure}
For $X_{\gamma,p},X_{\delta,p}\in T_pM$ and $\lambda\in\R$ the maps $X_{\gamma,p}\oplus X_{\delta,p}$ and $\lambda\odot X_{\gamma,p}$ are again elements of $T_pM$; hence $(T_pM,\oplus,\odot)$ is a real vector space \ts{9}{1}{23:30} \ts{9}{1}{23:52}.
\end{proposition}
\begin{proof}
\emph{The sum.} Employ a chart $(U,x)$ with $p\in U$ (the curves may be restricted to $\gamma^{-1}(U)$; ``let us not be so picky'') and define
\[
  \sigma(\lambda):=x^{-1}\Bigl(\bigl(x\circ\gamma\bigr)(\lambda)+\bigl(x\circ\delta\bigr)(\lambda)-x(p)\Bigr).
\]
This looks funny; first see whether it makes sense. The maps $x\circ\gamma$ and $x\circ\delta$ are the chart representations of the two curves, maps from (a part of) $\R$ to $\R^D$, $D=\dim M$; the third term is the constant map $\lambda\mapsto x(p)$. As maps into $\R^D$ they can be added
(addition in $\R^D$), and the term $-x(p)$ ensures that at $\lambda=0$ the resulting curve passes through $x(p)$: indeed
$\sigma(0)=x^{-1}(x(p)+x(p)-x(p))=p$ \ts{9}{1}{24:13} \ts{9}{1}{24:59} \ts{9}{1}{25:21} \ts{9}{1}{26:02} \ts{9}{1}{27:03} \ts{9}{1}{29:07}.
The picture: one wants a curve in the manifold whose tangent vector is the sum, but curves on the manifold cannot be added, since the manifold has not enough structure; so one goes to the chart, where the chart images of the two curves through $x(p)$ can be added, and lifts the result back up with $x^{-1}$
\ts{9}{1}{27:44} \ts{9}{1}{28:05} \ts{9}{1}{28:25} \ts{9}{1}{28:46}.

This seems to be ill-defined: the same curves could be represented in very different charts, and the sum curve in one chart looks very different from the sum in another. If one defined the sum of the curves as such it would be wildly chart dependent; the curve $\sigma$ certainly is. What saves the day is
that one is not interested in the whole sum curve but only in its derivative at $0$, and the derivative takes care of all that: the sum of the tangent vectors is well defined \ts{9}{1}{29:28} \ts{9}{1}{29:48} \ts{9}{1}{30:09} \ts{9}{1}{30:30}.

Write $\hat f:=f\circ x^{-1}\colon\R^D\to\R$ for the chart representative of $f$. Then
\[\begin{aligned}
X_{\sigma,p}(f)&=(f\circ\sigma)'(0)\\
&=\bigl(\hat f\circ(x\circ\gamma+x\circ\delta-x(p))\bigr)'(0)\\
&=\partial_a\hat f(x(p))\,\bigl[(x^a\circ\gamma)'(0)\\
&\qquad+(x^a\circ\delta)'(0)\bigr]\\
&=(\hat f\circ x\circ\gamma)'(0)+(\hat f\circ x\circ\delta)'(0)\\
&=X_{\gamma,p}(f)+X_{\delta,p}(f).
\end{aligned}\]
The third line is the multi-dimensional chain rule ($\hat f$ is the chart representative of $f$), applied to $x\circ\gamma+x\circ\delta-x(p)\colon\R\to\R^D$; the derivative of the constant is zero. The point at which $\partial_a(f\circ x^{-1})$ is evaluated is $x(\sigma(0))=x(p)$. The fifth line
uses the chain rule backwards, together with $x^{-1}\circ x=\mathrm{id}$ and the associativity of the composition. So $X_{\gamma,p}\oplus X_{\delta,p}=X_{\sigma,p}\in T_pM$: ``we pull this stunt of going to a chart and constructing a sum of curves, something terribly ill-defined; what is not ill-defined is its first derivative,
independent of the chart because the chart map cancels out again. That's the trick'' \ts{9}{1}{30:51} \ts{9}{1}{31:53} \ts{9}{1}{33:18} \ts{9}{1}{34:20} \ts{9}{1}{36:06} \ts{9}{1}{37:07}.

\emph{The scalar multiplication} is much simpler: one does the same with $\lambda$ times $x\circ\gamma$ and ``does not need all the rest'' \ts{9}{1}{37:07} \ts{9}{1}{37:28}.
\end{proof}
\begin{coursenote}[Details left out in class]
(i) The curve $\sigma$ is defined only for those $\lambda$ for which $x\circ\gamma(\lambda)+x\circ\delta(\lambda)-x(p)$ lies in $x(U)$; this is a neighbourhood of $0$, which is all that the derivative at $0$ needs (and one may reparametrise to the full line). (ii) For the scalar multiplication the cleanest curve is the
reparametrisation $\tau(\lambda):=\gamma(\mu\lambda)$ with $\mu$ the scalar: $X_{\tau,p}(f)=\mu\,(f\circ\gamma)'(0)$, exactly as in the doubling computation above (\cref{int:tv-velocity}); the lecture's version is $x^{-1}\bigl(x(p)+\mu\,(x\circ\gamma-x(p))\bigr)$,
which gives the same vector. (iii) The zero vector is the tangent vector of the constant curve. Not said in class.
\end{coursenote}
Looking back at the definition of $T_pM$ one sees that it is a very simple one, once the tangent vectors are accepted to be these linear maps: addition and scalar multiplication are those of maps, and one has shown that they close. Now that one knows, one never needs to worry again:
for any sum or scalar multiple there is an underlying curve. Note also that there are many curves producing the same tangent vector, for instance different curves through $p$ with the same speed give rise, although they are different curves, to the same tangent vector; it suffices to know that some curve exists
\ts{9}{1}{37:49} \ts{9}{1}{38:10} \ts{9}{1}{45:53}.

% ==================================================================
\section{Aside: algebras and derivations}\label{sec:tv-derivation}

Before looking at the properties of the tangent vectors, a little aside on \emph{algebras and derivations} (a technical term; not the same as derivatives) \ts{9}{1}{38:10} \ts{9}{1}{38:32}.

\begin{definition}[New concepts: algebra]\label{def:tv-algebra}
A vector space $(V,+,\cdot)$ over a field $K$ equipped with a \textbf{product}, a bilinear map
\[
  \bullet\colon V\times V\to V,
\]
is called an \textbf{algebra}; so it is a quadruple, where besides the scalar multiplication with respect to the underlying field one has a product between the elements of the vector space itself \ts{9}{1}{38:53} \ts{9}{1}{39:14} \ts{9}{1}{39:36} \ts{9}{1}{39:57}.
\end{definition}
The product can take various forms (``we will meet several in the future''); in particular, the \emph{Lie product} will be an antisymmetric one; no symmetry is assumed, only bilinearity. An example from high school: the cross product on $\R^3$ makes $\R^3$ an algebra (``you can't do this on $\R^4$'')
\ts{9}{1}{40:18} \ts{9}{1}{40:39}.
\begin{coursenote}[On the cross product in $\R^4$]
The lecture says that one cannot do this on $\R^4$. A bilinear antisymmetric product certainly exists on every vector space; what exists only in dimensions $3$ and $7$ is a cross product with the properties of the one on $\R^3$ (orthogonal to both factors, norm $|a||b|\sin\theta$). I cannot confirm that
this is what is meant; [?]. Not said in class.
\end{coursenote}
\begin{example}[$C^\infty(M)$ is an algebra]\label{ex:tv-cinf-alg}
On $C^\infty(M)$ define $\bullet\colon C^\infty(M)\times C^\infty(M)\to C^\infty(M)$ pointwise by $(f\bullet g)(p):=f(p)\,g(p)$, the multiplication of the real numbers. This is the multiplication of two functions; whereas the scalar multiplication is only the multiplication with a fixed constant: ``it's structurally totally
different, very important to realise this''. So $C^\infty(M)$ is not only a vector space but, with this product, an $\R$ \emph{algebra}. It seems formal overkill, since one always knew how to multiply two functions, but the point is to realise the structural level of the whole thing
\ts{9}{1}{40:39} \ts{9}{1}{41:00} \ts{9}{1}{41:21} \ts{9}{1}{41:41} \ts{9}{1}{42:02}.
\end{example}

\begin{definition}[New concepts: derivation]\label{def:tv-derivation}
Let $(A,+,\cdot,\bullet_A)$ and $(B,+,\cdot,\bullet_B)$ be algebras. A \textbf{derivation} is a linear map $D\colon A\to B$ which in addition satisfies the \textbf{Leibniz rule}
\[
  D(f\bullet_A g)=D(f)\bullet_B g+f\bullet_B D(g)\qquad\text{for all }f,g\in A.
\]
Linearity refers to the vector space underlying the algebras; the Leibniz rule is the compatibility with the products \ts{9}{1}{42:22} \ts{9}{1}{42:46} \ts{9}{1}{43:07} \ts{9}{1}{43:49} \ts{9}{1}{44:09} \ts{9}{1}{46:57}.
\end{definition}
(In class the definition was first written for a single algebra and then generalised, after some hesitation, to a map between two algebras.) The Leibniz rule terribly resembles the product rule of differentiation, and that is no accident: ordinary differentiation is a derivation, linear (the sum rule) and satisfying the Leibniz rule (the
product rule); the abstract version works in far more situations than the one used here \ts{9}{1}{44:30} \ts{9}{1}{44:51} \ts{9}{1}{45:11}.

\begin{example}[Tangent vectors are derivations]\label{ex:tv-der-tangent}
Take $A=C^\infty(M)$ with the pointwise product and $B=\R$ as an algebra over itself (all products are the multiplication in $\R$). Every $X\in T_pM$ is a derivation $C^\infty(M)\to\R$:
\[
  X(f\bullet g)=X(f)\,g(p)+f(p)\,X(g).
\]
One does not need to say to which curve and at which point $X$ corresponds: being in $T_pM$ it comes from some curve through $p$. This is quickly checked and goes to the problem sheet (\cref{ex:tv-leibniz}). So a tangent vector at $p$ is a particular example of a derivation
\ts{9}{1}{45:11} \ts{9}{1}{45:32} \ts{9}{1}{46:57} \ts{9}{1}{47:18} \ts{9}{1}{47:38} \ts{9}{1}{47:58}.
\end{example}
\begin{exercise}[Tangent vectors satisfy the Leibniz rule]\label{ex:tv-leibniz}
Show that $X_{\gamma,p}(f\bullet g)=X_{\gamma,p}(f)\,g(p)+f(p)\,X_{\gamma,p}(g)$ using the product rule for $(f\circ\gamma)\cdot(g\circ\gamma)\colon\R\to\R$. (Problem sheet.)
\end{exercise}
One may say this abstract view is heavy machinery for something simple; but ``we'll quite soon meet another type of algebra'' \ts{9}{1}{47:58} \ts{9}{1}{48:19}.

\begin{example}[The Lie algebra $\mathrm{End}(V)$ and the commutator as a derivation]\label{ex:tv-lie}
(After a false start with matrices: ``nonsense''.) Take $A=\mathrm{End}(V)$ for a vector space $V$ (\cref{def:te-end}); this is a vector space, not yet an algebra. Define the product
\[
  [\ ,\ ]\colon A\times A\to A,\qquad[\varphi,\psi]:=\varphi\circ\psi-\psi\circ\varphi,
\]
the \textbf{commutator}, the difference of the two compositions, quite abstractly, no matrices or matrix multiplication needed; it is again an endomorphism and it is bilinear. It is an \emph{antisymmetric} product (the bilinear map was not required to be symmetric), so $\mathrm{End}(V)$ is an algebra; one calls it a
\textbf{Lie algebra}, since the product satisfies the \textbf{Jacobi identity}
\[
  [\varphi,[\psi,\rho]]+[\psi,[\rho,\varphi]]+[\rho,[\varphi,\psi]]=0,
\]
``quickly calculated'' and following from the definition. In general any algebra whose product satisfies the Jacobi identity (and is antisymmetric) is a Lie algebra; this is one particular way to produce one \ts{9}{1}{48:19} \ts{9}{1}{49:01} \ts{9}{1}{49:42} \ts{9}{1}{50:44} \ts{9}{1}{51:25} \ts{9}{1}{52:06}.
Now the derivation: for a fixed $H\in A$ the map $D:=[H,\cdot\,]$ (a bilinear map with one slot left open is a linear map) is linear, and using the Jacobi identity one finds
\[
  D\bigl([a,b]\bigr)=[D(a),b]+[a,D(b)],
\]
which is the Leibniz rule with the product $\bullet=[\ ,\ ]$ (were the product written $\bullet$ one would see $D(a\bullet b)=D(a)\bullet b+a\bullet D(b)$). So $D=[H,\cdot\,]$ is a derivation of the Lie algebra $A$ into itself \ts{9}{1}{52:06} \ts{9}{1}{52:29} \ts{9}{1}{52:50} \ts{9}{1}{53:15} \ts{9}{1}{53:36} \ts{9}{1}{53:57}.
\end{example}
\begin{exercise}[Jacobi identity and the derivation property]\label{ex:tv-jacobi}
\emph{Not said in class.} Expand the three double commutators of \cref{ex:tv-lie} in terms of compositions, check that the twelve terms cancel in pairs, and deduce $[H,[a,b]]=[[H,a],b]+[a,[H,b]]$.
\end{exercise}
\begin{remark}[Where one has met derivations]\label{rem:tv-physics}
Where does one know this kind of derivation from? From classical mechanics: the \emph{Poisson bracket} is a derivation on the algebra of functions on phase space, the \emph{observables}; and only this is the reason why in classical mechanics one can identify the time derivative with the Poisson bracket with the Hamiltonian, since a
derivative must satisfy the product rule and so the bracket has to do so too. And from quantum mechanics: $\mathrm d/\mathrm dt$ is to be identified with $[H,\cdot\,]$ acting on observables, now self-adjoint operators. The idea of a derivation on an algebra satisfying the product rule has been met in classical mechanics, in quantum mechanics and now
in differential geometry: ``an extremely common concept, and maybe worth once and for all clarifying its algebraic status'' \ts{9}{1}{53:57} \ts{9}{1}{54:20} \ts{9}{1}{55:08} \ts{9}{1}{55:28} \ts{9}{1}{56:10} \ts{9}{1}{56:31}.
\end{remark}
\begin{coursenote}[Factors in the quantum case]
In quantum mechanics the Heisenberg equation is $\mathrm dA/\mathrm dt=\tfrac{i}{\hbar}[H,A]$ (and in classical mechanics $\mathrm df/\mathrm dt=\{f,H\}$, with the opposite order); the lecture identifies them only up to such factors and signs. Not said in class.
\end{coursenote}
The main point of the aside: the tangent vectors are derivations. Any questions so far? No \ts{9}{1}{56:31} \ts{9}{1}{56:51}. This ends the aside \ts{9}{1}{56:51}.

% ==================================================================
\section{The dimension of the tangent space}\label{sec:tv-dim}

An important theorem \ts{9}{1}{57:17}. The tangent space at another point is a vector space too, and as a vector space it has a dimension (the size of a basis, if finite).

\begin{theorem}[Dimension of the tangent space]\label{thm:tv-dim}
Let $M$ be a smooth manifold and $p\in M$. Then
\[
  \dim_{\text{vector space}}T_pM=\dim_{\text{manifold}}M.
\]
\ts{9}{1}{57:17} \ts{9}{1}{57:38} \ts{9}{1}{57:59}
\end{theorem}
These are two entirely different notions of dimension, and there is nothing obvious about the statement: the vector space dimension has technically nothing to do with the manifold dimension, which was the $\R^D$ of the charts, and the charts were not considered vector spaces, only topological spaces (\cref{ch:l06}). In principle
one could call them $\mathrm{Dim}$ and $\dim$; it needs a proof, and the strategy is to construct a vector space basis from a chart \ts{9}{1}{58:19} \ts{9}{1}{58:41} \ts{9}{1}{59:04}.

\paragraph{The chart-induced curves.} Choose a chart $(U,x)$ with $p\in U$, which one may take to be \textbf{centred at $p$}, i.e.\ $x(p)=0\in\R^D$ (without loss of generality: subtract $x(p)$ from the chart; terribly convenient). Consider $D=\dim M$ curves $\gamma_{(a)}\colon\R\to U\subseteq M$, $a=1,\dots,D$ (the index in brackets,
not to be confused with a component: it labels the curve) with the property that in the chart
\[
  x^b\bigl(\gamma_{(a)}(\lambda)\bigr)=\delta^b_a\,\lambda,\qquad b=1,\dots,D,
\]
i.e.\ in the chart the curve has all components $0$ except the $a$-th, which grows linearly with the parameter. ($\delta^b_a=1$ if $a=b$, $0$ otherwise, ``just a short notation for an array of equations''; this was asked by a student.) In $\R^D$, the Cartesian product of $D$ copies of $\R$, one has a clear notion of first, second, \dots\ component;
in the chart these curves are the coordinate axes through $0$, and $\gamma_{(a)}=x^{-1}\circ(\text{the $a$-th axis})$ are the \textbf{coordinate lines through $p$} drawn on the manifold. Another chart produces another coordinatisation. One could have drawn this picture much earlier, to get an idea of the back and forth between the charts and the manifold
\ts{9}{1}{59:27} \ts{9}{1}{1:00:09} \ts{9}{1}{1:01:11} \ts{9}{1}{1:02:35} \ts{9}{1}{1:03:38} \ts{9}{1}{1:06:26}.
Define the tangent vectors $E_a:=X_{\gamma_{(a)},p}$; all the curves pass through $p$ at $\lambda=0$ (this is why a chart centred at $p$ is used) \ts{9}{1}{1:04:40} \ts{9}{1}{1:05:02} \ts{9}{1}{1:05:22} \ts{9}{1}{1:05:43}.

Now see how the $a$-th tangent vector acts on a smooth function $f$:
\begin{align*}
E_a(f)&=(f\circ\gamma_{(a)})'(0)=\bigl((f\circ x^{-1})\circ(x\circ\gamma_{(a)})\bigr)'(0)\\
&=\partial_b\bigl(f\circ x^{-1}\bigr)\bigl(x(p)\bigr)\cdot\bigl(x^b\circ\gamma_{(a)}\bigr)'(0)=\partial_b\bigl(f\circ x^{-1}\bigr)\bigl(x(p)\bigr)\,\delta^b_a\\
&=\partial_a\bigl(f\circ x^{-1}\bigr)\bigl(x(p)\bigr),
\end{align*}
by inserting the identity $x^{-1}\circ x$, the associativity of the composition, the chain rule (the inner function is $x\circ\gamma_{(a)}\colon\R\to\R^D$, evaluated at $x(\gamma_{(a)}(0))=x(p)$) and $(x^b\circ\gamma_{(a)})'(0)=\delta^b_a$
\ts{9}{1}{1:06:47} \ts{9}{1}{1:07:28} \ts{9}{1}{1:07:49} \ts{9}{1}{1:08:30} \ts{9}{1}{1:08:51} \ts{9}{1}{1:09:33}.

\begin{definition}[New concepts: $\partial/\partial x^a$, the chart-induced tangent vectors]\label{def:tv-ddx}
For a chart $(U,x)$ with $p\in U$ and $a=1,\dots,\dim M$ the tangent vector $\bigl(\partial/\partial x^a\bigr)_p\in T_pM$ is defined by
\[
  \Bigl(\frac{\partial}{\partial x^a}\Bigr)_p(f):=\partial_a\bigl(f\circ x^{-1}\bigr)\bigl(x(p)\bigr)\qquad\text{for }f\in C^\infty(M).
\]
By the computation above $\bigl(\partial/\partial x^a\bigr)_p=X_{\gamma_{(a)},p}=E_a$ \ts{9}{1}{1:09:57} \ts{9}{1}{1:15:33} \ts{9}{1}{1:15:55}.
\end{definition}

\begin{intuition}[Two different objects: $\partial_a$ and $\partial/\partial x^a$]\label{int:tv-two-partials}
This ``at first mysterious'' remark is important and a very frequent confusion. $\partial_a$ is the $a$-th partial derivative of a function $\R^D\to\R$: it sends a smooth function $\R^D\to\R$ to a smooth function $\R^D\to\R$, ``it always did, it always will'', and it acts on functions on a Cartesian product, where it is
known which entry is the first, second, third. In physics one writes $\partial/\partial x^a$ for it. That makes no sense: it does not matter how one calls the entries when defining the function ($f(x,y,z)=x^2+y^2+z^2$ or with $b$ in place of $y$); the abstract derivative with respect to the $a$-th
entry has nothing to do with the name of the independent variable, as little as in one dimension it matters whether one writes $x$ or $\lambda$.\par
The symbol $\partial/\partial x^a$ is an \emph{entirely new symbol}, which has nothing to do with the former: it acts on functions $f$ on the manifold, $M\to\R$, for which one has no idea what the ``first entry'' would be (the manifold is a topological space); one first needs a chart map to coordinatise the manifold around $p$, obtain the chart
representative $f\circ x^{-1}$ on $\R^D$ and only then take the derivative in the $a$-th direction. So $\partial/\partial x^a$ takes a smooth function on the manifold to a real number, at the point $p$ (or the corresponding point $x(p)$); one could also produce the derivative function, but here only the value at the point is used. For the
manifold it makes sense, indeed it is necessary, that a chart has been picked. Whenever $\partial/\partial x^a$ is written in this course (and the lecturer recommends it for own calculations) the $x$ is the $x$ of the chart map: it is just a symbol, but it refers to the specific chart by which one can speak of the $a$-th component. Its value is
computed by the partial derivative of the chart representative. ``So far you never needed this; you always needed that. And here it makes sense because it refers to a chart. Very important: do not confuse the two'' \ts{9}{1}{1:09:57} \ts{9}{1}{1:10:59} \ts{9}{1}{1:12:02} \ts{9}{1}{1:13:04} \ts{9}{1}{1:14:09} \ts{9}{1}{1:15:11}.
\end{intuition}
So the tangent vectors to the coordinate-induced curves $\gamma_{(a)}$ are the chart-induced vectors $\partial/\partial x^a|_p$ \ts{9}{1}{1:15:33} \ts{9}{1}{1:15:55} \ts{9}{1}{1:16:18} \ts{9}{1}{1:16:39}. What is their meaning?

\begin{proposition}[The chart-induced vectors generate $T_pM$]\label{prop:tv-generate}
Every $X\in T_pM$ can be written as
\[
  X=X^a\,\frac{\partial}{\partial x^a}\Bigr|_p\qquad(\text{summation convention, \cref{def:te-einstein}}),
\]
with real numbers $X^a$; i.e.\ $\partial/\partial x^1|_p,\dots,\partial/\partial x^{\dim M}|_p$ is a \emph{generating system} of $T_pM$. In fact $X^a=(x^a\circ\mu)'(0)$ for any curve $\mu$ with $X=X_{\mu,p}$ \ts{9}{1}{1:16:39} \ts{9}{1}{1:16:59} \ts{9}{1}{1:17:21} \ts{9}{1}{1:17:41} \ts{9}{1}{1:23:34}.
\end{proposition}
\begin{proof}
Certainly there is a smooth curve $\mu$ through $p$ with $\mu(0)=p$ such that $X=X_{\mu,p}$: this is exactly what being an element of $T_pM$ means (there are many such curves; one is chosen) \ts{9}{1}{1:18:03} \ts{9}{1}{1:19:04} \ts{9}{1}{1:19:25} \ts{9}{1}{1:19:46}.
``How does one study a map? By letting it act on an arbitrary element of its domain'': for $f\in C^\infty(M)$,
\begin{align*}
X_{\mu,p}(f)&=(f\circ\mu)'(0)=\bigl((f\circ x^{-1})\circ(x\circ\mu)\bigr)'(0)\\
&=\partial_b\bigl(f\circ x^{-1}\bigr)\bigl(x(p)\bigr)\,(x^b\circ\mu)'(0)\\
&=(x^b\circ\mu)'(0)\,\frac{\partial}{\partial x^b}\Bigr|_p(f).
\end{align*}
The little diagram helps to see that $f\circ\mu=(f\circ x^{-1})\circ(x\circ\mu)$: restricted to the chart, $x\circ\mu$ is the representative of the curve and $f\circ x^{-1}$ that of the function, ``they come in these pairs'', and then the brackets can be dropped by associativity. The chain rule is the same calculation over and over again, and
$x(\mu(0))=x(p)$. Hence $X=X^b\,\partial/\partial x^b|_p$ with $X^b:=(x^b\circ\mu)'(0)$, a number for each $b$ \ts{9}{1}{1:19:46} \ts{9}{1}{1:20:27} \ts{9}{1}{1:21:09} \ts{9}{1}{1:21:29} \ts{9}{1}{1:22:10} \ts{9}{1}{1:22:53}.
\end{proof}
\begin{coursenote}[A slip in the board]
The board writes ``$x^b$ after $\mu$ inverse'' at one point; what is meant is $x^b\circ\mu$, not an inverse. The same remark: a few steps are written with $\gamma$ and $\mu$ interchanged. Not said in class.
\end{coursenote}
So every tangent vector is a linear combination, constructively found, of the chart-induced vectors; the numbers $X^b$ are the components with respect to this generating system (``I almost said basis, but we haven't shown that yet'') \ts{9}{1}{1:23:34} \ts{9}{1}{1:23:55}.
This is what one always did in physics: take a curve, not in the manifold abstractly but in a chart representation (the particle flies along the curve $\mu$, the first component is this function of time, the second is that function of time) and derive componentwise; the components of the tangent vector, physically the derivatives,
with respect to the chart-induced basis of the tangent space. But since one usually takes $\R^D$ as the basic space, every tangent space can be identified with $\R^D$ itself and everything is mixed up; this is the real logic behind it \ts{9}{1}{1:23:55} \ts{9}{1}{1:24:15} \ts{9}{1}{1:24:36} \ts{9}{1}{1:24:57}.
The generating system is not unique: building linear combinations of certain types gives other ones; this one is very special, induced from the coordinates, but with the advantage that derivatives of curves can be calculated in the chart picture. Later one will consider bases of the tangent space that are \emph{not} coordinate
induced, especially for vector fields on extended regions (at a point a basis can always be constructed from coordinates) \ts{9}{1}{1:25:17} \ts{9}{1}{1:25:38} \ts{9}{1}{1:25:58} \ts{9}{1}{1:26:18}.

\paragraph{Linear independence.} Suppose $\lambda^a\,\partial/\partial x^a|_p=0$ (summation over $a$; the $\lambda^a$ are real numbers, a different $\lambda$ from the curve parameter). If the vectors are linearly independent one must show that every $\lambda^a=0$. The left-hand side is an element of $T_pM$, hence an operator, so
the equation is an operator equation; it is handled by letting it act on an arbitrary element of the domain; and here one chooses very specific functions, the \emph{coordinate functions} $x^b$ \ts{9}{1}{1:26:39} \ts{9}{1}{1:27:00} \ts{9}{1}{1:27:21} \ts{9}{1}{1:27:42}.

\begin{remark}[Is it legal to act on $x^b$?]\label{rem:tv-legal}
A lawyer would say ``it depends''; the answer is yes, but not obviously. The $b$-th component $x^b\colon U\to\R$ of the chart map is continuous (the chart map is a homeomorphism onto its image, and a component of a homeomorphism to $\R^D$ is continuous; it cannot be a homeomorphism, since the dimensions differ). But $T_pM$ acts only on
\emph{smooth} functions $C^\infty(M)$, and by construction $x^b$ was only continuous on the manifold: one would not even know how to say it is differentiable. However, smoothness of a function $f$ on the manifold is \emph{decided} by looking at $f\circ x^{-1}$ in a chart: $f$ is smooth if $f\circ x^{-1}$ is. For $f=x^b$:
\[
  x^b\circ x^{-1}\colon x(U)\to\R,\qquad(a^1,\dots,a^D)\longmapsto a^b,
\]
the \emph{projection to the $b$-th component}, which is of course smooth. Since the definition of smoothness is well defined (\cref{prop:df-welldef}), the coordinate functions are smooth, and acting on them is legal; one could say that the definition of smooth function is chosen such that the chart component maps are smooth.
(This is a favourite for an oral exam: ``at first this can be very confusing''.) (A technical caveat: $x^b$ is defined on $U$ only; to act on it as an element of $C^\infty(M)$ one needs a smooth function on $M$ agreeing with it near $p$, or to let $T_pM$ act on functions defined near $p$; neither is discussed in the lecture; not said in class [?].)
\ts{9}{1}{1:27:42} \ts{9}{1}{1:28:46} \ts{9}{1}{1:29:50} \ts{9}{1}{1:30:34} \ts{9}{1}{1:31:36} \ts{9}{1}{1:32:39}.
\end{remark}
Now one can act: by the definition of the symbol, $\partial/\partial x^a|_p(x^b)=\partial_a(x^b\circ x^{-1})(x(p))$, and $x^b\circ x^{-1}$ is the projection to the $b$-th component, so the derivative in the $a$-th direction is
\[
  \frac{\partial}{\partial x^a}\Bigr|_p(x^b)=\delta^b_a.
\]
Hence $0=\lambda^a\,\partial/\partial x^a|_p(x^b)=\lambda^a\delta^b_a=\lambda^b$ for each $b$: ``you start on the assumption that this guy is zero, as an operator, so it must be zero if it acts on particular functions; we picked clever ones; we derive $\lambda^b=0$'' \ts{9}{1}{1:32:59} \ts{9}{1}{1:33:20} \ts{9}{1}{1:33:41} \ts{9}{1}{1:34:25} \ts{9}{1}{1:34:45} \ts{9}{1}{1:35:06}.
\begin{proof}[Proof of \cref{thm:tv-dim}]
$\partial/\partial x^1|_p,\dots,\partial/\partial x^{\dim M}|_p$ generate $T_pM$ (\cref{prop:tv-generate}) and are linearly independent (above); they are a basis, with $\dim M$ elements. Hence $\dim T_pM=\dim M$ \ts{9}{1}{1:35:06} \ts{9}{1}{1:35:27}.
\end{proof}
Questions? None. \ts{9}{1}{1:35:27} \ts{9}{1}{1:35:56}.
\begin{coursenote}[Students and the dots]
At \ts{9}{1}{34:44} there is a short exchange about whether something was said ``in between, like other dots''; a technical remark about the notation on the board, without content.
\end{coursenote}

% ==================================================================
\section{Components, change of chart and the myth of position vectors}\label{sec:tv-components}

\begin{definition}[New concepts: components of a tangent vector]\label{def:tv-components}
If a vector $X\in T_pM$ is written $X=X^a\,\partial/\partial x^a|_p$, the real numbers $X^1,\dots,X^{\dim M}$ are the \textbf{components} of $X$ with respect to the tangent space basis induced from the chart $(U,x)$ (where of course $p\in U$) \ts{9}{1}{1:35:56} \ts{9}{1}{1:36:17} \ts{9}{1}{1:36:39} \ts{9}{1}{1:37:03} \ts{9}{1}{1:37:23}.
\end{definition}
The problem sheet will ask for the change of components under a change of charts; a very obvious question. The key point: a chart transition map can be highly non-linear (Cartesian to polar coordinates is a highly non-linear function $\R^D\to\R^D$); but at any given point it induces a \emph{linear} transformation on the
vector space at that point, and it must, because vectors must transform linearly. If a new chart $y$ is given in terms of the old one by $y^a=y^a(x^1,\dots,x^{\dim M})$ (physicist notation), the Jacobian matrix $\partial y^a/\partial x^b$ evaluated at the point $x(p)\in\R^D$ governs the change of the basis of $T_pM$ (``kill me, I don't know whether it is the matrix $A$ or
$B$ of last time'': see \cref{prop:te-transform}); the Jacobian is a square array of functions, but it is evaluated at this one point \ts{9}{1}{1:37:43} \ts{9}{1}{1:38:04} \ts{9}{1}{1:38:47} \ts{9}{1}{1:39:07} \ts{9}{1}{1:39:50} \ts{9}{1}{1:40:10}.
\begin{exercise}[Change of components under a change of chart]\label{ex:tv-change}
Let $(U,x)$ and $(V,y)$ be charts around $p$. Show that $\dfrac{\partial}{\partial y^a}\Big|_p=\dfrac{\partial x^b}{\partial y^a}\Big|_{y(p)}\dfrac{\partial}{\partial x^b}\Big|_p$ by letting both sides act on an arbitrary $f$ (chain rule applied to $f\circ y^{-1}=(f\circ x^{-1})\circ(x\circ y^{-1})$), and find the transformation of the components $X^a$.
(Problem sheet; the long version with polar coordinates is in \cref{exam:ch09}.)
\end{exercise}

\begin{intuition}[There are no position vectors]\label{int:tv-position}
This is the full point of what is always said about there being no position vectors: the position is this. On the manifold a point is a point; what are the coordinates of Paris? There are none, unless one chooses a chart: then ``it's 7, 2'', in that chart. The chart can be changed in any diffeomorphic way ($C^\infty$ maps
$\R^D\to\R^D$), highly non-linearly. The coordinates $x^a(p)$, the components of the chart map, are certainly \emph{not} components of a vector: if they were, they would transform linearly under a change of chart, and they transform highly non-linearly. The tangent vectors to curves, on the other hand, are ``totally unimpressed'' by the non-linearity of the chart
change: at the point they only see the Jacobian matrix at that point as the linear transformation of the basis of the tangent space \ts{9}{1}{1:40:10} \ts{9}{1}{1:40:31} \ts{9}{1}{1:40:51} \ts{9}{1}{1:41:12} \ts{9}{1}{1:41:32}.
\end{intuition}
In the old books a vector is ``what transforms like this'' (with the Jacobian or its inverse, depending on the convention); but positions simply do not. Only if one restricts attention to \emph{linear} changes of coordinates does the transformation of the coordinates coincide with that of the vector components (the Jacobian of a linear map is the map itself, constant),
and then one can yield to the illusion that the coordinates constitute a vector transforming linearly. But that means one can only do Cartesian coordinates and never go to polar ones, let alone anything more interesting \ts{9}{1}{1:41:32} \ts{9}{1}{1:41:53} \ts{9}{1}{1:42:14} \ts{9}{1}{1:42:35}.
Much of special relativity is done with this confusion: the Lorentz transformations are said to act on all of position space, the $x^\mu$ ``are being transformed''; no, the $x^\mu$ can be transformed in any way one wishes, and only the derivatives, the tangent vectors at specific points, transform. If the chart transition is non-linear, the linear transformation in $T_pM$ is
not the same as that in $T_qM$ at another point $q$; it differs from point to point, from tangent space to tangent space, whereas for a linear change of chart it is everywhere the same. ``We should never abuse this by arbitrary restrictions of these functions only to save our concept of a position vector: an artificial concept, pointless, it causes confusion''
\ts{9}{1}{1:42:35} \ts{9}{1}{1:42:56} \ts{9}{1}{1:43:16} \ts{9}{1}{1:43:37} \ts{9}{1}{1:43:57}.

% ==================================================================
\section*{Key concepts}
\begin{keyconcepts}
\textbf{Retrieval questions.}
\begin{enumerate}
  \item What is the vector space $C^\infty(M)$, and which product makes it an algebra? Why is the scalar multiplication different from the product?
  \item Define the tangent vector $X_{\gamma,p}$ to a curve as an operator. Why is it the ``velocity'' of $\gamma$? What happens at double speed?
  \item Define $T_pM$. Why must one prove that the sum of two tangent vectors is again a tangent vector, and how does the curve $\sigma$ do it? Why is the chart dependence of $\sigma$ harmless?
  \item Define algebra, derivation, Lie algebra. Why is $T_pM\ni X$ a derivation, and why is $[H,\cdot\,]$ one?
  \item State and prove $\dim T_pM=\dim M$: which curves and which basis are used?
  \item What is the difference between $\partial_a$ and $\partial/\partial x^a|_p$? Why can $T_pM$ act on the coordinate functions $x^b$?
  \item How do the components of a tangent vector change under a change of chart? Why are the coordinates of a point not components of a vector?
\end{enumerate}
\textbf{Concepts.}
\begin{itemize}
  \item $C^\infty(M)$ (\cref{def:tv-cinf}); tangent vector as directional derivative operator (\cref{def:tv-vector,int:tv-married,int:tv-velocity}); the tangent space and its closure (\cref{def:tv-space,prop:tv-closure}).
  \item Algebra, derivation, Lie algebra, Poisson bracket (\cref{def:tv-algebra,def:tv-derivation,ex:tv-der-tangent,ex:tv-lie,rem:tv-physics}).
  \item $\dim T_pM=\dim M$ via the chart-induced basis (\cref{thm:tv-dim,def:tv-ddx,prop:tv-generate,rem:tv-legal}); $\partial_a$ versus $\partial/\partial x^a$ (\cref{int:tv-two-partials}).
  \item Components, change of chart, no position vectors (\cref{def:tv-components,ex:tv-change,int:tv-position}).
\end{itemize}
\end{keyconcepts}

\flushside
\end{paracol}

% !TEX root = ../main.tex
\chapter{Construction of the tangent bundle}\label{ch:l10}
\begin{paracol}{2}

\begin{sources}
\begin{tabularx}{\linewidth}{@{}l l L@{}}
\toprule
\textbf{Video} & \textbf{Length} & \textbf{Content}\\
\midrule
\seg{10}{1}{Lecture 10: Construction of the tangent bundle} & 1:48:50 & Cotangent spaces and tensor spaces at a point; canonical and non-canonical isomorphisms; the gradient $\mathrm d_pf$ as a co-vector and the level sets; the chart-induced dual basis $\mathrm dx^a$;
push-forward (the derivative of a smooth map) and pull-back; vectors are pushed forward, forms pulled back; immersions and embeddings with the circle as example; the Whitney embedding and immersion theorems and the Klein bottle;
the tangent bundle as a set bundle, a correction about a bundle of real and complex lines, the topology and the smooth atlas on $TM$ (compatibility left to the problem sheet).\\
\bottomrule
\end{tabularx}

\smallskip
\textbf{Transcript} \file{materials/transcripts/C10-V01.txt} (YouTube automatic captions).
\end{sources}

\begin{reading}[Where this lecture sits in the other notes]
\emph{Gravity and Light}, Lecture~6 (Fields); \emph{Theoretische Mechanik}, Lecture~5. The tangent bundle and the push-forward are the prerequisites of vector fields; the Whitney theorems, the immersion/embedding distinction and the discussion of the gradient
as a co-vector are, as far as the titles say, specific to this course; the modules announced here are the subject of \cref{ch:l11}.
(From the lecture titles and the plan of the course; not checked in detail.)
\end{reading}

% ==================================================================
\section{Plan of the lecture}\label{sec:tb-plan}

Today some technology is developed, collecting along the way a number of interesting results; the two most important ones \ts{10}{1}{0:04}. First the \textbf{Whitney embedding theorems}, which clarify whether there is really a difference between the intrinsic definition of a manifold and of its tangent spaces and the \emph{extrinsic} one, not
followed in this course (\cref{int:df-embedding}), where the manifold is thought of as embedded in some $\R^N$ and the tangent spaces are planes in the embedding space. The intrinsic point of view was argued for, since for physics, if the manifold is to represent the universe, it would be funny to think of an even bigger space into which the
universe is embedded; the question is whether something could be gained by taking the embedding approach, and the Whitney theorems say that the two are equivalent \ts{10}{1}{0:04} \ts{10}{1}{0:26} \ts{10}{1}{0:47} \ts{10}{1}{1:07}.
Second, the construction of the tangent bundle: vector fields will be defined, but the space of all vector fields is \emph{not} a vector space (the space of all vectors at a point was): it is merely a \emph{module}, a vector space over a commutative ring; modules are the subject of the next lecture (\cref{ch:l11}) \ts{10}{1}{1:28} \ts{10}{1}{1:49} \ts{10}{1}{2:10}.

% ==================================================================
\section{Cotangent spaces and tensor spaces at a point}\label{sec:tb-cotangent}

Having constructed the tangent spaces last time (\cref{ch:l09}) it is quite simple to come up with the cotangent space \ts{10}{1}{2:10} \ts{10}{1}{2:30}.
\begin{definition}[New concepts: cotangent space]\label{def:tb-cotangent}
Let $M$ be a smooth manifold (with its topology and smooth atlas). The \textbf{cotangent space} to $M$ at $p\in M$ is the vector space dual of the tangent space (\cref{def:te-hom}),
\[
  T^*_pM:=(T_pM)^*,
\]
with the inherited addition and scalar multiplication; the star ``jumps'' in the notation. It is a vector space, ``no secret here'' \ts{10}{1}{2:30} \ts{10}{1}{2:51} \ts{10}{1}{3:13} \ts{10}{1}{3:34}.
\end{definition}
\begin{remark}[Finite dimension, and a non-canonical isomorphism]\label{rem:tb-noncanonical}
The manifolds defined in this course are finite dimensional, since one always maps into some $\R^D$; the definition did not include infinite dimensional manifolds, which can be seen as a defect: with more trouble (``not much more, but some more'') one could. Since $\dim T_pM=\dim M<\infty$ (\cref{thm:tv-dim}),
$T_pM$ and $T^*_pM$ have the same dimension and are \emph{isomorphic as vector spaces}: there exists a bijective linear map between them. ``Which one? Many; but none of them can be constructed without extra structure''. The isomorphisms exist but are \textbf{non-canonical}: from the bare vector space structure
one cannot provide the map without picking, say, arbitrarily the entries of a matrix defining it. With extra structure (an inner product, for instance) one can construct one. This is quite different from $V\cong V^{**}$ (\cref{thm:te-dd}, on the current problem sheet): there a \textbf{canonical} isomorphism exists, one that needs no extra structure beyond the level
at which one is discussing, here the level of bare vector spaces \ts{10}{1}{3:55} \ts{10}{1}{4:36} \ts{10}{1}{4:58} \ts{10}{1}{5:38} \ts{10}{1}{5:59} \ts{10}{1}{6:40}.
\end{remark}
For the tangent space there was a nuts-and-bolts construction with curves and directional derivative operators, step by step; for the cotangent space this is not necessary, one just takes the vector space dual. Hence at a point $p$ one has the vector space $T_pM$ (the only thing that had to be built from scratch, over which everything is then pure
linear algebra) and its dual \ts{10}{1}{7:01} \ts{10}{1}{7:21}.
\begin{definition}[New concepts: tensors at a point, $T^r_s(T_pM)$]\label{def:tb-tensors}
Similarly one can construct the space of $(r,s)$ tensors at $p$,
\[
  T^r_s(T_pM):=\{\,T\colon\underbrace{T^*_pM\times\dots\times T^*_pM}_{r}\times\underbrace{T_pM\times\dots\times T_pM}_{s}\to\R\ \text{multilinear}\,\},
\]
in the notation of \cref{def:te-tspace}, with the tangent space $T_pM$ as the underlying vector space ($T$ the tensor, $T_pM$ the tangent space: different $T$'s; no smoothness condition is required, it is a pure vector space matter, the multilinear maps are on vector spaces at the same point). It is a vector space with the pointwise
addition and scalar multiplication, and every tensor at $p$ lives there \ts{10}{1}{7:21} \ts{10}{1}{7:42} \ts{10}{1}{8:23} \ts{10}{1}{8:44} \ts{10}{1}{9:27} \ts{10}{1}{9:47}.
\end{definition}
The only time one has to construct something from scratch is the tangent space itself with the curves; everything else is generated by pure linear algebra. ``Tensors have nothing to do with differential geometry, tensors have all to do with linear algebra and multilinear maps; the only thing we had to construct was an underlying vector space over which the tensor space is
built'', always at this one point $p$ \ts{10}{1}{10:07} \ts{10}{1}{10:29} \ts{10}{1}{10:50}.

% ==================================================================
\section{The gradient}\label{sec:tb-gradient}

\begin{definition}[New concepts: gradient operator, gradient]\label{def:tb-gradient}
Let $f$ be a smooth function on $M$. At every point $p\in M$ there is a linear map, written $\mathrm d_p$ (``if you want to emphasise that we look at this only at one point''; a different map for every point),
\[
  \mathrm d_p\colon C^\infty(M)\to T^*_pM,\qquad f\mapsto\mathrm d_pf,\qquad(\mathrm d_pf)(X):=X(f)\quad\text{for }X\in T_pM.
\]
The map $\mathrm d_p$ is called the \textbf{gradient operator at $p$} and $\mathrm d_pf$ the \textbf{gradient of $f$ at $p$} \ts{10}{1}{10:50} \ts{10}{1}{11:31} \ts{10}{1}{11:52} \ts{10}{1}{12:34} \ts{10}{1}{12:55} \ts{10}{1}{14:42}.
\end{definition}
This seems formalism taken to an extreme: one could simply write $X(f)$ instead of defining a map $\mathrm d_pf\colon T_pM\to\R$ acting on $X$. The reason is that one wants to speak of this object \emph{without providing the vector it acts on}; as in analysis one often wants to talk about the function as such and not about its value $f(x)$ at some point
\ts{10}{1}{13:18} \ts{10}{1}{13:39} \ts{10}{1}{14:00}.
\begin{remark}[The gradient is a co-vector, not a vector]\label{rem:tb-covector}
(a) $\mathrm d_pf\in T^*_pM$ is a \emph{co-vector}, not a vector \ts{10}{1}{14:42} \ts{10}{1}{15:05}. One often hears that the gradient is orthogonal to the level lines of $f$; but orthogonality needs an inner product, and with a bare vector space, which is the case here, there is none in sight.\par
(b) Let $X\in T_pM$ be tangent to the \textbf{level set} $N_c(f):=\{p\in M\mid f(p)=c\}$, $c\in\R$ (``tangent to a level set'' has not been defined formally yet, but the idea: a function on a three-dimensional manifold constant on a sphere, where the tangent to the level set is the tangent to a smooth submanifold; to be treated later). Then
automatically $\mathrm d_pf(X)=0$: the co-vector acts on the vector and annihilates it. The statement ``the gradient is orthogonal to the level sets'' needs no inner product \ts{10}{1}{15:05} \ts{10}{1}{15:25} \ts{10}{1}{16:08} \ts{10}{1}{16:31} \ts{10}{1}{17:12} \ts{10}{1}{17:32}.
In standard vector analysis one commits a ``perversion'': with a metric one constructs a non-canonical isomorphism to raise the index and make the gradient a vector, and then uses the metric again to dot this gradient vector into another vector; the inverse metric used for the first step and the metric used for the second cancel. ``It is totally perverse,
because the gradient naturally is a co-vector'' \ts{10}{1}{17:32} \ts{10}{1}{17:53} \ts{10}{1}{18:14} \ts{10}{1}{18:35}.
\end{remark}

% ==================================================================
\section{The chart-induced dual basis $\mathrm dx^a$}\label{sec:tb-dx}

The gradient is useful in itself, but it can also be used to construct a basis of the cotangent space \ts{10}{1}{19:03} \ts{10}{1}{19:24}.
\begin{definition}[New concepts: chart-induced co-vector basis]\label{def:tb-dx}
Let $p\in M$ and $(U,x)$ a chart with $p\in U$. Each component function $x^a\colon U\to\R$, $a=1,\dots,\dim M$, of the chart map is smooth (\cref{rem:tv-legal}); one can act on it with the gradient operator and obtains the $\dim M$ co-vectors
\[
  \mathrm d_px^a\in T^*_pM,\qquad a=1,\dots,\dim M,
\]
written $\mathrm dx^a|_p$ for short. They are called the \textbf{chart-induced co-vector basis} at $p$ (Theorem below shows that it is a basis, and indeed a dual one) \ts{10}{1}{19:24} \ts{10}{1}{20:05} \ts{10}{1}{20:25} \ts{10}{1}{21:09} \ts{10}{1}{21:29} \ts{10}{1}{22:11}.
\end{definition}
There are enough of them (the cotangent space has dimension $\dim M$ since the manifold is finite dimensional), so they are a candidate for a basis; this is one of those ``definition-theorems'' where one has to provide a proof.
\begin{theorem}[The dual basis]\label{thm:tb-dualbasis}
$\mathrm dx^a|_p(\partial/\partial x^b|_p)=\delta^a_b$. Hence $\mathrm dx^1|_p,\dots,\mathrm dx^{\dim M}|_p$ is a basis of $T^*_pM$, namely the \emph{dual basis} (\cref{def:te-dualbasis}) of the chart-induced basis $\partial/\partial x^b|_p$ of $T_pM$ \ts{10}{1}{22:11} \ts{10}{1}{27:27} \ts{10}{1}{27:48} \ts{10}{1}{28:08}.
\end{theorem}
\begin{proof}
Let the $a$-th co-vector act, as a linear map $T_pM\to\R$, on the $b$-th basis vector. By the definition of $\mathrm d_p$ and then of $\partial/\partial x^b|_p$ (\cref{def:tv-ddx}):
\[
  \mathrm dx^a|_p\Bigl(\frac{\partial}{\partial x^b}\Bigr|_p\Bigr)=\frac{\partial}{\partial x^b}\Bigr|_p(x^a)=\partial_b\bigl(x^a\circ x^{-1}\bigr)\bigl(x(p)\bigr)=\delta^a_b,
\]
since $x^a\circ x^{-1}$ is the projection to the $a$-th component of $\R^{\dim M}$, whose $b$-th partial derivative is $1$ if $b=a$ and $0$ otherwise (``it doesn't matter where you evaluate it''). ``It's also good practice to calculate this: a mere application of the definitions'' \ts{10}{1}{22:11} \ts{10}{1}{23:14} \ts{10}{1}{24:17} \ts{10}{1}{24:58} \ts{10}{1}{26:01} \ts{10}{1}{27:05}.
For linear independence (and the generating property): the equation $\mu_a\,\mathrm dx^a|_p=0$ evaluated on $\partial/\partial x^b|_p$ gives $\mu_b=0$; moreover for $\omega\in T^*_pM$ one has $\omega=\omega(\partial/\partial x^a|_p)\,\mathrm dx^a|_p$ (``for every basis element here one finds one that does not annihilate it and all the others annihilate it''),
and there are $\dim M$ of them: a basis, which by $\delta^a_b$ is the dual one. This proves more than was wanted: it is the dual basis of the dual of $T_pM$, ``I chose this more arcane formulation to use the word dual twice in one sentence, with slightly different meanings'': the dual of $T_pM$ is $T^*_pM$, and the dual basis refers to
applying one basis to the other and getting $\delta^a_b$ \ts{10}{1}{27:27} \ts{10}{1}{27:48} \ts{10}{1}{28:08} \ts{10}{1}{28:29} \ts{10}{1}{28:49}.
\end{proof}
\begin{intuition}[What $\mathrm dx^a$ is: not a small change in $x^a$]\label{int:tb-dx}
It is important not to look at $\mathrm dx^1$ as one does in undergraduate physics, as ``a small change in $x^1$''; it is precisely the co-vector defined above, and that makes full sense; one then has full technical control over what one is doing. Whenever one meets $\mathrm dx^1,\dots,\mathrm dx^{\dim M}$ in the future, for instance in the measure of multiple integrals $\mathrm dx^1\,\mathrm dx^2\,\mathrm dx^3$
(``there is a wedge product between them, we'll come to that''), this is exactly what appears under the integral, and it is not just a symbol: it can be dissected in this way, and it is very important to do so if one wants to do analysis on a smooth manifold \ts{10}{1}{28:49} \ts{10}{1}{29:11} \ts{10}{1}{29:32} \ts{10}{1}{29:53} \ts{10}{1}{30:13}.
\end{intuition}
\begin{exercise}[Components of $\mathrm d_pf$]\label{ex:tb-df}
\emph{Not said in class.} Show that for $f\in C^\infty(M)$ one has $\mathrm d_pf=\partial_a(f\circ x^{-1})(x(p))\,\mathrm dx^a|_p$, by evaluating both sides on $\partial/\partial x^b|_p$ (\cref{thm:tb-dualbasis,def:tv-ddx}).
\end{exercise}

% ==================================================================
\section{The push-forward}\label{sec:tb-pushforward}

The next short section is a very important one, with two important concepts: the \textbf{push-forward} and its cousin, the \textbf{pull-back} \ts{10}{1}{30:13} \ts{10}{1}{30:34} \ts{10}{1}{31:01}.
\begin{definition}[New concepts: push-forward]\label{def:tb-push}
Let $\phi\colon M\to N$ be a smooth map between smooth manifolds (each with its own topology and atlas, a priori unrelated; smoothness of $\phi$ is defined by pushing $\phi$ down to the chart representations, \cref{def:df-map}) \ts{10}{1}{31:01} \ts{10}{1}{31:24} \ts{10}{1}{31:45}.
The \textbf{push-forward} of $\phi$ at $p\in M$ is the linear map
\[
  \phi_{*,p}\colon T_pM\longrightarrow T_{\phi(p)}N,\qquad\bigl(\phi_{*,p}(X)\bigr)(f):=X(f\circ\phi)\quad\text{for }f\in C^\infty(N).
\]
\end{definition}
(In class the symbol was first written with the star up, $\phi^*$, and corrected to the standard convention $\phi_*$ with the star down; ``it's nothing to do with the dual star''. At one point also the target is said to be a ``cotangent space'', a slip: it is $T_{\phi(p)}N$.)
Why is the definition forced? The result must be an element of the tangent space of the target, hence it must act on a smooth function $f$ on $N$; but $f\circ\phi$ is a smooth function \emph{on $M$}, and $X$, a tangent vector to $M$ at $p$, can act on it. The only point of $N$ that is distinguished, given $p\in M$ and a map $\phi$, is $\phi(p)$; the only
linear map from $T_pM$ to $T_{\phi(p)}N$ one can build from the given data is this one. ``Don't memorise stuff you don't need to memorise; every time I construct this when I need it'' \ts{10}{1}{31:45} \ts{10}{1}{33:11} \ts{10}{1}{34:34} \ts{10}{1}{35:36} \ts{10}{1}{36:59} \ts{10}{1}{38:24}.
(Linearity of $\phi_{*,p}$ in $X$ follows from that of $X$; one may also check that $\phi_{*,p}(X)$ is linear and satisfies the Leibniz rule, i.e.\ is a derivation, \cref{def:tv-derivation}; not carried out in class.)

\begin{remark}[Names: derivative, push-forward]\label{rem:tb-names}
(a) As said, it is the only linear map $T_pM\to T_{\phi(p)}N$ that can be constructed from the given data. (b) A small $p$ may be written for each point (``a different such map at every point''): $\phi_{*,p}$. The push-forward is labelled by the star \emph{down} by convention. (c) $\phi_{*,p}$ is often called the \textbf{derivative} of $\phi$ at $p$: when somebody
talks about the derivative of a smooth map between smooth manifolds, this induced map from tangent space to tangent space is meant. (d) \emph{Push-forward} is a very good name: it pushes vectors from the tangent space in the domain to the tangent space in the target, in the direction in which the map goes. If $\phi$ is a diffeomorphism one could also go
back with the inverse, but maps are often not invertible and may nevertheless be smooth in one direction: even then every vector in the domain can be pushed to a vector in the target \ts{10}{1}{38:24} \ts{10}{1}{39:10} \ts{10}{1}{39:31} \ts{10}{1}{40:12} \ts{10}{1}{40:32} \ts{10}{1}{41:13}.
\end{remark}
\begin{proposition}[Push-forward of the tangent vector of a curve]\label{prop:tb-push-curve}
For a smooth map $\phi\colon M\to N$ and a smooth curve $\gamma$ through $p$ (with $\gamma(0)=p$)
\[
  \phi_{*,p}\bigl(X_{\gamma,p}\bigr)=X_{\phi\circ\gamma,\phi(p)}:
\]
the tangent vector of $\gamma$ at $p$ is pushed forward to the tangent vector of the curve $\phi\circ\gamma$ at $\phi(p)$ \ts{10}{1}{41:33} \ts{10}{1}{41:55} \ts{10}{1}{42:17} \ts{10}{1}{42:37}.
\end{proposition}
\emph{Illustration.} Two manifolds $M,N$; in $M$ a curve $\gamma$ through $p$ and its tangent vector (``I do something terrible: I draw a little arrow, but you know it is the directional derivative operator; in charts it corresponds to a $\dim M$-tuple of numbers, the components'', allowed for illustration only). The map takes every point to a point, in a smooth manner,
hence also the points of the curve: the image curve $\phi\circ\gamma$ in $N$, passing through $\phi(p)$; one can construct its tangent vector $Y$ directly, without pushing forward. The claim is that $Y=\phi_{*,p}(X)$ \ts{10}{1}{43:03} \ts{10}{1}{43:43} \ts{10}{1}{44:04} \ts{10}{1}{44:48} \ts{10}{1}{45:09} \ts{10}{1}{45:51}.
\begin{proof}
``It is good practice for the formalism.'' Let $f\in C^\infty(N)$ be arbitrary (a function on the target, since the result is a vector there). By the definition of the push-forward and then of the tangent vector of a curve (with $\gamma(0)=p$; otherwise put the parameter value where $\gamma$ passes through $p$),
\[\begin{aligned}
  \bigl(\phi_{*,p}X_{\gamma,p}\bigr)(f)&=X_{\gamma,p}(f\circ\phi)\\
  &=\bigl((f\circ\phi)\circ\gamma\bigr)'(0)\\
  &=\bigl(f\circ(\phi\circ\gamma)\bigr)'(0)\\
  &=X_{\phi\circ\gamma,\phi(p)}(f),
\end{aligned}\]
using the associativity of the composition and the definition of the tangent vector backwards: $\phi\circ\gamma$ is a curve in $N$ through $\phi(\gamma(0))=\phi(p)$ and $f$ a function on $N$. The two operators agree on an arbitrary function, hence they are equal \ts{10}{1}{45:51} \ts{10}{1}{46:33} \ts{10}{1}{47:35} \ts{10}{1}{48:17} \ts{10}{1}{49:20} \ts{10}{1}{50:01}.
\end{proof}
This is the geometric notion and intuition of the push-forward, ``or, as it's called, the derivative of the function at that particular chosen point'', and it can be done at every point \ts{10}{1}{50:01} \ts{10}{1}{50:22}.
\begin{exercise}[Chain rule for the push-forward]\label{ex:tb-chain}
\emph{Not said in class.} Show that for smooth $\phi\colon M\to N$, $\psi\colon N\to P$ one has $(\psi\circ\phi)_{*,p}=\psi_{*,\phi(p)}\circ\phi_{*,p}$ and $(\mathrm{id}_M)_{*,p}=\mathrm{id}_{T_pM}$ (hence for a diffeomorphism $\phi_{*,p}$ is invertible). Use \cref{def:tb-push}.
\end{exercise}

% ==================================================================
\section{The pull-back}\label{sec:tb-pullback}

\begin{definition}[New concepts: pull-back]\label{def:tb-pull}
Let $\phi\colon M\to N$ be again a smooth map. The \textbf{pull-back} of $\phi$ at the point $\phi(p)\in N$ is the linear map
\[
  \phi^{*}_{p}\colon T^*_{\phi(p)}N\longrightarrow T^*_pM,\qquad\bigl(\phi^*_p(\omega)\bigr)(X):=\omega\bigl(\phi_{*,p}(X)\bigr)\quad\text{for }X\in T_pM,
\]
this time with the star up (\cref{def:tb-push}) \ts{10}{1}{50:47} \ts{10}{1}{51:11} \ts{10}{1}{51:32} \ts{10}{1}{51:53} \ts{10}{1}{52:13} \ts{10}{1}{52:35}.
\end{definition}
It looks like the push-forward, only the cotangent spaces instead of the tangent spaces; the big issue is that the arrow goes in the \emph{opposite direction}, hence the name \emph{pull-back}: if $\phi$ is a diffeomorphism, one can go back and forth and write the pull-back of the inverse, which goes in the same direction; but this is not always the case: $\phi$
need not be surjective, for instance if $M$ has higher dimension than $N$ it cannot be injective and has no inverse. (In class the arrow was written ``for dramatic effect'' in the unusual direction; the point is as stated.) The map was a homework problem but ``it's too important not to do it right now'', and it was constructed together with the students
\ts{10}{1}{52:55} \ts{10}{1}{53:16} \ts{10}{1}{53:37} \ts{10}{1}{53:57}.
\emph{Construction, with the students.} Take $\omega\in T^*_{\phi(p)}N$, a co-vector on the target; $\phi^*_p(\omega)$ should be an element of $T^*_pM$, i.e.\ it is defined by its action on an arbitrary $X\in T_pM$ giving a real number, \emph{linearly} in $X$. The data are only $\phi$, $X$ and $\omega$. A student suggests: use the push-forward, which takes
$X\in T_pM$ to $\phi_{*,p}(X)\in T_{\phi(p)}N$. Then it remains to let $\omega$ act on it (``you said we let this act on $\omega$; no, the other way round: $\omega$ is an element of $T^*_{\phi(p)}N$ so it can act on this vector, it produces a real number''). It is linear because the push-forward is linear and $\omega$ is linear: sums and factors
pull out. Again it is the only way to define it, and one needs only to know from where to where it maps, and that it is linear \ts{10}{1}{53:57} \ts{10}{1}{54:58} \ts{10}{1}{55:40} \ts{10}{1}{56:42} \ts{10}{1}{57:25} \ts{10}{1}{58:27}.
\begin{intuition}[Motto: vectors are pushed forward, co-vectors are pulled back]\label{int:tb-motto}
Memorise: \emph{vectors are pushed forward, forms (co-vectors, and extensions, to come later) are pulled back}. Only under the special circumstance that $\phi$ is a diffeomorphism can one go back and forth, and then one can also pull vectors back, using the inverse map; but if the inverse does not exist, vectors can generically only be pushed forward and co-vectors
only be pulled back. This is very important and will come back once immersions and embeddings are introduced \ts{10}{1}{58:27} \ts{10}{1}{58:51} \ts{10}{1}{59:17} \ts{10}{1}{59:37}.
\end{intuition}
\begin{exercise}[Pull-back of a gradient]\label{ex:tb-pullback-df}
\emph{Not said in class.} Show that $\phi^*_p(\mathrm d_{\phi(p)}f)=\mathrm d_p(f\circ\phi)$ for $f\in C^\infty(N)$ (\cref{def:tb-gradient,def:tb-pull}).
\end{exercise}

% ==================================================================
\section{Immersions and embeddings}\label{sec:tb-immersion}

The push-forward can be used to decide the question: under which circumstances can a smooth manifold ``sit'' (quotation marks) inside some $\R^N$? \ts{10}{1}{1:00:00} \ts{10}{1}{1:00:28} \ts{10}{1}{1:00:49}. There are two notions of sitting, and $\R^N$ is chosen because it is immediately intuitive; one could generalise to one manifold sitting in another: \textbf{immersion} and
\textbf{embedding}; the embedding is the stronger one: every embedding is an immersion, but an immersion has fewer conditions, and sometimes one is happy with an immersion, sometimes one wants an embedding \ts{10}{1}{1:01:09} \ts{10}{1}{1:01:30} \ts{10}{1}{1:01:51}.
\begin{definition}[New concepts: immersion]\label{def:tb-immersion}
A smooth map $\phi\colon M\to\R^N$ is an \textbf{immersion} of $M$ into $\R^N$ if its derivative is injective at every point: the push-forward $\phi_{*,p}\colon T_pM\to T_{\phi(p)}\R^N$ is injective for all $p\in M$. (Injective: two different vectors are not mapped to the same target vector.)
\ts{10}{1}{1:02:13} \ts{10}{1}{1:02:33} \ts{10}{1}{1:02:56} \ts{10}{1}{1:03:18} \ts{10}{1}{1:03:41} \ts{10}{1}{1:04:02}.
\end{definition}
\begin{example}[Two immersions of the circle]\label{ex:tb-circle}
The technical checks are left as homework; intuitively:
\begin{itemize}
  \item \emph{(a)} $\phi\colon S^1\to\R^2$ takes the circle to a distorted circle in the plane. It is smooth (no corners are produced) and its push-forward is injective: the tangent space of the circle is one dimensional and its vectors, pointing in this or that direction, are mapped to vectors that do not coincide.
        (One could also immerse the circle into $\R^3$, as a loop in space, as long as there are no corners.)
  \item \emph{(b)} $\tilde\phi\colon S^1\to\R^2$ takes the circle to a figure-eight, a Lissajous-like curve (a sine in one direction and a sine with double frequency in the other). It is smooth, and \emph{not injective} as a map of sets: two points of the circle are mapped to the crossing point of the figure, one on each branch. But injectivity of
        the map was not required: the derivative $\tilde\phi_{*,p}$ is injective at every point, since the vectors at the two preimages of the crossing are mapped to different vectors, and no vector is mapped to the same as another. So $\tilde\phi$ is an immersion; ``an immersion is not such a strong condition''
\end{itemize}
\ts{10}{1}{1:04:22} \ts{10}{1}{1:05:25} \ts{10}{1}{1:06:08} \ts{10}{1}{1:07:09} \ts{10}{1}{1:07:51} \ts{10}{1}{1:08:54}.
\end{example}
\begin{definition}[New concepts: embedding]\label{def:tb-embedding}
A smooth map $\phi\colon M\to N$ is an \textbf{embedding} if (a) $\phi$ is an immersion and (b) $M$ is homeomorphic to its image $\phi(M)$, which is required to be a submanifold of the target $N$. (The lecture says: if one maps the entire domain, the image should be topologically isomorphic to the manifold; ``if it is topologically isomorphic and smooth,
it immediately implies that it is also diffeomorphic'', but one only needs to check that the image is again a manifold, a submanifold of the target.) It is then not difficult to see that every embedding is an immersion, by condition (a) \ts{10}{1}{1:08:54} \ts{10}{1}{1:09:19} \ts{10}{1}{1:09:40} \ts{10}{1}{1:10:01} \ts{10}{1}{1:10:22}.
\end{definition}
\begin{coursenote}[Wording of the embedding condition]
The usual definition asks that $\phi$ be a homeomorphism onto its image with the subspace topology, equivalently a diffeomorphism onto an embedded submanifold; the lecture's phrasing (``homeomorphic to the target'', corrected to the image) is loose, and the claim that smooth plus one-to-one implies diffeomorphic is not a general fact
(think of $t\mapsto t^3$, a smooth homeomorphism of $\R$ that is not a diffeomorphism), but for an immersion that is a homeomorphism onto its image it is true. Not said in class.
\end{coursenote}
So example (a) is an embedding: the image, the distorted circle, is again a manifold (a smooth one, since $\phi$ is smooth), ``the idea of an embedding: one can think of the smooth manifold as an immersed submanifold of the bigger space''; whereas example (b) is an immersion that is \emph{not} an embedding: at the crossing point the image is not locally homeomorphic to $\R^1$, it fails to be a
manifold. Sometimes one insists on this being a manifold \ts{10}{1}{1:10:22} \ts{10}{1}{1:10:43} \ts{10}{1}{1:11:04} \ts{10}{1}{1:11:26} \ts{10}{1}{1:11:47}.

% ==================================================================
\section{The Whitney theorems}\label{sec:tb-whitney}

So what is the answer to the original question: under which circumstances can a given manifold sit in $\R^N$ for a suitable $N$? \ts{10}{1}{1:11:47} \ts{10}{1}{1:12:08}
\begin{theorem}[Whitney embedding and immersion theorem]\label{thm:tb-whitney}
Any smooth manifold $M$ can be embedded in $\R^{2\dim M}$ and immersed in $\R^{2\dim M-1}$ \ts{10}{1}{1:12:08} \ts{10}{1}{1:12:39} \ts{10}{1}{1:12:59} \ts{10}{1}{1:13:22} \ts{10}{1}{1:13:43} \ts{10}{1}{1:15:25}.
\end{theorem}
(In class the order was confused: first ``it can be embedded in $\R^{2\dim M}$'', then the immersion ``in $2\dim M-1$'' after the Klein bottle example, once the stronger notion was recalled; the statement above is the one the lecture arrives at, and is the correct one.) The good news: there is no smooth manifold that cannot sit in some $\R^N$. The embedding space must
have in the worst case up to twice the dimension of the manifold; some manifolds can be embedded in $\R^{2\dim M-3}$ or the like, it depends on the manifold, but the worst case, ``I think provided by projective spaces'' [?], needs the double dimension \ts{10}{1}{1:13:43} \ts{10}{1}{1:14:03}.
\begin{example}[Circle, surfaces, Klein bottle]\label{ex:tb-klein}
The circle is one dimensional, so it can be embedded in $\R^2$: no surprise. Any two-dimensional manifold (sphere, torus, \dots) can be embedded in $\R^4$. A two-dimensional manifold that cannot be embedded in $\R^3$, which was on the blackboard before: the \textbf{Klein bottle}; but it can be immersed in $\R^{2\cdot2-1}=\R^3$. The picture: a throat
coming in, passing through the wall of the bottle and joining the other side (``like taking a torus and not gluing it like this but going in somewhere else and joining it from the other side''); in $\R^3$ it intersects itself, like the figure-eight before, which prevents the image from being a manifold, hence from being homeomorphic to the Klein bottle: an immersion, not an embedding.
So Whitney guarantees an embedding in $\R^4$, and the Klein bottle is an example in which one cannot beat that limit; for other manifolds one can obviously beat it: the two-sphere can be embedded in $\R^3$ (in the lecture ``immersed into $\R^4$ for sure, but also into $\R^3$'')
\ts{10}{1}{1:14:03} \ts{10}{1}{1:15:05} \ts{10}{1}{1:16:07} \ts{10}{1}{1:17:11} \ts{10}{1}{1:18:14} \ts{10}{1}{1:19:17}.
\end{example}
\begin{coursenote}[Slips in the statement]
At \ts{10}{1}{1:15:25} the lecture first gives ``embedded in $2\dim M-1$'' for the Klein bottle and then corrects to the immersion; and at \ts{10}{1}{1:22:48}--\ts{10}{1}{1:23:08} ``the strong Whitney theorem says a three-dimensional manifold can be embedded in $\R^5$'' where $2\cdot3-1=5$ is the immersion bound and the embedding bound is $\R^6$. The two-sphere is also in the lecture
described as immersed where embedded is meant. Not said in class.
\end{coursenote}
\begin{secondary}[Weaker, stronger and sharper versions]
The statement above is the \emph{strong} Whitney theorem, which takes some effort to prove; one can prove weaker versions more quickly, for example that any manifold can be embedded in \emph{some} $\R^N$ (``also that is not obvious from the beginning'') \ts{10}{1}{1:19:17} \ts{10}{1}{1:19:37} \ts{10}{1}{1:19:58}.
There are even stronger versions that improve the dimension bound: any smooth manifold $M$ can be \emph{immersed} into $\R^{2\dim M-\alpha(\dim M)}$, where $\alpha(n)$ is the number of ones in the binary expansion of the integer $n\ne0$. Example: $\dim M=3$ is $11$ in binary ($1\cdot2^1+1\cdot2^0$), so $\alpha(3)=2$ and any three-dimensional manifold can be immersed in $\R^{6-2}=\R^4$,
whereas the strong Whitney theorem gives $\R^5$. ``I find it quite amusing.'' In a particular case with even better estimates needed one should go to the literature: there are other bounds under particular additional circumstances which may or may not apply
\ts{10}{1}{1:19:58} \ts{10}{1}{1:20:40} \ts{10}{1}{1:21:22} \ts{10}{1}{1:22:26} \ts{10}{1}{1:23:08} \ts{10}{1}{1:23:53}.
\end{secondary}
\begin{history}[Whitney and the immersion bound]
\emph{Not said in class.} \emph{Checked.} H.~Whitney proved the embedding theorem in $\R^{2n}$ and the immersion theorem in $\R^{2n-1}$ in 1936 (\emph{Differentiable manifolds}, \emph{Ann.\ Math.}\ 37) and the easy embedding in $\R^{2n+1}$ earlier; the refinement to $\R^{2n-\alpha(n)}$ for immersions is the theorem of R.~Cohen (1985), the Whitney--Massey
immersion conjecture.\par
\emph{References.} H.~Whitney, \emph{Ann.\ Math.}\ 37 (1936) 645--680; R.~L.~Cohen, \emph{Ann.\ Math.}\ 122 (1985) 237--328.
\end{history}
The conclusion of the lecture: the Whitney theorem \emph{clarifies} whether the intrinsic approach is more powerful than the extrinsic one: it is not, they are the same. One could now check, although it was not proved, that the tangent spaces constructed intrinsically are vector space isomorphic to the tangent planes (higher dimensional planes, not hyperplanes) one would construct for the embedded manifold. For
this one needs an embedding (an immersion alone would not tell how to attach the tangent plane), but one goes one dimension higher at worst. The two approaches are equivalent; this was argued for a while in the history of differential geometry; the intrinsic approach is more aesthetic, and it also justifies drawing the tangent spaces as one intuitively draws them
\ts{10}{1}{1:23:53} \ts{10}{1}{1:24:14} \ts{10}{1}{1:24:35} \ts{10}{1}{1:24:55} \ts{10}{1}{1:25:16} \ts{10}{1}{1:25:36}.
\begin{recall}[Why intrinsic]
The reasons for the intrinsic construction were given in \cref{int:df-embedding}: no reference to an ambient space.
\end{recall}

% ==================================================================
\section{The tangent bundle}\label{sec:tb-bundle}

Another technically highly important topic: the \textbf{tangent bundle}, associated with which one looks at vector fields \ts{10}{1}{1:25:36} \ts{10}{1}{1:25:56}. So far one looked at a manifold and at the tangent space at one particular point; ``after the Whitney embedding theorem I feel more justified in drawing such childish pictures'' \ts{10}{1}{1:25:56} \ts{10}{1}{1:26:17}.
\begin{definition}[New concepts: tangent bundle, bundle projection]\label{def:tb-bundle}
Let $M$ be a smooth manifold. The \textbf{tangent bundle} of $M$ is the set
\[
  TM:=\bigsqcup_{p\in M}T_pM,
\]
the \emph{disjoint} union of all the tangent spaces: every tangent vector at every point of the manifold, as a set. The \textbf{bundle projection} is the map
\[
  \pi\colon TM\to M,\qquad X\mapsto p\qquad\text{where }X\in T_pM,
\]
sending a tangent vector to its base point \ts{10}{1}{1:26:38} \ts{10}{1}{1:26:59} \ts{10}{1}{1:29:05} \ts{10}{1}{1:29:25} \ts{10}{1}{1:30:07} \ts{10}{1}{1:30:29}.
\end{definition}
In pictures the tangent spaces $T_pM$ and $T_qM$ at different points seem to intersect, but one is no longer fooled: that is an artifact of the embedding; intrinsically no element of $T_qM$ is an element of $T_pM$. The dot on the union sign is just an emphasis. The fact that $\pi$ is a \emph{map} (one base point for each $X$) is a consequence of the disjointness
\ts{10}{1}{1:27:41} \ts{10}{1}{1:28:01} \ts{10}{1}{1:28:22} \ts{10}{1}{1:28:43} \ts{10}{1}{1:30:49} \ts{10}{1}{1:31:10}.
So one has a set $TM$ and a map $TM\to M$; this already looks like a bundle (\cref{def:mf-bundle}), but at this point only a \emph{bundle of sets}, not even a topological space bundle, let alone a topological or smooth manifold bundle \ts{10}{1}{1:31:10} \ts{10}{1}{1:31:31} \ts{10}{1}{1:31:52}.

\begin{remark}[A correction on the example of a non-fibre bundle]\label{rem:tb-correction}
In an earlier lecture the lecturer was asked for an example of a bundle that is not a fibre bundle, i.e.\ where not every fibre is the same as the others, and came up quickly with a manifold with a complex line attached at some points and real lines at the others. As a bundle of \emph{sets}, with the projection $\pi$ on top, it is clearly a bundle (no extra
requirement on $\pi$: one cannot talk about continuity without a topology on both sides) and not a fibre bundle: the fibres $\pi^{-1}(p)$ and $\pi^{-1}(q)$ are not isomorphic as sets, there is no bijection between $\R$ and $\mathbb C$ [?: ``that's tricky, there are proofs that you can, well certainly you can get some fibres here and some fibres there'', the lecturer
realises the cardinalities of $\R$ and $\mathbb C$ are equal]. The emphasis is on the way bundles were defined later, as bundles of \emph{topological manifolds} (\cref{def:mf-fibre-bundle}), where continuity is needed; and the biggest problem is that the total space is not a manifold: if $M$ is one dimensional, at one point the total space would be real two dimensional and at another real
three dimensional. So the example is one of a bundle of sets, but not of manifolds; it is really tricky to construct a non-fibre bundle for manifolds, ``but I think it's possible''. ``That was my sin there''; one should not take the example too seriously, although as a bundle of sets it can be held up \ts{10}{1}{1:32:14} \ts{10}{1}{1:32:56} \ts{10}{1}{1:33:16} \ts{10}{1}{1:33:58} \ts{10}{1}{1:34:19} \ts{10}{1}{1:34:59}.
\end{remark}
\begin{coursenote}[Sets of the same cardinality]
Indeed $\R$ and $\mathbb C$ are in bijection as sets (both have the cardinality of the continuum, \cref{ch:l03}); the example is therefore an example of a bundle whose fibres have different \emph{vector space} dimension, not different cardinality. Not said in class.
\end{coursenote}

\paragraph{From a bundle of sets to a smooth manifold.} How to make this into a bundle of topological manifolds, or of smooth manifolds? The base $M$ is smooth. The simplest would be to take the smooth atlas on $M$ and construct from it a smooth atlas on $TM$: then $TM$ is a manifold. A topology on $TM$: one wants that $\pi$ be continuous; the topology to establish is the
\textbf{initial topology} with respect to $\pi$ (it was on the last problem sheet; recall \cref{ch:l04}); in the picture, that already gives open sets of the form $\pi^{-1}(U)$ \ts{10}{1}{1:34:59} \ts{10}{1}{1:35:20} \ts{10}{1}{1:35:41} \ts{10}{1}{1:36:22} \ts{10}{1}{1:36:43} \ts{10}{1}{1:39:51}.
Idea: construct a smooth atlas on $TM$ from a given smooth atlas $\mathcal A_M$ on $M$. The picture: the manifold with the tangent space attached at every point, which looks funny since the tangent space should be tangent, but one does not need to subscribe to the embedded notion, so one can make a nice picture by putting the tangent spaces up; the whole is $TM$
\ts{10}{1}{1:37:06} \ts{10}{1}{1:37:26} \ts{10}{1}{1:37:47}.
Take a chart $(U,x)\in\mathcal A_M$. The new chart domain is the \emph{pre-image} $\pi^{-1}(U)$ (written $\pi^{-1}$ as usual, ``but that has always this taste of $\pi$ being invertible, which it is not at all''), an open subset of $TM$ in the initial topology \ts{10}{1}{1:37:47} \ts{10}{1}{1:38:07} \ts{10}{1}{1:38:28} \ts{10}{1}{1:39:10} \ts{10}{1}{1:39:30} \ts{10}{1}{1:40:12}.
The new chart map is $\xi\colon\pi^{-1}(U)\to\R^{2\dim M}$ (in the lecture the symbol sounds like ``chi''), a subset of $\R^{2\dim M}$ as image; it is intuitively clear that the dimension doubles, since the tangent space has the same dimension as $M$ and ``locally it is trivial'' \ts{10}{1}{1:40:12} \ts{10}{1}{1:40:32} \ts{10}{1}{1:40:52} \ts{10}{1}{1:41:12} \ts{10}{1}{1:41:33}.
\begin{definition}[New concepts: the induced chart on $TM$]\label{def:tb-chart}
For $X\in\pi^{-1}(U)$ (a point of the tangent bundle, an element of some tangent space) define
\[
  \xi(X):=\Bigl(x^1(\pi(X)),\dots,x^{\dim M}(\pi(X)),\ X^1,\dots,X^{\dim M}\Bigr)\in\R^{2\dim M},
\]
where the $X^a$ are the components in the expansion $X=X^a\,\dfrac{\partial}{\partial x^a}\Big|_{\pi(X)}$ with respect to the basis of the tangent space $T_{\pi(X)}M$ induced by the chart $(U,x)$. The first $\dim M$ coordinates are the coordinates of the base point $\pi(X)\in U$ (project down and use the
original chart: they say in which fibre $X$ lies); the other $\dim M$ say where in the fibre $X$ lies. This is possible because $X$ lies in the tangent space over its own base point, and $\partial/\partial x^a|_{\pi(X)}$ is a basis of it; the $X^a$ are real numbers
\ts{10}{1}{1:41:33} \ts{10}{1}{1:42:35} \ts{10}{1}{1:43:59} \ts{10}{1}{1:45:03} \ts{10}{1}{1:46:25} \ts{10}{1}{1:47:28}.
\end{definition}
This is a chart that assigns to an arbitrary point of the set $\pi^{-1}(U)$ a collection of $2\dim M$ numbers. What remains to be shown: that any two such charts $(\pi^{-1}(U),\xi)$ and $(\pi^{-1}(\tilde U),\tilde\xi)$ are $C^\infty$ compatible; this goes on the problem sheet (``quite simple''). Once this is shown, a smooth structure on $M$ is inherited by $TM$, which
is then a smooth manifold; one can speak of the projection $\pi$ being smooth, ``and then you have a smooth bundle'' \ts{10}{1}{1:47:28} \ts{10}{1}{1:47:50} \ts{10}{1}{1:48:12} \ts{10}{1}{1:48:32}.
\begin{exercise}[The charts of $TM$ are compatible]\label{ex:tb-compat}
Show that the transition map $\tilde\xi\circ\xi^{-1}$ between two induced charts is $C^\infty$: the first $\dim M$ components are $\tilde x\circ x^{-1}$, and the last ones are the Jacobian of $\tilde x\circ x^{-1}$ at the base point applied to $(X^a)$ (\cref{ex:tv-change}), linear in $X^a$. (Problem sheet; the full computation, not given in class, is in
\cref{exam:ch10}.)
\end{exercise}

% ==================================================================
\section*{Key concepts}
\begin{keyconcepts}
\textbf{Retrieval questions.}
\begin{enumerate}
  \item Define $T^*_pM$ and $T^r_s(T_pM)$. Why is $T_pM\cong T^*_pM$ non-canonical, while $V\cong V^{**}$ is canonical?
  \item Define the gradient $\mathrm d_pf$. Why is it a co-vector, and why does ``the gradient is orthogonal to the level sets'' not need an inner product?
  \item Prove that $\mathrm dx^a|_p$ is the dual basis of $\partial/\partial x^b|_p$. What does $\mathrm dx^a$ \emph{not} mean?
  \item Define the push-forward $\phi_{*,p}$ and show that it is forced by the data; prove $\phi_{*,p}X_{\gamma,p}=X_{\phi\circ\gamma,\phi(p)}$.
  \item Define the pull-back. Why does it go in the opposite direction? State the motto.
  \item Define immersion and embedding; give an immersion of the circle that is not an embedding.
  \item State the Whitney theorems, with the Klein bottle as an example and the improved bound $2\dim M-\alpha(\dim M)$. What do they say about the intrinsic versus the extrinsic definition of tangent spaces?
  \item Define the tangent bundle, the bundle projection and the chart on $TM$ induced by a chart on $M$. What remains to be proven?
\end{enumerate}
\textbf{Concepts.}
\begin{itemize}
  \item Cotangent space, tensors at a point (\cref{def:tb-cotangent,def:tb-tensors,rem:tb-noncanonical}); gradient (\cref{def:tb-gradient,rem:tb-covector}); chart-induced dual basis (\cref{def:tb-dx,thm:tb-dualbasis,int:tb-dx}).
  \item Push-forward and pull-back (\cref{def:tb-push,prop:tb-push-curve,def:tb-pull,int:tb-motto}).
  \item Immersion, embedding, Whitney (\cref{def:tb-immersion,def:tb-embedding,ex:tb-circle,thm:tb-whitney,ex:tb-klein}).
  \item Tangent bundle and its smooth structure (\cref{def:tb-bundle,rem:tb-correction,def:tb-chart,ex:tb-compat}).
\end{itemize}
\end{keyconcepts}

\flushside
\end{paracol}

% !TEX root = ../main.tex
\chapter{Tensor space theory II: over a ring}\label{ch:l11}
\begin{paracol}{2}

\begin{sources}
\begin{tabularx}{\linewidth}{@{}l l L@{}}
\toprule
\textbf{Video} & \textbf{Length} & \textbf{Content}\\
\midrule
\seg{11}{1}{Lecture 11: Tensor space theory II: over a ring} & 1:55:10 & Vector fields as smooth sections of $TM$; the pointwise addition and the multiplication by functions; rings and the ring $C^\infty(M)$; modules over a ring; modules over a division ring have a basis (Zorn's lemma, with the student question on upper bounds),
the hairy ball as an example without a basis; free and projective modules; the theorem on sections of vector bundles; $\mathrm{Hom}$ of projective modules over a commutative ring; the textbook definition of tensor fields.\\
\bottomrule
\end{tabularx}

\smallskip
\textbf{Transcript} \file{materials/transcripts/C11-V01.txt} (YouTube automatic captions).
\end{sources}

\begin{reading}[Where this lecture sits in the other notes]
\emph{Gravity and Light}, Lecture~6 (Fields): vector fields and the $C^\infty(M)$-module structure of the set of fields; \emph{Theoretische Mechanik}, Chapter~5 (tensor fields). Both treat vector and tensor fields on a smooth manifold; the emphasis here on rings, modules, Zorn's lemma and the absence of a global basis
is, as far as the titles say, special to this course. (From the lecture titles and the plan of the course; not checked in detail.)
\end{reading}

% ==================================================================
\section{Vector fields and the algebraic structure they carry}\label{sec:rm-fields}

Until now (\cref{ch:l09,ch:l10}) the tangent space of a manifold was considered only at one point; last time the tangent bundle, the disjoint union of all the tangent spaces, was constructed, and it is a \emph{smooth} bundle. That now allows one to define vector fields and, in general, tensor fields. Over one point the tangent space is a vector
space; but for the \emph{collection} of vector fields the structure will turn out to be only a \textbf{module}, not a vector space. ``Most people who apply mathematics consider rings and modules a somewhat esoteric subject'': today it will be seen that the funny properties of modules, compared with vector spaces, are of direct geometric
relevance and make one understand the subject better \ts{11}{1}{0:04} \ts{11}{1}{1:06} \ts{11}{1}{1:31} \ts{11}{1}{1:51}.

\begin{definition}[New concepts: vector field, $\Gamma(TM)$]\label{def:rm-field}
Let $M$ be a smooth manifold and $\pi\colon TM\to M$ its tangent bundle, a smooth bundle (that is what ``the tangent bundle is a smooth bundle'' means: $\pi$ is smooth; \cref{def:tb-bundle}). A \textbf{vector field} on $M$ is a \emph{smooth section} of the tangent bundle: a smooth map $\sigma\colon M\to TM$ with
\[
  \pi\circ\sigma=\mathrm{id}_M
\]
(\cref{def:mf-section}). The set of all vector fields is denoted $\Gamma(TM)$, the standard notation for the set of all sections of a bundle (strictly one should write $\Gamma(TM,\pi,M)$, since all three data enter; this is short notation) \ts{11}{1}{2:11} \ts{11}{1}{2:59} \ts{11}{1}{3:19} \ts{11}{1}{3:40} \ts{11}{1}{6:06} \ts{11}{1}{6:27} \ts{11}{1}{7:08} \ts{11}{1}{7:29} \ts{11}{1}{8:18} \ts{11}{1}{8:59}.
\end{definition}
A vector field assigns to each point $p\in M$ an element $\sigma(p)\in T_pM$: because $\pi\circ\sigma=\mathrm{id}_M$ the vector provided at $p$ lies in the tangent space \emph{at $p$}, not in the tangent space at another point $q$; at another point another vector is provided. One could not simply write ``a map $M\to T_pM$'': one does not know at which
point one is, and one would have to say what it means for a vector field to be smooth, that is how one differentiates vector fields, which would need a new notion of differentiation. Instead, $M$ and $TM$ are smooth manifolds and one knows what a smooth map between smooth manifolds is
(\cref{ch:l07}): this is the whole idea of the definition \ts{11}{1}{4:22} \ts{11}{1}{4:43} \ts{11}{1}{5:04} \ts{11}{1}{5:24} \ts{11}{1}{5:45}.
From now on ``vector field'' always means smooth vector field; on a topological manifold one could only speak of continuous ones (a continuous map into the then topological manifold $TM$), and at whatever level of structure one is, the corresponding notion of field is obtained in the same way \ts{11}{1}{6:27} \ts{11}{1}{6:47}.

\begin{coursenote}[Notation $\Gamma$ asked in class]
A student asked whether $\Gamma(TM)$ is ``the collection of all $\sigma$'': yes, the set of all sections, those maps with $\pi\circ\sigma=\mathrm{id}_M$; the symbol carries the total space only because the bundle map and the base are understood. That the collection is a set at all (not a class) is clear since
sections are defined as maps, which are subsets of Cartesian products. \ts{11}{1}{8:18}
\end{coursenote}

\paragraph{Two operations on $\Gamma(TM)$.} The first one is harmless. For the time being call it $\oplus$:
\[
  \oplus\colon\Gamma(TM)\times\Gamma(TM)\to\Gamma(TM),\qquad(\sigma\oplus\tau)(p):=\sigma(p)+\tau(p),
\]
defined \textbf{pointwise}: at each $p\in M$ one adds the two vectors in the vector space $T_pM$ (the $+$ on the right is the one of $T_pM$, which makes sense because of the section property). No secret here: take two vector fields, look at every point and add the two vectors there; the collection of results is again a vector field
\ts{11}{1}{9:20} \ts{11}{1}{9:41} \ts{11}{1}{10:01} \ts{11}{1}{10:21}.

The second one is more interesting. It is a kind of scalar multiplication, but one may ask by what. One can of course multiply a vector field by a real number, but ``let us not be so childish'': one can multiply it by an entire \emph{function} on the manifold,
\[
  \odot\colon C^\infty(M)\times\Gamma(TM)\to\Gamma(TM),\qquad(f\odot\sigma)(p):=f(p)\cdot\sigma(p),
\]
again pointwise, with the scalar multiplication of $T_pM$ on the right. The function has different values at different points and scales the vector of the field at each point by the value it has \emph{there}; there is no need for one global real number that doubles or triples the whole field. ``It is very innocent, and it has
dramatic effects in what follows'' \ts{11}{1}{10:42} \ts{11}{1}{11:03} \ts{11}{1}{11:24} \ts{11}{1}{11:45} \ts{11}{1}{12:05} \ts{11}{1}{12:26}.
\begin{coursenote}[Smoothness of the results]
That $\sigma\oplus\tau$ and $f\odot\sigma$ are again \emph{smooth} sections is not discussed in class; in the induced charts of $TM$ (\cref{def:tb-chart}) the components of $\sigma\oplus\tau$ and $f\odot\sigma$ are $\sigma^a+\tau^a$ and $f\sigma^a$, smooth if $f,\sigma^a,\tau^a$ are. Not said in class.
\end{coursenote}

The set $C^\infty(M)$ has appeared in several forms (\cref{def:tv-cinf}). As a real vector space, multiplied by real numbers (the small dot). But here something else is meant: the multiplication of two functions, pointwise, which (\cref{ex:tv-cinf-alg}) makes $C^\infty(M)$ a \emph{ring}; two functions can also be added pointwise. So the structure on $\Gamma(TM)$ is an
addition and a multiplication, but the scalars, the elements of $C^\infty(M)$, are \emph{not the elements of a field, only of a ring}. That could seem harmless, and it is not \ts{11}{1}{12:26} \ts{11}{1}{12:46} \ts{11}{1}{13:07} \ts{11}{1}{13:28} \ts{11}{1}{13:48} \ts{11}{1}{14:29} \ts{11}{1}{14:51}.

% ==================================================================
\section{Rings and modules}\label{sec:rm-rings}

\begin{recall}[Rings]
The definition of ring and the examples were given in \cref{def:te-ring,ex:te-rings}; here only the distinctions are needed again, in the form of the lecture: a ring has $+$ and $\cdot$ (both maps $R\times R\to R$), $(R,+)$ satisfies the commutative, associative, neutral and inverse laws, and the distributive law links the two.
\end{recall}
The \textbf{distributive law} was not written down last time (``a self-correcting mistake''): it links the two operations, $a\cdot(b+c)=a\cdot b+a\cdot c$ and the like. For a ring, in contrast to a field, the following may be missing, quite independently of each other (almost: see the third point) \ts{11}{1}{17:58} \ts{11}{1}{18:19} \ts{11}{1}{18:42}:
\begin{enumerate}
  \item the commutativity of $\cdot$: if it holds, the ring is a \textbf{commutative ring};
  \item the neutral element of $\cdot$: if it exists, the ring is a \textbf{unital} ring; only unital rings will be considered here, but there are rings without a unit;
  \item the multiplicative inverse of every non-zero element (``zero'' being the neutral element of $+$): if it exists, the ring is a \textbf{division ring}; a non-commutative structure which is otherwise a field is a division ring (a ``skew field''; the lecture recalls the term only vaguely: [?]).
\end{enumerate}
A division ring is necessarily unital, since the inverse is defined through the neutral element of the product: $a\cdot b=1$. A field is a commutative division ring \ts{11}{1}{19:02} \ts{11}{1}{19:23} \ts{11}{1}{19:43} \ts{11}{1}{20:04} \ts{11}{1}{20:25} \ts{11}{1}{20:45} \ts{11}{1}{21:06}.

\begin{proposition}[$C^\infty(M)$ is a commutative unital ring, not a division ring]\label{prop:rm-cinf-ring}
$\bigl(C^\infty(M),+,\bullet\bigr)$ with the pointwise operations is a commutative unital ring that is not a division ring \ts{11}{1}{21:26}.
\end{proposition}
\begin{proof}
\emph{Commutative}: the pointwise product inherits the commutativity of the product in $\R$.
\emph{Unital}: the neutral element is the constant function $1$ (every point goes to the number $1$), which is smooth, and $1\bullet f=f$.
\emph{Not a division ring}: let $f\ne0$ vanish at one point $p$ (without being the zero function, which is the neutral element of $+$). There is then no $g$ with $f\bullet g=1$, since $f(p)g(p)=0\ne1$:
the pointwise multiplication cannot ``get rid of the zero'' at $p$ \ts{11}{1}{21:47} \ts{11}{1}{22:08} \ts{11}{1}{22:29} \ts{11}{1}{22:49} \ts{11}{1}{23:09}.
\end{proof}
So one is immediately dealing with a commutative and unital ring, ``but not, and that is of key importance, a so-called division ring'' \ts{11}{1}{23:09} \ts{11}{1}{23:30}.
\begin{exercise}[Zero and non-zero functions]\label{ex:rm-zero}
Show that on $M=\R$ the function $f(x)=x$ has no inverse in $C^\infty(\R)$, although it is non-zero. Is $f$ invertible on $M=\R\setminus\{0\}$? Which functions in $C^\infty(M)$ \emph{are} invertible, for any $M$? \emph{Not said in class.}
\end{exercise}

\begin{definition}[New concepts: $R$-module]\label{def:rm-module}
Let $R$ be a ring. An \textbf{$R$-module} is a set $\mathcal M$ (in class written $M$, ``M like module'', ``I could choose a different letter'', although $M$ is the manifold; here $\mathcal M$) with an addition and a scalar multiplication
\[
  \oplus\colon\mathcal M\times\mathcal M\to\mathcal M,\qquad\odot\colon R\times\mathcal M\to\mathcal M,
\]
which satisfy the vector space axioms for $\oplus$ (commutative, associative, neutral element, inverse), the axioms for the interplay of the two (associativity $(r\cdot s)\odot m=r\odot(s\odot m)$ and the distributive laws), and where the scalars are taken from the ring $R$ instead of from a field \ts{11}{1}{23:52} \ts{11}{1}{24:14} \ts{11}{1}{24:35} \ts{11}{1}{24:55} \ts{11}{1}{25:16} \ts{11}{1}{25:37}.
\end{definition}
So it is a ``vector space over a ring'', which would be perfect terminology; but modules behave so differently from vector spaces that they received their own name \ts{11}{1}{25:37} \ts{11}{1}{25:58}.
\begin{coursenote}[Left and right modules, and the unit]
With the multiplication from the left one speaks of a \emph{left} $R$-module; one can define \emph{right} modules, multiplying from the right, and a bimodule if both are given compatibly: ``let us not go into this, it is not necessary'', and over a commutative ring such as $C^\infty(M)$ the distinction disappears. For a unital ring one also asks $1\odot m=m$ for all $m$; this axiom is used implicitly below
(``it is a unital ring, $v=1\cdot v$''). \ts{11}{1}{24:55} \ts{11}{1}{1:12:42}
\end{coursenote}
The motivating example: the set $\Gamma(TM)$ with the pointwise operations $\oplus$ and $\odot$ is a $C^\infty(M)$-module. ``The tricky thing in the module is not the module structure, it is the ring structure that underlies it'' \ts{11}{1}{26:19} \ts{11}{1}{26:40}.
Why not restrict to real numbers, which would give back a real vector space? One could, but then a basis would always be (at best) \emph{uncountable} even for simple examples, as shown below; roughly, one wants to provide the coefficients of a vector field as \emph{functions}, which change from point to point, as fields in the physicist's sense do \ts{11}{1}{16:36} \ts{11}{1}{16:57} \ts{11}{1}{17:17} \ts{11}{1}{17:37}.

% ==================================================================
\section{Modules need not have a basis}\label{sec:rm-basis}

The key fact: unlike a vector space, a module generically does \emph{not} have a basis \ts{11}{1}{27:00} \ts{11}{1}{27:23}. If the underlying ring is a division ring there is not much difference: then a module is, in all its structural properties, like a vector space (the lack of commutativity plays a role in some calculations, but not in the
structure) and this will be proven. But $C^\infty(M)$ is not a division ring, so $\Gamma(TM)$ in general has no basis, and there are geometric examples \ts{11}{1}{27:43} \ts{11}{1}{28:04} \ts{11}{1}{28:24} \ts{11}{1}{28:45} \ts{11}{1}{29:05} \ts{11}{1}{29:26}.

\begin{theorem}[Modules over a division ring have a basis]\label{thm:rm-basis}
Let $D$ be a division ring and $V$ a $D$-module. Then $V$ has a basis \ts{11}{1}{30:51} \ts{11}{1}{31:17} \ts{11}{1}{31:39}.
\end{theorem}
\begin{corollary}[Every vector space has a basis]\label{cor:rm-vs}
Every field is a division ring, hence every vector space has a basis (a Hamel basis, \cref{def:te-basis}) \ts{11}{1}{31:39} \ts{11}{1}{32:00}.
\end{corollary}
Here a \textbf{basis} of a module $V$ over a ring $R$ is, as for vector spaces, a subset $B\subseteq V$ such that every element of $V$ can be written \emph{uniquely} as a finite linear combination $\sum r^ie_i$ with $e_i\in B$, $r^i\in R$ \ts{11}{1}{37:39}. The theorem will be proven below; first two examples, one with and one without a basis.

\begin{example}[Vector fields on $\R^2$: a basis exists]\label{ex:rm-r2}
Let $M=\R^2$. The set $\Gamma(T\R^2)$ is a $C^\infty(\R^2)$-module; the ring is not a division ring, yet it \emph{has} a basis: ``I never said that without a division ring there is no basis, only that with one you are guaranteed to have one''. Namely there exist two vector \emph{fields} $E_1,E_2$ (not vectors), for example the constant fields
along the two axes, such that for every vector field $\sigma$ there are two coefficient functions with $\sigma=\sigma^1\odot E_1\oplus\sigma^2\odot E_2$: at each point one scales $E_1$ and $E_2$ differently to produce the value of $\sigma$ there. This is the standard intuition about vector fields and it is fine
\ts{11}{1}{32:41} \ts{11}{1}{33:03} \ts{11}{1}{33:26} \ts{11}{1}{33:47} \ts{11}{1}{34:07} \ts{11}{1}{34:28} \ts{11}{1}{34:49} \ts{11}{1}{35:11} \ts{11}{1}{35:32} \ts{11}{1}{35:54} \ts{11}{1}{36:15} \ts{11}{1}{36:36} \ts{11}{1}{36:56}.
\end{example}

\begin{example}[Vector fields on $S^2$: no basis]\label{ex:rm-s2}
Let $M=S^2$. A basis means that a decomposition with \emph{unique} coefficient functions exists (otherwise one only has a generating system). The claim is that $\Gamma(TS^2)$, as a $C^\infty(S^2)$-module, has \textbf{no basis}. It is due to a famous and very intuitive result: \emph{you cannot comb a hairy sphere}. Imagine the sphere with a hair at
every point, none of length zero (``I am an ape, fully haired''); try to comb it smoothly so that the hair gives a smooth vector field, nowhere zero: there is always at some point a curl (a parting), and there the field is not smooth, or has to vanish; with a second curl, the same. Formally: there is no everywhere non-zero and everywhere
smooth vector field on $S^2$ (the hairy ball theorem) \ts{11}{1}{36:56} \ts{11}{1}{37:16} \ts{11}{1}{37:39} \ts{11}{1}{37:59} \ts{11}{1}{38:20} \ts{11}{1}{38:41} \ts{11}{1}{39:02} \ts{11}{1}{39:22} \ts{11}{1}{39:43} \ts{11}{1}{40:03} \ts{11}{1}{40:24}.
How does this relate to a basis? If there is not even one everywhere non-zero smooth vector field, much less are there two that are linearly independent at every point, such that every vector field on the sphere could be uniquely decomposed in terms of them. A basis element cannot vanish; and if a candidate $E_1$ vanishes at some
point while $E_2$ happens to be non-zero there, a vector field with a component along a ``second direction'' at that point cannot be produced by $E_1$, since $E_1$ is zero there, and the coefficient function does not help. So the module over the non-division ring $C^\infty(S^2)$ has no basis: it is no surprise that here the module has no basis, whereas on $\R^2$, thanks to the special structure
of $\R^2$, everything was fine \ts{11}{1}{41:01} \ts{11}{1}{41:22} \ts{11}{1}{41:42} \ts{11}{1}{42:03} \ts{11}{1}{42:24} \ts{11}{1}{42:45} \ts{11}{1}{43:05}.
\end{example}
\begin{secondary}[A complete argument for the sphere]
\emph{Not said in class.} The lecture's argument is an intuition; here is a way to make it a proof, using the hairy ball theorem. Suppose $\Gamma(TS^2)$ had a basis $E_1,\dots,E_n$. Then every vector $v\in T_pS^2$ is a combination of the $E_i(p)$ (extend $v$ to a field with a bump function), so $E_i(p)$ span $T_pS^2$ for every $p$, and
$\Phi\colon S^2\times\R^n\to TS^2$, $(p,c)\mapsto\sum c^iE_i(p)$ is a surjective bundle map of constant rank. If $n>2$ its kernel is a vector bundle of positive rank, with a non-zero smooth section $c(p)$ (bump functions); then $\sum c^i\odot E_i=0$ with $c\ne0$ contradicts the uniqueness of the coefficients. So $n=2$,
$\Phi$ is an isomorphism, $TS^2$ is trivial and $E_1$ is an everywhere non-zero vector field, contradicting the hairy ball theorem.
\end{secondary}
\begin{coursenote}[Sets and sequences: bases]
At the end of the proof below the lecturer notes that a basis is really not a subset but more like a \emph{sequence} (a family): if the same element appears twice the elements are linearly dependent, whereas as a set $\{v,v\}=\{v\}$ is independent; and if one speaks of ``the first, second component'' an order is used, which a subset does not have. ``Don't get confused by this''; it
is a notational issue \ts{11}{1}{25:51} \ts{11}{1}{26:12} \ts{11}{1}{26:32}.
\end{coursenote}

\paragraph{Why not use real coefficients only.} One might retreat and consider $\Gamma(TM)$ simply as a real vector space, which it is; then it has a basis. But the coefficients in a decomposition would be real numbers. On $\R^2$, the pair $E_1,E_2$ would no longer be a basis: only fields with the \emph{same} combination of $E_1$ and $E_2$ at every point could be expressed, such as constant
fields. One would need uncountably many basis elements, which is not tractable. Therefore one wants coefficient \emph{functions} for tensor fields, and has to deal with this situation \ts{11}{1}{43:49} \ts{11}{1}{44:10} \ts{11}{1}{44:30} \ts{11}{1}{44:50}.
\medskip
Now the situation is understood algebraically; this amounts to the proof of \cref{thm:rm-basis}. The proof needs Zorn's lemma, the incarnation in which the \emph{axiom of choice} (\cref{sec:se-choice}) will be used. Everything that now comes is of real importance, the techniques and notions in it being ubiquitous in mathematics \ts{11}{1}{45:11} \ts{11}{1}{45:33} \ts{11}{1}{45:53} \ts{11}{1}{46:16}.

% ==================================================================
\section{Zorn's lemma}\label{sec:rm-zorn}

Zorn's lemma is in fact \emph{equivalent} to the axiom of choice in Zermelo--Fraenkel set theory (without choice): it is just another way to express it, but in a form that is very often useful. It is stated first, and the terms are explained afterwards \ts{11}{1}{46:16} \ts{11}{1}{46:37} \ts{11}{1}{46:57}.

\begin{theorem}[Zorn's lemma]\label{thm:rm-zorn}
Let $(P,\le)$ be a partially ordered set in which every totally ordered subset $T\subseteq P$ has an upper bound in $P$. Then $P$ contains a maximal element \ts{11}{1}{47:18} \ts{11}{1}{47:41} \ts{11}{1}{48:01} \ts{11}{1}{48:22}.
\end{theorem}
The proof is not given; in ZF set theory (without choice) the axiom of choice is equivalent to Zorn's lemma \ts{11}{1}{59:36} \ts{11}{1}{59:58}.
(``Apart from four undefined terms it is very enlightening'': partially ordered set, totally ordered subset, upper bound, maximal element.)

\begin{definition}[New concepts: partially ordered set]\label{def:rm-poset}
A \textbf{partially ordered set} is a set $P$ with a relation $\le$ (``keep an open mind about what this less-or-equal is'': it is not necessarily about numbers) which is
\begin{enumerate}
  \item \emph{reflexive}: $a\le a$ for all $a\in P$;
  \item \emph{antisymmetric} (not symmetric, as for equivalence relations): $a\le b$ and $b\le a$ imply $a=b$;
  \item \emph{transitive}: $a\le b$ and $b\le c$ imply $a\le c$ \ts{11}{1}{48:43} \ts{11}{1}{49:05} \ts{11}{1}{49:26} \ts{11}{1}{49:47} \ts{11}{1}{50:07}.
\end{enumerate}
\end{definition}
\begin{definition}[New concepts: totally ordered subset]\label{def:rm-total}
A subset $T\subseteq P$ is \textbf{totally ordered} if the restriction of $\le$ to $T$ satisfies: (a) \emph{totality}: for all $a,b\in T$, $a\le b$ or $b\le a$ (inclusive ``or''); (b) antisymmetry; (c) transitivity. Reflexivity is not required separately: it follows from totality
(take $a=b$) \ts{11}{1}{50:28} \ts{11}{1}{50:51} \ts{11}{1}{51:11} \ts{11}{1}{51:35} \ts{11}{1}{51:56} \ts{11}{1}{52:17}.
\end{definition}
In a partially ordered set two elements can stand in no relation to each other; in a totally ordered one any two elements are related one way or the other \ts{11}{1}{52:37} \ts{11}{1}{52:58}. A student asks whether the complex numbers are partially ordered. A partial order on $\C$ can certainly be defined: relate only objects of the real line by the standard $\le$ (nothing else is related). That is reflexive,
antisymmetric and transitive, hence a partial order, though not a total one on $\C$; restricted to $\R$ it is a total order since any two real numbers are related \ts{11}{1}{53:25} \ts{11}{1}{53:47} \ts{11}{1}{54:07} \ts{11}{1}{54:28} \ts{11}{1}{54:49}.

\begin{definition}[New concepts: upper bound, maximal element]\label{def:rm-upper}
Let $T\subseteq P$. An element $u\in P$ is an \textbf{upper bound} of $T$ if $t\le u$ for all $t\in T$ (the relation is that of $P$; $u$ comes from $P$, not necessarily from $T$). An element $m\in P$ is \textbf{maximal} if there is no $x\in P$ with $m\le x$
\ts{11}{1}{55:10} \ts{11}{1}{55:31} \ts{11}{1}{55:52} \ts{11}{1}{56:16} \ts{11}{1}{56:42} \ts{11}{1}{57:06} \ts{11}{1}{57:31}.
\end{definition}
\begin{coursenote}[Maximal element: a slip]
As written, with $x$ arbitrary, the definition of maximal is wrong: reflexivity gives $m\le m$, so no element would be maximal; the right statement is ``there is no $x\in P$ with $m\le x$ and $x\ne m$'' (equivalently: $m\le x$ implies $x=m$). Not said in class. The remark that follows in class is correct: the tempting alternative ``$u\le m$ for all $u\in P$'' is a different notion
(a \emph{greatest} element) and in a partial order is stronger: it moves the negation incorrectly, since a negation in front of $\le$ can be flipped only when any two elements are related, which holds for a total, not for a partial order. In the class the lecturer found the same error in his own preparation
(``very honest''): with a partial order there can be elements of $P$ unrelated to $m$ \ts{11}{1}{57:52} \ts{11}{1}{58:12} \ts{11}{1}{58:33} \ts{11}{1}{58:54} \ts{11}{1}{59:15}.
\end{coursenote}

\paragraph{The student question on upper bounds.} A student asks whether the upper bound $u$ has to be in $T$ or may be anywhere in $P$ (and raises the case that $u$ may then be unrelated to the elements of $T$). The lecturer first tries the modified definition ($u\in T$, to play it safe), checks whether it suits the application below (where the upper bound is the union of the chain, and he is not sure it lies in $T$ [?]), draws a diagram and thereupon \emph{returns to the original definition}. The reasoning: in a poset $P$ with a chain $T$ and some element $u\in P$ unrelated to $T$, then $u$ is simply not an upper bound of $T$ (the condition $t\le u$ for all $t$ fails), but
another element which is related to every $t\in T$ is, and that is no problem; the requirement ``$T$ totally ordered'' ensures that elements of $T$ can be compared with the bound. The original definition (with $u\in P$) is the \emph{weaker} hypothesis (more chains have upper bounds), hence Zorn's lemma is stronger and in the application one wants the freedom to
have the upper bound outside $T$ \ts{11}{1}{1:00:22} \ts{11}{1}{1:00:45} \ts{11}{1}{1:01:06} \ts{11}{1}{1:01:33} \ts{11}{1}{1:01:59} \ts{11}{1}{1:02:19} \ts{11}{1}{1:02:40} \ts{11}{1}{1:03:01} \ts{11}{1}{1:03:22} \ts{11}{1}{1:03:43} \ts{11}{1}{1:04:07} \ts{11}{1}{1:04:30} \ts{11}{1}{1:04:51} \ts{11}{1}{1:05:15} \ts{11}{1}{1:05:36}
\ts{11}{1}{1:05:57} \ts{11}{1}{1:06:18} \ts{11}{1}{1:06:38} \ts{11}{1}{1:06:59} \ts{11}{1}{1:07:20} \ts{11}{1}{1:07:44} \ts{11}{1}{1:08:04} \ts{11}{1}{1:08:25} \ts{11}{1}{1:08:46} \ts{11}{1}{1:09:11} \ts{11}{1}{1:09:31} \ts{11}{1}{1:09:53}.

% ==================================================================
\section{Proof that modules over a division ring have a basis}\label{sec:rm-proof}

\begin{proof}[Proof of \cref{thm:rm-basis}]
\emph{Step 1: a generating system.} Let $S\subseteq V$ be a \textbf{generating system}: for every $v\in V$ there are finitely many $e_1,\dots,e_n\in S$ and $v^1,\dots,v^n\in D$ with
\[
  v=v^ie_i.
\]
This is no basis yet: nothing is said about uniqueness. A generating system exists: take $S=V$, since $v=1\cdot v$ in the unital ring $D$. It may be ``fantastically huge'', but the question is whether a basis can be extracted from it \ts{11}{1}{1:10:16} \ts{11}{1}{1:10:37} \ts{11}{1}{1:10:59} \ts{11}{1}{1:11:20} \ts{11}{1}{1:11:40} \ts{11}{1}{1:12:00} \ts{11}{1}{1:12:20} \ts{11}{1}{1:12:42} \ts{11}{1}{1:13:02} \ts{11}{1}{1:13:24}.

\emph{Step 2: a partially ordered set.} Let
\[
  P:=\{U\in\mathcal P(S)\mid U\text{ linearly independent}\},
\]
partially ordered by inclusion $\subseteq$ (two subsets either are contained in one another or neither is; no problem, only a partial order is needed). Linear independence means that every finite subset $\{e_1,\dots,e_n\}$ of $U$ satisfies $a^ie_i=0\Rightarrow a^1=\dots=a^n=0$ \ts{11}{1}{1:13:45} \ts{11}{1}{1:14:08} \ts{11}{1}{1:14:29} \ts{11}{1}{1:14:50} \ts{11}{1}{1:15:12} \ts{11}{1}{1:15:33} \ts{11}{1}{1:15:53}.

\emph{Step 3: Zorn.} Let $T\subseteq P$ be an \emph{arbitrary} totally ordered subset (no further specification, to be able to apply Zorn's lemma). Then the union $\bigcup T$ of its elements is an upper bound of $T$ (every member of $T$ is contained in it) and is linearly independent: a finite subset of $\bigcup T$ lies in one member of $T$ (the members are totally
ordered by inclusion, so the largest of the finitely many members containing the elements is one of them), which is independent. So $\bigcup T\in P$. By Zorn's lemma $P$ has a maximal element $B$ (possibly several; one is picked: here the axiom of choice enters). By construction $B$ is a \emph{maximal linearly independent subset} of $S$ in the sense of inclusion: nothing can be
added without losing independence \ts{11}{1}{1:16:14} \ts{11}{1}{1:16:36} \ts{11}{1}{1:16:57} \ts{11}{1}{1:17:17} \ts{11}{1}{1:17:38} \ts{11}{1}{1:17:58} \ts{11}{1}{1:18:19} \ts{11}{1}{1:18:40} \ts{11}{1}{1:19:02} \ts{11}{1}{1:19:23} \ts{11}{1}{1:19:43} \ts{11}{1}{1:20:05} \ts{11}{1}{1:20:26} \ts{11}{1}{1:20:47}.

\emph{Step 4: $B$ spans $S$.} The claim is $S\subseteq\mathrm{span}_DB$ (finite linear combinations with coefficients in $D$) \ts{11}{1}{1:21:08} \ts{11}{1}{1:21:30} \ts{11}{1}{1:21:50}. Let $v\in S$ be arbitrary. If $v\in B$ there is nothing to show. If not, $B\cup\{v\}$ strictly contains the maximal $B$, so it is linearly dependent: there are
$e_1,\dots,e_n\in B$ and $a^1,\dots,a^n,b\in D$, not all zero, with
\[
  a^ie_i+b\,v=0.
\]
Here $b\ne0$: if $b=0$ then $a^ie_i=0$ with not all $a^i$ zero, contradicting the linear independence of $B$. Since $D$ is a division ring, $b$ has a multiplicative inverse, and
\[
  v=-\bigl(b^{-1}a^i\bigr)e_i,
\]
a finite linear combination of elements of $B$ (``this is where we use the division ring''; for non-commutative $D$ the order of the factors matters and the scalars act from the left) \ts{11}{1}{1:22:14} \ts{11}{1}{1:22:36} \ts{11}{1}{1:22:57} \ts{11}{1}{1:23:17} \ts{11}{1}{1:23:41} \ts{11}{1}{1:24:03} \ts{11}{1}{1:24:26} \ts{11}{1}{1:24:46} \ts{11}{1}{1:25:07} \ts{11}{1}{1:25:31} \ts{11}{1}{1:27:12} \ts{11}{1}{1:27:33} \ts{11}{1}{1:27:54} \ts{11}{1}{1:28:16} \ts{11}{1}{1:28:36} \ts{11}{1}{1:28:58} \ts{11}{1}{1:29:20} \ts{11}{1}{1:29:45} \ts{11}{1}{1:30:06} \ts{11}{1}{1:30:26} \ts{11}{1}{1:30:49}.

\emph{Step 5: conclusion.} $V$ is spanned by its generating system $S$, and $S$ is spanned by $B$, so $B$ generates $V$; it is linearly independent, hence unique decompositions exist (the difference of two decompositions is a vanishing combination). Thus $B$ is a basis of $V$, ``a basis, not the unique basis''; Zorn's lemma only guarantees that some maximal element exists
\ts{11}{1}{1:31:09} \ts{11}{1}{1:31:31} \ts{11}{1}{1:31:51}.
\end{proof}
\begin{coursenote}[Details left out in class]
(i) The empty chain $T=\emptyset$ has $\emptyset\in P$ as an upper bound: this is why Zorn's lemma in the form above needs $P\ne\emptyset$, which is guaranteed since $\emptyset\in P$. (ii) The lecturer's ``union of independent sets'' is exactly the finite-subset argument given in Step 3. (iii) The case $v\in B$ in Step 4 was discussed in class in
the form that $B\cup\{v\}=B$ is still independent, whence the remark on sets and sequences above. Not said in class. \ts{11}{1}{1:25:31}
\end{coursenote}
So every vector space has a Hamel basis (the construction may use infinite sets, but one speaks of finite linear combinations; again for vector spaces \emph{without} further structure such as a topology or a norm), and more generally every module over a division ring has a basis. But $C^\infty(M)$ is not a division ring, hence $\Gamma(TM)$ regarded as $C^\infty(M)$-module is
\emph{not guaranteed} to have a basis: it has one for $\R^2$, none for $S^2$ \ts{11}{1}{1:32:13} \ts{11}{1}{1:32:34} \ts{11}{1}{1:32:58} \ts{11}{1}{1:33:22} \ts{11}{1}{1:33:44}.

\begin{intuition}[No global component functions]\label{int:rm-nocomp}
If one thinks of vectors and vector fields always in terms of their components, one has the illusion that there are always components; but components refer to a basis, and if the module of vector fields has no basis there are \emph{no component functions}. One must stop thinking of tensor fields, in their most general form, in terms of component functions: ``there
are no such component functions in general, full stop''. Locally this is possible: on a patch of the sphere there are component functions, one changes chart and gets component functions on every chart, and that is what is done all the time; but there is not necessarily a global one, and this has nothing to do with choosing charts: abstractly no such thing exists. The formalism, which looked like formalism for
formalism's sake, has a geometric meaning \ts{11}{1}{1:34:04} \ts{11}{1}{1:34:25} \ts{11}{1}{1:34:45} \ts{11}{1}{1:35:07}.
\end{intuition}

% ==================================================================
\section{Free and projective modules}\label{sec:rm-projective}

Some terminology, which emerges here \ts{11}{1}{1:35:07}.

\begin{definition}[New concepts: free module]\label{def:rm-free}
An $R$-module is called \textbf{free} if it has a basis. This does not require $R$ to be a division ring: for example $\Gamma(T\R^2)$ is free, $\Gamma(TS^2)$ is not free. Unless $R$ is a division ring a module over it is not guaranteed to be free; sometimes it is, sometimes it is not \ts{11}{1}{1:35:28} \ts{11}{1}{1:35:50} \ts{11}{1}{1:36:11}.
\end{definition}
Free modules can be described via exact and split exact sequences, which will not be gone into here (``we may come back to it later in terms of cohomology''): a free module has a certain property in terms of such sequences, but there is another type of module that shares exactly this property without being free \ts{11}{1}{1:36:11} \ts{11}{1}{1:36:32} \ts{11}{1}{1:36:53}.
\begin{definition}[New concepts: projective module]\label{def:rm-proj}
An $R$-module $\Gamma$ is called \textbf{projective} if it is a direct summand of a free module $F$: there is an $R$-module $Q$ such that
\[
  \Gamma\oplus Q\ \text{ is free}\ (=F).
\]
(The name is not explained in the lecture; it comes from the exact-sequence property just mentioned, not said in class.) Clearly every free module is projective (take $Q=\{0\}$): \emph{free implies projective} \ts{11}{1}{1:36:53} \ts{11}{1}{1:37:15} \ts{11}{1}{1:37:37} \ts{11}{1}{1:37:58} \ts{11}{1}{1:38:18}.
\end{definition}
Projective modules generalise free modules. A finitely generated free $R$-module $F$ is isomorphic to a direct sum of copies of the ring, $F\cong R\oplus\dots\oplus R$, exactly as for vector spaces: a finite dimensional vector space is the direct sum of $\dim$ many copies of the field. (As a set the direct sum is a Cartesian product, with the linear structure working
componentwise, as known from the problem sheet.) A projective module fails to be free when the summand $Q$ cannot be chosen trivial: only after adding another module is it free \ts{11}{1}{1:38:38} \ts{11}{1}{1:39:01} \ts{11}{1}{1:39:21} \ts{11}{1}{1:39:42} \ts{11}{1}{1:40:03}.

The reason for all of this is a theorem, quoted without proof; the lecture attributes it to Swan and others, with extensions by others (the \emph{Serre--Swan theorem}).
\begin{theorem}[Serre--Swan, in the form quoted]\label{thm:rm-serre-swan}
The set of all smooth sections of a vector bundle $E\to M$ over a smooth manifold $M$ (a bundle whose fibres are vector spaces) is a \emph{finitely generated projective} $C^\infty(M)$-module, denoted $\Gamma(E)$ \ts{11}{1}{1:40:27} \ts{11}{1}{1:40:47} \ts{11}{1}{1:41:11} \ts{11}{1}{1:41:31} \ts{11}{1}{1:41:52} \ts{11}{1}{1:42:19}.
\end{theorem}
Examples of vector bundles: the tangent bundle, the cotangent bundle (every fibre the cotangent space), a tensor bundle (at every point a tensor space, itself a vector space): ``it is very general''. So $\Gamma(E)\oplus Q$ is free for some $C^\infty(M)$-module $Q$; the module $\Gamma(E)$ may have no basis, but behaves as a summand of a module that very much behaves like a vector space.
If $Q$ could be chosen $\{0\}$ (the module with only the additive zero) it would be free: for $E=T\R^2$ it can, for $E=TS^2$ it cannot \ts{11}{1}{1:42:41} \ts{11}{1}{1:43:01} \ts{11}{1}{1:43:22}.
\begin{coursenote}[Hypotheses of the theorem]
As said in class (not checked): the theorem is attributed to Swan. R.~G.~Swan, \emph{Vector bundles and projective modules}, Trans.\ Amer.\ Math.\ Soc.\ 105 (1962), proved it for vector bundles over a compact Hausdorff space and continuous sections; the version for smooth manifolds with $C^\infty(M)$ is the one quoted here. The
finite generation uses that a vector bundle over a manifold is covered by finitely many trivialising open sets (``of finite type''), true for manifolds; the precise hypotheses are not discussed in the lecture [?]. Not said in class.
\end{coursenote}
So $Q$ quantifies how much $\Gamma(E)$ fails to have a basis (``corollary, if you want''). Since a module over a non-division ring is not expected to have a basis, when one has one it is an accident, a coincidence; whoever is surprised that the vector fields on the sphere have no basis can say: ``why would it? It is just a projective module, not a free one'': the algebraic expression of the geometric fact,
which is of course not special to the sphere \ts{11}{1}{1:43:22} \ts{11}{1}{1:43:43} \ts{11}{1}{1:44:06} \ts{11}{1}{1:44:28} \ts{11}{1}{1:44:48}.

\paragraph{Homomorphisms of projective modules.} Two more remarks to wrap up \ts{11}{1}{1:44:48}.
\begin{theorem}[Hom of projective modules]\label{thm:rm-hom}
Let $P,Q$ be projective $R$-modules over a \emph{commutative} ring $R$. Then $\mathrm{Hom}_R(P,Q)$, the set of all $R$-linear maps $P\to Q$ (linear as for vector spaces, with respect to the linear structure of the modules and the ring), is again projective; if $P$ and $Q$ are finitely generated, so is $\mathrm{Hom}_R(P,Q)$ \ts{11}{1}{1:45:09} \ts{11}{1}{1:45:29} \ts{11}{1}{1:45:50} \ts{11}{1}{1:46:11} \ts{11}{1}{1:46:31}.
\end{theorem}
The proof is not given; the only surprise is the commutativity. Equip $\mathrm{Hom}_R(P,Q)$ with the pointwise addition and the multiplication by ring elements, $(r\odot\varphi)(v):=r\odot\varphi(v)$. For this to be again $R$-linear one needs $r\odot\varphi(s\odot v)=s\odot(r\odot\varphi(v))$: by linearity of $\varphi$ the
$s$ comes out, $\varphi(s\odot v)=s\odot\varphi(v)$, but then to bring $r$ and $s$ in the right order one needs $rs=sr$: ``only over a commutative ring does one have this actually expected structure'', homomorphisms are again a structure of the same type, and projectivity and finite generation are preserved
\ts{11}{1}{1:46:52} \ts{11}{1}{1:47:12} \ts{11}{1}{1:47:36} \ts{11}{1}{1:47:57} \ts{11}{1}{1:48:17} \ts{11}{1}{1:48:38} \ts{11}{1}{1:48:59} \ts{11}{1}{1:49:20}.
\begin{coursenote}[Dots in the lecture's computation]
The board computation of this commutativity argument contains several dots ($\odot$ in the ring module, the ring product, the scalar multiplication in $Q$) and the lecturer himself says ``oh dear, we have various dots'' and ``I wrote the same thing twice, not very meaningful''; the above is the reconstructed content. Not said in class. \ts{11}{1}{1:47:57}
\end{coursenote}

In particular one can apply the theorem with $R=C^\infty(M)$, which is commutative, and send vector fields to $C^\infty(M)$ itself (trivially a module over itself). The result
\[
  \Gamma(TM)^*:=\mathrm{Hom}_{C^\infty(M)}\bigl(\Gamma(TM),C^\infty(M)\bigr)
\]
is again a finitely generated module over $C^\infty(M)$, the \textbf{dual} of the module of vector fields, and one can show that it is the module of sections of the cotangent bundle, $\Gamma(TM)^*=\Gamma(T^*M)$ (not proven). That closes the loop \ts{11}{1}{1:49:20} \ts{11}{1}{1:49:41} \ts{11}{1}{1:50:02} \ts{11}{1}{1:50:24}.

% ==================================================================
\section{Tensor fields}\label{sec:rm-tensorfields}

One is now fully prepared for the standard textbook definition \ts{11}{1}{1:50:24} \ts{11}{1}{1:50:45}.
\begin{definition}[New concepts: tensor field]\label{def:rm-tensorfield}
A \textbf{$(r,s)$ tensor field} $T$ on a smooth manifold $M$ (in the lecture $(r,s)$; $(p,q)$ in \cref{def:te-tensor}) is a $C^\infty(M)$-\emph{multilinear} map
\[
  T\colon\underbrace{\Gamma(T^*M)\times\dots\times\Gamma(T^*M)}_{r\text{ copies}}\times\underbrace{\Gamma(TM)\times\dots\times\Gamma(TM)}_{s\text{ copies}}\longrightarrow C^\infty(M).
\]
It eats $r$ covector fields and $s$ vector fields and makes of them a smooth function on the manifold \ts{11}{1}{1:51:07} \ts{11}{1}{1:51:31} \ts{11}{1}{1:51:52} \ts{11}{1}{1:52:12}.
\end{definition}
Pointwise: at a given point the tensor field takes $r$ covectors, those given by the covector fields, and $s$ vectors, and returns a real number; doing this at every point gives a function, and that everything is smooth, and that multilinearity holds, is encoded in the fact that one has a $C^\infty(M)$-multilinear map. The set of all tensor fields is again
a module over $C^\infty(M)$ \ts{11}{1}{1:52:12} \ts{11}{1}{1:52:34} \ts{11}{1}{1:52:54} \ts{11}{1}{1:53:16}.
\begin{remark}[$C^\infty(M)$-multilinear, not $\R$-multilinear]\label{rem:rm-multilinear}
Textbooks write ``multilinear'' and sometimes say it is not $\R$-multilinear but $C^\infty(M)$-multilinear, and \emph{that is the key thing to check}: since the objects one feeds in are elements of a module, multilinearity must refer to the ring structure, otherwise the definition makes no sense
(compare \cref{ex:rm-leibniz}). One now sees the module structure ``lurking underneath'' the textbook definition \ts{11}{1}{1:53:16} \ts{11}{1}{1:53:37} \ts{11}{1}{1:53:57} \ts{11}{1}{1:54:18}.
\end{remark}
\begin{exercise}[Linear over $\R$ but not over $C^\infty(M)$]\label{ex:rm-leibniz}
Show that $\Gamma(TM)\to C^\infty(M)$, $X\mapsto X(g)$ for a fixed $g\in C^\infty(M)$ is $C^\infty(M)$-linear, whereas $f\mapsto X(f)$ (with fixed $X$, defined on the \emph{ring} $C^\infty(M)$) is only $\R$-linear, not $C^\infty(M)$-linear: the Leibniz rule (\cref{def:tv-derivation}) prevents it.
\emph{Not said in class.}
\end{exercise}
Hence one would never get the idea that the set of tensor fields should itself have a basis: not even the objects that one feeds in have necessarily a basis, since everything is a module structure and not a vector space structure. Keeping that in mind one can of course work \emph{locally} with a local basis, but not with a global one. This is the relationship of modules,
vector fields and tensor fields \ts{11}{1}{1:54:18} \ts{11}{1}{1:54:39} \ts{11}{1}{1:54:59}.

% ==================================================================
\section*{Key concepts}
\begin{keyconcepts}
\textbf{Retrieval questions.}
\begin{enumerate}
  \item What is a vector field, and how do the pointwise $\oplus$ and the multiplication by a function work on $\Gamma(TM)$? Why is a function (not only a real number) the right scalar?
  \item Which of the ring properties (commutative, unital, division) does $C^\infty(M)$ have? Why is it not a division ring?
  \item Define an $R$-module. What is the difference from a vector space, and what is the key fact about bases?
  \item State Zorn's lemma, with partial order, totally ordered subset, upper bound and maximal element. Which formulation of ``maximal'' is correct for a partial order?
  \item Prove that a module over a division ring has a basis: which poset, which chain, why is the union an upper bound, where is the division ring used?
  \item Why does $\Gamma(T\R^2)$ have a basis and $\Gamma(TS^2)$ not? Why can one not think of tensor fields globally in terms of components?
  \item What are free and projective modules? What does the theorem on sections of vector bundles say, and what does $Q$ measure?
  \item Why must a tensor field be $C^\infty(M)$-multilinear? What is $\Gamma(TM)^*$?
\end{enumerate}
\textbf{Concepts.}
\begin{itemize}
  \item Vector field, $\Gamma(TM)$, pointwise operations (\cref{def:rm-field}); ring $C^\infty(M)$ (\cref{prop:rm-cinf-ring}); module (\cref{def:rm-module}).
  \item Basis of a module over a division ring, Zorn's lemma (\cref{thm:rm-basis,cor:rm-vs,thm:rm-zorn,def:rm-poset,def:rm-total,def:rm-upper}); $\R^2$ versus $S^2$ (\cref{ex:rm-r2,ex:rm-s2,int:rm-nocomp}).
  \item Free and projective modules, sections of vector bundles, $\mathrm{Hom}$ of projective modules (\cref{def:rm-free,def:rm-proj,thm:rm-serre-swan,thm:rm-hom}).
  \item Tensor fields as $C^\infty(M)$-multilinear maps (\cref{def:rm-tensorfield,rem:rm-multilinear}).
\end{itemize}
\end{keyconcepts}

\flushside
\end{paracol}

% !TEX root = ../main.tex
% STUB: generated by tools/stubs.py; write the chapter from materials/transcripts/C12-V01.txt
\chapter{Grassmann algebra and de Rham cohomology}\label{ch:l12}
\begin{paracol}{2}

\begin{sources}
\begin{tabularx}{\linewidth}{@{}l l L@{}}
\toprule
\textbf{Video} & \textbf{Length} & \textbf{Content}\\
\midrule
\seg{12}{1}{Lecture 12: Grassmann algebra and de Rham cohomology} & 1:47:19 & (to be written)\\
\bottomrule
\end{tabularx}

\smallskip
\textbf{Transcript} \file{materials/transcripts/C12-V01.txt}.
\end{sources}

\begin{reading}[Where this lecture sits in the other notes]
\emph{Gravity and Light}, Lecture 12 (Integration on manifolds). (From the lecture titles; to be checked when the chapter is written.)
\end{reading}

\emph{This chapter has not been written yet.}

\flushside
\end{paracol}

% !TEX root = ../main.tex
% STUB: generated by tools/stubs.py; write the chapter from materials/transcripts/C13-V01.txt
\chapter{Lie groups and their Lie algebras}\label{ch:l13}
\begin{paracol}{2}

\begin{sources}
\begin{tabularx}{\linewidth}{@{}l l L@{}}
\toprule
\textbf{Video} & \textbf{Length} & \textbf{Content}\\
\midrule
\seg{13}{1}{Lecture 13: Lie groups and their Lie algebras} & 1:43:12 & (to be written)\\
\bottomrule
\end{tabularx}

\smallskip
\textbf{Transcript} \file{materials/transcripts/C13-V01.txt}.
\end{sources}

\begin{reading}[Where this lecture sits in the other notes]
\emph{Gravity and Light}, Lecture 11 (Symmetry); \emph{Quantum Theory}, Lecture 13 (Spin). (From the lecture titles; to be checked when the chapter is written.)
\end{reading}

\emph{This chapter has not been written yet.}

\flushside
\end{paracol}

% !TEX root = ../main.tex
% STUB: generated by tools/stubs.py; write the chapter from materials/transcripts/C14-V01.txt
\chapter{Classification of Lie algebras and Dynkin diagrams}\label{ch:l14}
\begin{paracol}{2}

\begin{sources}
\begin{tabularx}{\linewidth}{@{}l l L@{}}
\toprule
\textbf{Video} & \textbf{Length} & \textbf{Content}\\
\midrule
\seg{14}{1}{Lecture 14: Classification of Lie algebras and Dynkin diagrams} & 1:46:44 & (to be written)\\
\bottomrule
\end{tabularx}

\smallskip
\textbf{Transcript} \file{materials/transcripts/C14-V01.txt}.
\end{sources}

\begin{reading}[Where this lecture sits in the other notes]
Not treated in the other courses (only here). (From the lecture titles; to be checked when the chapter is written.)
\end{reading}

\emph{This chapter has not been written yet.}

\flushside
\end{paracol}

% !TEX root = ../main.tex
% STUB: generated by tools/stubs.py; write the chapter from materials/transcripts/C15-V01.txt
\chapter{The Lie group SL(2,C) and its Lie algebra sl(2,C)}\label{ch:l15}
\begin{paracol}{2}

\begin{sources}
\begin{tabularx}{\linewidth}{@{}l l L@{}}
\toprule
\textbf{Video} & \textbf{Length} & \textbf{Content}\\
\midrule
\seg{15}{1}{Lecture 15: The Lie group SL(2,C) and its Lie algebra sl(2,C)} & 1:51:13 & (to be written)\\
\bottomrule
\end{tabularx}

\smallskip
\textbf{Transcript} \file{materials/transcripts/C15-V01.txt}.
\end{sources}

\begin{reading}[Where this lecture sits in the other notes]
\emph{Quantum Theory}, Lectures 13 and 15 (spin) [?]. (From the lecture titles; to be checked when the chapter is written.)
\end{reading}

\emph{This chapter has not been written yet.}

\flushside
\end{paracol}

% !TEX root = ../main.tex
% STUB: generated by tools/stubs.py; write the chapter from materials/transcripts/C16-V01.txt
\chapter{Dynkin diagrams from Lie algebras, and vice versa}\label{ch:l16}
\begin{paracol}{2}

\begin{sources}
\begin{tabularx}{\linewidth}{@{}l l L@{}}
\toprule
\textbf{Video} & \textbf{Length} & \textbf{Content}\\
\midrule
\seg{16}{1}{Lecture 16: Dynkin diagrams from Lie algebras, and vice versa} & 1:40:32 & (to be written)\\
\bottomrule
\end{tabularx}

\smallskip
\textbf{Transcript} \file{materials/transcripts/C16-V01.txt}.
\end{sources}

\begin{reading}[Where this lecture sits in the other notes]
(to be written)
\end{reading}

\emph{This chapter has not been written yet.}

\flushside
\end{paracol}

% !TEX root = ../main.tex
\chapter{Representation theory of Lie groups and Lie algebras}\label{ch:l17}
\begin{paracol}{2}

\begin{sources}
\begin{tabularx}{\linewidth}{@{}l l L@{}}
\toprule
\textbf{Video} & \textbf{Length} & \textbf{Content}\\
\midrule
\seg{17}{1}{Lecture 17: Representation theory of Lie groups and Lie algebras} & 1:32:39 & A short section added on request (``mostly you meet Lie algebras and groups in physics in terms of their representations''), without proofs: linear representations of a Lie algebra and the representation space; examples ($\mathfrak{sl}(2,\C)$ on $\C^2$, the rotation algebra on $\R^3$ and on $\C^2$, the trivial and the adjoint representation); reducible, irreducible and faithful representations; the Killing form of a representation, the Casimir operator, Schur's lemma and the value $\dim L/\dim V$, checked on the vector and spin representations (with a long digression on a factor $2$); representations of a Lie group, the example $SO(2)$ and the adjoint representation of a Lie group.\\
\bottomrule
\end{tabularx}

\smallskip
\textbf{Transcript} \file{materials/transcripts/C17-V01.txt} (YouTube automatic captions; the blackboard matrices are reconstructed from the verbal description, several signs are unsure already in the lecture and are checked in the course notes).
\end{sources}

\begin{reading}[Where this lecture sits in the other notes]
\emph{Quantum Theory}, Lectures 13--15; \emph{Theoretische Quantenmechanik}, Lectures 14--16 (angular momentum, Clebsch--Gordan): forward pointers, since the spin representations of $\mathfrak{so}(3)$ and the addition of angular momenta are the physics of the examples below [check later]. It uses the algebras of \cref{ch:l13,ch:l14,ch:l16}; the next lecture goes back from the algebra to the group (\cref{ch:l18}).
\end{reading}

% ==================================================================
\section{Representations of a Lie algebra}\label{sec:rp-algebra}

After the previous lecture some listeners asked about representations of Lie algebras and groups, and the lecturer decided to add a brief section, because in physics one mostly meets Lie algebras and Lie groups \emph{through} their representations, and very often they are even defined by them \ts{17}{1}{0:00} \ts{17}{1}{0:21}. In this course a Lie group was defined as a smooth manifold with a group law compatible with the smooth structure, and its Lie algebra as the left-invariant vector fields (which one can pull back to $T_{\mathbb 1}G$ with an induced bracket, \cref{thm:ll-iso}). ``But what on earth are these representations?'' \ts{17}{1}{0:41}. A half-semester course could be given on the subject, so the lecture presents \emph{no proofs}, only the important concepts \ts{17}{1}{1:02} \ts{17}{1}{1:28} \ts{17}{1}{1:49}.

\begin{definition}[New concepts: representation of a Lie algebra]\label{def:rp-rep}
Let $(L,[\cdot,\cdot])$ be a Lie algebra. A \textbf{representation} of $L$ is a linear map
\[
  \rho\colon L\to\mathrm{End}(V)
\]
into the endomorphisms of some vector space $V$ over $\R$ or $\C$ (``linear'': $L$ is a vector space, and $\mathrm{End}(V)$ is a vector space), such that for all $a,b\in L$
\[
  \rho\bigl([a,b]\bigr)=\bigl[\rho(a),\rho(b)\bigr]_{\mathrm{End}(V)} .
\]
The vector space $V$ is the \textbf{representation space}; if $V$ is finite-dimensional the representation is finite-dimensional. Since the maps $\rho(a)$ are linear maps on $V$ (not vector fields), the bracket on the right is not the differential-geometric one but the commutator, $[\phi,\psi]:=\phi\circ\psi-\psi\circ\phi$. This is bilinear, antisymmetric by construction and satisfies the Jacobi identity: this is why one maps into $\mathrm{End}(V)$, which thanks to its natural composition is a Lie algebra at once. So: one may first take the bracket in $L$ and then map, or map first and then take the commutator, with the same result \ts{17}{1}{2:10} \ts{17}{1}{2:52} \ts{17}{1}{3:34} \ts{17}{1}{4:17} \ts{17}{1}{4:59} \ts{17}{1}{5:43}.
\end{definition}
The terminology can confuse: one might think that $\mathrm{End}(V)$ is the ``space where $L$ is represented'', but the name \emph{representation space} is reserved for $V$. The idea is that, once $\rho(a)$ is known, it can \emph{act on} the vectors of $V$: an endomorphism takes a vector and produces a new vector \ts{17}{1}{6:03} \ts{17}{1}{6:25} \ts{17}{1}{7:07}. (``Linear'' in ``linear representation'' refers to the fact that the elements are mapped to endomorphisms of a vector space.) The name comes from this: the abstract elements of $L$ are \emph{represented} by concrete linear maps (matrices) \ts{17}{1}{7:28}.
\begin{coursenote}[Not said in class: real algebra, complex space]
$L$ and $V$ need not be over the same field: a real Lie algebra such as $\mathfrak{so}(3,\R)$ may be represented on a complex space such as $\C^2$ (then $\rho$ is $\R$-linear), as is done below. \ts{17}{1}{19:37}
\end{coursenote}

% ------------------------------------------------------------------
\section{Examples}\label{sec:rp-examples}

\subsection{$\mathfrak{sl}(2,\C)$ on $\C^2$}\label{sec:rp-sl2}
The algebra of \cref{ch:l15,ch:l16} has three basis vectors with
\[
  [X_1,X_2]=2X_2,\qquad[X_1,X_3]=-2X_3,\qquad[X_2,X_3]=X_1 ,
\]
derived from the left-invariant vector fields; so far $X_1,X_2,X_3$ were only the representatives at the identity of those fields, ``we didn't know more'' \ts{17}{1}{7:49} \ts{17}{1}{8:31}.
\begin{recall}[The three brackets]
The relations are those of \cref{thm:sa-bracket}.
\end{recall}
A representation $\rho\colon\mathfrak{sl}(2,\C)\to\mathrm{End}(\C^2)$ is obtained by prescribing it on a basis (linearity does the rest) \ts{17}{1}{8:52} \ts{17}{1}{9:14}:
\[
  \rho(X_1)=\begin{pmatrix}1&0\\0&-1\end{pmatrix},\quad
  \rho(X_2)=\begin{pmatrix}0&1\\0&0\end{pmatrix},\quad
  \rho(X_3)=\begin{pmatrix}0&0\\1&0\end{pmatrix}.
\]
Here $\mathrm{End}(\C^2)$ is written by matrices with respect to the canonical basis, and the composition of endomorphisms is matrix multiplication. The compatibility condition has to be checked on the pairs of basis vectors, because both sides are bilinear in $a,b$. For $(X_1,X_2)$:
\[
  \rho(X_1)\rho(X_2)-\rho(X_2)\rho(X_1)=\begin{pmatrix}0&1\\0&0\end{pmatrix}-\begin{pmatrix}0&-1\\0&0\end{pmatrix}=2\begin{pmatrix}0&1\\0&0\end{pmatrix}=2\rho(X_2)=\rho(2X_2)=\rho([X_1,X_2]),
\]
by linearity of $\rho$ \ts{17}{1}{10:15} \ts{17}{1}{10:35} \ts{17}{1}{11:18} \ts{17}{1}{12:00}. The pairs $(X_1,X_3)$ and $(X_2,X_3)$ are checked in the same way (the lecture does not carry this out; one finds $-2\rho(X_3)$ and $\rho(X_1)$) \ts{17}{1}{12:21}. So this is a representation, and it is how $\mathfrak{sl}(2,\C)$ is usually presented.

The image $\rho(\mathfrak{sl}(2,\C))\subset\mathrm{End}(\C^2)$ is the span of the three matrices, that is, the matrices $\begin{pmatrix}\alpha&\beta\\\gamma&\delta\end{pmatrix}$ with $\alpha+\delta=0$: the \emph{trace-free} matrices. This is the origin of the usual sentence ``the group $SL(2,\C)$ consists of the matrices of determinant $1$, and its algebra of the trace-free $2\times2$ matrices'': it is just a statement about \emph{one} representation. Here one went the abstract way, from the group to the left-invariant fields, and the algebra is not a priori a space of matrices. And the point of the whole business is that a given Lie algebra has, under suitable conditions, \emph{several} representations; this is only one of them \ts{17}{1}{12:42} \ts{17}{1}{13:45} \ts{17}{1}{14:05}.

\subsection{The rotation algebra on $\R^3$ and on $\C^2$}\label{sec:rp-so3}
The algebra of the rotation group appears in physics all the time \ts{17}{1}{14:27}. The algebra $\mathfrak{so}(3,\R)$ is three-dimensional and is defined, like any Lie algebra, by the choice of a basis $X_1,X_2,X_3$ and of the structure constants (\cref{def:cl-structure}). The lecture does it as follows: first the Killing form $K_{ij}=C^m{}_{in}C^n{}_{jm}$ (\cref{def:cl-killing}); then the first index of the structure constants is pulled down with it, $C_{kij}:=K_{km}C^m{}_{ij}$ (``one could say it basis-free, but that is just a way of saying this''); and the algebra $\mathfrak{so}(3,\R)$ is the one for which
\[
  C_{kij}=\varepsilon_{ijk}\in\{+1,-1,0\}
\]
(even permutation of $123$, odd permutation, no permutation: for instance $112\mapsto0$) \ts{17}{1}{14:48} \ts{17}{1}{15:08} \ts{17}{1}{15:50} \ts{17}{1}{16:31} \ts{17}{1}{16:52}. These are totally antisymmetric structure coefficients, and to be able to say \emph{totally} antisymmetric one needs the Killing form to pull the index down; conversely, from $\varepsilon_{ijk}$ one reconstructs the $C^k{}_{ij}$ \ts{17}{1}{17:12}.
\begin{coursenote}[Not said in class: the normalisation]
Taken literally, the two conditions are not compatible with unit constants: for $C^k{}_{ij}=c\,\varepsilon_{ijk}$ one has $K_{ij}=-2c^2\delta_{ij}$, hence $K_{km}C^m{}_{ij}=-2c^3\varepsilon_{ijk}$. What is meant, and what is used afterwards, is the standard algebra $[X_i,X_j]=\varepsilon_{ijk}X_k$, whose Killing form is $K=-2\delta$ (negative definite). \ts{17}{1}{16:52}
\end{coursenote}
\textbf{On $\R^3$} (the vector representation) one takes $\rho_V\colon\mathfrak{so}(3,\R)\to\mathrm{End}(\R^3)$ with three antisymmetric matrices, each with a pair of $\pm1$ entries and zeros elsewhere \ts{17}{1}{17:33} \ts{17}{1}{17:53} \ts{17}{1}{18:14}. The lecture does not trust its signs (``up to signs this is right; please check for yourself'') and says one has to check the compatibility condition; the consistent choice is
\[
  \bigl(\rho_V(X_i)\bigr)_{jk}=-\varepsilon_{ijk},\quad\text{e.g. }\rho_V(X_1)=\begin{pmatrix}0&0&0\\0&0&-1\\0&1&0\end{pmatrix}
\]
(\emph{not said in class}: with this choice $[\rho_V(X_i),\rho_V(X_j)]=\varepsilon_{ijk}\rho_V(X_k)$.) \ts{17}{1}{18:35} \ts{17}{1}{18:55}. This is what one learns first: the rotation algebra acts on $\R^3$.

\textbf{On $\C^2$} (the spin representation): famously from quantum mechanics there is another one, $\rho_{\mathrm{spin}}\colon\mathfrak{so}(3,\R)\to\mathrm{End}(\C^2)$, which acts on two-component complex objects, the spinors of Schr\"odinger theory, and the matrices are $\frac12$ of the Pauli matrices $\sigma_1,\sigma_2,\sigma_3$ ``with or without a factor $\mathrm{i}$; please figure it out'' \ts{17}{1}{19:17} \ts{17}{1}{19:57} \ts{17}{1}{20:18}. The consistent choice is worked out later in the lecture (\cref{sec:rp-casimir-spin}): $\rho_{\mathrm{spin}}(X_k)=-\frac{\mathrm i}{2}\sigma_k$. So the \emph{same} rotation algebra can be represented on $\R^3$ (rotating in real three-space) or on $\C^2$ (spinors), and these are different representations \ts{17}{1}{20:39} \ts{17}{1}{21:00}.

\subsection{The trivial and the adjoint representation}\label{sec:rp-triv}
\textbf{Trivial representation} \ts{17}{1}{21:21} \ts{17}{1}{21:41}: for every Lie algebra and on every vector space $V$ there is $\rho_{\mathrm{triv}}\colon L\to\mathrm{End}(V)$, $a\mapsto0$ (the zero map, ``the matrix with only zeros''). It is a representation: $\rho([a,b])=0$ and $[\rho(a),\rho(b)]=[0,0]=0$, the same zero of $\mathrm{End}(V)$, so compatibility always holds \ts{17}{1}{22:09} \ts{17}{1}{22:29}. Physicists call it the \emph{scalar} representation (nothing changes under the algebra) \ts{17}{1}{23:11}.

\textbf{Adjoint representation} \ts{17}{1}{23:31}. Does every Lie algebra have a \emph{non-trivial} representation? Yes. A representation has two choices: the representation space, and, once the space is fixed, still the linear map, so different representations may live on the same space \ts{17}{1}{23:52} \ts{17}{1}{24:13} \ts{17}{1}{24:54}. If one is short of vector spaces, one is already given one: the vector space underlying $L$. Take $V=L$ and
\[
  \mathrm{ad}\colon L\to\mathrm{End}(L),\qquad a\mapsto\mathrm{ad}_a=[a,\,\cdot\,],
\]
that is, $\mathrm{ad}_a(\ell)=[a,\ell]\in L$ (\cref{def:cl-ad}). Compatibility is the Jacobi identity. This is the \textbf{adjoint representation} \ts{17}{1}{25:14} \ts{17}{1}{26:05} \ts{17}{1}{26:27} \ts{17}{1}{26:49}. The whole classification of semisimple and simple Lie algebras (\cref{ch:l14}) made essential use of it: the Cartan subalgebra is where $\mathrm{ad}$ is diagonalised, and all the other directions are eigenvectors \ts{17}{1}{27:10}.

% ------------------------------------------------------------------
\section{Reducible, irreducible and faithful representations}\label{sec:rp-irred}

To study representations one wants \emph{building blocks}: if a representation can be understood through two others, there is no need to study it separately \ts{17}{1}{27:30}.
\begin{definition}[New concepts: reducible and irreducible representations]\label{def:rp-irred}
A representation $\rho\colon L\to\mathrm{End}(V)$ is \textbf{reducible} if there exists a vector subspace $U\subseteq V$, with $U\neq\{0\}$ and $U\neq V$, such that
\[
  \rho(a)u\in U\qquad\text{for all }a\in L,\ u\in U .
\]
Such a $U$ is an \textbf{invariant subspace}. Equivalently $\rho$ restricts to a representation $\rho_U\colon L\to\mathrm{End}(U)$. Otherwise $\rho$ is called \textbf{irreducible}. (``Reducible'': it can be reduced to the smaller pieces $U$; ``irreducible'': it cannot.) \ts{17}{1}{27:51} \ts{17}{1}{28:12} \ts{17}{1}{28:53} \ts{17}{1}{29:13} \ts{17}{1}{29:54}
\end{definition}
\begin{answer}{Should the invariant subspace be non-trivial?}[Proper invariant subspace]
A student asks about the condition ``$U$ non-zero'': it does not suffice to exclude $\{0\}$, which is always invariant; the whole space $V$ must be excluded too. The lecturer agrees (``I should have written non-trivial''), and the definition above contains both conditions. \ts{17}{1}{36:52}
\end{answer}
\textbf{Example} \ts{17}{1}{30:15}. Take the two representations $\rho_V\colon\mathfrak{so}(3,\R)\to\mathrm{End}(\R^3)$ and $\rho_{\mathrm{spin}}\colon\mathfrak{so}(3,\R)\to\mathrm{End}(\C^2)$ and build a new one on the direct sum, $V=\R^3\oplus\C^2$ (the real $\R^3$ is considered as inside $\C^3$, so the whole is a complex matrix representation): to $a$ one associates the block matrix
\[
  \rho(a)=\begin{pmatrix}\rho_V(a)&0\\0&\rho_{\mathrm{spin}}(a)\end{pmatrix}
\]
with zeros off the diagonal blocks \ts{17}{1}{30:56} \ts{17}{1}{31:17}. Matrix multiplication can be carried out blockwise, so the commutator of two such matrices has the same block form with the commutators of the blocks; since each block is a representation, $\rho$ is again one \ts{17}{1}{32:20} \ts{17}{1}{32:41}. Now $U=\R^3$ (the vectors with the last two entries zero) and $U'=\C^2$ (the first three entries zero) are vector subspaces and \emph{invariant}: $\rho(a)$ acts on a vector with entries only in the upper block and, because of the block form, nothing is produced in the lower one \ts{17}{1}{33:04} \ts{17}{1}{33:46}. So a representation with a block form is reducible. One can also show (not done in class) that $\rho_V$ and $\rho_{\mathrm{spin}}$ themselves are not reducible \ts{17}{1}{34:07}.
\begin{coursenote}[Not said in class: building blocks]
``One can always understand a reducible representation through irreducible building blocks'' holds for finite-dimensional representations of \emph{semisimple} Lie algebras over $\C$ (complete reducibility, Weyl's theorem): every such representation is a direct sum of irreducible ones. For a general Lie algebra a reducible representation need not be a direct sum, only an invariant subspace with a quotient (a block-\emph{triangular} form). \ts{17}{1}{34:28}
\end{coursenote}
\begin{aside}[Elementary particles]
Physics: ``particles are irreducible representations of the symmetry'' \ts{17}{1}{34:48}. For the particle of classical mechanics the symmetry is the Galilei group (rotations, Galilei boosts, translations in space-time); to keep it simple, one traps the particle at a point: what remains is rotational invariance, with the group $SO(3)$ and the algebra above. The particle may then be represented by a vector, by a scalar (the trivial representation) or, in quantum mechanics, by a two-component spinor. If one took a reducible representation one could always make \emph{two} particles out of it (two particles also transform in some representation of the symmetry, and one wants to study each on its own); an \emph{elementary} particle is an irreducible one. This is just the reason why the term comes up in particle physics \ts{17}{1}{35:08} \ts{17}{1}{35:50} \ts{17}{1}{36:10}.
\end{aside}
\begin{definition}[New concepts: faithful representation]\label{def:rp-faithful}
A representation $\rho\colon L\to\mathrm{End}(V)$ is \textbf{faithful} if $\rho$ is injective: no two different elements of $L$ are mapped to the same endomorphism. (The representation then ``represents $L$ without loss'': $L$ can be read off from its image.) \ts{17}{1}{37:13} \ts{17}{1}{37:35} \ts{17}{1}{38:33}
\end{definition}
\begin{exercise}[Adjoint representation: irreducible, faithful?]\label{ex:rp-ad}
Check irreducibility and faithfulness of the adjoint representation $\rho=\mathrm{ad}$ (the lecture: ``let us formulate this as a homework problem''). \emph{Hint, not said in class:} an invariant subspace of $\mathrm{ad}$ is exactly an ideal, and $\ker\mathrm{ad}$ is the centre of $L$; see \cref{def:cl-simple,thm:sa-simple}. (\cref{exam:ch17}.) \ts{17}{1}{38:54}
\end{exercise}

% ------------------------------------------------------------------
\section{The Casimir operator}\label{sec:rp-casimir}

A very important concept in the study of representations is the \emph{Casimir operator} (the name is not explained in class; not said in class: it is named after the Dutch physicist Hendrik Casimir) \ts{17}{1}{39:40}. In physics one is very often handed the Lie algebra already in a representation, so there is some confusion: a Casimir operator cannot be constructed for a Lie algebra, but for \textbf{any given representation of any given Lie algebra} \ts{17}{1}{40:00} \ts{17}{1}{40:23}. This takes several steps.

\textbf{Step A: the Killing form of a representation.} \ts{17}{1}{40:44} The space $\mathrm{End}(V)$ with the commutator is a Lie algebra, so it has a Killing form, but the lecture wants something different, a form on $L$ itself \ts{17}{1}{41:05} \ts{17}{1}{41:47}.
\begin{definition}[New concepts: the Killing form of a representation]\label{def:rp-killing}
For a representation $\rho\colon L\to\mathrm{End}(V)$ the \textbf{Killing form with respect to $\rho$} is
\[
  K_\rho\colon L\times L\to\C\ (\text{or }\R),\qquad K_\rho(a,b):=\mathrm{tr}_V\bigl(\rho(a)\circ\rho(b)\bigr).
\]
For $\rho=\mathrm{ad}$ this is the Killing form $K(a,b)=\mathrm{tr}(\mathrm{ad}_a\circ\mathrm{ad}_b)$ used so far (\cref{def:cl-killing}): that one was the Killing form of the \emph{adjoint} representation. \ts{17}{1}{42:07} \ts{17}{1}{42:50} \ts{17}{1}{43:12}
\end{definition}
\textbf{Step B: the Casimir operator} \ts{17}{1}{43:36} \ts{17}{1}{43:57}.
\begin{definition}[New concepts: the Casimir operator]\label{def:rp-casimir}
Let $\rho\colon L\to\mathrm{End}(V)$ be a representation. Choose a basis $X_1,\dots,X_{\dim L}$ of $L$ and let $\tilde X_1,\dots,\tilde X_{\dim L}$ be the \emph{dual basis} of $L$ with respect to $K_\rho$, that is
\[
  K_\rho\bigl(X_i,\tilde X_j\bigr)=\delta_{ij}.
\]
The \textbf{Casimir operator} of $\rho$ is
\[
  \Omega_\rho:=\sum_{i=1}^{\dim L}\rho(X_i)\circ\rho(\tilde X_i)\ \in\mathrm{End}(V).
\]
It is independent of the choice of the basis $X_i$. \ts{17}{1}{44:18} \ts{17}{1}{45:02} \ts{17}{1}{45:24} \ts{17}{1}{45:44} \ts{17}{1}{46:05}
\end{definition}
\begin{coursenote}[Not said in class: existence of the dual basis, signs]
The dual basis exists iff $K_\rho$ is non-degenerate. The lecture first writes $\tilde X_i$ ``for an orthonormal basis'', $\tilde X_i=X_i$, and says this works only if the group is compact (``otherwise I do not find such a basis''). For a compact group $K$ is \emph{negative} definite, so an orthonormal basis has $K_\rho(X_i,X_j)=-\delta_{ij}$ and one writes $\Omega_\rho=-\sum_i\rho(X_i)\rho(X_i)$; the lecture corrects this sign later \ts{17}{1}{1:00:08} \ts{17}{1}{1:00:28}. With the dual basis of \cref{def:rp-casimir} no compactness is needed, and the signs take care of themselves (see the examples). Basis independence: under $X_i'=M_i{}^jX_j$ the dual basis transforms with $(M^{-1})^T$ and the sum is invariant. \ts{17}{1}{46:25} \ts{17}{1}{46:50}
\end{coursenote}
The Casimir operator can be applied to a vector of the representation space and gives a vector of the same space \ts{17}{1}{47:11}. The first property:
\begin{theorem}[The Casimir operator commutes with the representation]\label{thm:rp-commute}
For every $a\in L$: $[\Omega_\rho,\rho(a)]=0$. Informally: it commutes with everything in the image $\rho(L)$.
\end{theorem}
The lecture does not prove it and asks for the proof on the problem sheet (\cref{exam:ch17}) \ts{17}{1}{47:32} \ts{17}{1}{47:52} \ts{17}{1}{48:13}. This is how the operator is usually introduced \ts{17}{1}{48:36}. The second statement is \textbf{Schur's lemma} (the transcript has ``shortlemma''; named after Issai Schur, not said in class) \ts{17}{1}{49:00}:
\begin{theorem}[Schur's lemma for the Casimir operator]\label{thm:rp-schur}
If $\rho$ is irreducible, then there is a number $c_\rho\in\C$ (or $\R$) such that
\[
  \Omega_\rho=c_\rho\,\mathrm{id}_V ,
\]
a diagonal matrix with $c_\rho$ on the diagonal. (The lecture gives no proof.) \ts{17}{1}{49:21}\ts{17}{1}{50:03}
\end{theorem}
So the Casimir operator of an irreducible representation is a constant multiple of the identity. Not said in class: over $\R$ the lemma can fail (an irreducible real representation may have commutant larger than the multiples of the identity, e.g.\ the rotation $\begin{pmatrix}0&-1\\1&0\end{pmatrix}$ commutes with the irreducible $SO(2)$ action on $\R^2$); the proof needs $\C$ (an eigenvalue exists, and its eigenspace is invariant).

The lecture then tries to show that this number is always given by the ratio of the dimensions \ts{17}{1}{50:25} \ts{17}{1}{50:45}:
\begin{theorem}[The value of the Casimir operator]\label{thm:rp-value}
If $\rho\colon L\to\mathrm{End}(V)$ is irreducible, $\dim V\ne0$, and $K_\rho$ is non-degenerate, then
\[
  c_\rho=\frac{\dim L}{\dim V}.
\]
\end{theorem}
\begin{proof}
By \cref{thm:rp-schur}, $\Omega_\rho=c_\rho\,\mathrm{id}_V$, so taking the trace (linear, and $\mathrm{tr}\,\mathrm{id}_V=\dim V$):
\[
  \mathrm{tr}\,\Omega_\rho=c_\rho\,\dim V .
\]
On the other hand, by linearity of the trace and the definitions,
\[
  \begin{aligned}
  \mathrm{tr}\,\Omega_\rho&=\sum_i\mathrm{tr}\bigl(\rho(X_i)\rho(\tilde X_i)\bigr)\\
  &=\sum_iK_\rho(X_i,\tilde X_i)\\
  &=\sum_{i=1}^{\dim L}1=\dim L .
  \end{aligned}
\]
Hence $c_\rho\dim V=\dim L$. \ts{17}{1}{51:15} \ts{17}{1}{51:35} \ts{17}{1}{52:16} \ts{17}{1}{52:37} \ts{17}{1}{53:01} \ts{17}{1}{53:43} \ts{17}{1}{54:04}
\end{proof}
So different irreducible representations of the same Lie algebra on different representation spaces $V_1,V_2$ are labelled by this number, and one can act with the Casimir operator on any vector of the representation space of one particular irreducible representation and always get the same number times the vector \ts{17}{1}{54:29} \ts{17}{1}{54:50}.

\subsection{Check: the vector representation}\label{sec:rp-casimir-vec}
In physics the rotation algebra reads $[J_i,J_j]=\varepsilon_{ijk}J_k$ (``something of that type''; here the abstract $X_i$) and $\vec J^2$ is not literally a square of a vector but $\sum_i\rho(J_i)\rho(J_i)$ \ts{17}{1}{55:10} \ts{17}{1}{55:31} \ts{17}{1}{56:12}. Take the vector representation. Each $\rho_V(X_i)^2$ is a diagonal matrix with $-1$ at the two positions in its block, so the three squares add up to $-2\cdot\mathbb 1_3$ (the lecture phrases it with the rotation blocks of the $\pm1$ entries and gets the pattern $(1,1,0)$, $(1,0,1)$, $(0,1,1)$ for the three squares, so that the sum is $2\cdot\mathbb 1$) \ts{17}{1}{56:33} \ts{17}{1}{57:16} \ts{17}{1}{57:37} \ts{17}{1}{57:57}. With the prefactor $\frac12$ this gives $c_{\rho_V}=1$, and \cref{thm:rp-value} agrees: $\dim\mathfrak{so}(3)/\dim\R^3=3/3=1$ \ts{17}{1}{58:18} \ts{17}{1}{58:39} \ts{17}{1}{59:00} \ts{17}{1}{59:25}. The lecture stresses that in physics the factor $\frac12$ is absent, and everything is then off by a factor $2$.
\begin{answer}{Where is the sign?}[The minus sign in the Casimir operator]
A student remarks on a sign in the computation (the question is inaudible in the transcript [?]; it concerns the minus signs coming from the squares of the blocks). The lecturer answers that this is the earlier problem with the sign: for a compact group the Killing form is negative definite and the definition of $\Omega_\rho$ needs the overall minus sign, ``then the signs are right''. See the course note after \cref{def:rp-casimir}. \ts{17}{1}{59:48} \ts{17}{1}{1:00:08} \ts{17}{1}{1:00:48}
\end{answer}

\subsection{Check: the spin representation, and a factor $2$}\label{sec:rp-casimir-spin}
Now the second example, the same abstract algebra $[X_i,X_j]=\varepsilon_{ijk}X_k$ in the spin representation $\rho_{\mathrm{spin}}$ \ts{17}{1}{1:01:09}. The Pauli matrices are (the lecturer needs reminding)
\[
  \sigma_1=\begin{pmatrix}0&1\\1&0\end{pmatrix},\quad\sigma_2=\begin{pmatrix}0&-\mathrm i\\\mathrm i&0\end{pmatrix},\quad\sigma_3=\begin{pmatrix}1&0\\0&-1\end{pmatrix}.
\]
Check of $[\sigma_1,\sigma_2]$: $\sigma_1\sigma_2=\mathrm i\sigma_3$ and $\sigma_2\sigma_1=-\mathrm i\sigma_3$, so $[\sigma_1,\sigma_2]=2\mathrm i\sigma_3$ \ts{17}{1}{1:01:30} \ts{17}{1}{1:01:50} \ts{17}{1}{1:02:14} \ts{17}{1}{1:02:37}. The factor $2$ is removed by taking $\tau_i=\frac12\sigma_i$, for which $[\tau_1,\tau_2]=\mathrm i\tau_3$; the factor $\mathrm i$ by putting $-\mathrm i$ in front, since $(-\mathrm i)^2[\tau_1,\tau_2]=-\mathrm i\tau_3$ is exactly the representative of $X_3$ (the lecturer counts the minus signs aloud and arrives at this choice). In summary \ts{17}{1}{1:02:57} \ts{17}{1}{1:03:19} \ts{17}{1}{1:03:40} \ts{17}{1}{1:04:02} \ts{17}{1}{1:04:22} \ts{17}{1}{1:04:42}:
\[
  \rho_{\mathrm{spin}}(X_k)=-\frac{\mathrm i}{2}\sigma_k,\qquad[\rho_{\mathrm{spin}}(X_i),\rho_{\mathrm{spin}}(X_j)]=\varepsilon_{ijk}\rho_{\mathrm{spin}}(X_k).
\]
(In quantum mechanics one wants hermitian generators and puts an extra $\mathrm i$ in.) The Casimir operator is then computed with the same recipe, $\Omega=-\frac12\sum_i\rho(X_i)\rho(X_i)$ \ts{17}{1}{1:05:02} \ts{17}{1}{1:05:23}. The lecture first evaluates it with the matrices $-\mathrm i\sigma_i$ (the factor $\frac12$ inside is forgotten): each square is $-\mathbb 1$, ``the minuses all go out and there is a three everywhere'', so $\Omega=\frac32\mathbb 1$, in agreement with $\dim L/\dim V=\frac32$ (the lecture also writes it as $\frac12(\frac12+1)$, which is $\frac34$: a slip, see the course note below) \ts{17}{1}{1:05:44} \ts{17}{1}{1:06:26} \ts{17}{1}{1:07:07} \ts{17}{1}{1:07:30} \ts{17}{1}{1:07:50}.

Then something is wrong: with the correct matrices $-\frac{\mathrm i}2\sigma_i$ each square is $-\frac14\mathbb 1$ and one gets $\frac12\cdot3\cdot\frac14=\frac38$, not $\frac32$ \ts{17}{1}{1:08:12} \ts{17}{1}{1:08:32} \ts{17}{1}{1:08:55}. And $\frac38$ is not $j(j+1)$ either (``in the vector case it wasn't either'': for the vector representation one has $j=1$, $j(j+1)=2$ but the number found is $1$). So the lecture's first diagnosis is a mismatch of conventions between physics and mathematics, by an overall factor $2$: the physical $\vec J^2$ is \emph{twice} the Casimir operator defined here, ``because one wants the dual basis with respect to the Killing form on the same space, to make everything basis independent'' \ts{17}{1}{1:09:16} \ts{17}{1}{1:09:36} \ts{17}{1}{1:09:57} \ts{17}{1}{1:10:18} \ts{17}{1}{1:10:38} \ts{17}{1}{1:10:59} \ts{17}{1}{1:11:20}.
\begin{answer}{Why does the number not come out as $3/2$?}[A student: the basis depends on $\rho$]
After a pause for thinking (``I do not get the algebra straight'') a student points out the origin of the trouble: the dual vectors $\tilde X_i$ were defined by $K_\rho(X_i,\tilde X_j)=\delta_{ij}$ (the lecture: by $-\delta_{ij}$ on a compact group), and $K_\rho(X_i,X_j)=\mathrm{tr}(\rho(X_i)\rho(X_j))$ \emph{depends on the representation}. In the spin representation the proportionality constant between $\tilde X_i$ and $X_i$ is therefore different from the one in the vector representation, and the earlier computation, done with the constant of the vector representation, does not apply. The lecturer thanks him (``it's an excellent problem-sheet question'': the dual basis is constructed with respect to $K_\rho$, so it has to be adapted to every $\rho$), and says that it is now clear that the missing factor $2$ should be there, so that one gets $\frac32$, but ``I should make no further guesses, just calculate''. \ts{17}{1}{1:11:50} \ts{17}{1}{1:13:46} \ts{17}{1}{1:14:07} \ts{17}{1}{1:14:50} \ts{17}{1}{1:15:32} \ts{17}{1}{1:15:52} \ts{17}{1}{1:16:13}
\end{answer}
\begin{coursenote}[Not said in class: the computation, and both factors of 2]
With $\rho_{\mathrm{spin}}(X_k)=-\frac{\mathrm i}{2}\sigma_k$ one has $K_{\rho_{\mathrm{spin}}}(X_i,X_j)=\mathrm{tr}\bigl(-\tfrac14\sigma_i\sigma_j\bigr)=-\tfrac12\delta_{ij}$. The dual basis is therefore $\tilde X_i=-2X_i$ and
\[
  \Omega_{\mathrm{spin}}=\sum_i\rho(X_i)\rho(\tilde X_i)=-2\sum_i\Bigl(-\tfrac14\Bigr)\mathbb 1=\tfrac32\,\mathbb 1=\frac{\dim L}{\dim V}\mathbb 1 .
\]
For the vector representation $K_{\rho_V}(X_i,X_j)=-2\delta_{ij}$, $\tilde X_i=-\frac12X_i$ and $\Omega_V=\mathbb 1=\frac33\mathbb 1$ (as above). So the class has the right numbers when the dual basis is taken \emph{with respect to $K_\rho$}. If instead one uses always the \emph{same} basis, dual with respect to the adjoint Killing form $K=-2\delta$ ($\tilde X_i=-\frac12X_i$ for every $\rho$), then $-\frac12\sum_i\rho(X_i)^2=\frac38$ in the spin and $1$ in the vector representation, that is $\frac12j(j+1)$ with $j=\frac12,1$: this is the lecture's ``factor $2$ with respect to physics'', where $\vec J^2=-\sum_i\rho(X_i)\rho(X_i)=j(j+1)$ ($\frac34$ and $2$). So both statements of the lecture are right, they refer to two different normalisations. (The first evaluation, $\frac32=\frac12(\frac12+1)$, was a slip: $\frac12(\frac12+1)=\frac34$ and $\frac32=2j(j+1)$.) \ts{17}{1}{1:09:36}
\end{coursenote}
So different representation spaces, at least when they differ in dimension, are labelled by different numbers in front of the identity in the Casimir operator, and one says: the particle has spin $\frac12$ or spin $1$, and the Casimir operator gives the same value on any representing vector. The lecturer still maintains the factor $2$ between mathematics and physics, but the other comment is important \ts{17}{1}{1:16:35} \ts{17}{1}{1:16:57} \ts{17}{1}{1:17:17}.

% ------------------------------------------------------------------
\section{Representations of a Lie group}\label{sec:rp-group}

Now Lie groups \ts{17}{1}{1:17:44}.
\begin{definition}[New concepts: representation of a Lie group]\label{def:rp-grouprep}
Let $(G,\bullet)$ be a finite-dimensional Lie group. A \textbf{representation} of $G$ on a vector space $V$ is a map
\[
  R\colon G\to GL(V),
\]
where $GL(V)\subseteq\mathrm{End}(V)$ is the set of linear maps $\varphi\colon V\to V$ with $\det\varphi\ne0$, that is, the \emph{invertible} endomorphisms (over $\C^n$ or $\R^n$ this is the general linear group). Since the composition of invertible maps is invertible, $\circ$ is a group law on $GL(V)$, and one requires compatibility with the group operations:
\[
  R(g_1\bullet g_2)=R(g_1)\circ R(g_2)\qquad\text{for all }g_1,g_2\in G ,
\]
that is, $R$ is a group homomorphism. The map $R$ should also be smooth (the lecture does not say it explicitly [?]). \ts{17}{1}{1:18:08} \ts{17}{1}{1:18:29} \ts{17}{1}{1:18:50} \ts{17}{1}{1:19:10} \ts{17}{1}{1:19:31} \ts{17}{1}{1:19:51} \ts{17}{1}{1:20:12} \ts{17}{1}{1:20:33}
\end{definition}
It looks very similar to the algebra case, but $GL(V)$ is ``a little less'' than $\mathrm{End}(V)$: the requirement $\det\ne0$ is needed so that the images of group elements, which have inverses, are themselves invertible.

\textbf{Example: $SO(2)$ on $\R^2$} \ts{17}{1}{1:20:55}. As a manifold $SO(2)$ is $S^1$, and the group operation moves points along the circle \ts{17}{1}{1:21:17}. To write down a representation one uses a chart $\varphi\colon U\to\R$, the angle map, defined on the circle with one point removed (otherwise the domain would not be open and one cannot cover the circle by one chart) \ts{17}{1}{1:21:37} \ts{17}{1}{1:21:58} \ts{17}{1}{1:22:18}. The group element with coordinate $\phi$ is $\varphi^{-1}(\phi)$, and it is represented by the rotation matrix
\[
  R\bigl(\varphi^{-1}(\phi)\bigr)=\begin{pmatrix}\cos\phi&\sin\phi\\-\sin\phi&\cos\phi\end{pmatrix}.
\]
By the addition theorems for cosine and sine, adding angles equals composing matrices:
\[
  R\bigl(\varphi^{-1}(\phi_1+\phi_2)\bigr)=R\bigl(\varphi^{-1}(\phi_1)\bigr)\circ R\bigl(\varphi^{-1}(\phi_2)\bigr),
\]
so this is a representation of the Lie group $SO(2)$ on $\R^2$, which is usually how one first meets it \ts{17}{1}{1:22:39} \ts{17}{1}{1:23:21} \ts{17}{1}{1:23:42} \ts{17}{1}{1:24:02} \ts{17}{1}{1:24:23} \ts{17}{1}{1:24:44}. The inverse chart is only needed to say which group elements are represented: on groups this is just more complicated than on $\R^n$.
\begin{coursenote}[Not said in class: the whole group, and the group law]
The lecture says one has a representation ``not on the entire group, because one chooses coordinates to write it down, but on $U$''. In fact the matrix is $2\pi$-periodic in $\phi$, so it is well defined also at the removed point and $R$ is defined on all of $SO(2)$. The relation $\phi_1+\phi_2$ on the left is in $\R$ only as long as the sum stays in the chart; globally it is addition mod $2\pi$. \ts{17}{1}{1:22:39}
\end{coursenote}
\textbf{The adjoint representation of a Lie group} \ts{17}{1}{1:25:07}. Are there non-trivial representations of every Lie group? Yes, and we can construct one that goes by the same name as the one for algebras, written with a capital $\mathrm{Ad}$ (for the group), versus the small $\mathrm{ad}$ (for the algebra): they are related, but a priori different. It is built in two steps \ts{17}{1}{1:25:28} \ts{17}{1}{1:25:48}.
\begin{definition}[New concepts: the adjoint map and the adjoint representation of a Lie group]\label{def:rp-Ad}
For $g\in G$ the \textbf{adjoint map} is
\[
  \mathrm{Ad}_g\colon G\to G,\qquad \mathrm{Ad}_g(h):=g\bullet h\bullet g^{-1},
\]
a left translation with $g$ followed by a right translation with $g^{-1}$ (not a linear map: there is no linearity of any kind here). \ts{17}{1}{1:26:10} \ts{17}{1}{1:26:31} \ts{17}{1}{1:26:53} It fixes the identity: $\mathrm{Ad}_g(e)=g\bullet e\bullet g^{-1}=e$ \ts{17}{1}{1:27:13} \ts{17}{1}{1:27:34}. So its push-forward (\cref{def:tb-push}) at the identity goes from the tangent space at $e$ to the tangent space at $\mathrm{Ad}_g(e)=e$ again \ts{17}{1}{1:27:54} \ts{17}{1}{1:28:15} \ts{17}{1}{1:28:36}:
\[
  (\mathrm{Ad}_g)_*\colon T_eG\to T_eG ,
\]
a linear map from the space isomorphic to the Lie algebra (\cref{thm:ll-iso}) to itself \ts{17}{1}{1:28:56}. It is invertible, since its inverse is $(\mathrm{Ad}_{g^{-1}})_*$, because $\mathrm{Ad}_{g^{-1}}$ is the inverse of $\mathrm{Ad}_g$. The \textbf{adjoint representation of the Lie group} is
\[
  \mathrm{Ad}\colon G\to GL(T_eG),\qquad g\mapsto(\mathrm{Ad}_g)_* .
\]
\ts{17}{1}{1:29:17} \ts{17}{1}{1:29:39} \ts{17}{1}{1:30:00} \ts{17}{1}{1:30:21}
\end{definition}
So a group has always a representation on the vector space $T_eG\cong L(G)$ that comes with it: one can find a representation of the Lie algebra on the Lie algebra (\cref{sec:rp-triv}), and a representation of the Lie group on the Lie algebra. If one has no imagination and is just given a Lie group and knows about the Lie algebra, ``with a group I always have a vector space that comes with it'', and the map that does the job is $\mathrm{Ad}$ \ts{17}{1}{1:30:41} \ts{17}{1}{1:31:02} \ts{17}{1}{1:31:22}.
\begin{coursenote}[Not said in class: why it is a representation, and how $\mathrm{Ad}$ and $\mathrm{ad}$ are related]
Homomorphism: $\mathrm{Ad}_{g_1g_2}=\mathrm{Ad}_{g_1}\circ\mathrm{Ad}_{g_2}$, so by the chain rule for push-forwards $(\mathrm{Ad}_{g_1g_2})_*=(\mathrm{Ad}_{g_1})_*\circ(\mathrm{Ad}_{g_2})_*$; and $g\mapsto(\mathrm{Ad}_g)_*$ is smooth because the group operations are. The relation between the two: the derivative of $\mathrm{Ad}$ at the identity is $\mathrm{ad}$, that is, $\mathrm{ad}_a=\frac{\mathrm d}{\mathrm dt}\big|_{t=0}\mathrm{Ad}_{\exp(ta)}$ (``a representation of the group differentiates to a representation of the algebra'', a statement for \cref{ch:l18}; no proof given). \ts{17}{1}{1:25:28}
\end{coursenote}
The lecture closes: finite-dimensional Lie algebras and Lie groups have been discussed; the next time one finally goes back from the algebra up to the group (and, in general, one \emph{cannot} hit the entire group), through the \emph{exponential map}, a concept from differential geometry which on Lie groups has even nicer properties than on general smooth manifolds \ts{17}{1}{1:31:43} \ts{17}{1}{1:32:03} \ts{17}{1}{1:32:23}. This is \cref{ch:l18}.

% ==================================================================
\section*{Key concepts}
\begin{keyconcepts}
\textbf{Retrieval questions.}
\begin{enumerate}
  \item State the definition of a representation of a Lie algebra. Which bracket appears on each side?
  \item Give three different representations of Lie algebras from the lecture, and say on which spaces they live.
  \item When is a representation reducible? Why must the invariant subspace be non-trivial? What is a faithful representation?
  \item Define $K_\rho$ and the Casimir operator $\Omega_\rho$. What do they depend on?
  \item Prove that $c_\rho=\dim L/\dim V$ for an irreducible $\rho$. Where is Schur's lemma used?
  \item Compute $K_\rho$ and $\Omega_\rho$ for $\rho_{\mathrm{spin}}$ of $\mathfrak{so}(3)$. Where does the factor $2$ with respect to physics come from?
  \item Define a representation of a Lie group and the group $GL(V)$. Give the example of $SO(2)$ and the adjoint representation $\mathrm{Ad}$.
\end{enumerate}
\textbf{Concepts.}
\begin{itemize}
  \item Representation of a Lie algebra, representation space (\cref{def:rp-rep}); examples (\cref{sec:rp-examples}); trivial and adjoint representation (\cref{sec:rp-triv}).
  \item Reducible, irreducible, faithful (\cref{def:rp-irred,def:rp-faithful}).
  \item Killing form of a representation, Casimir operator, Schur's lemma, $c_\rho=\dim L/\dim V$ (\cref{def:rp-killing,def:rp-casimir,thm:rp-schur,thm:rp-value}).
  \item Representation of a Lie group and $\mathrm{Ad}$ (\cref{def:rp-grouprep,def:rp-Ad}).
\end{itemize}
\end{keyconcepts}

\flushside
\end{paracol}

% !TEX root = ../main.tex
% STUB: generated by tools/stubs.py; write the chapter from materials/transcripts/C18-V01.txt
\chapter{Reconstruction of a Lie group from its algebra}\label{ch:l18}
\begin{paracol}{2}

\begin{sources}
\begin{tabularx}{\linewidth}{@{}l l L@{}}
\toprule
\textbf{Video} & \textbf{Length} & \textbf{Content}\\
\midrule
\seg{18}{1}{Lecture 18: Reconstruction of a Lie group from its algebra} & 0:44:10 & (to be written)\\
\bottomrule
\end{tabularx}

\smallskip
\textbf{Transcript} \file{materials/transcripts/C18-V01.txt}.
\end{sources}

\begin{reading}[Where this lecture sits in the other notes]
(to be written)
\end{reading}

\emph{This chapter has not been written yet.}

\flushside
\end{paracol}

% !TEX root = ../main.tex
\chapter{Principal fibre bundles}\label{ch:l19}
\providecommand{\lact}{\mathbin{\triangleright}}
\providecommand{\ract}{\mathbin{\triangleleft}}
\providecommand{\lactf}{\mathbin{\blacktriangleright}}
\providecommand{\ractf}{\mathbin{\blacktriangleleft}}
\begin{paracol}{2}

\begin{sources}
\begin{tabularx}{\linewidth}{@{}l l L@{}}
\toprule
\textbf{Video} & \textbf{Length} & \textbf{Content}\\
\midrule
\seg{19}{1}{Lecture 19: Principal fibre bundles} & 2:33:32 & Motivation (frame bundle, gauge theory, densities, spinors); left and right Lie group actions, the right action defined from a left one; equivariant maps; orbits, orbit space, stabilisers, free actions; the definition of a principal $G$-bundle; the frame bundle $LM$ as a principal $GL(d,\R)$-bundle; principal bundle maps (same group, and with a group homomorphism $\rho$; the spin frame bundle); the lemma that a principal bundle map over the identity is an isomorphism; trivial principal bundles and the theorem ``trivial $\Leftrightarrow$ there is a section'', with both directions of the proof; the frame bundle of $S^2$ is not trivial.\\
\bottomrule
\end{tabularx}

\smallskip
\textbf{Transcript} \file{materials/transcripts/C19-V01.txt} (YouTube automatic captions; the blackboard symbols for the two actions are reconstructed from the verbal description: ``arrow to the right'' for the left action, ``arrow to the left'' for the right action, a filled and an empty triangle for two different actions).
\end{sources}

\begin{reading}[Where this lecture sits in the other notes]
Not treated in the other courses (only here). Bundles were introduced in \cref{ch:l06} (bundle, bundle morphism, isomorphism, section, trivial bundle: \cref{def:mf-bundle,def:mf-morphism,def:mf-iso,def:mf-section,def:mf-trivial}); the lecturer asks to revise pull-backs, sub-bundles and restrictions (\cref{sec:mf-morph,sec:mf-pullback}). It uses Lie groups (\cref{ch:l13}), representations (\cref{ch:l17}), the exponential map (\cref{ch:l18}) and the hairy ball theorem (\cref{ch:l11}). Next: the associated bundles (\cref{ch:l20}). (From the lecture titles; not checked in detail.)
\end{reading}

% ==================================================================
\section{Motivation}\label{sec:pb-motivation}

``We have come a long way from manifolds over Lie groups, Lie algebras and back, and we are now fully prepared to come to the at least technical climax of this course'': the chapter on principal fibre bundles \ts{19}{1}{0:00} \ts{19}{1}{0:25}. What is a principal fibre bundle and why is it so important a notion? Since the technical development will be heavy at some points, it is good to know first what it is all about \ts{19}{1}{0:45}. Very roughly speaking, \emph{a principal fibre bundle is a bundle whose fibre is a Lie group}; to make this precise is exactly the point of the technical developments. This is why Lie groups are so important: they provide the fibres of fibre bundles \ts{19}{1}{1:05} \ts{19}{1}{1:26}.

Why is it so immensely important? Because principal bundles allow one to understand \emph{any} fibre bundle with a fibre $F$ on which the Lie group $G$ acts (the lecturer says ``Lie algebra'' by a slip \ts{19}{1}{2:14}); the fibre bundles on which $G$ acts and that are \emph{associated} with the principal bundle are called \textbf{associated bundles} \ts{19}{1}{1:47} \ts{19}{1}{2:35} \ts{19}{1}{2:57} (\cref{ch:l20}).

Fibre bundles are ubiquitous in physics, with a few examples \ts{19}{1}{3:17}:
\begin{itemize}
  \item In \emph{general relativity} they appear with the Lorentz group $SO(1,3)$ as the group, or, if one considers spinors, with $SL(2,\C)$, the double cover of $SO(1,3)$ \ts{19}{1}{3:38}.
  \item In \emph{Yang--Mills theory} and its generalisations (non-abelian gauge theories), the co-vector potential of ``electrodynamics'' all of a sudden becomes Lie-algebra valued; in the standard model the groups are $SU(2)$ (roughly, the electroweak theory) or $SU(3)$ (the strong interaction). These theories, and the subtleties in them, can only be properly understood by considering principal fibre bundles \ts{19}{1}{3:59} \ts{19}{1}{4:20} \ts{19}{1}{4:40}.
  \item In \emph{mathematics}: studying how the fibres of a principal bundle \emph{twist} (the bundle need not be trivial) teaches a lot about the manifold; the field of K-theory is mentioned, where one studies the bundles one can construct over a manifold $M$ \ts{19}{1}{5:01} \ts{19}{1}{5:22} \ts{19}{1}{5:43} \ts{19}{1}{6:04}.
  \item In \emph{differential geometry}: the definition of a tensor density is tricky without bundles. One looks at the principal bundle whose fibres are the general linear group of the appropriate dimension (the \emph{frame bundle}), and then the corresponding associated bundle on which $GL(d,\R)$ acts on tensor fibres with an additional determinant factor: these are the densities. The heavily used tensor densities of general relativity are thus properly understood through principal bundles and their associated bundles \ts{19}{1}{6:25} \ts{19}{1}{6:46} \ts{19}{1}{7:06}.
\end{itemize}
\begin{coursenote}[K-theory and cohomology]
Not said in class: in the lecture the cohomology groups of $M$ are said to be ``essentially defined as the various bundles one can construct over $M$''. This is a loose statement: K-theory groups are built from the \emph{vector bundles} over $M$ (isomorphism classes with a sum), not the de Rham cohomology groups of \cref{ch:l12}; the two are related (for instance through Chern classes) but not the same. \ts{19}{1}{5:22}
\end{coursenote}
``So this is lots of motivation which at the moment you do not understand, but what you do understand is: there is a Lie group, a Lie group, a Lie group \ldots{} and that is the principal fibre bundle that we want to develop'' (and a joke asking to mention the motivation in the course evaluation) \ts{19}{1}{7:26}.

% ==================================================================
\section{Lie group actions on a manifold}\label{sec:pb-actions}

We need to prepare the whole business: Lie group actions on a manifold, a topic that could have been discussed with Lie groups in \cref{ch:l13} but fits here so well that it is a preparation (section 5.1 of the lecture) \ts{19}{1}{7:46} \ts{19}{1}{8:07}.

\begin{definition}[New concepts: left action of a Lie group]\label{def:pb-left}
Let $(G,\bullet)$ be a Lie group and $M$ a smooth manifold. A smooth map
\[
  \lact\colon G\times M\to M,\qquad(g,p)\mapsto g\lact p ,
\]
(smooth, because $G\times M$ is made into a smooth manifold in the usual way) satisfying
\begin{enumerate}
  \item[(i)] $\mathbb 1\lact p=p$ for all $p\in M$ (the identity does nothing),
  \item[(ii)] $g_2\lact(g_1\lact p)=(g_2\bullet g_1)\lact p$ for all $g_1,g_2\in G$ and $p\in M$,
\end{enumerate}
is called a \textbf{left $G$-action} on $M$. \ts{19}{1}{8:29} \ts{19}{1}{8:50} \ts{19}{1}{9:11} \ts{19}{1}{9:31} \ts{19}{1}{9:51} \ts{19}{1}{10:14} \ts{19}{1}{10:35}
\end{definition}
(A map defined on a Cartesian product has its symbol put in between the two arguments, ``like the addition you know from elementary school''; in (ii) the brackets are for comfort only, there is just one way to read the expression; the order $g_2,g_1$ is kept, one can act one after the other or at once with the product \ts{19}{1}{10:14} \ts{19}{1}{10:35}.) The lecture says ``left action of the Lie algebra of the Lie group $G$'' at \ts{19}{1}{10:56} by a slip: it is the group that acts. This is ``a central notion we are going to use'' \ts{19}{1}{10:56} \ts{19}{1}{11:17}.

\begin{example}[Representations are left actions]\label{ex:pb-rep-action}
Let $G$ be a Lie group, $V$ a representation space with a linear representation $R\colon G\to GL(V)$ (\cref{ch:l17}; one needs $GL(V)$ rather than $\mathrm{End}(V)$ so that $R(g)$ is invertible). Take $M=V$ (a smooth manifold) and define $g\lact v:=R(g)v$. This is a left action. \ts{19}{1}{11:17} \ts{19}{1}{11:40} \ts{19}{1}{12:03} \ts{19}{1}{12:24} \ts{19}{1}{12:44}
\end{example}
The representations of $G$ on a vector space, as usually found in the introductory lectures of physics, are a very special case of left actions \ts{19}{1}{13:08} \ts{19}{1}{13:30}. One is almost tempted to say that left actions are ``non-linear representations'', but ``non-linear'' would suggest that it could be linear: on a general manifold there is no vector space structure, so one would not even know what linear means. The name is misleading; the idea is that one loosens the restriction \ts{19}{1}{13:30} \ts{19}{1}{13:51}.

\begin{definition}[New concepts: right action of a Lie group]\label{def:pb-right}
A smooth map $\ract\colon M\times G\to M$, $(p,g)\mapsto p\ract g$, is a \textbf{right $G$-action} if
\begin{enumerate}
  \item[(i)] $p\ract\mathbb 1=p$ for all $p\in M$,
  \item[(ii)] $(p\ract g_1)\ract g_2=p\ract(g_1\bullet g_2)$ for all $g_1,g_2\in G$ and $p\in M$.
\end{enumerate}
\ts{19}{1}{14:11} \ts{19}{1}{14:31} \ts{19}{1}{14:52} \ts{19}{1}{15:12} \ts{19}{1}{15:33} \ts{19}{1}{15:54}
\end{definition}
Whether to let a group act from the left or from the right is in principle a matter of convention. In \emph{principal} bundle theory we shall always use \emph{right} actions; for the fibre bundles on which the same group acts (the associated bundles) we shall use \emph{left} actions, and the two play nicely together \ts{19}{1}{14:11} \ts{19}{1}{14:31} \ts{19}{1}{14:52}. The point of writing the group element on the right: in (ii) the order $g_1\bullet g_2$ is the order in which one reads the acting elements, whereas for a left action the first one acting ($g_1$) stands on the right of the product, so it is convenient to write the symbol on the side of the action: ``you can always keep the order, you do not need to think about it''. It is nevertheless something really different \ts{19}{1}{16:14} \ts{19}{1}{16:35}.

\begin{proposition}[A right action from a left action]\label{prop:pb-right-from-left}
If $\lact$ is a left $G$-action on $M$, then
\[
  p\ract g:=g^{-1}\lact p
\]
is a right $G$-action on $M$. \ts{19}{1}{16:35} \ts{19}{1}{16:55} \ts{19}{1}{17:16} \ts{19}{1}{17:38}
\end{proposition}
\begin{proof}
It is smooth: the inversion of a Lie group is smooth (``one of the points where this property is important; it was implicitly important before'') \ts{19}{1}{18:40}. Property (i): $p\ract\mathbb 1=\mathbb 1^{-1}\lact p=\mathbb 1\lact p=p$ \ts{19}{1}{17:58} \ts{19}{1}{18:20}. Property (ii):
\begin{align*}
  (p\ract g_1)\ract g_2&=g_2^{-1}\lact(g_1^{-1}\lact p)\\
  &=(g_2^{-1}\bullet g_1^{-1})\lact p\\
  &=(g_1\bullet g_2)^{-1}\lact p=p\ract(g_1\bullet g_2),
\end{align*}
by the left action property and because the inverse of a product is the product of the inverses \emph{in the opposite order}. \ts{19}{1}{19:01} \ts{19}{1}{19:21} \ts{19}{1}{19:42} \ts{19}{1}{20:03} \ts{19}{1}{20:23}
\end{proof}
It is crucial that one takes the inverse: it flips the order. Without the inverse the order would not be flipped and one would not get a right action. If left and right actions were the same, this would be mere notation, but ``because of these properties a right action is a different thing'' \ts{19}{1}{20:44} \ts{19}{1}{21:04}.

\subsection{The picture to keep in mind: bases and components}\label{sec:pb-bases}
Once everything is expressed in terms of principal and associated bundles, we shall see the following recipe \ts{19}{1}{21:28}. Recall that a basis $e_1,\dots,e_{\dim M}$ of a tangent space is labelled by \emph{lower} indices, the components of a vector $X\in T_pM$ by \emph{upper} indices (``this is just a convention'') \ts{19}{1}{21:51} \ts{19}{1}{22:15} \ts{19}{1}{22:38}. If one defines a new basis by a matrix $A$, $\tilde e_m=A^n{}_me_n$, then the new components are given by precisely the inverse transformation, $\tilde X^m=(A^{-1})^m{}_nX^n$, because the linear combination of the new components with the new basis must still be the same vector: the vector is the vector, only expressed differently \ts{19}{1}{23:07} \ts{19}{1}{23:27} \ts{19}{1}{23:48}. This $A$ and $A^{-1}$ will be understood as a \emph{right} action of the general linear group $GL(d,\R)$ on the basis and a \emph{left} action of the same group on the components: ``one must be the inverse of the other'' \ts{19}{1}{24:11} \ts{19}{1}{24:32}. The bundle of bases will be the \emph{frame bundle}, the principal bundle we are going to define, and the bundle of components (say an $\R^4$ bundle) is just a bundle on which the group also acts \ts{19}{1}{24:32} \ts{19}{1}{24:53}.

This is how ``a tensor is what transforms as a tensor does'' is often defined in physics; to say precisely what it means (``tensor components are things that transform as tensor components would'') one needs principal fibre bundles \ts{19}{1}{25:14} \ts{19}{1}{25:34}. This was just a remark to show where we are going.

\subsection{Equivariant maps}\label{sec:pb-equivariant}
\begin{definition}[New concepts: $\rho$-equivariant map]\label{def:pb-equivariant}
Let $(G,\bullet)$ and $(H,\star)$ be Lie groups, $\rho\colon G\to H$ a Lie group homomorphism (smooth, and $\rho(g_1\bullet g_2)=\rho(g_1)\star\rho(g_2)$; not a representation here, but a map between groups), $M$ and $N$ smooth manifolds with a left $G$-action $\lact$ on $M$ and a left $H$-action $\lactf$ on $N$. A smooth map $f\colon M\to N$ is called \textbf{$\rho$-equivariant} if the diagram
\begin{center}
\begin{tikzpicture}[>=Stealth,node distance=3.4cm,baseline=(current bounding box.center)]
\node (GM) {$G\times M$};\node (HN) [right=of GM] {$H\times N$};
\node (M) [below=2.2cm of GM] {$M$};\node (N) [below=2.2cm of HN] {$N$};
\draw[->] (GM)--node[above,font=\small]{$\rho\times f$}(HN);
\draw[->] (M)--node[below,font=\small]{$f$}(N);
\draw[->] (GM)--node[left,font=\small]{$\lact$}(M);
\draw[->] (HN)--node[right,font=\small]{$\lactf$}(N);
\end{tikzpicture}
\end{center}
commutes, where $\rho\times f$ means that $\rho$ eats the first factor and $f$ the second; that is,
\[
  f(g\lact m)=\rho(g)\lactf f(m)\qquad\text{for all }g\in G,\ m\in M .
\]
\ts{19}{1}{25:56} \ts{19}{1}{26:18} \ts{19}{1}{26:40} \ts{19}{1}{27:00} \ts{19}{1}{27:22} \ts{19}{1}{27:43} \ts{19}{1}{28:04} \ts{19}{1}{28:28} \ts{19}{1}{28:52} \ts{19}{1}{29:14} \ts{19}{1}{29:35} \ts{19}{1}{29:56} \ts{19}{1}{30:17} \ts{19}{1}{30:38} \ts{19}{1}{30:58}
\end{definition}
Commuting means: one can first act with $g$ on $m$ by the left action on $M$ and then apply $f$, or first map $g$ by $\rho$ and $m$ by $f$ and then use the action on $N$ \ts{19}{1}{30:17} \ts{19}{1}{30:38}. So equivariant maps are the structure preserving maps between manifolds when the structure is a left action of a Lie group on each manifold, the groups being related by a homomorphism: it is again the same principle as always, structure preserving maps between objects \ts{19}{1}{31:21} \ts{19}{1}{31:42} \ts{19}{1}{32:02}. One can have $H=G$, but in physics the case $G=SO(1,3)$, $H=SL(2,\C)$ occurs: if one transforms the measurement apparatus (a Lorentz transformation, ``I rotate it or I put it into a moving system'') and the apparatus measures the spin of an object, one needs to know what the Lorentz transformation does to a spin particle, so one needs a $2:1$ homomorphism relating the two groups \ts{19}{1}{32:24} \ts{19}{1}{32:45} \ts{19}{1}{33:05} \ts{19}{1}{33:27}. (The lecture states the homomorphism in the direction $SO(1,3)\to SL(2,\C)$ here; the double covering map goes the other way, from $SL(2,\C)$ to the (orthochronous) Lorentz group: see \cref{sec:pb-pbmap}.) ``It looks excessively mathematical but it is real life, at least in theoretical physics.''

\subsection{Orbits, orbit space, stabilisers, free actions}\label{sec:pb-orbits}
\begin{definition}[New concepts: orbit]\label{def:pb-orbit}
Let $\lact\colon G\times M\to M$ be a left action. The \textbf{orbit} of a point $p\in M$ is the subset of $M$
\[
  \mathcal O_p:=\{\,q\in M\mid\exists g\in G:\ q=g\lact p\,\}
\]
of all points that can be reached from $p$ by acting with some group element. \ts{19}{1}{33:49} \ts{19}{1}{34:09} \ts{19}{1}{34:31} \ts{19}{1}{34:52}
\end{definition}
(The curly $\mathcal O$ is not a topology, it is an orbit.) \textbf{Example} \ts{19}{1}{35:16}: $M=\R^2$, $G=SO(2)$, and the left action is the vector representation,
\[
  \varphi\lact(p_1,p_2)=\begin{pmatrix}\cos\varphi&\sin\varphi\\-\sin\varphi&\cos\varphi\end{pmatrix}\begin{pmatrix}p_1\\p_2\end{pmatrix}
\]
(in quotation marks, since one has to choose coordinates to write this), the usual rotations of the plane. The orbit of a point $p\neq0$ is the circle through $p$ around the origin; the origin cannot be reached from there. The orbits are subsets of the manifold: all the circles, and also one point, the origin, which is the orbit of itself \ts{19}{1}{35:38} \ts{19}{1}{36:00} \ts{19}{1}{36:20} \ts{19}{1}{36:42} \ts{19}{1}{37:02} \ts{19}{1}{37:24}.

\begin{definition}[New concepts: orbit space]\label{def:pb-orbitspace}
For $p,q\in M$ let $p\sim q$ if there exists $g\in G$ with $q=g\lact p$ (equivalently $p=g^{-1}\lact q$): two points are equivalent if they lie on the same orbit. This is an equivalence relation, and the quotient set (\cref{def:cl-quotient}) is the \textbf{orbit space}, written
\[
  M/G:=M/\!\sim .
\]
One element of the orbit space is exactly the set of points that makes up an orbit. \ts{19}{1}{37:46} \ts{19}{1}{38:07} \ts{19}{1}{38:27} \ts{19}{1}{38:47} \ts{19}{1}{39:08} \ts{19}{1}{39:30} \ts{19}{1}{39:52}
\end{definition}
Why it is an equivalence relation: reflexive, because the identity is in the group; symmetric, because every element has an inverse (if $q=g\lact p$, then $g^{-1}\lact q=g^{-1}\lact(g\lact p)=(g^{-1}\bullet g)\lact p=\mathbb 1\lact p=p$); transitivity ``we do not do'' \ts{19}{1}{38:27} \ts{19}{1}{38:47} \ts{19}{1}{39:08}. (Not said in class: if $q=g\lact p$ and $r=h\lact q$, then $r=(h\bullet g)\lact p$.) Because the quotient is defined only using the group and the action, one writes $M/G$ and does not carry the construction around. In the example the equivalence classes are all the circles and the origin: a foliation of the space into orbits. ``This orbit space will play a very important role in principal fibre bundle theory'' \ts{19}{1}{39:52} \ts{19}{1}{40:13} \ts{19}{1}{40:34}.

\begin{definition}[New concepts: stabiliser]\label{def:pb-stabiliser}
The \textbf{stabiliser} of a point $p\in M$ is the subset of the \emph{group}
\[
  S_p:=\{\,g\in G\mid g\lact p=p\,\}
\]
of all group elements that keep the point invariant. \ts{19}{1}{40:55} \ts{19}{1}{41:16} \ts{19}{1}{41:37}
\end{definition}
Always keep in mind: the orbit is a subset of the space, the stabiliser a subset of the group \ts{19}{1}{41:37}. In the same example the stabiliser of $p\ne0$ under the rotations is just the identity, whereas the stabiliser of the origin is the entire group \ts{19}{1}{41:58} \ts{19}{1}{42:18} \ts{19}{1}{42:41}.

\begin{definition}[New concepts: free action]\label{def:pb-free}
A left (or right) action is called \textbf{free} if for \emph{every} point $p\in M$ the stabiliser consists of the identity only, $S_p=\{\mathbb 1\}$. \ts{19}{1}{43:03} \ts{19}{1}{43:25} \ts{19}{1}{43:46}
\end{definition}
So the $SO(2)$-action on $\R^2$ above is \emph{not} free (the origin is fixed by everything). A free action is one where every point is stabilised only by the identity.
\begin{answer}{Isn't $\varphi=2\pi$ a non-trivial element fixing every point?}[Free action and the circle]
A student objects that, for the rotations, one can choose another element of the group, $\varphi=2\pi$, that also leaves the point where it is. The lecturer: it is a very good example of what was said before about the exponential map \cref{ch:l18}. The rotation is $\exp(\varphi E)$ with the generator $E$; the group $SO(2)$ is compact, the group manifold is $S^1$, hence $\exp$ cannot be injective: one hits the same group element for $\varphi=0,2\pi,4\pi,\dots$. But this speaks about the \emph{parameter}, not about the group: every time one gets the identity \emph{element}. So there are no other elements than the identity that fix the point; the feeling that there must be more comes from the non-injective generation of the group elements by $\exp$ \ts{19}{1}{44:06} \ts{19}{1}{44:28} \ts{19}{1}{44:49} \ts{19}{1}{45:09} \ts{19}{1}{45:30} \ts{19}{1}{45:51} \ts{19}{1}{46:11}
\end{answer}
Can an action defined by a linear representation ever be free? No: a linear map always sends the zero vector to the zero vector, so such an action is never free \ts{19}{1}{46:33} \ts{19}{1}{46:53}. What about the translations of the plane? A translation moves every point, apart from the trivial one: so the translation action is free \ts{19}{1}{46:53} \ts{19}{1}{47:15} \ts{19}{1}{47:36}.

There are other notions, \emph{transitive} actions, \emph{effective} actions and so on; they are defined along similar lines and ``nothing to be afraid of'' if one stumbles across them, but they are not needed now: for principal bundles we only need the free actions \ts{19}{1}{47:57} \ts{19}{1}{48:18} \ts{19}{1}{48:38}.
\begin{proposition}[Orbits of a free action]\label{prop:pb-orbit-G}
If the action is free, then each orbit $\mathcal O_p$ is, as a smooth manifold, diffeomorphic to the group $G$. \ts{19}{1}{49:03} \ts{19}{1}{49:27}
\end{proposition}
The lecture only states this (``I will ask you to show that on a problem sheet''). \emph{Not said in class}: the map $G\to\mathcal O_p$, $g\mapsto g\lact p$ is surjective by definition of the orbit and injective by freeness (if $g_1\lact p=g_2\lact p$ then $g_2^{-1}\bullet g_1\in S_p=\{\mathbb 1\}$); smoothness of the inverse needs an argument and is the exercise \cref{ex:pb-orbit-G}.
\begin{example}[The punctured plane]\label{ex:pb-punctured}
Let now $M=\R^2\setminus\{0\}$ and again $SO(2)$ acting by the vector representation. This is no longer a linear representation, since $M$ is not a vector space, but still a left action (``a very trivial modification of a linear representation''). The only point with non-trivial stabiliser has been removed, so the action is \emph{free}. The orbits are the circles, and each of them is diffeomorphic to $SO(2)$, because $SO(2)$ is the Lie group built on the circle $S^1$ as underlying manifold: ``in a sense it is $SO(2)$ seen purely as a smooth manifold''. \ts{19}{1}{49:50} \ts{19}{1}{50:13} \ts{19}{1}{50:33} \ts{19}{1}{50:54} \ts{19}{1}{51:14} \ts{19}{1}{51:34} \ts{19}{1}{51:56} \ts{19}{1}{52:16}
\end{example}
This is one of the purposes of free actions: one knows immediately that the orbits are isomorphic to the group \ts{19}{1}{52:16} \ts{19}{1}{52:37}.
\begin{exercise}[Orbits of a free action]\label{ex:pb-orbit-G}
Prove \cref{prop:pb-orbit-G}. (The long version is in \cref{exam:ch19}, \cref{ex:pb-full}.)
\end{exercise}

% ==================================================================
\section{Principal fibre bundles}\label{sec:pb-def}

After these preparations the central definition \ts{19}{1}{52:58}. Whenever we talk about bundles from now on we assume everything smooth: the spaces, the projection maps, all the maps; ``if I forget to mention smooth at some point you may insert it'' \ts{19}{1}{53:24}. Recall that a bundle consists of a total space, a projection map and a base space (\cref{def:mf-bundle}).

If $E$ is equipped with a right $G$-action one says that $E$ is a \emph{right $G$-space}; this is just terminology which one does not really need, but it is convenient \ts{19}{1}{54:27} \ts{19}{1}{54:51}.
\begin{definition}[New concepts: principal $G$-bundle]\label{def:pb-principal}
Let $G$ be a Lie group. A bundle $(E,\pi,M)$ is called a \textbf{principal $G$-bundle} if
\begin{enumerate}
  \item[(i)] $E$ is a right $G$-space, i.e.\ it carries a right $G$-action $\ract\colon E\times G\to E$;
  \item[(ii)] this right $G$-action is \emph{free};
  \item[(iii)] the bundle $(E,\pi,M)$ is isomorphic, as a bundle, to the bundle $(E,\rho,E/G)$, where $E/G$ is the orbit space of the right action and $\rho\colon E\to E/G$ the canonical projection, which sends a point $\varepsilon\in E$ to its equivalence class (its orbit).
\end{enumerate}
\ts{19}{1}{53:44} \ts{19}{1}{54:07} \ts{19}{1}{55:15} \ts{19}{1}{55:37} \ts{19}{1}{55:59} \ts{19}{1}{56:20} \ts{19}{1}{56:40} \ts{19}{1}{57:02} \ts{19}{1}{57:22}
\end{definition}
Isomorphism as bundles means (\cref{def:mf-iso}): there exist diffeomorphisms $u\colon E\to E$ and $f\colon M\to E/G$ such that the diagram
\begin{center}
\begin{tikzpicture}[>=Stealth,node distance=2.6cm,baseline=(current bounding box.center)]
\node (E) {$E$};\node (E2) [right=of E] {$E$};
\node (M) [below=2cm of E] {$M$};\node (EG) [below=2cm of E2] {$E/G$};
\draw[->] (E)--node[above,font=\small]{$u$}(E2);
\draw[->] (M)--node[below,font=\small]{$f$}(EG);
\draw[->] (E)--node[left,font=\small]{$\pi$}(M);
\draw[->] (E2)--node[right,font=\small]{$\rho$}(EG);
\end{tikzpicture}
\end{center}
commutes, $\rho\circ u=f\circ\pi$. ``One always has an object that consists of three elements (total space, projection, base space); an isomorphism must map each of them, and since there is additional structure, the projection, it has to be preserved: one does not need to remember this, one can always reconstruct it by thinking about what is structure preserving'' \ts{19}{1}{58:48} \ts{19}{1}{59:08} \ts{19}{1}{59:30} \ts{19}{1}{59:51} \ts{19}{1}{1:00:12} \ts{19}{1}{1:00:33} \ts{19}{1}{1:00:53} \ts{19}{1}{1:01:14}.

Why this is the right notion: one has a manifold $E$ with a free right $G$-action; every orbit is the group itself (\cref{prop:pb-orbit-G}); $E/G$ is the space of all orbits; and the fibres of $\rho$, the preimages $\rho^{-1}([\varepsilon])$ of a point of the orbit space, are exactly the orbits, hence diffeomorphic to the group. The condition says that the bundle $(E,\pi,M)$ looks like this one \ts{19}{1}{56:40} \ts{19}{1}{57:02} \ts{19}{1}{57:22} \ts{19}{1}{57:43}.
\begin{remark}[Observation: the fibres are the group]\label{rem:pb-fibres}
Since the right action is free, the fibres $\pi^{-1}(x)$ of a principal $G$-bundle are diffeomorphic to the Lie group $G$: if $\varepsilon\in\pi^{-1}(x)$, the fibre is the orbit of $\varepsilon$. And two points of the same fibre are related by a \emph{unique} group element. This is just an observation, not part of the definition. \ts{19}{1}{57:43} \ts{19}{1}{58:06} \ts{19}{1}{58:28}
\end{remark}
In this sense a principal fibre bundle is, as announced in the beginning, a bundle whose fibre is the Lie group: it is isomorphic to a bundle which is a $G$-fibre bundle. But in that isomorphism the $G$ is not immediately visible: we only know that $G$ can act on $E$. ``This is the simplest definition one can give that does what you want'' \ts{19}{1}{1:01:36} \ts{19}{1}{1:01:57}.
\begin{coursenote}[What the definition leaves implicit]
Not said in class: (1) taking $u=\mathrm{id}_E$ in the isomorphism is the usual (and the natural) choice: then the fibres of $\pi$ are exactly the orbits; for an arbitrary diffeomorphism $u$ the diagram says only that $\pi$ and $\rho$ have fibres of the same kind. (2) For $E/G$ to be a smooth manifold and $\rho$ a smooth bundle one needs more than freeness (for instance that the action is also \emph{proper}, as it is whenever $G$ is compact); the lecture does not discuss this. (3) The standard definition adds local triviality of the bundle explicitly; it is hidden here in the requirement that the bundle be isomorphic to $E\to E/G$. \ts{19}{1}{1:01:36}
\end{coursenote}
The lecturer asks to revise, until the next lecture, the definitions of pull-backs of bundles, sub-bundles and restrictions of bundles (\cref{sec:mf-morph,sec:mf-pullback}), because they will be needed \ts{19}{1}{1:01:57} \ts{19}{1}{1:02:19}.

% ------------------------------------------------------------------
\subsection{The frame bundle}\label{sec:pb-frame}
We now consider a particular example for a principal $G$-bundle, ``of the utmost importance in differential geometry'': the \textbf{frame bundle} $LM$ \ts{19}{1}{1:02:19} \ts{19}{1}{1:02:39}.
\begin{definition}[New concepts: frame, frame bundle $LM$]\label{def:pb-frame}
Let $M$ be a smooth manifold of dimension $d$. At a point $x\in M$ let
\[
  L_xM:=\{\,(e_1,\dots,e_d)\mid e_1,\dots,e_d\text{ is a basis of }T_xM\,\}
\]
the set of ordered $d$-tuples of vectors that form a basis of the tangent space: the \textbf{frames} at $x$. (It could be called the basis bundle, but it is tradition to call it frame bundle.) The \textbf{frame bundle} is the disjoint union
\[
  LM:=\bigsqcup_{x\in M}L_xM .
\]
\ts{19}{1}{1:03:02} \ts{19}{1}{1:03:24} \ts{19}{1}{1:03:45} \ts{19}{1}{1:04:49}
\end{definition}
Each $L_xM$ is in bijection with $GL(d,\R)$: if one expresses the vectors of a basis in terms of the canonical basis of $\R^d$ one gets column vectors, which, put next to each other, form an invertible matrix \ts{19}{1}{1:04:07} \ts{19}{1}{1:04:27}. (Not said in class: this uses a chosen reference frame at $x$ and is not canonical; in particular it is not a group isomorphism, $L_xM$ has no group structure of itself.)
\begin{aside}[The disjoint union: apples and pears]
Some had asked about this kind of disjoint union \ts{19}{1}{1:05:09}. The tangent bundle was defined as the disjoint union of all the $T_xM$. Think of the spaces $T_xM$ and $T_{x'}M$ at two points as entirely different spaces: the elements in the first are apples, in the second pears; put all the apples and pears into one bag, close the eyes and pull out an element of the union: ``in my hand I have an apple. Is there any doubt where the apple comes from?'' So if one takes a vector $X$ from the tangent bundle, one always knows that there exists a unique $x$ with $X\in T_xM$, and can define the projection that assigns to $X$ the base point $x$. A construction of the type $\bigcup_xT_xM\times\{x\}$ (pairs of a point and a tangent vector) makes the point explicit, but does not solve the problem of the dimension: it is first of all just a set, and a set has no dimension; one needs to make it into a topological space and then a manifold. \ts{19}{1}{1:05:30} \ts{19}{1}{1:05:50} \ts{19}{1}{1:06:11} \ts{19}{1}{1:06:32} \ts{19}{1}{1:06:52} \ts{19}{1}{1:07:13} \ts{19}{1}{1:07:34} \ts{19}{1}{1:07:55} \ts{19}{1}{1:08:15}
\end{aside}
The manifold structure is built in complete analogy to the tangent bundle (\cref{def:tb-bundle,def:tb-chart}): $LM$ is equipped with a smooth atlas inherited from the smooth atlas of $M$. One takes an element of $LM$, expresses it in a chart-induced basis, decomposes each of its vectors, and gets a square array of numbers $e^a{}_b$ which coordinatise the element, together with the $d$ chart coordinates of the base point. The dimension of the smooth manifold $LM$ is therefore
\[
  \dim LM=\dim M+(\dim M)^2=\dim M\,(1+\dim M) .
\]
\ts{19}{1}{1:08:37} \ts{19}{1}{1:08:58} \ts{19}{1}{1:09:19} \ts{19}{1}{1:09:40} \ts{19}{1}{1:10:01} \ts{19}{1}{1:10:22}

What is still missing is the projection map. The \textbf{projection} $\pi\colon LM\to M$ is the ``apples and pears'' map: an element $(e_1,\dots,e_d)$ of $LM$ lies in $L_xM$ for a \emph{unique} $x\in M$ (the union is disjoint), and $\pi(e_1,\dots,e_d):=x$. So $(LM,\pi,M)$ is a smooth bundle; ``in order to really show this in detail you have to check a lot of things, but this is the sketch'' \ts{19}{1}{1:10:42} \ts{19}{1}{1:11:02} \ts{19}{1}{1:11:23} \ts{19}{1}{1:11:46} \ts{19}{1}{1:12:06}.

Now we want to make this into a principal $G$-bundle with $G=GL(d,\R)$: first it is a bundle (as just seen), second we establish a right $GL(d,\R)$-action on the total space \ts{19}{1}{1:12:06} \ts{19}{1}{1:12:26} \ts{19}{1}{1:12:47}. For $g\in GL(d,\R)$ with matrix entries $g^m{}_n\in\R$ ($m,n=1,\dots,d$, $\det g\ne0$ ``that is our definition of $GL$'') define
\[
  (e_1,\dots,e_d)\ract g:=\bigl(g^m{}_1e_m,\ g^m{}_2e_m,\ \dots,\ g^m{}_de_m\bigr)
\]
(summation over $m$). This is just what one does to change a basis by an invertible transformation: the new first frame vector is a linear combination of all the old frame vectors, the first collection of coefficients; the second new vector is the second collection of coefficients, and so on: ``this is precisely how you transform from one basis to another'' \ts{19}{1}{1:13:09} \ts{19}{1}{1:13:31} \ts{19}{1}{1:13:52} \ts{19}{1}{1:14:13} \ts{19}{1}{1:14:33} \ts{19}{1}{1:14:54}. (It is an action of the group on the right: not said in class, with $e'_n=g_1{}^m{}_ne_m$ one gets $(e\ract g_1)\ract g_2=(g_2{}^n{}_kg_1{}^m{}_n e_m)_k=e\ract(g_1g_2)$, the matrix product $(g_1g_2)^m{}_k=g_1{}^m{}_ng_2{}^n{}_k$.)

\textbf{Is it a free action?} A free action is one in which only the identity (the Kronecker $\delta^m{}_n$) keeps the point unchanged \ts{19}{1}{1:15:35} \ts{19}{1}{1:15:55}. For a single vector one finds non-trivial $GL$ transformations that keep it invariant, and for two vectors still; in three dimensions it gets trickier; ``but if one has an entire frame, there is nothing one can do to keep every vector of the frame invariant'': so it is free \ts{19}{1}{1:15:55} \ts{19}{1}{1:16:16} \ts{19}{1}{1:16:37}. (That is the intuitive argument; the formal proof, not said in class: if $e\ract g=e$ then $g^m{}_ne_m=e_n$ for all $n$, and since the $e_m$ are linearly independent, $g^m{}_n=\delta^m{}_n$.)

\begin{proposition}[The frame bundle is a principal $GL(d,\R)$-bundle]\label{prop:pb-frame-principal}
The bundle $(LM,\pi,M)$ with the right action above is a principal $GL(d,\R)$-bundle. \ts{19}{1}{1:16:37} \ts{19}{1}{1:16:57} \ts{19}{1}{1:17:19}
\end{proposition}
What has to be checked is that there is a bundle isomorphism $(u,f)$ between $(LM,\pi,M)$ and $(LM,\rho,LM/GL(d,\R))$, an isomorphism that in particular matches the base spaces \ts{19}{1}{1:17:40} \ts{19}{1}{1:18:02} \ts{19}{1}{1:18:22}. The intuition: the orbit space consists of orbits; take a frame at a point of the manifold and act with the whole group $GL(d,\R)$ on it: one obtains \emph{any other frame at the same point}, and only at that point (a change of basis keeps the vectors in the same $T_xM$). So the orbit of a frame is the whole $L_xM$, and if one forgets the orbit (collapses it), all that remains is the point $x$. This is precisely the identification $LM/GL(d,\R)\cong M$ that we want: ``from any one frame one can go to any other frame, and that is a free action'' \ts{19}{1}{1:18:43} \ts{19}{1}{1:19:04} \ts{19}{1}{1:19:24} \ts{19}{1}{1:19:45} \ts{19}{1}{1:20:05}. The check is to be made first intuitively (is it reasonable that it could be one?) and then written down. Another check for freeness: one must be able to reach any point of a fibre from any other point of the same fibre by some action, and by a \emph{unique} group element \ts{19}{1}{1:20:05} \ts{19}{1}{1:20:26} \ts{19}{1}{1:20:47}. ``So indeed this is a principal $GL(d,\R)$-bundle, and that is very important'' \ts{19}{1}{1:20:47} \ts{19}{1}{1:21:08}.

Why important: this is how physicists say a basis is transformed; if a basis is transformed, then also the components of a vector, a tensor or a tensor density are transformed. These objects are not elements of this bundle; they will be elements of an \emph{associated} bundle, on which there will be a \emph{left} $GL(d,\R)$-action (\cref{ch:l20}). The example cannot be taken further now, but later it will nicely become the formal setting for what has been told in the first semester of physics about tensors in classical mechanics, the inertia tensor and so on \ts{19}{1}{1:21:08} \ts{19}{1}{1:21:29} \ts{19}{1}{1:21:51}.

% ==================================================================
\section{Principal bundle maps}\label{sec:pb-pbmap}

At the beginning of the course the lecturer kept emphasising that one always studies the structure preserving maps between objects, and this remains true \ts{19}{1}{1:21:51}. Recall (\cref{def:mf-morphism}) that for two bundles $(E,\pi,M)$ and $(E',\pi',M')$ a \emph{bundle map} is a pair $(u,f)$, $u\colon E\to E'$, $f\colon M\to M'$, with $\pi'\circ u=f\circ\pi$, i.e.\ the diagram commutes \ts{19}{1}{1:22:11} \ts{19}{1}{1:22:34} \ts{19}{1}{1:22:54}. A principal bundle is certainly a bundle, but it has additional structure, the right action (``sometimes suppressed in the notation, but it should not be''), which has to be taken into account \ts{19}{1}{1:23:16} \ts{19}{1}{1:24:19} \ts{19}{1}{1:24:40}.

\begin{definition}[New concepts: principal bundle map]\label{def:pb-pbmap}
Let $(P,\pi,M)$ and $(P',\pi',M')$ be principal $G$-bundles with right actions $\ract$ and $\ract'$ (the \emph{same} group, but the action on the two total spaces may be different). A pair $(u,f)$, $u\colon P\to P'$, $f\colon M\to M'$, is a \textbf{principal bundle map} if the diagram
\begin{center}
\begin{tikzpicture}[>=Stealth,node distance=2.6cm,baseline=(current bounding box.center)]
\node (PG) {$P\times G$};\node (PG2) [right=3.4cm of PG] {$P'\times G$};
\node (P) [below=2cm of PG] {$P$};\node (P2) [below=2cm of PG2] {$P'$};
\node (M) [below=2cm of P] {$M$};\node (M2) [below=2cm of P2] {$M'$};
\draw[->] (PG)--node[above,font=\small]{$u\times\mathrm{id}_G$}(PG2);
\draw[->] (P)--node[above,font=\small]{$u$}(P2);
\draw[->] (M)--node[below,font=\small]{$f$}(M2);
\draw[->] (PG)--node[left,font=\small]{$\ract$}(P);
\draw[->] (PG2)--node[right,font=\small]{$\ract'$}(P2);
\draw[->] (P)--node[left,font=\small]{$\pi$}(M);
\draw[->] (P2)--node[right,font=\small]{$\pi'$}(M2);
\end{tikzpicture}
\end{center}
commutes: the lower square is the bundle map condition, $\pi'\circ u=f\circ\pi$; the upper square is the equivariance condition,
\[
  u(p\ract g)=u(p)\ract'g\qquad\text{for all }p\in P,\ g\in G .
\]
\ts{19}{1}{1:22:34} \ts{19}{1}{1:23:37} \ts{19}{1}{1:23:59} \ts{19}{1}{1:25:01} \ts{19}{1}{1:25:22} \ts{19}{1}{1:25:42} \ts{19}{1}{1:26:03} \ts{19}{1}{1:26:25} \ts{19}{1}{1:26:47}
\end{definition}
It is only the natural thing to require that the entire diagram commutes: it seems a subtlety, but a principal bundle map is different from a bundle map, it is a stronger notion \ts{19}{1}{1:25:22} \ts{19}{1}{1:27:10}. Questions? ``Straightforward enough'' \ts{19}{1}{1:27:10}.

\textbf{The maximal generalisation} \ts{19}{1}{1:27:34}. Take a principal $G$-bundle $(P,\pi,M)$ with right action $\ract$ and another principal $G'$-bundle $(P',\pi',M')$ with right action $\ract'$ of a possibly \emph{different} group $G'$. To talk about a map one needs an extra ingredient: a way to relate the elements of $G$ to elements of $G'$, a Lie group homomorphism $\rho\colon G\to G'$ (smooth and $\rho(g_1\bullet g_2)=\rho(g_1)\bullet'\rho(g_2)$) \ts{19}{1}{1:27:55} \ts{19}{1}{1:28:15} \ts{19}{1}{1:28:36} \ts{19}{1}{1:28:56} \ts{19}{1}{1:29:17}. Given $\rho$, $u$ and $f$, the triple is a \textbf{principal bundle map} (a bundle morphism of principal bundles) if
\[
  \pi'\circ u=f\circ\pi\qquad\text{and}\qquad u(p\ract g)=u(p)\ract'\rho(g)\ \ \text{for all }p\in P,\,g\in G ,
\]
that is, if the diagram with $u\times\rho$ in the upper row commutes. One first applies $\rho$ to turn an element of $G$ into an element of $G'$, so that it can act from the right on $P'$: ``it is totally obvious once you draw the diagrams, everything becomes child's play'' \ts{19}{1}{1:29:37} \ts{19}{1}{1:29:59} \ts{19}{1}{1:30:19} \ts{19}{1}{1:30:40} \ts{19}{1}{1:31:00} \ts{19}{1}{1:31:24}. (A bundle morphism is a bundle map: ``morph is form, a morphism is a structure preserving map, that is the game we have been playing all along'' \ts{19}{1}{1:31:24}.)

\begin{application}[General relativity: spin frame bundle]
Is this generalisation only for its own sake? There is an important application \ts{19}{1}{1:31:44} \ts{19}{1}{1:32:04}. In general relativity one often looks at the frame bundle of the space-time manifold with group $GL(d,\R)$; if there is a metric one can restrict to the orthogonal group, and for Lorentzian signature to the Lorentz group $SO(1,3)$: this is equivalent to the introduction of a metric, because then only frames related by a Lorentz transformation make up the orbits, and one has to restrict the bundle to such frames \ts{19}{1}{1:32:04} \ts{19}{1}{1:32:25} \ts{19}{1}{1:32:45} \ts{19}{1}{1:33:06}. Then one can define a second principal bundle, the \textbf{spin frame bundle}, with group $SL(2,\C)$, and a homomorphism $\rho$ between $SL(2,\C)$ and $SO(1,3)$ which is $2\!:\!1$ \ts{19}{1}{1:33:06} \ts{19}{1}{1:33:26} \ts{19}{1}{1:33:47}. A principal bundle map between them then answers: if I rotate my measurement apparatus (a Lorentz transformation) what is the corresponding spin transformation on the spinors? The details need associated bundles (\cref{ch:l20}) \ts{19}{1}{1:34:08} \ts{19}{1}{1:34:28}.

Why spinors need this: there is no (finite-dimensional, unitary, irreducible) spin representation of $GL(4,\R)$, so one has to restrict to $SO(1,3)$ (or its cover), and for that one first needs the metric; tensor fields do not need a metric, they exist because the manifold is smooth, but finite-dimensional spinor representations require a geometry before one can even say what a spinor is. The principal bundle map then tells physically that, once one transforms a local frame in the frame bundle, one has to use the corresponding transformation in the spin frame bundle \ts{19}{1}{1:34:28} \ts{19}{1}{1:34:49} \ts{19}{1}{1:35:10} \ts{19}{1}{1:35:30} \ts{19}{1}{1:35:51}. Should these remarks remain mysterious: ``have patience''.
\end{application}
\begin{coursenote}[Direction of $\rho$ in the example]
By the definition just given $\rho$ goes from the group of the source bundle to the group of the target bundle, $\rho\colon G\to G'$ with $u(p\ract g)=u(p)\ract'\rho(g)$. The covering homomorphism goes from $SL(2,\C)$ to the Lorentz group, so in the application the \emph{spin frame bundle} (group $SL(2,\C)$) is the source $P$ and the (Lorentz) frame bundle the target $P'$; the lecture's labels are not consistent on this point [?]. \ts{19}{1}{1:33:47}
\end{coursenote}
\begin{answer}{Does $u$ determine $\rho$?}[The role of $\rho$]
A student asks (inaudible) about the connection between $\rho$ and $u$. The lecturer: $\rho$ is an extra piece of data. The connection is only established by the requirement that $u$ and $\rho$ be chosen so that they satisfy the condition. One cannot say that, given $u$, the $\rho$ is determined: there might be several $\rho$ making $u$ into a principal bundle map of this type; for any Lie group homomorphism one can impose the condition. The data are all prescribed and the condition is imposed, as in the example where one already knows how $SL(2,\C)$ is mapped into $SO(1,3)$; in particular circumstances $\rho$ may be determined by $u$, but in full generality not \ts{19}{1}{1:36:20} \ts{19}{1}{1:36:43} \ts{19}{1}{1:37:03} \ts{19}{1}{1:37:24} \ts{19}{1}{1:37:44} \ts{19}{1}{1:38:06}
\end{answer}

\subsection{A principal bundle map over the identity is an isomorphism}\label{sec:pb-lemma}
Principal bundle maps are a stronger notion than bundle maps, and a consequence of this is the following lemma \ts{19}{1}{1:38:06} \ts{19}{1}{1:38:28}.
\begin{lemma}[Principal bundle maps over the identity]\label{lem:pb-iso}
Let $(P,\pi,M)$ and $(P',\pi',M)$ be principal $G$-bundles over the \emph{same} base manifold $M$, and $u\colon P\to P'$ a principal bundle map with $\pi'\circ u=\pi$ (base map the identity) and $u(p\ract g)=u(p)\ract'g$. Then $u$ is bijective (and a diffeomorphism): it is a bundle isomorphism. \ts{19}{1}{1:38:49} \ts{19}{1}{1:39:09} \ts{19}{1}{1:39:30} \ts{19}{1}{1:39:52} \ts{19}{1}{1:40:12}
\end{lemma}
\begin{proof}
If one has a map from $M$ to $M$ that commutes with the projections, it is the identity (``we go to the same point''), which is why no $f$ appears \ts{19}{1}{1:42:38}. We prove injectivity and surjectivity; then $u$ is invertible, and the rest (smoothness) follows \ts{19}{1}{1:40:32} \ts{19}{1}{1:40:54}.\par
\emph{Injective.} Let $u(p_1)=u(p_2)$. Then $\pi(p_1)=\pi'(u(p_1))=\pi'(u(p_2))=\pi(p_2)$, so $p_1,p_2$ lie in the same fibre; hence there is a \emph{unique} $g$ with $p_2=p_1\ract g$ (freeness, and principal bundle: the fibre is an orbit). Then
\[
  u(p_1)=u(p_2)=u(p_1\ract g)=u(p_1)\ract'g ,
\]
by equivariance. So $g$ fixes the point $u(p_1)$, and since the action is free, $g=\mathbb 1$. Hence $p_2=p_1$. (Both the bundle map property and the extra equivariance property were used.) \ts{19}{1}{1:41:15} \ts{19}{1}{1:41:35} \ts{19}{1}{1:43:18} \ts{19}{1}{1:43:42} \ts{19}{1}{1:44:03} \ts{19}{1}{1:44:24} \ts{19}{1}{1:44:45} \ts{19}{1}{1:45:06} \ts{19}{1}{1:45:26} \ts{19}{1}{1:45:47} \ts{19}{1}{1:46:07} \ts{19}{1}{1:46:28} \ts{19}{1}{1:46:49} \ts{19}{1}{1:47:09} \ts{19}{1}{1:47:30}\par
\emph{Surjective.} Let $p'\in P'$. Choose some $p\in P$ in the fibre $\pi^{-1}(\pi'(p'))$ over the same base point; it exists because $P$ is a bundle over $M$ with non-empty fibres. Then $\pi'(u(p))=\pi(p)=\pi'(p')$ (bundle map aspect only), so $u(p)$ and $p'$ lie in the same fibre of $P'$ and there is a unique $g\in G$ with $u(p)\ract'g=p'$. By equivariance $u(p\ract g)=u(p)\ract'g=p'$, so $\tilde p:=p\ract g$ is the point we were looking for. \ts{19}{1}{1:47:52} \ts{19}{1}{1:48:13} \ts{19}{1}{1:48:33} \ts{19}{1}{1:48:57} \ts{19}{1}{1:49:17} \ts{19}{1}{1:49:38} \ts{19}{1}{1:49:59} \ts{19}{1}{1:50:19} \ts{19}{1}{1:50:40} \ts{19}{1}{1:51:00} \ts{19}{1}{1:51:21} \ts{19}{1}{1:51:44} \ts{19}{1}{1:52:06} \ts{19}{1}{1:52:27} \ts{19}{1}{1:52:47} \ts{19}{1}{1:53:09} \ts{19}{1}{1:53:29}
\end{proof}
Without the extra principal-bundle condition beyond the pure bundle map, neither injectivity nor surjectivity could be proven. ``It is really a strengthening of what a bundle map is on a principal bundle: a principal bundle map is more than a bundle map, and it makes this automatically a diffeomorphism'' \ts{19}{1}{1:53:50} \ts{19}{1}{1:54:12}.
\begin{coursenote}[Smoothness of the inverse]
Not said in class: the lecture proves bijectivity only and says that the smoothness ``also will be the case''. One way: over a small open set $U\subseteq M$ choose local sections $\sigma$ of $P$ and $\sigma'$ of $P'$ (they exist by local triviality) and write $u(\sigma(x))=\sigma'(x)\ract'\gamma(x)$ with a smooth $\gamma\colon U\to G$; then $u$ acts as $\sigma(x)\ract g\mapsto\sigma'(x)\ract'(\gamma(x)\bullet g)$, whose inverse is smooth since $G$ is a Lie group. \ts{19}{1}{1:54:12}
\end{coursenote}
In the proof we used several times the fact that two points in the same fibre of a principal bundle differ by a unique group element (\cref{rem:pb-fibres}). The lemma is the basis of the next definition and theorem \ts{19}{1}{1:54:32} \ts{19}{1}{1:54:55}.

% ==================================================================
\section{Trivial principal bundles and sections}\label{sec:pb-trivial}

\begin{definition}[New concepts: trivial principal bundle]\label{def:pb-trivial}
The \textbf{trivial principal $G$-bundle} over $M$ is the bundle $(M\times G,\pi_1,M)$, with $\pi_1(x,h):=x$ the projection onto the first factor, and the right action
\[
  (x,h)\ractf g:=(x,h\bullet g)
\]
(the base point untouched; one acts from the right by $g$ on the group factor). It is a principal $G$-bundle: the action is free because in the group there is always the inverse ($h\bullet g=h$ forces $g=\mathbb 1$). A principal $G$-bundle $(P,\pi,M)$ is called \textbf{trivial} if there exists a principal bundle map $u\colon P\to M\times G$ over $\mathrm{id}_M$ (i.e.\ $\pi_1\circ u=\pi$ and $u(p\ract g)=u(p)\ractf g$). \ts{19}{1}{1:54:55} \ts{19}{1}{1:55:24} \ts{19}{1}{1:55:44} \ts{19}{1}{1:56:05} \ts{19}{1}{1:56:25} \ts{19}{1}{1:56:45} \ts{19}{1}{1:57:06} \ts{19}{1}{1:57:26} \ts{19}{1}{1:57:47}
\end{definition}
Since a principal bundle map is automatically a diffeomorphism (\cref{lem:pb-iso}), ``trivial'' can be phrased either as: diffeomorphic as a principal $G$-bundle to $M\times G$, or as: there exists a principal bundle map to $M\times G$ \ts{19}{1}{1:55:24} \ts{19}{1}{1:55:44} \ts{19}{1}{1:59:52} \ts{19}{1}{2:00:13}. One should insist that it is \emph{not} enough to have a bundle map to the trivial bundle (one would be tempted to define it like that): one only calls it trivial if there exists a \emph{principal} bundle map, which is the stronger notion \ts{19}{1}{2:00:13} \ts{19}{1}{2:00:33}.

At the beginning of the lecture we said that a principal $G$-bundle very roughly is a bundle with fibre $G$. This trivial bundle is the very trivial way to make $G$ a fibre: one has a manifold and attaches the group at every point. A principal $G$-bundle is also, very precisely, a bundle with fibre $G$ (because of the isomorphism with $P\to P/G$), but it is certainly not always trivial like this \ts{19}{1}{1:57:47} \ts{19}{1}{1:58:07} \ts{19}{1}{1:58:28} \ts{19}{1}{1:58:48} \ts{19}{1}{1:59:12} \ts{19}{1}{1:59:32}.

Why are we interested in trivial principal bundles? One can think of them globally as having the group just attached as a fibre: very simple. The question is how common that is, and what it has to do with what we did before \ts{19}{1}{2:00:55} \ts{19}{1}{2:01:15}.
\begin{answer}{What exactly is called trivial here?}[Trivial bundle and trivial principal bundle]
A question (partly inaudible) about the colours of the diagram. The lecturer: the right-hand side $M\times G$ with the projection $\pi_1$ is certainly what one calls a trivial \emph{bundle}, because it is a Cartesian product. Any principal bundle that is diffeomorphic to it \emph{as a principal bundle}, i.e.\ by a diffeomorphism of the total spaces that preserves all the group actions, is called trivial: all the structure may look more complicated, but after an appropriate diffeomorphism it takes this simple form. It was the same with bundles before: a bundle is trivial if the total space is the Cartesian product of the base with a fibre. In the extended picture, the principal bundle $P\to M$ is, by definition, bundle-isomorphic to $P\to P/G$ with its free $G$-action, so a $G$-fibre bundle with the same fibre at every point; hence it makes sense to ask the question. Not every principal bundle is trivial; otherwise the notion would not be worth a definition \ts{19}{1}{2:01:36} \ts{19}{1}{2:01:57} \ts{19}{1}{2:02:17} \ts{19}{1}{2:02:37} \ts{19}{1}{2:02:58} \ts{19}{1}{2:03:19} \ts{19}{1}{2:03:39} \ts{19}{1}{2:04:02} \ts{19}{1}{2:04:27} \ts{19}{1}{2:04:47} \ts{19}{1}{2:05:09} \ts{19}{1}{2:05:29} \ts{19}{1}{2:05:50}
\end{answer}

\begin{theorem}[Triviality and sections]\label{thm:pb-trivial-section}
A principal $G$-bundle $(P,\pi,M)$ is trivial if and only if it admits a smooth section, i.e.\ a smooth map $\sigma\colon M\to P$ with $\pi\circ\sigma=\mathrm{id}_M$ (\cref{def:mf-section}). \ts{19}{1}{2:06:12} \ts{19}{1}{2:06:32} \ts{19}{1}{2:06:52} \ts{19}{1}{2:07:15} \ts{19}{1}{2:07:36}
\end{theorem}
The point of the theorem: finding a diffeomorphism or a principal bundle map can be tricky, and it is hard to show that it does not exist; so it is good to have a criterion, here a necessary and sufficient one \ts{19}{1}{2:06:12} \ts{19}{1}{2:06:32}.
\begin{example}[The frame bundle of the 2-sphere is not trivial]\label{ex:pb-S2}
Let $M=S^2$ and $P=LS^2$, the frame bundle, a principal $GL(2,\R)$-bundle (\cref{prop:pb-frame-principal}). A section of the frame bundle would be a frame at every point, that is a basis of the tangent space smoothly depending on the point: two linearly independent, everywhere non-vanishing smooth vector fields. But there is not even one everywhere non-vanishing smooth vector field on $S^2$: ``you cannot comb a sphere'' (the hairy ball theorem of \cref{ch:l11}). Hence there is no smooth section, and by \cref{thm:pb-trivial-section} the principal $GL(2,\R)$-bundle $LS^2$ is not trivial. So there are principal bundles that are non-trivial, and others, like $M\times G$, that are trivial. \ts{19}{1}{2:07:58} \ts{19}{1}{2:08:19} \ts{19}{1}{2:08:40} \ts{19}{1}{2:09:01} \ts{19}{1}{2:09:24} \ts{19}{1}{2:09:48} \ts{19}{1}{2:10:08} \ts{19}{1}{2:10:29} \ts{19}{1}{2:10:49}
\end{example}
(At the end of the lecture the lecturer relates this to the previous lectures: there is no basis of the module of vector fields, the $C^\infty(M)$-module of vector fields is not free, only projective, the theorem of Serre and Swan (\cref{ch:l11}): ``all these things belong together'' \ts{19}{1}{2:31:13} \ts{19}{1}{2:31:34} \ts{19}{1}{2:31:55}.) Now we want to understand the theorem better by proving it \ts{19}{1}{2:10:49} \ts{19}{1}{2:11:10}.
\begin{proof}
\emph{($\Rightarrow$)} Suppose the bundle is trivial: there is a principal bundle map $u\colon P\to M\times G$ over the identity, a diffeomorphism by \cref{lem:pb-iso}. Define
\[
  \sigma(x):=u^{-1}(x,\mathbb 1) ,
\]
that is, put the identity of the group in the second slot and push it back. It is smooth, and a section because $\pi\bigl(\sigma(x)\bigr)=\pi_1\bigl(u(\sigma(x))\bigr)=\pi_1(x,\mathbb 1)=x$, since $u$ is a bundle map. \ts{19}{1}{2:11:32} \ts{19}{1}{2:11:53} \ts{19}{1}{2:12:13} \ts{19}{1}{2:12:34} \ts{19}{1}{2:12:54}\par
\emph{($\Leftarrow$)} Suppose a smooth section $\sigma$ is given \ts{19}{1}{2:13:16} \ts{19}{1}{2:13:36}. Let $p\in P$ be any point; it lies in some fibre. The point $\sigma(\pi(p))$ lies in the \emph{same} fibre, so there is a unique group element $\chi_\sigma(p)\in G$ with
\[
  \sigma(\pi(p))\ract\chi_\sigma(p)=p .
\]
It exists, and is unique, because the action is free and the fibre is an orbit. It depends on the section $\sigma$ (another section would need another element to bridge the gap) and on the point $p$, so it is a map $\chi_\sigma\colon P\to G$ \ts{19}{1}{2:13:57} \ts{19}{1}{2:14:20} \ts{19}{1}{2:14:40} \ts{19}{1}{2:15:01} \ts{19}{1}{2:15:22} \ts{19}{1}{2:15:42} \ts{19}{1}{2:16:03} \ts{19}{1}{2:16:24} \ts{19}{1}{2:16:45} \ts{19}{1}{2:17:06}.\par
\emph{Claim:} $\chi_\sigma(p\ract g)=\chi_\sigma(p)\bullet g$ for all $g\in G$. Indeed, $\pi(p\ract g)=\pi(p)$ (the action stays in the fibre), hence $\sigma(\pi(p\ract g))=\sigma(\pi(p))$ and, by definition of $\chi_\sigma$, $\sigma(\pi(p))\ract\chi_\sigma(p\ract g)=p\ract g$. On the other hand, acting on the defining equation with $g$,
\[
  \sigma(\pi(p))\ract\bigl(\chi_\sigma(p)\bullet g\bigr)=\bigl(\sigma(\pi(p))\ract\chi_\sigma(p)\bigr)\ract g=p\ract g ,
\]
because it is a right action. So the same point is acted upon by the two elements $\chi_\sigma(p\ract g)$ and $\chi_\sigma(p)\bullet g$ and the results coincide; since the action is free, the two group elements are equal. \ts{19}{1}{2:17:27} \ts{19}{1}{2:17:48} \ts{19}{1}{2:18:10} \ts{19}{1}{2:18:32} \ts{19}{1}{2:18:53} \ts{19}{1}{2:20:00} \ts{19}{1}{2:20:22} \ts{19}{1}{2:21:03} \ts{19}{1}{2:21:24} \ts{19}{1}{2:22:09} \ts{19}{1}{2:22:30} \ts{19}{1}{2:22:52} \ts{19}{1}{2:23:13} \ts{19}{1}{2:23:33} \ts{19}{1}{2:23:53} \ts{19}{1}{2:24:14}\par
Now define
\[
  u_\sigma\colon P\to M\times G,\qquad u_\sigma(p):=\bigl(\pi(p),\chi_\sigma(p)\bigr).
\]
It is a bundle map: $\pi_1\circ u_\sigma=\pi$, obvious since $\pi_1$ projects out the first component. It is equivariant:
\[
  u_\sigma(p\ract g)=\bigl(\pi(p),\chi_\sigma(p)\bullet g\bigr)=u_\sigma(p)\ractf g ,
\]
where we used the claim (and $\pi(p\ract g)=\pi(p)$): ``this is why we had to show that if one first acts on $p$ and then uses $\chi_\sigma$, one can also first use $\chi_\sigma$ and then blob with $g$''. So $u_\sigma$ is a principal bundle map into the trivial bundle: the bundle is trivial \ts{19}{1}{2:24:36} \ts{19}{1}{2:24:56} \ts{19}{1}{2:25:17} \ts{19}{1}{2:25:37} \ts{19}{1}{2:25:58} \ts{19}{1}{2:26:20} \ts{19}{1}{2:26:41} \ts{19}{1}{2:27:01} \ts{19}{1}{2:27:22} \ts{19}{1}{2:27:43} \ts{19}{1}{2:28:04} \ts{19}{1}{2:28:25} \ts{19}{1}{2:28:46} \ts{19}{1}{2:29:07} \ts{19}{1}{2:29:28} \ts{19}{1}{2:29:48} \ts{19}{1}{2:30:11} \ts{19}{1}{2:30:32} \ts{19}{1}{2:30:53}.
\end{proof}
So: having a section means that the bundle is trivial, and conversely, if the bundle is trivial, one can construct a section. The lecture does not discuss the smoothness of $\chi_\sigma$.
\begin{coursenote}[Smoothness of $\chi_\sigma$]
Not said in class: $\chi_\sigma$ is smooth; locally one has a local section $\sigma'$ of the principal bundle with $p=\sigma'(\pi(p))\ract\chi'(p)$ for a smooth $\chi'\colon\pi^{-1}(U)\to G$, and $\sigma(x)=\sigma'(x)\ract\gamma(x)$ with $\gamma$ smooth, so $\chi_\sigma(p)=\gamma(\pi(p))^{-1}\bullet\chi'(p)$. The trial in the lecture of writing the equation for $\chi_\sigma(p\ract g)$ had a few false starts at the board (``hang on, can that be the thing I want to do?''), cleaned up here. \ts{19}{1}{2:19:13} \ts{19}{1}{2:19:34}
\end{coursenote}
\begin{aside}[Looking ahead]
The lecture ends here: ``see you after Christmas, and we will continue with associated bundles, and then you will finally understand why we are so interested in principal bundles, because they have associated bundles, and in the associated bundles live the objects we like so much: spinors, tensors, tensor densities''. If in physics we say we rotate something, we do an $SO(3)$ transformation, how do the different objects, which are sections of different associated bundles, transform under the same transformation? Understood in terms of the principal bundle: ``in this sense the principal bundle really thrones over all of the others. If you understand the principal bundle, you understand all the associated bundles'' \ts{19}{1}{2:32:17} \ts{19}{1}{2:32:38} \ts{19}{1}{2:32:58} \ts{19}{1}{2:33:19}
\end{aside}

% ==================================================================
\section*{Key concepts}
\begin{keyconcepts}
\textbf{Retrieval questions.}
\begin{enumerate}
  \item Define a left and a right $G$-action. How does one get a right action from a left one, and why is the inverse needed?
  \item What is a $\rho$-equivariant map? Write the commuting diagram and the equation.
  \item Define orbit, orbit space $M/G$ and stabiliser. Which is a subset of the space, which of the group? When is an action free?
  \item Why is the rotation action on $\R^2$ not free, and why is it free on $\R^2\setminus\{0\}$? Why is a linear action never free?
  \item State the three conditions of a principal $G$-bundle. Why are its fibres diffeomorphic to $G$?
  \item Define the frame bundle $LM$, its dimension, the projection and the right $GL(d,\R)$-action. Why is the action free?
  \item What is a principal bundle map (same group; with $\rho$)? What does the lemma say and how does the proof use freeness and equivariance?
  \item Define the trivial principal bundle and state the section theorem. Sketch both directions, including the role of $\chi_\sigma$ and the identity $\chi_\sigma(p\ract g)=\chi_\sigma(p)\bullet g$.
  \item Why is $LS^2$ not trivial?
\end{enumerate}
\textbf{Concepts.}
\begin{itemize}
  \item Left and right actions, equivariance (\cref{def:pb-left,def:pb-right,prop:pb-right-from-left,def:pb-equivariant}).
  \item Orbits, orbit space, stabilisers, free actions (\cref{def:pb-orbit,def:pb-orbitspace,def:pb-stabiliser,def:pb-free,prop:pb-orbit-G}).
  \item Principal bundles and the frame bundle (\cref{def:pb-principal,def:pb-frame,prop:pb-frame-principal}).
  \item Principal bundle maps and the isomorphism lemma (\cref{def:pb-pbmap,lem:pb-iso}).
  \item Trivial principal bundles and sections (\cref{def:pb-trivial,thm:pb-trivial-section,ex:pb-S2}).
\end{itemize}
\end{keyconcepts}

\flushside
\end{paracol}

% !TEX root = ../main.tex
% STUB: generated by tools/stubs.py; write the chapter from materials/transcripts/C20-V01.txt
\chapter{Associated fibre bundles}\label{ch:l20}
\begin{paracol}{2}

\begin{sources}
\begin{tabularx}{\linewidth}{@{}l l L@{}}
\toprule
\textbf{Video} & \textbf{Length} & \textbf{Content}\\
\midrule
\seg{20}{1}{Lecture 20: Associated fibre bundles} & 1:42:31 & (to be written)\\
\bottomrule
\end{tabularx}

\smallskip
\textbf{Transcript} \file{materials/transcripts/C20-V01.txt}.
\end{sources}

\begin{reading}[Where this lecture sits in the other notes]
Not treated in the other courses (only here). (From the lecture titles; to be checked when the chapter is written.)
\end{reading}

\emph{This chapter has not been written yet.}

\flushside
\end{paracol}

% !TEX root = ../main.tex
% STUB: generated by tools/stubs.py; write the chapter from materials/transcripts/C21-V01.txt
\chapter{Connections and connection 1-forms}\label{ch:l21}
\begin{paracol}{2}

\begin{sources}
\begin{tabularx}{\linewidth}{@{}l l L@{}}
\toprule
\textbf{Video} & \textbf{Length} & \textbf{Content}\\
\midrule
\seg{21}{1}{Lecture 21: Connections and connection 1-forms} & 1:05:01 & (to be written)\\
\bottomrule
\end{tabularx}

\smallskip
\textbf{Transcript} \file{materials/transcripts/C21-V01.txt}.
\end{sources}

\begin{reading}[Where this lecture sits in the other notes]
(to be written)
\end{reading}

\emph{This chapter has not been written yet.}

\flushside
\end{paracol}

% !TEX root = ../main.tex
\chapter{Local representations of a connection on the base manifold: Yang-Mills fields}\label{ch:l22}
\providecommand{\lact}{\mathbin{\triangleright}}
\providecommand{\ract}{\mathbin{\triangleleft}}
\begin{paracol}{2}

\begin{sources}
\begin{tabularx}{\linewidth}{@{}l l L@{}}
\toprule
\textbf{Video} & \textbf{Length} & \textbf{Content}\\
\midrule
\seg{22}{1}{Lecture 22: Local representations of a connection on the base manifold: Yang-Mills fields} & 1:29:33 & Recall of the connection 1-form; what ``Lie-algebra valued 1-form'' means and that pull-backs do not care about the target; a local section gives the Yang--Mills field $\omega_U$ and a local trivialisation; the Maurer--Cartan form; the frame bundle and the connection coefficients $\Gamma^i{}_{j\mu}$; the Maurer--Cartan form of $GL(d,\R)$; the gauge map between two sections and the transformation law of Yang--Mills fields; the $GL$ case and the transformation law of the $\Gamma$ symbols.\\
\bottomrule
\end{tabularx}

\smallskip
\textbf{Transcript} \file{materials/transcripts/C22-V01.txt} (YouTube automatic captions; the blackboard symbols, in particular the several Greek letters, are reconstructed from the verbal description and from \cref{ch:l21}; the captions render Maurer--Cartan as ``market on'' and ``Rialto prevail'' for ``Lie-algebra valued'').
\end{sources}

\begin{reading}[Where this lecture sits in the other notes]
\emph{Gravity and Light}, lecture 7 (connection coefficients $\Gamma^i{}_{jk}$ and their transformation law on the tangent bundle); gauge fields and the transformation law of the Yang--Mills field are not treated in the other courses (only here). It continues \cref{ch:l21} (connection, connection 1-form $\omega$ and its properties), uses the frame bundle (\cref{ch:l19}), left-invariant fields and the exponential map (\cref{ch:l13}, \cref{ch:l18}), the adjoint map (\cref{ch:l17}) and pull-backs and push-forwards (\cref{ch:l10}). Next: parallel transport (\cref{ch:l23}). (From the lecture titles; not checked in detail.)
\end{reading}

% ==================================================================
\section{Recall, and what a Lie-algebra valued 1-form is}\label{sec:ym-recall}

Last time we studied connections on a principal bundle: a connection is a distribution of horizontal subspaces, a choice (there is no canonical one), and it can be expressed conveniently as a Lie-algebra valued 1-form on the principal bundle. Today we study how this 1-form can be expressed \emph{locally} as a Lie-algebra valued 1-form on the base manifold: this is how it is usually met in physics, and depending on context it is called the local connection coefficients $\Gamma$ (general relativity) or the Yang--Mills field (elementary particle physics). The lecture is section 5.5, ``local representation of a connection''. The terminology \emph{Yang--Mills field} is used invariantly, and put in quotation marks, since it may apply to situations that a physicist would not call a Yang--Mills field in the first place \ts{22}{1}{0:07} \ts{22}{1}{0:28} \ts{22}{1}{0:49} \ts{22}{1}{1:11} \ts{22}{1}{1:32}.
\begin{recall}[Connection 1-form]
Principal $G$-bundle $(P,\pi,M)$ with right action $\ract$ and connection 1-form $\omega\colon T_pP\to T_eG$ (\cref{def:cn-omega}), satisfying (\cref{thm:cn-omega}): (a) $\omega(X^A)=A$ for the vector field $X^A$ generated by $A\in T_eG$; (b) $r_g^*\omega=(\mathrm{Ad}_{g^{-1}})_*\omega$; (c) $\omega$ smooth.
\end{recall}
\ts{22}{1}{1:53} \ts{22}{1}{2:14} \ts{22}{1}{2:34} \ts{22}{1}{2:55} \ts{22}{1}{3:16} \ts{22}{1}{3:37} \ts{22}{1}{3:58}. So much for the recall (the lecturer first writes the brackets in (b) the wrong way and corrects them: first one applies $\omega$ to $X$, the result is a Lie algebra element, and on it the push-forward of the adjoint map at $g^{-1}$ can act).

\textbf{What precisely is a Lie-algebra valued 1-form?} Somebody asked this last time \ts{22}{1}{3:58}. The simple answer: exactly what the definition says, a map from the sections of the tangent bundle, linear, not into $C^\infty(P)$, not into the real numbers or some number field, but into the Lie algebra; so one may call it a Lie-algebra valued 1-form. ``Nice, maybe suggestive name'': in what sense can one still treat it as a 1-form, and in what sense not? \ts{22}{1}{4:18} \ts{22}{1}{4:39}
\begin{answer}{In which sense is $\omega$ still a 1-form? Do pull-backs still work?}[Lie-algebra valued 1-forms]
One particular worry: all things defined for ordinary 1-forms and heavily used, like pull-backs, do they generalise? Take a map $u\colon P\to P$ such that $(u,\mathrm{id}_M)$ [?: the lecture draws a diagram with a map $f$ on the base] is a principal bundle morphism; over the same $P$ it is already an isomorphism, an \emph{automorphism} of the principal bundle (the usual two compatibility conditions, that the diagram commutes and that $u$ is equivariant, \cref{def:pb-pbmap}). If $\omega$ lives on $P$, we should be able to pull it back by $u$. But we defined the pull-back for forms whose value is in $C^\infty(P)$, or in the number field. The pull-back is, at a point,
\[
  (u^*\omega)_p(X):=\omega_{u(p)}\bigl(u_{*,p}X\bigr),
\]
defined by the push-forward of the \emph{vector} $X$, which is an ordinary vector: the pull-back does not care about the target of $\omega$, it is as Lie-algebra valued as $\omega$ is. So the one-form aspect is not harmed by changing the destination; if in doubt, go back to the definitions and check. \ts{22}{1}{4:59} \ts{22}{1}{5:19} \ts{22}{1}{5:41} \ts{22}{1}{6:02} \ts{22}{1}{6:22} \ts{22}{1}{6:43} \ts{22}{1}{7:04} \ts{22}{1}{7:24} \ts{22}{1}{7:45} \ts{22}{1}{8:06} \ts{22}{1}{8:27}
\end{answer}

% ==================================================================
\section{A local section: the Yang--Mills field and the local trivialisation}\label{sec:ym-local}

We have the bundle and the connection 1-form $\omega$, globally defined, which is no problem. In practice however we often want to restrict to part of the bundle, over a subset $U\subseteq M$ of the base manifold, for instance for calculations: $\pi\colon\pi^{-1}(U)\to U$ (the restriction is clear). Physics also takes place locally: if the base is space-time or the physical space and $P$ the principal bundle of the transformations one considers, then one does local experiments and writes down laws in a local patch; the globalisation is another issue \ts{22}{1}{8:27} \ts{22}{1}{8:48} \ts{22}{1}{9:16} \ts{22}{1}{9:40} \ts{22}{1}{10:00} \ts{22}{1}{10:20}.

So choose a \emph{local section} $\sigma\colon U\to P$, $\pi\circ\sigma=\mathrm{id}_U$ \ts{22}{1}{10:40} \ts{22}{1}{11:02}. It does two things.

\textbf{First, the Yang--Mills field} \ts{22}{1}{11:23}.
\begin{definition}[New concepts: the Yang--Mills field $\omega_U$ of a section]\label{def:ym-field}
Let $\sigma\colon U\to P$ be a local section and $\omega$ a connection 1-form on $P$. The \textbf{``Yang--Mills field''} (with respect to $\sigma$) is the Lie-algebra valued 1-form on $U\subseteq M$
\[
  \omega_U:=\sigma^*\omega ,\qquad(\omega_U)_x(v)=\omega_{\sigma(x)}\bigl(\sigma_{*,x}v\bigr)\in T_eG ,\quad v\in T_xM .
\]
\ts{22}{1}{11:46} \ts{22}{1}{12:09} \ts{22}{1}{12:30}
\end{definition}
It lives on an entirely different level than $\omega$: $\omega$ is a 1-form on the principal bundle, $\omega_U$ is a (still Lie-algebra valued) 1-form only on a local patch of the base manifold \ts{22}{1}{12:09} \ts{22}{1}{12:30} \ts{22}{1}{12:52}.

\textbf{Second, a local trivialisation} \ts{22}{1}{13:13}. Recall that a principal bundle of the form $M\times G$ is trivial, and that a principal bundle is trivial if and only if it has a global section (\cref{thm:pb-trivial-section}); in general there is none, which is why one has to restrict to a local section in the first place. But a local section gives a local trivialisation:
\begin{definition}[New concepts: the local trivialisation $h_\sigma$]\label{def:ym-triv}
The section $\sigma$ defines the map
\[
  h\colon U\times G\to P,\qquad h(m,g):=\sigma(m)\ract g ,
\]
a point $m$ of $U$ is lifted by $\sigma$ to a point of $P$, and $g$ moves one in the fibre. It is a diffeomorphism of $U\times G$ onto $\pi^{-1}(U)$, compatible with the projections and equivariant: $h(m,g\bullet g')=h(m,g)\ract g'$. The \textbf{local representation of the connection 1-form} is the pull-back $h^*\omega$, a Lie-algebra valued 1-form on $U\times G$. \ts{22}{1}{13:38} \ts{22}{1}{13:58} \ts{22}{1}{14:19} \ts{22}{1}{14:40} \ts{22}{1}{15:02} \ts{22}{1}{15:24}
\end{definition}
The lecture draws the full diagram: $U\times G\xrightarrow{\ h\ }P$, over it the projection to the first factor (not a principal bundle map, since only part of $P$ is involved, there is the right action of $G$ on the second factor with the primed action, and underneath the identity on $U$). We can pull $\omega$ back along $h$, which gives $h^*\omega$ on the local trivialisation, and along $\sigma$, which gives $\omega_U$ on the base \ts{22}{1}{15:45} \ts{22}{1}{16:06} \ts{22}{1}{16:27}. Which of the two contains all the information of $\omega$ over $U$? One would guess $h^*\omega$, and that is true; but also the Yang--Mills field $\omega_U$, living down on the base, contains the same information over $U$: both contain less than $\omega$ only because they are defined over a local patch \ts{22}{1}{16:52} \ts{22}{1}{17:13} \ts{22}{1}{17:34} \ts{22}{1}{17:55}. In physics one usually writes down only the local Yang--Mills field, in coordinates: the potential $A_\mu$ of electrodynamics, the (Lie-algebra valued) gauge potential of non-abelian gauge theories, the connection coefficients $\Gamma^\mu{}_{\nu\sigma}$ of differential geometry and general relativity: all are local shadows, local representatives, of the abstract everywhere defined connection 1-form, and they are important because they come up all the time \ts{22}{1}{17:55} \ts{22}{1}{18:15} \ts{22}{1}{18:36} \ts{22}{1}{18:57}.

\textbf{The relation} between $h^*\omega$ and $\omega_U$ \ts{22}{1}{19:19}. At a point $(m,g)$ the pull-back $h^*\omega$ takes a tangent vector of $U\times G$, that is (since $T_{(m,g)}(U\times G)=T_mU\oplus T_gG$) a pair $(v,\xi)$ of a tangent vector to the base and a tangent vector to the group, to the Lie algebra. For this one needs the Maurer--Cartan form.
\begin{coursenote}[A slip in the lecture]
While setting up the pull-back $h^*\omega$ the lecturer first writes a push-forward, then a cotangent space, then ``it takes a vector, it takes an element of $U\times G$'' and corrects himself several times (``I am sorry, it is the pull-back, I wrote down some push-forward''); what is meant is: $h_*$ maps $T_{(m,g)}(U\times G)\to T_{h(m,g)}P$, and $\omega$ eats the result. \ts{22}{1}{19:19}
\end{coursenote}
\ts{22}{1}{19:40} \ts{22}{1}{20:12} \ts{22}{1}{20:33} \ts{22}{1}{20:54}
\begin{definition}[New concepts: Maurer--Cartan form]\label{def:ym-mc}
Let $G$ be a Lie group. The \textbf{Maurer--Cartan form} is the Lie-algebra valued 1-form $\Xi$ on $G$ which at $g\in G$ is the map
\[
  \Xi_g\colon T_gG\to T_eG,\qquad\Xi_g(X):=(L_{g^{-1}})_*X ,
\]
that is: every $X\in T_gG$ is the value at $g$ of exactly one left-invariant vector field $L^A$, $L^A_g=(L_g)_*A=X$ (\cref{thm:ll-iso}), and $\Xi_g(X):=A$. Since $\Xi$ eats a tangent vector at a point of the group and yields a Lie algebra element, it is a Lie-algebra valued 1-form \emph{on the group}. \ts{22}{1}{21:15} \ts{22}{1}{21:36} \ts{22}{1}{21:56} \ts{22}{1}{22:19} \ts{22}{1}{23:22} \ts{22}{1}{23:43} \ts{22}{1}{24:03}
\end{definition}
The lecture draws the picture of the group as a sphere: at the identity $e$ the tangent space, at $g$ the tangent space; take an arbitrary vector at $g$; there is always a specific vector $A$ at the identity such that the left-invariant field generated by $A$ reproduces it at $g$. $\Xi_g$ is the inverse of this: ``you take the vector here and you find from which vector it has been produced'' \ts{22}{1}{24:24} \ts{22}{1}{24:44} \ts{22}{1}{25:05} \ts{22}{1}{25:25} \ts{22}{1}{25:47}. Here $L_g$ denotes the left multiplication $\tilde g\mapsto g\bullet\tilde g$, and its push-forward goes $T_eG\to T_gG$, as in \cref{def:ll-push-field} \ts{22}{1}{27:45} \ts{22}{1}{28:06} \ts{22}{1}{28:27}. So we now have three Lie-algebra valued 1-forms at three levels: $\omega$ on the total space of the principal bundle, $\omega_U$ on (part of) the base, and the Maurer--Cartan form on the group; but the last is fixed, a property of the group alone \ts{22}{1}{26:09} \ts{22}{1}{26:29} \ts{22}{1}{26:50}. The Maurer--Cartan form of $GL(d,\R)$ is computed in section~\ref{sec:ym-mc}; all this machinery is needed only because we want to represent the whole thing on the base space \ts{22}{1}{27:23}.
\begin{theorem}[Local form of the connection 1-form]\label{thm:ym-local}
For $(m,g)\in U\times G$, $v\in T_mU$ and $\xi\in T_gG$,
\[
  (h^*\omega)_{(m,g)}(v,\xi)=(\mathrm{Ad}_{g^{-1}})_*\bigl(\omega_U(v)\bigr)+\Xi_g(\xi) .
\]
\ts{22}{1}{21:36} \ts{22}{1}{21:56} \ts{22}{1}{22:19} \ts{22}{1}{22:41} \ts{22}{1}{23:01}
\end{theorem}
The Yang--Mills field $\omega_U$ appears and eats the tangent vector $v$, producing a Lie algebra element on which $(\mathrm{Ad}_{g^{-1}})_*$ is then applied; and an extra term which is the Maurer--Cartan form acting on $\xi$ \ts{22}{1}{22:41} \ts{22}{1}{23:01}.
\begin{proof}
\emph{(The lecture states the relation and explains the pieces, but gives no proof; the argument is added, not said in class.)} Let $A:=\Xi_g(\xi)$, so that $\xi=(L_g)_*A$. Choose a curve $c(t)=g\,\exp(tA)$ with $\dot c(0)=\xi$, and a curve $m(t)$ with $\dot m(0)=v$. Then
\[
  \begin{aligned}
  h(m(t),c(t))&=\sigma(m(t))\ract g\ract\exp(tA)\\
  \tfrac{\mathrm d}{\mathrm dt}\big|_0&=(r_g)_*\sigma_*v+X^A_{h(m,g)} .
  \end{aligned}
\]
Apply $\omega$: the first term gives $(r_g^*\omega)(\sigma_*v)=(\mathrm{Ad}_{g^{-1}})_*\omega_U(v)$ by (b); the second gives $A$ by (a). (For a general $(v,\xi)$ use linearity in $T_mU\oplus T_gG$.)
\end{proof}
(This also shows that $h^*\omega$, hence $\omega$ over $U$, is determined by $\omega_U$ alone: the two contain the same information, as said above.)

% ==================================================================
\section{Example 1: the frame bundle and the connection coefficients}\label{sec:ym-frame}

An example to motivate the choice of a section as ``localisation'' \ts{22}{1}{28:27} \ts{22}{1}{28:48} \ts{22}{1}{29:08}. Consider the frame bundle $LM$ (\cref{def:pb-frame}), where to every point all bases of the tangent space are attached. Every choice of a coordinate chart $(U,x)$ on $M$ induces a section of $LM$,
\[
  \sigma(m):=\Bigl(\frac{\partial}{\partial x^1}\Big|_m,\dots,\frac{\partial}{\partial x^{d}}\Big|_m\Bigr),\qquad d=\dim M,
\]
the frame provided by the coordinates at the point $m$ of the chart domain: choosing coordinates on the base immediately induces a section of the frame bundle \ts{22}{1}{29:28} \ts{22}{1}{29:51} \ts{22}{1}{30:12} \ts{22}{1}{30:32} \ts{22}{1}{30:53} \ts{22}{1}{31:15}. The Yang--Mills field $\omega_U=\sigma^*\omega$ is a 1-form on $U$ with values in the Lie algebra of $GL(d,\R)$, that is in the $d\times d$ matrices \ts{22}{1}{31:45} \ts{22}{1}{32:06}. In components
\[
  (\omega_U)^i{}_{j\mu}=:\Gamma^i{}_{j\mu},
\]
where $\mu=1,\dots,d$ is the 1-form index on the manifold, and $i$ (up), $j$ (down) are indices of the matrix Lie algebra of $GL$, ``the Lie algebra representation, not a $GL$ matrix'' \ts{22}{1}{32:30} \ts{22}{1}{32:51} \ts{22}{1}{33:12}. These are the connection coefficients of differential geometry (of general relativity, if one chooses an orthogonal frame, an extra restriction). All three indices run from $1$ to $d$, and very often one writes $\Gamma^\alpha{}_{\beta\gamma}$ pretending that all three have to do with the base manifold; but they do not: two of them label matrix elements of the Lie algebra, only $\mu$ is a proper 1-form index; \emph{these indices have a different logical status even if they run over the same range} \ts{22}{1}{33:32} \ts{22}{1}{33:53} \ts{22}{1}{34:13}. So this is an example of how Yang--Mills fields come up, by the choice of a section, which for the frame bundle is induced by the choice of coordinates \ts{22}{1}{34:34} \ts{22}{1}{34:55}.

% ==================================================================
\section{Example 2: the Maurer--Cartan form of $GL(d,\R)$}\label{sec:ym-mc}

The Maurer--Cartan form depends only on the group; we construct it for $GL(d,\R)$ (an important case because the Lie groups one meets are very often $GL$ or a subgroup, so that the same procedure works essentially in the same way) \ts{22}{1}{34:55} \ts{22}{1}{35:16} \ts{22}{1}{35:37} \ts{22}{1}{35:57}.

Choose coordinates on a chart domain of the group containing $g$ (the lecturer: on the connected component $GL^+$ of the identity; in fact $GL(d,\R)$ is an open subset of $\R^{d\times d}$, the determinant condition being open, so there is one chart on the whole group): $x^i{}_j\colon GL(d,\R)\to\R$, the $(i,j)$ entry of the matrix $g$ \ts{22}{1}{36:20} \ts{22}{1}{36:41} \ts{22}{1}{37:02} \ts{22}{1}{37:25} \ts{22}{1}{37:46}. A student corrects the target of the chart (it is $\R^{d\times d}$, the single coordinate functions map to $\R$) \ts{22}{1}{38:29} \ts{22}{1}{38:55}.

Consider the left-invariant vector field $L^A$ generated by $A\in T_eG$ (matrices); it exists abstractly (\cref{thm:ll-iso}), and we want to compute it. A vector field acts on functions, in particular on the coordinate functions, and by definition, at the point $g$, $L^A_g\,f=\frac{\mathrm d}{\mathrm dt}f\bigl(g\,\exp(tA)\bigr)\big|_{t=0}$: the curve $t\mapsto g\,\exp(tA)$ through $g$ is the one-parameter subgroup generated by $A$ moved by $g$ (it is the integral curve of the left-invariant field, which is how $\exp$ was defined, \cref{def:rc-exp}) \ts{22}{1}{38:08} \ts{22}{1}{39:16} \ts{22}{1}{39:36} \ts{22}{1}{39:56} \ts{22}{1}{40:17} \ts{22}{1}{40:38} \ts{22}{1}{40:59} \ts{22}{1}{41:19} \ts{22}{1}{41:40}. Then, in the matrix group, the multiplication is matrix multiplication and $\exp(tA)$ is the matrix exponential (\cref{ch:l18}), so
\[
  \begin{aligned}
  L^A_g\,x^i{}_j&=\tfrac{\mathrm d}{\mathrm dt}\big|_0\bigl(g\,e^{tA}\bigr)^i{}_j\\
  &=g^i{}_k\,A^k{}_j ,
  \end{aligned}
\]
since $g$ is a fixed matrix and $\frac{\mathrm d}{\mathrm dt}e^{tA}|_0=A$ \ts{22}{1}{42:00} \ts{22}{1}{42:21} \ts{22}{1}{42:42} \ts{22}{1}{43:02} \ts{22}{1}{43:24} \ts{22}{1}{43:46} \ts{22}{1}{44:06}. A quick check: the result is a real number depending on $i,j$, the $(i,j)$ entry of the matrix product $gA$, as it must be. Hence the components of $L^A_g$ are
\[
  L^A_g=g^i{}_k\,A^k{}_j\ \frac{\partial}{\partial x^i{}_j}\Big|_g ,
\]
all indices summed, in the basis $\partial/\partial x^i{}_j$ of $T_gG$ \ts{22}{1}{46:43} \ts{22}{1}{47:04} \ts{22}{1}{47:24} \ts{22}{1}{47:45}.
\begin{coursenote}[A moment of confusion about the definition of a vector]
At \ts{22}{1}{44:06} the lecturer doubts that the derivative can be taken ``without'' the function and says ``no no, I did not mean it''; what is meant is the definition used from the beginning: the vector tangent to a curve $\gamma$ at $p$ acts on $f$ as $(f\circ\gamma)'(0)$, with $\gamma(0)=p$, here $\gamma(0)=g$ (the brackets belong outside). With this the computation above is correct. \ts{22}{1}{44:27} \ts{22}{1}{44:48} \ts{22}{1}{45:09} \ts{22}{1}{45:42} \ts{22}{1}{46:03} \ts{22}{1}{46:23}
\end{coursenote}
\begin{proposition}[Maurer--Cartan form of $GL(d,\R)$]\label{prop:ym-mc-gl}
In the coordinates $x^i{}_j$,
\[
  \Xi_g=(g^{-1})^i{}_k\,\mathrm dx^k{}_j\ \ \text{(component $(i,j)$)}\qquad\text{at the point }g .
\]
\ts{22}{1}{48:06} \ts{22}{1}{48:27} \ts{22}{1}{48:55} \ts{22}{1}{49:15}
\end{proposition}
\ts{22}{1}{47:45} It is a form, so its components are $\mathrm dx^k{}_j$, and the result lies in the Lie algebra, with indices $(i,j)$ \ts{22}{1}{49:15}.
\begin{proof}
It suffices to check $\Xi_g(L^A_g)=A$ for all $A$, since $\Xi_g$ is defined as the inverse of $A\mapsto L^A_g$. Use the dual bases $\mathrm dx^k{}_j(\partial/\partial x^p{}_q)=\delta^k_p\delta_j^q$ and write the field with the dummy indices $p,r,q$:
\[
  \begin{aligned}
  \Xi_g(L^A_g)^i{}_j&=(g^{-1})^i{}_k\,g^p{}_rA^r{}_q\ \mathrm dx^k{}_j\Bigl(\tfrac{\partial}{\partial x^p{}_q}\Bigr)\\
  &=(g^{-1})^i{}_p\,g^p{}_r\,A^r{}_j=A^i{}_j .
  \end{aligned}
\]
\ts{22}{1}{49:37} \ts{22}{1}{49:59} \ts{22}{1}{50:20} \ts{22}{1}{50:41} \ts{22}{1}{51:02}
\end{proof}
The constants pull out, the $\mathrm dx$ are dual to the $\partial/\partial x$ by construction, giving Kronecker deltas, and the matrix multiplication $g^{-1}g$ gives the identity: ``fantastic, you plug in the left-invariant vector field generated from $A$ and you get back the $A^i{}_j$'' \ts{22}{1}{50:41} \ts{22}{1}{51:02} \ts{22}{1}{51:22}. \emph{The lesson}: to compute the Maurer--Cartan form, compute how the left-invariant field acts on chosen coordinate functions, and read off, killing the factor $g$ in front (``the $L^A$ is $gA$, so you have to get rid of the $g$ in front in order to get the $A$'') \ts{22}{1}{51:43} \ts{22}{1}{52:04}.

So, summarising \ts{22}{1}{52:04} \ts{22}{1}{52:26}: from a connection 1-form on a principal bundle, looking at part of the base manifold and choosing a section $\sigma$, we pull back $\omega$ to the Yang--Mills field $\sigma^*\omega$, a Lie-algebra valued 1-form on that domain \ts{22}{1}{52:47}.

% ==================================================================
\section{Change of section: the gauge map and the transformation of Yang--Mills fields}\label{sec:ym-gauge}

In physics we are often led to write down a Yang--Mills field because we understand the situation locally, and we can then speculate about whether a global connection 1-form can be constructed out of various local pieces: two local pieces $U_1$, $U_2$ with $U_1\cap U_2\ne\emptyset$ and a Yang--Mills field on each. There must be a compatibility condition on the overlap, and this is what we study now \ts{22}{1}{52:47} \ts{22}{1}{53:09} \ts{22}{1}{53:29} \ts{22}{1}{53:50} \ts{22}{1}{54:11}.

So let $\sigma_1,\sigma_2$ be local sections over $U_1,U_2$ (the lecturer first draws non-intersecting domains and then the intersection), $\omega$ a connection 1-form on $P$ with the two Yang--Mills fields $\omega_{U_1}=\sigma_1^*\omega$ and $\omega_{U_2}=\sigma_2^*\omega$ \ts{22}{1}{54:34} \ts{22}{1}{54:54} \ts{22}{1}{55:15}. We want the transition between the two on $U_1\cap U_2$ \ts{22}{1}{55:36}.
\begin{definition}[New concepts: the gauge map $\Omega$]\label{def:ym-gauge}
The \textbf{gauge map} (``gauge function'') is the map $\Omega\colon U_1\cap U_2\to G$ such that, for all $m\in U_1\cap U_2$, $\Omega(m)$ is the unique group element with
\[
  \sigma_2(m)=\sigma_1(m)\ract\Omega(m).
\]
\ts{22}{1}{55:56} \ts{22}{1}{56:17} \ts{22}{1}{56:38} \ts{22}{1}{56:59} \ts{22}{1}{57:20}
\end{definition}
It always exists and is uniquely defined: two points in the same fibre are related by a unique group element, because the right action is free (\cref{def:pb-free}) and the fibres are orbits \ts{22}{1}{56:59} \ts{22}{1}{57:20} (not said in class: $\Omega$ is smooth, because in a local trivialisation it is $\Omega(m)=\chi_{\sigma_1}(\sigma_2(m))$, with the map $\chi$ of \cref{thm:pb-trivial-section}). It intertwines the two choices of sections.
\begin{theorem}[Transformation of Yang--Mills fields]\label{thm:ym-gauge}
Under the above assumptions, on $U_1\cap U_2$,
\[
  \omega_{U_2}=(\mathrm{Ad}_{\Omega^{-1}})_*\,\omega_{U_1}+\Omega^*\,\Xi ,
\]
that is, for $v\in T_mM$ with $m\in U_1\cap U_2$: $(\omega_{U_2})_m(v)=(\mathrm{Ad}_{\Omega(m)^{-1}})_*\bigl(\omega_{U_1}(v)\bigr)+\Xi_{\Omega(m)}\bigl(\Omega_{*}v\bigr)$. \ts{22}{1}{57:41} \ts{22}{1}{58:03} \ts{22}{1}{58:24} \ts{22}{1}{58:47} \ts{22}{1}{59:09}
\end{theorem}
``It looks similar to the other formulas, in an entirely different formula'': it says how one connection, think of the connection coefficients, is expressed in the other frames; the first term is the adjoint action with $\Omega^{-1}$ on the first Yang--Mills field, pushed forward, the second the pull-back of the Maurer--Cartan form along $\Omega$ \ts{22}{1}{59:09} \ts{22}{1}{59:31}. If one wants one can make explicit that everything is a 1-form at the point $m$ \ts{22}{1}{59:09}.
\begin{coursenote}[The proof of \cref{thm:ym-gauge}]
The lecturer announces ``before we come to a proof of this, let us work it out again for the $GL$ case'' (\cref{sec:ym-gl}) and does not come back to it. The proof is given here, not said in class. \ts{22}{1}{59:55}
\end{coursenote}
\begin{proof}
\emph{(Added, not said in class.)} Let $v=\dot m(0)$. By the definition of $\Omega$, $\sigma_2(m(t))=\sigma_1(m(t))\ract\Omega(m(t))$, so
\[
  \sigma_{2*}v=(r_{\Omega(m)})_*\sigma_{1*}v+X^{A}_{\sigma_2(m)}
\]
with $A=\Xi_{\Omega(m)}(\Omega_*v)$ (the same computation as in the proof of \cref{thm:ym-local}). Apply $\omega$: by (b) and (a)
\[
  \omega_{U_2}(v)=(\mathrm{Ad}_{\Omega^{-1}})_*\omega_{U_1}(v)+A .
\]
\end{proof}
The extra term does not depend on $\omega$ at all, only on the group and the gauge map. Since the Maurer--Cartan form of an \emph{abelian} group is just $\mathrm d$ of the coordinates (and $\mathrm{Ad}$ is trivial), for $U(1)$ this reads $\omega_{U_2}=\omega_{U_1}+\Omega^*\Xi$, the familiar gauge transformation of the potential in electrodynamics (not said in class; the lecture mentions only that the expression ``must have been seen before with gauge transformations in physics'').

% ==================================================================
\section{The case $G=GL(d,\R)$: the transformation law of the $\Gamma$ symbols}\label{sec:ym-gl}

``Before we come to a proof, let us work it out for the $GL$ case, because that is usually the best way to understand this, since it is rather formalistic'' \ts{22}{1}{59:55} \ts{22}{1}{1:00:23}. We need two things: the pull-back $\Omega^*\Xi$ and $(\mathrm{Ad}_{\Omega^{-1}})_*$ acting on $\omega_{U_1}$, for $G=GL(d,\R)$ and the frame bundle \ts{22}{1}{1:00:43} \ts{22}{1}{1:01:04}.

\textbf{The pull-back of the Maurer--Cartan form} \ts{22}{1}{1:01:04} \ts{22}{1}{1:01:24}. $\Omega$ goes from $U_1\cap U_2$ to the group, and $\Xi$ takes $T_gG$ to $T_eG$, so the pull-back $\Omega^*\Xi$ is a Lie-algebra valued 1-form on $U_1\cap U_2$ (the lecturer first begins to say that $\Xi$ acts on $\partial/\partial x^{pq}$ ``at the point $p$'' and calls it nonsense: it acts on tangent vectors of the group, and the pull-back acts on tangent vectors of the base) \ts{22}{1}{1:01:48} \ts{22}{1}{1:02:09} \ts{22}{1}{1:02:29} \ts{22}{1}{1:02:49} \ts{22}{1}{1:03:11}. It suffices to know its action on every tangent vector $\partial/\partial x^\mu$ of the base: by the definition of the pull-back
\[
  (\Omega^*\Xi)^i{}_j\Bigl(\tfrac{\partial}{\partial x^\mu}\Bigr)=\Xi^i{}_j\Bigl(\Omega_*\tfrac{\partial}{\partial x^\mu}\Bigr)
\]
at the point $\Omega(m)$ \ts{22}{1}{1:03:32} \ts{22}{1}{1:03:53} \ts{22}{1}{1:04:15}. By \cref{prop:ym-mc-gl}, $\Xi=(\Omega^{-1})^i{}_k\,\mathrm dx^k{}_j$ at that point, and $\mathrm dx^k{}_j$ acts on the vector $\Omega_*\partial_\mu$ (a tangent vector to the group, since $\Omega$ lands in the group) by the action of the vector on the function $x^k{}_j$; by the definition of the push-forward this is $\partial_\mu(x^k{}_j\circ\Omega)$, and $x^k{}_j\circ\Omega$ is just the $(k,j)$ entry of the matrix $\Omega$ \ts{22}{1}{1:04:35} \ts{22}{1}{1:04:56} \ts{22}{1}{1:05:16} \ts{22}{1}{1:05:45} \ts{22}{1}{1:06:06} \ts{22}{1}{1:06:27} \ts{22}{1}{1:06:54} \ts{22}{1}{1:07:14} \ts{22}{1}{1:07:35} \ts{22}{1}{1:07:56} \ts{22}{1}{1:08:16}. Therefore
\[
  (\Omega^*\Xi)^i{}_{j\mu}=(\Omega^{-1})^i{}_k\ \partial_\mu\Omega^k{}_j ,
\]
\ts{22}{1}{1:08:38} \ts{22}{1}{1:08:59} \ts{22}{1}{1:09:21} \ts{22}{1}{1:09:41} \ts{22}{1}{1:10:03} and, in matrix notation (the product is matrix multiplication; $\Omega$ is a matrix, and the second factor is $\mathrm d$ of the matrix),
\[
  \Omega^*\Xi=\Omega^{-1}\,\mathrm d\Omega .
\]
``You must have seen this before with gauge transformations in physics'' \ts{22}{1}{1:10:24} \ts{22}{1}{1:10:45}. It looks as if there is no more trace of the Maurer--Cartan form, but this is the very special structure: $\Omega$ is just the gauge map interpolating between the two chosen sections \ts{22}{1}{1:11:05}. (Student question: all indices run to $d$; the lecturer: yes, $\Omega$ is a matrix, one takes the inverse, a matrix product, and the result has indices $i,j$ \ts{22}{1}{1:11:25} \ts{22}{1}{1:11:47}.)

\textbf{The adjoint term} \ts{22}{1}{1:11:47} \ts{22}{1}{1:12:08} \ts{22}{1}{1:12:31}. Recall $\mathrm{Ad}_g\colon G\to G$, $h\mapsto g\bullet h\bullet g^{-1}$, an action on the group, which fixes $e$, so that its push-forward sends $T_eG$ to $T_eG$, the Lie algebra to itself \ts{22}{1}{1:12:52} \ts{22}{1}{1:13:14} \ts{22}{1}{1:13:38}. In practice, for a matrix group,
\[
  (\mathrm{Ad}_g)_*A=g\,A\,g^{-1}
\]
with matrix multiplication on the right, a tangent vector being the derivative of a curve $t\mapsto e^{tA}$ in the group, and the push-forward of the vector being the derivative of the image curve $g\,e^{tA}g^{-1}$: the only things affected are the factors that stay in front and behind \ts{22}{1}{1:13:58} \ts{22}{1}{1:14:18} \ts{22}{1}{1:14:40} \ts{22}{1}{1:15:03} \ts{22}{1}{1:15:24} \ts{22}{1}{1:15:44}.
\begin{aside}[A warning about the star]
The lecturer insists: on the left the product is the group multiplication, on the right it is matrix multiplication; it only looks as if one could plug in another matrix, an algebra element rather than a group element, but formally $\mathrm{Ad}_g$ acts on the group and what acts on the algebra is its \emph{push-forward}. In many physics books the star is forgotten, writing $\mathrm{Ad}_gA=gAg^{-1}$, which is ``terribly misleading'', and one starts confusing the two ``only because of this numerical coincidence''. \ts{22}{1}{1:16:05} \ts{22}{1}{1:16:26}
\end{aside}
\textbf{Together} \ts{22}{1}{1:16:26} \ts{22}{1}{1:16:49}: for the frame bundle, $G=GL(d,\R)$, on $U_1\cap U_2$ the transition rule between two Yang--Mills fields on the same domain is
\[
  \omega_{U_2}=\Omega^{-1}\,\omega_{U_1}\,\Omega+\Omega^{-1}\,\mathrm d\Omega ,
\]
in components (the matrix indices $i,j$, the form index $\mu$)
\[
  \begin{aligned}
  (\omega_{U_2})^i{}_{j\mu}&=(\Omega^{-1})^i{}_k\,(\omega_{U_1})^k{}_{l\mu}\,\Omega^l{}_j\\
  &\quad+(\Omega^{-1})^i{}_k\,\partial_\mu\Omega^k{}_j .
  \end{aligned}
\]
\ts{22}{1}{1:17:12} \ts{22}{1}{1:17:32} \ts{22}{1}{1:17:52} \ts{22}{1}{1:18:13} \ts{22}{1}{1:18:34} \ts{22}{1}{1:18:54} \ts{22}{1}{1:19:15}
\begin{coursenote}[Index placement]
The captions give the placement of the $\Omega$'s ambiguously (``$\Omega^{ik}$ acting on the first Yang--Mills field with its tangent space index and its algebra index $kl$, then again $\Omega$, with respect to $\Omega^{-1}$ [and so the inverse is there]''); the version above is the one that follows from $(\mathrm{Ad}_{\Omega^{-1}})_*A=\Omega^{-1}A\,\Omega$ and from the Maurer--Cartan form of \cref{prop:ym-mc-gl} [?: reconstructed]. \ts{22}{1}{1:18:13}
\end{coursenote}
\textbf{How one uses this} \ts{22}{1}{1:19:36}. This is how one Yang--Mills field must transform into the other on the same domain. In practice one will rather have a Yang--Mills field here, here and here, and if they are to stem from the same connection 1-form on the principal bundle one needs to \emph{require} this as a transition function. This is the coordinate approach to connections: one has lots of connection coefficients, going from one to the other they must transform as above, and then there is one global connection on the principal bundle which corresponds to them \ts{22}{1}{1:19:36} \ts{22}{1}{1:19:57} \ts{22}{1}{1:20:18}. (The formula is entirely open as to which sections $\sigma_1,\sigma_2$ were chosen: it holds for any such sections, each Yang--Mills field being a pull-back of the global $\omega$.)

% ==================================================================
\section{Sections induced by coordinates: the $\Gamma$ symbols}\label{sec:ym-gamma}

The formula known for the transformation of connection coefficients is the one where the sections are induced by coordinates \ts{22}{1}{1:20:40} \ts{22}{1}{1:21:00}. So specialise: $\sigma_1\colon U_1\to LM$ and $\sigma_2\colon U_2\to LM$ the sections induced by two charts, coordinates $x^1,\dots,x^d$ on $U_1$ and $y^1,\dots,y^d$ on $U_2$, that is the frames $\partial/\partial x^i$ and $\partial/\partial y^j$ \ts{22}{1}{1:21:20} \ts{22}{1}{1:21:45} \ts{22}{1}{1:22:07}. Then the gauge map $\Omega$ is the Jacobian matrix: since $\frac{\partial}{\partial y^j}=\frac{\partial x^i}{\partial y^j}\frac{\partial}{\partial x^i}$ and $(e\ract g)_j=g^i{}_je_i$ (\cref{prop:pb-frame-principal}),
\[
  \Omega^i{}_j=\frac{\partial x^i}{\partial y^j},\qquad(\Omega^{-1})^i{}_j=\frac{\partial y^i}{\partial x^j}
\]
(the lecturer first tries the other way around, ``we need to check whether it is the right way around'', and settles on the one above) \ts{22}{1}{1:22:28} \ts{22}{1}{1:22:49} \ts{22}{1}{1:23:11}. In the transition formula above both sides still have the same $\mu$, a component with respect to the same coordinate basis: one wants to express the second Yang--Mills field also in terms of the new coordinates, so one needs yet another Jacobian, for the form index, $\partial x^\mu/\partial y^\nu$, in front of the remaining formula \ts{22}{1}{1:23:11} \ts{22}{1}{1:23:32} \ts{22}{1}{1:24:00}. The result (the lecturer checks the placement of $x$ and $y$ several times and corrects it) is
\[
  \begin{aligned}
  (\omega_{U_2})^i{}_{j\nu}&=\frac{\partial y^i}{\partial x^k}\frac{\partial x^l}{\partial y^j}\frac{\partial x^\mu}{\partial y^\nu}(\omega_{U_1})^k{}_{l\mu}\\
  &\quad+\frac{\partial y^i}{\partial x^k}\frac{\partial^2x^k}{\partial y^\nu\,\partial y^j} .
  \end{aligned}
\]
\ts{22}{1}{1:24:22} \ts{22}{1}{1:24:44} \ts{22}{1}{1:25:05} \ts{22}{1}{1:25:26} \ts{22}{1}{1:25:47} With $(\omega_U)^i{}_{j\mu}=\Gamma^i{}_{j\mu}$ this is precisely the well-known transformation formula of the $\Gamma$ symbols: each index transforms like a tensor, plus an extra (inhomogeneous) term \ts{22}{1}{1:26:07}. The apparent absence of one transformation matrix in the general formula (three indices, two Jacobians) is because there one did not assume that the $\Omega$'s have anything to do with the coordinate transformation; in the special case everything is linked and the third Jacobian is the one for the form index \ts{22}{1}{1:26:28} \ts{22}{1}{1:26:48}.

\textbf{What do we see?} \ts{22}{1}{1:26:48} \ts{22}{1}{1:27:08} The $\Gamma$ symbols of standard differential geometry (for a connection, whether one comes from a covariant derivative or not) transform in a funny, arcane way, and it is always emphasised that they do not transform as a $(1,2)$ tensor (one index up, two down, not just because of that). But this complicated object, the Yang--Mills field, comes from a very simple geometric object: a global 1-form, which is tensorial, not a 1-form on the base manifold like the Yang--Mills field but a 1-form on the principal bundle \ts{22}{1}{1:27:29} \ts{22}{1}{1:27:49}. On the principal bundle a connection is something very simple, a Lie-algebra valued 1-form with two properties; if one pushes it down to the base manifold by force, it acquires this complicated transformation behaviour. That may be where physics starts, because one looks at a patch and glues the patches obeying this rule; but then one grows up to the principal bundle and everything becomes simple in terms of the connection 1-form of \cref{ch:l21}, itself just a convenient encoding of a choice of horizontal subspaces, ``which is the actual idea of connections'' \ts{22}{1}{1:28:10} \ts{22}{1}{1:28:30}.

``In principle everything we did today we only did because we had to pay tribute to what we usually write down in physics.'' For certain calculations it may be useful, and it must be done; but next time, when we discuss parallel transport on the principal bundle and the associated bundles, and then the covariant derivative on vector bundles, we will not need this any more: we can use it, but we will go back to the nicely defined Lie-algebra valued 1-form on the principal bundle \ts{22}{1}{1:28:50} \ts{22}{1}{1:29:11}.

% ==================================================================
\section*{Key concepts}
\begin{keyconcepts}
\textbf{Retrieval questions.}
\begin{enumerate}
  \item Why does the pull-back of a Lie-algebra valued 1-form make sense? What does it not depend on?
  \item Define the Yang--Mills field $\omega_U=\sigma^*\omega$ and the local trivialisation $h(m,g)=\sigma(m)\ract g$. State the relation between $h^*\omega$ and $\omega_U$ (with a proof).
  \item What is the Maurer--Cartan form? Compute it for $GL(d,\R)$ and check it on $L^A$.
  \item Which section of the frame bundle does a chart induce? Which indices of $\Gamma^i{}_{j\mu}$ are manifold indices?
  \item Define the gauge map $\Omega$ and state the transformation law of the Yang--Mills fields (with a proof). Why does it hold for any two sections?
  \item Derive, for $GL(d,\R)$, $\Omega^*\Xi=\Omega^{-1}\mathrm d\Omega$ and $(\mathrm{Ad}_g)_*A=gAg^{-1}$; why does the star matter?
  \item Derive the transformation of the $\Gamma$ symbols from the gauge formula with $\Omega=\partial x/\partial y$.
\end{enumerate}
\textbf{Concepts.}
\begin{itemize}
  \item Yang--Mills field and local trivialisation (\cref{def:ym-field,def:ym-triv,thm:ym-local}).
  \item Maurer--Cartan form (\cref{def:ym-mc,prop:ym-mc-gl}).
  \item Gauge map and transformation law (\cref{def:ym-gauge,thm:ym-gauge}), the $GL$ case and the $\Gamma$ symbols (\cref{sec:ym-gl,sec:ym-gamma}).
\end{itemize}
\end{keyconcepts}

\flushside
\end{paracol}

% !TEX root = ../main.tex
% STUB: generated by tools/stubs.py; write the chapter from materials/transcripts/C23-V01.txt
\chapter{Parallel transport}\label{ch:l23}
\begin{paracol}{2}

\begin{sources}
\begin{tabularx}{\linewidth}{@{}l l L@{}}
\toprule
\textbf{Video} & \textbf{Length} & \textbf{Content}\\
\midrule
\seg{23}{1}{Lecture 23: Parallel transport} & 1:44:32 & (to be written)\\
\bottomrule
\end{tabularx}

\smallskip
\textbf{Transcript} \file{materials/transcripts/C23-V01.txt}.
\end{sources}

\begin{reading}[Where this lecture sits in the other notes]
(to be written)
\end{reading}

\emph{This chapter has not been written yet.}

\flushside
\end{paracol}

% !TEX root = ../main.tex
% STUB: generated by tools/stubs.py; write the chapter from materials/transcripts/C24-V01.txt
\chapter{Curvature and torsion on principal bundles}\label{ch:l24}
\begin{paracol}{2}

\begin{sources}
\begin{tabularx}{\linewidth}{@{}l l L@{}}
\toprule
\textbf{Video} & \textbf{Length} & \textbf{Content}\\
\midrule
\seg{24}{1}{Lecture 24: Curvature and torsion on principal bundles} & 1:16:10 & (to be written)\\
\bottomrule
\end{tabularx}

\smallskip
\textbf{Transcript} \file{materials/transcripts/C24-V01.txt}.
\end{sources}

\begin{reading}[Where this lecture sits in the other notes]
\emph{Gravity and Light}, Lecture 8. (From the lecture titles; to be checked when the chapter is written.)
\end{reading}

\emph{This chapter has not been written yet.}

\flushside
\end{paracol}

% !TEX root = ../main.tex
% STUB: generated by tools/stubs.py; write the chapter from materials/transcripts/C25-V01.txt
\chapter{Covariant derivatives}\label{ch:l25}
\begin{paracol}{2}

\begin{sources}
\begin{tabularx}{\linewidth}{@{}l l L@{}}
\toprule
\textbf{Video} & \textbf{Length} & \textbf{Content}\\
\midrule
\seg{25}{1}{Lecture 25: Covariant derivatives} & 1:16:36 & (to be written)\\
\bottomrule
\end{tabularx}

\smallskip
\textbf{Transcript} \file{materials/transcripts/C25-V01.txt}.
\end{sources}

\begin{reading}[Where this lecture sits in the other notes]
\emph{Gravity and Light}, Lectures 7--8; \emph{Theoretische Mechanik}, Lecture 6. (From the lecture titles; to be checked when the chapter is written.)
\end{reading}

\emph{This chapter has not been written yet.}

\flushside
\end{paracol}

% !TEX root = ../main.tex
% STUB: generated by tools/stubs.py; write the chapter from materials/transcripts/C26-V01.txt
\chapter{Application: Quantum mechanics on curved spaces}\label{ch:l26}
\begin{paracol}{2}

\begin{sources}
\begin{tabularx}{\linewidth}{@{}l l L@{}}
\toprule
\textbf{Video} & \textbf{Length} & \textbf{Content}\\
\midrule
\seg{26}{1}{Lecture 26: Application: Quantum mechanics on curved spaces} & 1:32:16 & (to be written)\\
\bottomrule
\end{tabularx}

\smallskip
\textbf{Transcript} \file{materials/transcripts/C26-V01.txt}.
\end{sources}

\begin{reading}[Where this lecture sits in the other notes]
(to be written)
\end{reading}

\emph{This chapter has not been written yet.}

\flushside
\end{paracol}

% !TEX root = ../main.tex
% STUB: generated by tools/stubs.py; write the chapter from materials/transcripts/C27-V01.txt
\chapter{Application: Spin structures}\label{ch:l27}
\begin{paracol}{2}

\begin{sources}
\begin{tabularx}{\linewidth}{@{}l l L@{}}
\toprule
\textbf{Video} & \textbf{Length} & \textbf{Content}\\
\midrule
\seg{27}{1}{Lecture 27: Application: Spin structures} & 1:39:15 & (to be written)\\
\bottomrule
\end{tabularx}

\smallskip
\textbf{Transcript} \file{materials/transcripts/C27-V01.txt}.
\end{sources}

\begin{reading}[Where this lecture sits in the other notes]
(to be written)
\end{reading}

\emph{This chapter has not been written yet.}

\flushside
\end{paracol}

% !TEX root = ../main.tex
% STUB: generated by tools/stubs.py; write the chapter from materials/transcripts/C28-V01.txt
\chapter{Application: Kinematical and dynamical symmetries}\label{ch:l28}
\begin{paracol}{2}

\begin{sources}
\begin{tabularx}{\linewidth}{@{}l l L@{}}
\toprule
\textbf{Video} & \textbf{Length} & \textbf{Content}\\
\midrule
\seg{28}{1}{Lecture 28: Application: Kinematical and dynamical symmetries} & 1:32:48 & (to be written)\\
\bottomrule
\end{tabularx}

\smallskip
\textbf{Transcript} \file{materials/transcripts/C28-V01.txt}.
\end{sources}

\begin{reading}[Where this lecture sits in the other notes]
(to be written)
\end{reading}

\emph{This chapter has not been written yet.}

\flushside
\end{paracol}

% END CHAPTERS

\appendix
% !TEX root = ../main.tex
\chapter{Pseudo-exam and second reading}\label{app:exam}
\begin{paracol}{2}

The long exercises of each chapter, moved here verbatim so that a first reading of the chapters is not held
up by them. Use this appendix as a second reading or, once the course is finished, as a cumulative
self-test: retrieval spaced weeks after first study is among the best-supported ways to make material stick
(Cepeda et al.\ 2006; Adesope et al.\ 2017).

% ------------------------------------------------------------------
\section{Logic of propositions and predicates (\cref{ch:logic})}\label{exam:ch01}

\begin{exercise}[NAND is functionally complete]\label{ex:lo-nand-full}
Show that every unary and every binary operator of propositional logic can be written using only the
NAND operator $\uparrow$. Start with $\neg p$, then $p\wedge q$, $p\vee q$, $p\Rightarrow q$ and
$p\Leftrightarrow q$; then argue that every operator of any order $n$ can be built from these.
\end{exercise}

% ------------------------------------------------------------------
\section{Classification of sets (\cref{ch:l03})}\label{exam:ch03}

\begin{exercise}[Well-definedness of the addition on $\mathbb Z$]\label{ex:cl-wd-Z}
On $\N\times\N$ let $(m,n)\sim(p,q):\Leftrightarrow m+q=p+n$. Show that $\sim$ is an equivalence relation (transitivity
included) and that $[(m,n)]+_{\mathbb Z}[(p,q)]:=[(m+p,\,n+q)]$ does not depend on the chosen representatives. Check that
$[(0,0)]$ is neutral and that $[(n,0)]+_{\mathbb Z}[(0,n)]=[(0,0)]$ for every $n\in\N$. Define a multiplication on $\mathbb Z$ by
inheriting the one of $\N$, and show that it is well defined (see \cref{def:cl-Z,def:cl-addZ}).
\end{exercise}

\begin{exercise}[The rational numbers and the fool's addition]\label{ex:cl-wd-Q}
On $\mathbb Z\times\mathbb Z^*$ let $(x,y)\sim(u,v):\Leftrightarrow xv=uy$. Show that $\sim$ is an equivalence relation, and that
$[(x,y)]+_\Q[(u,v)]:=[(xv+yu,\,yv)]$ and $[(x,y)]\cdot_\Q[(u,v)]:=[(xu,\,yv)]$ are well defined (note that $yv\ne0$). Show that
$[(x,y)]\,\tilde+\,[(u,v)]:=[(x+u,\,y+v)]$ is \emph{not} well defined by finding representatives that give different results
(see \cref{def:cl-Q,def:cl-opsQ,rem:cl-fool}). Show that $x\mapsto[(x,1)]$ is injective and compatible with the operations.
\end{exercise}

% ------------------------------------------------------------------
\section{Topological spaces: construction and purpose (\cref{ch:l04})}\label{exam:ch04}

\begin{exercise}[The topology inherited by a quotient set]\label{ex:to-quotient-full}
Let $(M,\mathcal O)$ be a topological space and $\sim$ an equivalence relation on $M$ (see \cref{def:cl-equiv,def:cl-quotient}). Define on the
quotient set $M/{\sim}$ a topology inherited from $\mathcal O$ (the problem-sheet exercise of the lecture does not give the definition: it
asks to find it), and show that it is a topology. Apply it to $M=\R$ with the standard topology and $x\sim y:\Leftrightarrow y-x\in2\pi\mathbb Z$, the
circle as a quotient set (see \cref{ex:to-quotient}).
\end{exercise}

% ------------------------------------------------------------------
\section{Topological spaces: some heavily used invariants (\cref{ch:l05})}\label{exam:ch05}

\begin{exercise}[Limits in Hausdorff spaces and the Zariski topology]\label{ex:iv-limit-full}
Prove that in a Hausdorff space every convergent sequence has a unique limit point (\cref{thm:iv-limit}). Show that the chaotic topology on a set with at least two
points is not $T_1$. Look up the Zariski topology and show that it is $T_1$ but not $T_2$ (\cref{ex:iv-zariski}).
\end{exercise}

\begin{exercise}[{$[0,1]$} is connected]\label{ex:iv-interval-full}
Prove that $[0,1]$ with the topology induced from the standard topology of $\R$ is connected (\cref{thm:iv-interval}); not proved in class, ``a little work''.
\end{exercise}

\begin{exercise}[The fundamental group is a group]\label{ex:iv-group-full}
Show that homotopy is an equivalence relation on curves with fixed end points (transitivity included), that $[\gamma]\bullet[\delta]:=[\gamma*\delta]$ is well defined
(independent of the representatives), associative, has the class of the constant loop as neutral element and $[\gamma\circ(1-\cdot)]$ as inverse
(\cref{def:iv-pi1,thm:iv-group,ex:iv-wd}).
\end{exercise}

% ------------------------------------------------------------------
\section{Topological manifolds and manifold bundles (\cref{ch:l06})}\label{exam:ch06}

\begin{exercise}[The Möbius strip is locally trivial but not trivial]\label{ex:mf-moebius-full}
For the Möbius strip as a bundle over $S^1$ with fibre $[-1,1]$ (\cref{ex:mf-moebius}) find open sets $U\subseteq S^1$ over which the restricted bundle is isomorphic to a product bundle, so that it is locally trivial, and show that the bundle
is not trivial (\cref{def:mf-trivial,ex:mf-cyl-moeb}). Write a section locally as a map $U\to[-1,1]$ (\cref{ex:mf-loctriv}).
\end{exercise}

\begin{exercise}[Pulled-back sections]\label{ex:mf-pullback-full}
Let $(E,\pi,M)$ be a bundle with a section $\sigma$ and $f\colon M'\to M$ a map. Show that $\sigma'(m'):=(m',\sigma(f(m')))$ is a section of the pull-back bundle, that it is continuous if $f$ and $\sigma$ are, and that it is the only
section $\sigma'$ with $u\circ\sigma'=\sigma\circ f$ (\cref{def:mf-pullback,prop:mf-pullback-section,ex:mf-pullback-sec}).
\end{exercise}

\begin{exercise}[Continuity of maps and curves through charts]\label{ex:mf-cont-charts-full}
Let $M,N$ be topological manifolds and $\phi\colon M\to N$ a map. Show that if $\phi$ is continuous, then $y\circ\phi\circ x^{-1}$ is continuous for all charts $(U,x)$ of $M$ and $(V,y)$ of $N$ with $\phi(U)\subseteq V$; and conversely that $\phi$ is continuous if every point of $M$ lies in the domain
of a chart $(U,x)$ for which there is a chart $(V,y)$ of $N$ with $\phi(U)\subseteq V$ and $y\circ\phi\circ x^{-1}$ continuous [?: the precise statement of the problem sheet is not in the lecture]. Deduce the criterion for a curve $\gamma\colon\R\to M$ of \cref{sec:mf-curve}
(\cref{ex:mf-cont-charts}).
\end{exercise}

\flushside
\end{paracol}


\end{document}
